\documentclass[11pt,reqno]{amsart}

\input{preamble/packages}
\input{preamble/macros}
\input{preamble/theorem}
\renewcommand{\thepart}{\Roman{part}}

\title{The topology of Gromov--Hausdorff space}
\author{Yoshito Ishiki}
\address{Department of Mathematical Sciences\\
Tokyo Metropolitan University\\
Minami-osawa, Hachioji, Tokyo 192-0397, Japan}
\email{ishiki-yoshito@tmu.ac.jp}

\date{}
\keywords{Gromov--Hausdorff space, Hilbert space, absolute retract,
probability measure, hyperspace}

\begin{document}
\raggedbottom

\begin{abstract}
\paraid{p-9ef91cc774d8}
We prove that the  Gromov--Hausdorff space is homeomorphic to the  Hilbert space. We construct a continuous assignment of full-support probability measures that is equivariant under isometries and finite-dimensional local approximations that control all pairwise distances. These approximations yield the absolute retract property for all metrizable spaces. We also prove that any countable family of continuous maps from compact metrizable spaces can be approximated, with respect to a prescribed open cover, by maps whose images form a discrete family.
\end{abstract}

\maketitle
\tableofcontents

\section{Introduction}\label{sec:introduction}
\paraid{p-a26667c35d82}
The Gromov--Hausdorff distance
$\GHDistance$
quantifies the difference in shape between
two compact metric spaces.
Let
$\GHSpace$
denote the space of isometry classes of nonempty compact metric spaces,
equipped with the Gromov--Hausdorff distance
$\GHDistance$.
We call
the space
$(\GHSpace, \GHDistance)$
the \emph{Gromov--Hausdorff space}.
Antonyan
\cite[p.~2]{Antonyan2020Euclidean}
asked whether
$\GHSpace$
 is homeomorphic to
$\SeparableHilbertSpace$.
We answer this question affirmatively by proving the following classification.

\begin{theorem}[\Cref{thm:part-iv-main}]\label{thm:main}
\paraid{p-f4cb7088c681}
The space
$\GHSpace$
is homeomorphic to the Hilbert space
$\SeparableHilbertSpace$.
\end{theorem}

\paraid{p-414311e7f43d}
Zava
\cite[arXiv v3, Theorem~B]{Zava2025Coarse}
proves that the subspace of
$\GHSpace$
consisting of finite metric spaces cannot be coarsely embedded into any
Hilbert space.
In particular,
no homeomorphism in \Cref{thm:main} can be bi-Lipschitz.

\medskip
\noindent\textbf{Related work.}
\paraid{p-fb73239bb9f0}
For comparison with \Cref{thm:main},
we first recall a classification of subspaces of
$\GHSpace$.
For each positive integer
$\FiniteSize $,
Antonyan identifies the subspace of
$\GHSpace$
represented by compact subsets of
$\RealNumbers^\FiniteSize $
with the Hilbert cube minus a point
\cite[Corollary~5.3]{Antonyan2021Euclidean}.

\paraid{p-31fbde650eca}
Geodesics and optimal correspondences describe the metric structure
underlying our topological classification of
$\GHSpace$.
Ivanov,
Nikolaeva,
and Tuzhilin
\cite[arXiv v1, Theorem~1]{IvanovNikolaevaTuzhilin2015}
prove that any two points of
$\GHSpace$
can be joined by a geodesic.
Ivanov,
Iliadis,
and Tuzhilin
\cite[arXiv v1, Theorem~2.6 and Corollary~2.7]{IvanovIliadisTuzhilin2016}
prove the existence of optimal correspondences between compact metric spaces
and use them to attain the Gromov--Hausdorff distance as the Hausdorff distance
between isometric copies in a common metric space.
A further point of comparison for our classification by homeomorphisms
is rigidity under surjective isometries.
Ivanov and Tuzhilin
\cite[Main Theorem]{IvanovTuzhilin2019Isometries}
prove that every surjective isometry of
$\GHSpace$
is the identity.
Tuzhilin surveys Hausdorff and Gromov--Hausdorff distance geometry in
\cite{Tuzhilin2020}.

\paraid{p-9f8f911356f3}
Earlier work studies the topology of
$\GHSpace$
through continuous families and embedded compact spaces.
Our result identifies the topological type of
$\GHSpace$
itself.
The author constructed families of branching Gromov--Hausdorff geodesics
continuously parametrized by the Hilbert cube
\cite[Theorem~1.3]{Ishiki2022Branching}.
The author also constructed topological embeddings of arbitrary compact metrizable
spaces into the subspaces of
$\GHSpace$
consisting of continua
\cite[Theorem~1.1]{Ishiki2022Continua},
spaces with prescribed dimensions
\cite[Theorem~1.3]{Ishiki2023FractalDimensions},
and compact metric trees
\cite[Theorem~1.1]{Ishiki2023MetricTrees}.
In each case,
the embeddings can be chosen with finitely many distinct prescribed values.
Byakuno
\cite[arXiv v1, Corollary~1.4]{Byakuno2026Embedding}
proves that countable products of closed intervals with positive summable lengths,
equipped with the supremum metric,
admit isometric embeddings into
$\GHSpace$.

\paraid{p-6f8cb7cfa910}
For comparison with our classification of the full space
$\GHSpace$,
we recall local descriptions and dimension formulas for its subspaces
of finite metric spaces.
Iliadis,
Ivanov,
and Tuzhilin
\cite[arXiv v1, Theorem~4.1]{IliadisIvanovTuzhilin2017Local}
prove that sufficiently small neighborhoods of generic finite metric spaces,
within the subspace of spaces with the same number of points,
are isometric to open subsets of a finite-dimensional space with the maximum norm.
Here generic means that the nonzero distances are pairwise distinct
and all triangle inequalities involving three distinct points are strict.
Nakajima,
Yamauchi,
and Zava
\cite[arXiv v1, Theorem~A]{NakajimaYamauchiZava2025}
prove that the subspace of spaces with at most
$\FiniteSize$
points has topological dimension
$\FiniteSize(\FiniteSize-1)/2$.
A complementary approach to the global topology of
$\GHSpace$
studies its compactifications.
Nakajima and Shioya
\cite[arXiv v1, Main Theorem~1.2]{NakajimaShioyaCompactification}
construct a compact metrizable space containing a dense topological copy of
$\GHSpace$.

\paraid{p-d8629ae38ad9}
Part~\ref{part:measures} selects probability measures continuously over
$\GHSpace$.
Topological results for metric measure spaces provide a comparison when
the measure is included as part of the data.
Kazukawa,
Nakajima,
and Shioya
\cite[arXiv v1, Theorems~1.4--1.6]{KazukawaNakajimaShioya2024}
prove that the space of metric measure spaces is contractible and locally
path connected both in the box topology and in the concentration topology.
Our classification concerns compact metric spaces.
An extension of the Gromov--Hausdorff framework treats noncomplete spaces
together with their boundary.
Shibahara
\cite[arXiv v1, Definition~3.2 and Theorem~1.1]{Shibahara2021Boundary}
introduces a Gromov--Hausdorff metric with boundary on the isometry classes
of noncomplete precompact locally compact metric spaces.
The boundary is the complement of the space in its metric completion.
Our local models in Part~\ref{part:spectral} keep track of isometries
between representatives.
Isometries between individual spaces are also retained by the moduli stack
used by
Yuji
\cite[arXiv v2, Definition~5.1]{Yuji2026ModuliStack}
to study families of compact metric spaces.


\paraid{p-86969a2d1c01}
Part~\ref{part:topology} studies
$\GHSpace$
as a global analogue of a hyperspace with a fixed ambient metric space.
The link between these two settings also appears in the study of geodesics.
M\'emoli and Wan
\cite[arXiv v2, Theorem~1]{MemoliWan2023}
prove that every Gromov--Hausdorff geodesic can be realized as a Hausdorff
geodesic of compact subsets of one compact metric space.
For comparison with our absolute retract theorem,
we recall the case of compact subsets of a fixed Banach space.
Curtis
\cite[Theorem~1.6]{Curtis1980}
proves,
in particular,
that the hyperspace of nonempty compact subsets of a Banach space,
equipped with the Hausdorff metric of its norm,
is an absolute retract.
We construct approximations to the identity of
$\GHSpace$
using quotients of hyperspaces of nonempty compact subsets of Banach spaces
by compact groups.

\paraid{p-5f89de56c669}
Our use of Toru\'nczyk's characterization in Part~\ref{part:hilbert}
has a precedent in the study of the Urysohn universal metric space.
In the early twenty-first century,
Uspenskij
\cite[arXiv v1, Theorem~2.1]{Uspenskij2004Urysohn}
proved that this space is homeomorphic to
$\SeparableHilbertSpace$.
The Urysohn space is the complete separable metric space that contains
an isometric copy of every separable metric space and in which every isometry
between finite subsets extends to a surjective isometry of the whole space.


\medskip
\noindent\textbf{Proof strategy and the four parts.}
\paraid{p-ce4b777e5ddb}
We first prove that
$\GHSpace$
is an absolute retract for metrizable spaces.
We approximate a sequence of maps from compact metrizable domains so that
the images of the approximating maps form a discrete family.
Since
$\GHSpace$
is complete and separable,
the absolute retract property and the discrete approximation property
imply \Cref{thm:main} by Toru\'nczyk's characterization
\cite[p.~248, assertion~(i)]{Torunczyk1981},
using the correction in
\cite[Section~C]{Torunczyk1985}.
In Parts~\ref{part:measures}--\ref{part:topology},
we construct the approximations needed for the absolute retract property.
The discrete approximation argument in Part~\ref{part:hilbert}
uses only the common preliminaries and its own metric constructions.
Figure~\ref{fig:proof-strategy}
shows how these two arguments combine to prove the classification.

\paraid{p-69b86b1ef543}
A difficulty in the absolute retract argument is that the isometry group
depends on the compact metric space.
Rouyer's results on generic compact metric spaces
\cite[arXiv v1, Theorems~2 and~4]{Rouyer2011}
imply that
\[
 \overline{
 \{\MetricSpace{\BaseCarrier}{\MetricSymbol}\in\GHSpace
   \mid\Isom\MetricSpace{\BaseCarrier}{\MetricSymbol}
      =\{\IdentityMap_\BaseCarrier\}\}
 }=\GHSpace.
\]
This also implies that the assignment of isometry groups
with their uniform metrics is discontinuous as a map from
$\GHSpace$
to itself.
Our construction must also handle spaces with nontrivial isometry groups.
On each model neighborhood,
we construct homomorphisms from the varying isometry groups into one fixed
orthogonal group so that the local point maps are equivariant.
We then use a partition of unity to combine the local point maps in a Banach space.
A countable product of orthogonal groups acts on this Banach space.
Passing to the orbit space of nonempty compact subsets of this Banach space
makes the resulting approximation independent of all coordinate choices.

\medskip
\noindent\textbf{Part~\ref{part:measures}.
Continuous invariant probability measures.}
\paraid{p-b0e31260869e}
In the first part,
we construct measures for integration on varying compact spaces.
The main result of this part,
\Cref{thm:invariant-measures},
assigns a Borel probability measure
$\SelectedMeasure{\BaseCarrier}{\MetricSymbol}$
to every nonempty compact metric space
$\MetricSpace{\BaseCarrier}{\MetricSymbol}$.
Each assigned measure has full support,
\[
 \supp(\SelectedMeasure{\BaseCarrier}{\MetricSymbol})=\BaseCarrier.
\]
For every isometry
$\IdentifyingIsometry\colon
\MetricSpace{\BaseCarrier}{\MetricSymbol}\to
\MetricSpace{\ComparisonCarrier}{\ComparisonMetric}$,
the assignment satisfies
\[
 \IdentifyingIsometry\Pushforward\SelectedMeasure{\BaseCarrier}{\MetricSymbol}
 =\SelectedMeasure{\ComparisonCarrier}{\ComparisonMetric}.
\]
The assignment of measures is also continuous in the following sense.
If compact subspaces
$\BaseCarrier_\SequenceIndex$
and
$\BaseCarrier$
of a compact metric space
$\MetricSpace{\AmbientCarrier}{\AmbientMetric}$
carry the restricted metrics
$\MetricSymbol_\SequenceIndex$
and
$\MetricSymbol$,
then
\[
 \HausdorffDistance{\AmbientMetric}(\BaseCarrier_\SequenceIndex,\BaseCarrier)\to0
 \quad\Longrightarrow\quad
 \SelectedMeasure{\BaseCarrier_\SequenceIndex}{\MetricSymbol_\SequenceIndex}
 \to\SelectedMeasure{\BaseCarrier}{\MetricSymbol}
 \quad\text{weakly in }\ProbabilityMeasures(\AmbientCarrier).
\]
Here the measures are regarded as probabilities on the common ambient space.

\paraid{p-968906a0f31d}
For a fixed compact space,
averaging over its isometry group produces an invariant probability measure.
Our construction must also preserve continuity when the space varies.
To construct such measures,
we prove that the space of invariant full-support measured compact spaces
is Polish and that the map that forgets the measure is an open surjection
onto
$\GHSpace$.
By Valov's theorem (\Cref{thm:valov}),
we obtain continuous probability laws on its fibers.
Averaging these laws produces the selected measures.
We use full support and continuity in the spectral approximation of
Part~\ref{part:spectral}.

\medskip
\noindent\textbf{Part~\ref{part:spectral}.
Continuous finite-dimensional approximation.}
\paraid{p-16a5923a2a64}
In the second part,
we replace distance functions by finite-dimensional data while controlling
all pairwise distances.
The main result of this part is \Cref{thm:local-models}.
For every
$\MetricSpace{\BaseCarrier}{\MetricSymbol}\in\GHSpace$
and every
$\LocalTolerance>0$,
we obtain a neighborhood
$\ModelNeighborhood$
of
$\MetricSpace{\BaseCarrier}{\MetricSymbol}$
and an integer
$\ModelDimension\geq1$
with the following property.
For each compact metric space
$\MetricSpace{\ComparisonCarrier}{\ComparisonMetric}$
whose isometry class belongs to
$\ModelNeighborhood$,
we construct a norm
$\NormSymbol_\ComparisonCarrier$
on
$\RealNumbers^\ModelDimension$
and a continuous map
$\FeatureCoordinates_\ComparisonCarrier\colon
\ComparisonCarrier\to\RealNumbers^\ModelDimension$.
For every
$\ComparisonPoint,\IntegrationPoint\in\ComparisonCarrier$,
define the induced pseudometric by
\[
 \Pseudometric_\ComparisonCarrier(\ComparisonPoint,\IntegrationPoint)
 =
 \NormSymbol_\ComparisonCarrier
 \bigl(\FeatureCoordinates_\ComparisonCarrier(\ComparisonPoint)
       -\FeatureCoordinates_\ComparisonCarrier(\IntegrationPoint)\bigr).
\]
After shrinking
$\ModelNeighborhood$,
we obtain
\[
 \sup_{\ComparisonPoint,\IntegrationPoint\in\ComparisonCarrier}
 \bigl|\Pseudometric_\ComparisonCarrier(\ComparisonPoint,\IntegrationPoint)
       -\ComparisonMetric(\ComparisonPoint,\IntegrationPoint)\bigr|
 <\LocalTolerance.
\]
The dimension
$\ModelDimension$
is fixed on
$\ModelNeighborhood$,
but may depend on the center
$\MetricSpace{\BaseCarrier}{\MetricSymbol}$
and the error scale
$\LocalTolerance$.
The norms
$\NormSymbol_\ComparisonCarrier$
and point maps
$\FeatureCoordinates_\ComparisonCarrier\colon
\ComparisonCarrier\to\RealNumbers^\ModelDimension$
vary continuously up to changes of coordinates in
$\OrthogonalGroup(\ModelDimension)$
in the sense of
\Cref{thm:local-models}\ref{item:model-continuity}.
Each isometry group
$\Isom\MetricSpace{\ComparisonCarrier}{\ComparisonMetric}$
acts on the coordinate image
$\FeatureCoordinates_\ComparisonCarrier(\ComparisonCarrier)
\subset\RealNumbers^\ModelDimension$
through a continuous homomorphism
\[
 \IsometryRepresentation{\ComparisonCarrier}{\ComparisonMetric}
 \colon\Isom\MetricSpace{\ComparisonCarrier}{\ComparisonMetric}
 \to\OrthogonalGroup(\ModelDimension).
\]
For every
$\FirstIsometry\in\Isom\MetricSpace{\ComparisonCarrier}{\ComparisonMetric}$,
it satisfies
\[
 \FeatureCoordinates_\ComparisonCarrier\circ\FirstIsometry
 =\IsometryRepresentation{\ComparisonCarrier}{\ComparisonMetric}(\FirstIsometry)
  \circ\FeatureCoordinates_\ComparisonCarrier,
 \qquad
 \NormSymbol_\ComparisonCarrier
  \circ\IsometryRepresentation{\ComparisonCarrier}{\ComparisonMetric}(\FirstIsometry)
 =\NormSymbol_\ComparisonCarrier
\]
by Remark~\ref{rem:local-isometry-representation}.
Thus the induced pseudometrics
$\Pseudometric_\ComparisonCarrier$
respect isometries while the coordinate changes take place in the fixed
compact group
$\OrthogonalGroup(\ModelDimension)$.

\paraid{p-458636bfa314}
To construct $\ModelNeighborhood$, we use the measures of Part~\ref{part:measures}
to define the distance operator
$\DistanceOperator{\BaseCarrier}{\MetricSymbol}$
on
$\LebesgueSpace^2(\BaseCarrier,\SelectedMeasure{\BaseCarrier}{\MetricSymbol})$.
For every
$f\in\LebesgueSpace^2(\BaseCarrier,\SelectedMeasure{\BaseCarrier}{\MetricSymbol})$
and
$\BasePoint\in\BaseCarrier$,
set
\[
 (\DistanceOperator{\BaseCarrier}{\MetricSymbol}f)(\BasePoint)
 =\int_\BaseCarrier
   \MetricSymbol(\BasePoint,\IntegrationPoint)f(\IntegrationPoint)
   \,\IntegrationDifferential
   \SelectedMeasure{\BaseCarrier}{\MetricSymbol}(\IntegrationPoint).
\]
These operators and their spectral coordinates were studied by Maria,
Oudot,
and Solomon
\cite[Section~3]{MariaOudotSolomon2020}.
Finite spectral subspaces of the distance operators approximate the distance
functions.
For our local models,
we choose the unique best approximations in
$\LebesgueSpace^\IntegrabilityExponent$
for a finite exponent
$2\leq\IntegrabilityExponent<\infty$
chosen to achieve the error bound at the center in
\ref{item:model-center-error} of \Cref{thm:local-models}.
We keep whole eigenspaces to handle eigenvalue multiplicities by orthogonal
changes of basis.
We retain both the compact coordinate image and its norm for the
construction in Part~\ref{part:topology}.
We expect this approach to provide a bridge between the study of the
Gromov--Hausdorff space and geometric analysis or spectral geometry.

\paraid{p-24f2b5c3820e}
By \cite[Theorem~1.1]{Ishiki2026Interpolation},
we obtain a uniform error bound when extending prescribed metrics from a
discrete family of closed subsets of a fixed metrizable space.
The absolute retract argument here uses approximations of metrics on varying
compact spaces whose coordinate pairs satisfy the convergence conditions in
\ref{item:model-continuity} of \Cref{thm:local-models}.

\medskip
\noindent\textbf{Part~\ref{part:topology}.
The absolute retract property.}
\paraid{p-159ea0252542}
In \Cref{thm:part-iii-main},
the main result of the third part,
we prove that
$\GHSpace$
is an absolute retract for all metrizable spaces.
In \Cref{thm:absolute-extensor},
we prove that
$\GHSpace$
is an absolute extensor for all metrizable spaces.

\paraid{p-02e057f0dca4}
Fix a continuous function
$\ErrorControl\colon\GHSpace\to(0,\infty)$.
We choose a countable locally finite cover by local models whose errors
are less than
$\ErrorControl$
throughout their domains.
For each model index
$\CoverIndex\in\NonnegativeIntegers$,
let
$\ModelDimension_\CoverIndex$
be the model dimension.
We represent the varying norms isometrically in fixed Banach spaces
and combine the point maps by a partition of unity in their
$\SummableNorm$-sum,
a separable Banach space
$\BanachAmbientSpace$.
The norm representation respects orthogonal changes of coordinates
(\Cref{lem:dual-representation,lem:dual-representation-equivariance}).
The compact group
\[
 \ActingGroup
 =\prod_{\CoverIndex\in\NonnegativeIntegers}
   \OrthogonalGroup(\ModelDimension_\CoverIndex)
\]
acts continuously on
$\BanachAmbientSpace$
by linear isometries.
Passing to the orbit of each compact coordinate image removes the
choices of bases and isometric representatives.
The orbit space
$\AuxiliaryAR=\CompactHyperspace(\BanachAmbientSpace)/\ActingGroup$
is an absolute retract by Antonyan's results
(\Cref{lem:hyperspace-orbit-ar}).
By \Cref{thm:global-domination},
we obtain continuous maps
\[
 \GHSpace\xrightarrow{\ \ApproximationMap\ }
 \AuxiliaryAR\xrightarrow{\ \RealizationMap\ }\GHSpace
\]
and a continuous homotopy
$\ApproximationHomotopy\colon\GHSpace\times\UnitInterval\to\GHSpace$.
For every
$\MetricSpace{\ComparisonCarrier}{\ComparisonMetric}\in\GHSpace$,
this homotopy satisfies
\begin{equation}\label{eq:intro-domination-1}
 \ApproximationHomotopy(\MetricSpace{\ComparisonCarrier}{\ComparisonMetric},0)
 =\MetricSpace{\ComparisonCarrier}{\ComparisonMetric},
\end{equation}
and
\begin{equation}\label{eq:intro-domination-2}
 \ApproximationHomotopy(\MetricSpace{\ComparisonCarrier}{\ComparisonMetric},1)
 =\RealizationMap\ApproximationMap\MetricSpace{\ComparisonCarrier}{\ComparisonMetric}.
\end{equation}
The homotopy tracks satisfy
\begin{equation}\label{eq:intro-domination-3}
 \diam_{\GHDistance}
 \bigl\{\ApproximationHomotopy(\MetricSpace{\ComparisonCarrier}{\ComparisonMetric},t)
       \mid t\in\UnitInterval\bigr\}
 <\ErrorControl\MetricSpace{\ComparisonCarrier}{\ComparisonMetric}.
\end{equation}
Applying Hanner's domination theorem (\Cref{thm:hanner-domination})
to the small homotopy tracks in \eqref{eq:intro-domination-3},
we obtain the ANR property.
Scaling all distances to
$0$
contracts
$\GHSpace$
to the one-point space.
Using the ANR/ANE equivalence in \Cref{thm:retract-extensor},
we apply \Cref{thm:contractible-ane} to conclude that
$\GHSpace$
is an AE for metrizable spaces.
By the AE/AR equivalence in \Cref{thm:retract-extensor},
we obtain the absolute retract property.

\begin{figure}[H]
\centering
\begin{tikzpicture}[
 >=Stealth,
 font=\small,
 partbox/.style={
  draw=black,
  fill=white,
  rounded corners=3pt,
  line width=.65pt,
  text width=4.9cm,
  minimum height=2.35cm,
  align=center,
  inner sep=7pt
 },
 route/.style={->,line width=.85pt,draw=black}
]
\node[partbox] (measures) at (-2.95,0) {
 \textbf{Part~\ref{part:measures}}\\
 Continuous invariant\\
 full-support probabilities\\
 \Cref{thm:invariant-measures}
};
\node[partbox] (models) at (2.95,0) {
 \textbf{Part~\ref{part:spectral}}\\
 Finite-dimensional local models\\
 $\Isom(Y,e)\to\OrthogonalGroup(n)$\\
 \Cref{thm:local-models}
};
\node[partbox,minimum height=2.95cm] (ar) at (2.95,-3.55) {
 \textbf{Part~\ref{part:topology}}\\
 $\ActingGroup=\prod_{i\in\NonnegativeIntegers}\OrthogonalGroup(n_i)$\\
 $\CompactHyperspace(\BanachAmbientSpace)/\ActingGroup$ is an AR\\
 Small tracks $\Longrightarrow$ ANR\\
 Scaling contraction $\Longrightarrow$ AR\\
 \Cref{thm:part-iii-main}
};
\node[
 partbox,
 minimum height=2.95cm
] (discrete) at (-2.95,-3.55) {
 \textbf{Part~\ref{part:hilbert}}\\
 Products with finite equilateral spaces\\
 Discrete approximation\\
 \Cref{thm:discrete-approximation}
};
\node[
 partbox,
 text width=10.8cm,
 minimum height=1.9cm
] (hilbert) at (0,-7) {
 \textbf{Toru\'nczyk's characterization}\\
 AR and discrete approximation,
 with completeness and separability\\
 $\GHSpace$ is homeomorphic to real $\SeparableHilbertSpace$\\
 \textbf{Part~\ref{part:hilbert}.
 \Cref{thm:part-iv-main}}
};
\draw[route] (measures.east)--(models.west);
\draw[route] (models.south)--(ar.north)
 node[midway,left=5pt,font=\footnotesize,align=right]
 {Partition of unity\\in the Banach space $\BanachAmbientSpace$};
\draw[route] (ar.south)--(2.95,-5.5)--(0,-5.5)--(hilbert.north);
\draw[line width=.85pt,draw=black]
 (discrete.south)--(-2.95,-5.5)--(0,-5.5);
\fill[black] (0,-5.5) circle (1.5pt);
\end{tikzpicture}
\caption{The two arguments used to identify the topology of
$\GHSpace$.
Arrows indicate dependence of the proofs.
In Parts~\ref{part:measures}--\ref{part:topology},
we prove the absolute retract property.
The discrete approximation theorem in Part~\ref{part:hilbert}
uses the common preliminaries and its own metric constructions.
These two conclusions,
together with completeness and separability,
imply the Hilbert space classification.}
\label{fig:proof-strategy}
\end{figure}


\medskip
\noindent\textbf{Part~\ref{part:hilbert}.
Discrete approximation and Hilbert space topology.}
\paraid{p-64cc2f159292}
In the fourth part,
we prove the remaining approximation condition and then
obtain the classification in \Cref{thm:part-iv-main}.
\Cref{thm:discrete-approximation},
the discrete approximation theorem of this part,
applies to any sequence
$\{\ApproximationDomain_\FamilyIndex\}_{\FamilyIndex\in\NonnegativeIntegers}$
of compact metrizable spaces.
If
$\InputMap_\FamilyIndex\colon\ApproximationDomain_\FamilyIndex\to\GHSpace$
are continuous and
$\OpenCover$
is an open cover of
$\GHSpace$,
then there are continuous maps
$\ApproximatingMap_\FamilyIndex\colon\ApproximationDomain_\FamilyIndex\to\GHSpace$
such that
\[
 \forall\FamilyIndex\in\NonnegativeIntegers\
 \forall x\in\ApproximationDomain_\FamilyIndex\
 \exists U\in\OpenCover,\qquad
 \{\InputMap_\FamilyIndex(x),\ApproximatingMap_\FamilyIndex(x)\}\subset U,
\]
and
\[
 \{\ApproximatingMap_\FamilyIndex(\ApproximationDomain_\FamilyIndex)
   \}_{\FamilyIndex\in\NonnegativeIntegers}
 \text{ is discrete in }\GHSpace.
\]
The last condition means that every point of
$\GHSpace$
has a neighborhood that intersects at most one of these image sets.

\paraid{p-66ca83693c4d}
Taking products with sufficiently small finite equilateral spaces changes each
input by a controlled amount in Gromov--Hausdorff distance.
We choose the cardinalities recursively to exceed the bounds on separated subsets
in the previously constructed compact images.
This separates the images,
and letting the cardinalities tend to infinity prevents accumulation in
$\GHSpace$.
This construction is independent of the results of Parts~\ref{part:measures}--\ref{part:topology}.
If all domains are the Hilbert cube
$\HilbertCube=\UnitInterval^{\NonnegativeIntegers}$,
then we obtain the discrete approximation condition in Toru\'nczyk's theorem.
We combine this condition with completeness,
separability,
and the absolute retract property from Part~\ref{part:topology}
to prove \Cref{thm:main}.

\paraid{p-59ed21a6dca9}
The main theorem also answers the local contractibility question.
We use local contractibility in the sense that every neighborhood
$\OpenNeighborhood $
of a point contains a neighborhood
$\SmallerNeighborhood $
whose inclusion into
$\OpenNeighborhood $
is null-homotopic.
By \Cref{thm:main},
every point has a neighborhood basis consisting of open sets that
strongly deformation retract onto that point.
This conclusion is topological and does not prescribe the neighborhoods
as balls for
$\GHDistance$.

\medskip
\noindent\textbf{Organization.}
\paraid{p-10cf165b0a9f}
The introduction and
\Cref{sec:preliminaries}
provide the common background for four parts.
Each part also has its own preliminaries,
in Sections~\ref{sec:prelim-measures},
\ref{sec:prelim-spectral},
\ref{sec:prelim-topology},
and~\ref{sec:prelim-hilbert},
respectively.
Part~\ref{part:measures},
introduced in \Cref{sec:intro-measures} and developed in \Cref{sec:measures},
constructs continuous invariant probabilities.
Part~\ref{part:spectral},
introduced in \Cref{sec:intro-spectral} and developed in \Cref{sec:spectral},
uses these probabilities to construct finite-dimensional local models.
Part~\ref{part:topology},
introduced in \Cref{sec:intro-topology} and developed in \Cref{sec:ar},
combines the models into global approximations
and proves the absolute retract property.
Part~\ref{part:hilbert},
introduced in \Cref{sec:intro-hilbert} and developed in \Cref{sec:approximation},
proves the discrete approximation property and completes the proof of
\Cref{thm:main}.
\Cref{sec:questions} concludes with questions about related spaces
and the geometry of GH balls.

\medskip
\noindent\textbf{Conventions and notation.}
\paraid{p-67846154e305}
For a set
$\BaseCarrier$,
we denote by
$\Card(\BaseCarrier)$
its cardinality.
Every metric space
$\MetricSpace{\BaseCarrier }{\MetricSymbol }$
representing a point of
$\GHSpace$
is nonempty and compact.
We identify such a pair with its isometry class when it occurs in
$\GHSpace$.
The symbol
$\BaseCarrier $
denotes the underlying set,
equipped with the topology induced by
$\MetricSymbol $.
Subscripts on constructions such as
$\SelectedMeasure{\BaseCarrier}{\MetricSymbol}$
refer to the specified pair
$\MetricSpace{\BaseCarrier }{\MetricSymbol }$.
When isometric embeddings into
$\MetricSpace{\AmbientCarrier }{\AmbientMetric }$
are specified,
we use the same symbols for the embedded subsets and the restricted
metrics.
The notation
$\IdentityMap_{\BaseCarrier}$
denotes the identity map on
$\BaseCarrier$.
We omit the subscript when the domain is specified.
Banach and Hilbert spaces are real unless complex scalars are explicitly indicated.
We denote by
$\CompactHyperspace(\BaseCarrier )$
the set of all nonempty compact subsets of
$\BaseCarrier $.
For a specified metric
$\MetricSymbol $,
we equip it with the Hausdorff metric
$\HausdorffDistance{\MetricSymbol }$.
For a metrizable space
$\DomainSpace$,
we denote by
$\ProbabilityMeasures(\DomainSpace )$
the set of Radon probability measures on
$\DomainSpace$
and equip it with the topology of weak convergence against bounded continuous
functions.
An isometric embedding
$\IsometricEmbeddingMap \colon\MetricSpace{\BaseCarrier }{\MetricSymbol }\to\MetricSpace{\ComparisonCarrier }{\ComparisonMetric }$
is a map satisfying
$\ComparisonMetric (\IsometricEmbeddingMap (\BasePoint ),\IsometricEmbeddingMap (\ComparisonPoint ))=\MetricSymbol (\BasePoint ,\ComparisonPoint )$
for every
$\BasePoint ,\ComparisonPoint \in \BaseCarrier $.
Throughout this paper,
an isometry is a surjective isometric embedding.
We use
$\NonnegativeIntegers=\{0,1,2,\ldots\}$
as the index set of sequences unless stated otherwise.

\begin{useofai}
\paraid{p-64859c0d8f24}
OpenAI Codex,
based on GPT-6,
was used to explore and develop the mathematical constructions and proofs,
search the literature,
and prepare the exposition and \LaTeX{} source.
The research and writing workflow used customized versions of
\texttt{math-research-harness}
(\url{https://github.com/haruhisa-enomoto/math-research-harness})
and
\texttt{math-paper-skills}
(\url{https://github.com/haruhisa-enomoto/math-paper-skills}),
both made publicly available by Haruhisa Enomoto.
The author takes full responsibility for the mathematical content and the final
manuscript.
\end{useofai}

\begin{acknowledgements}
\paraid{p-e60be87ab1f9}
The author was supported by JSPS KAKENHI Grant Number JP24KJ0182.
\end{acknowledgements}


\section{Common preliminaries}\label{sec:preliminaries}
\paraid{p-faeeb28244f8}
We recall metric convergence and compactness,
comparison of metric quotients,
and weak convergence of probability measures.
The final subsection fixes the contraction and absolute extension terminology.

\subsection{Gromov--Hausdorff convergence and compactness}\label{subsec:preliminaries-metric}
\paraid{p-9622f56720d9}
We begin with the correspondence formula for the distance and compactness criteria.

\paraid{p-f7c8be172ef5}
For a metric space
$\MetricSpace{\BaseCarrier}{\MetricSymbol}$,
a point
$\BasePoint\in\BaseCarrier$,
and a radius
$\Radius>0$,
we write
\[
 \MetricBall(\BasePoint,\Radius;\MetricSymbol)
 =\{\ComparisonPoint\in\BaseCarrier\mid\MetricSymbol(\BasePoint,\ComparisonPoint)<\Radius\}.
\]
For a subset
$\BorelSubset\subset\BaseCarrier$,
we put
\[
 \dist_{\MetricSymbol}(\BasePoint,\BorelSubset)
 =\inf_{\ComparisonPoint\in\BorelSubset}\MetricSymbol(\BasePoint,\ComparisonPoint),
 \qquad \dist_{\MetricSymbol}(\BasePoint,\emptyset)=\infty.
\]
We denote by
$\CompactHyperspace(\BaseCarrier)$
the set of all nonempty compact subsets of
$\BaseCarrier$.
The \emph{Hausdorff distance} on this set is defined by
\[
 \HausdorffDistance{\MetricSymbol}\colon
 \CompactHyperspace(\BaseCarrier)\times\CompactHyperspace(\BaseCarrier)\to[0,\infty).
\]
For every
$\FirstCompactSet,\SecondCompactSet\in\CompactHyperspace(\BaseCarrier)$,
put
\[
 \HausdorffDistance{\MetricSymbol}(\FirstCompactSet,\SecondCompactSet)
 =\max\left\{
   \sup_{\BasePoint\in\FirstCompactSet}\dist_{\MetricSymbol}(\BasePoint,\SecondCompactSet),
   \sup_{\ComparisonPoint\in\SecondCompactSet}\dist_{\MetricSymbol}(\ComparisonPoint,\FirstCompactSet)
 \right\}.
\]
The \emph{Vietoris topology} on
$\CompactHyperspace(\BaseCarrier)$
has a basis consisting of the sets obtained as follows.
For an integer
$m\geq1$
and open subsets
$U_1,\ldots,U_m$
of
$\BaseCarrier$,
take
\[
 \left\{\FirstCompactSet\in\CompactHyperspace(\BaseCarrier)\;\middle|\;
 \FirstCompactSet\subset\bigcup_{i=1}^{m}U_i,
 \quad \FirstCompactSet\cap U_i\ne\emptyset\ (1\leq i\leq m)\right\}.
\]

\begin{theorem}[{\cite[Exercises~2.2.11(a),(b) and~3.2.9]{Beer1993}}]\label{thm:compact-hyperspace-topology}
\paraid{p-974030144cf7}
Let
$\MetricSpace{\BaseCarrier}{\MetricSymbol}$
be a metric space.
Then the Hausdorff distance induces the Vietoris topology on
$\CompactHyperspace(\BaseCarrier)$.
If
$\BaseCarrier$
is separable,
then
$\CompactHyperspace(\BaseCarrier)$
is separable.
\end{theorem}

\begin{lemma}[{\cite[Chapter~2, Proposition~5.3, p.~40]{Wicks1991}}]\label{lem:hausdorff-lipschitz-images}
\paraid{p-eb3ed5b8a76b}
Let
$\MetricSpace{\BaseCarrier}{\MetricSymbol}$
and
$\MetricSpace{\ComparisonCarrier}{\ComparisonMetric}$
be metric spaces,
let
$L\geq0$,
and let
$\Mapping\colon\BaseCarrier\to\ComparisonCarrier$
be
$L$-Lipschitz.
Then for all
$\FirstCompactSet,\SecondCompactSet\in\CompactHyperspace(\BaseCarrier)$,
we have
$\HausdorffDistance{\ComparisonMetric}
 (\Mapping(\FirstCompactSet),\Mapping(\SecondCompactSet))
 \leq L\cdot \HausdorffDistance{\MetricSymbol}(\FirstCompactSet,\SecondCompactSet)$.
\end{lemma}

% \paraid{p-109cc338209b}

\paraid{p-0a1ce4229ada}
For
$\ErrorTolerance>0$,
a subset
$\NetSet\subset\BaseCarrier$
is an
\emph{$\ErrorTolerance$-net}
if
\[
 \forall\BasePoint\in\BaseCarrier\quad
 \exists\NetPoint\in\NetSet\quad
 \MetricSymbol(\BasePoint,\NetPoint)<\ErrorTolerance.
\]
Every compact metric space admits a finite
$\ErrorTolerance$-net
for every
$\ErrorTolerance>0$.

\paraid{p-d7c66dfce89f}
A \emph{correspondence} between nonempty sets
$\BaseCarrier $
and
$\ComparisonCarrier $
is a subset of
$\BaseCarrier \times \ComparisonCarrier $
whose two coordinate projections are surjective.
We denote by
$\Cor(\BaseCarrier,\ComparisonCarrier)$
the set of all correspondences between
$\BaseCarrier$
and
$\ComparisonCarrier$.
For metric spaces
$\MetricSpace{\BaseCarrier }{\MetricSymbol }$
and
$\MetricSpace{\ComparisonCarrier }{\ComparisonMetric }$,
the distortion of a correspondence
$\Correspondence\in\Cor(\BaseCarrier,\ComparisonCarrier)$
is
\[
 \dis(\Correspondence )=
 \sup\left\{
   \abs{\MetricSymbol (\BasePoint _0,\BasePoint _1)-\ComparisonMetric (\ComparisonPoint _0,\ComparisonPoint _1)}
   \,\middle|\,
   (\BasePoint _0,\ComparisonPoint _0),(\BasePoint _1,\ComparisonPoint _1)\in \Correspondence
 \right\}.
\]
\begin{theorem}[{\cite[Theorem~7.3.25]{BuragoBuragoIvanov2001} and \cite[arXiv v1, Theorem~6.12]{Tuzhilin2020}}]\label{thm:gh-correspondence}
\paraid{p-fc061a808065}
For nonempty compact metric spaces
$\MetricSpace{\BaseCarrier }{\MetricSymbol }$
and
$\MetricSpace{\ComparisonCarrier }{\ComparisonMetric }$,
the Gromov--Hausdorff distance satisfies
\begin{equation}\label{eq:correspondence}
 \GHDistance(\MetricSpace{\BaseCarrier }{\MetricSymbol },\MetricSpace{\ComparisonCarrier }{\ComparisonMetric })=\frac12\inf_{\Correspondence\in\Cor(\BaseCarrier,\ComparisonCarrier)}\dis(\Correspondence).
\end{equation}
Equivalently,
it is the infimum of the Hausdorff distances between isometric copies
of
$\MetricSpace{\BaseCarrier }{\MetricSymbol }$
and
$\MetricSpace{\ComparisonCarrier }{\ComparisonMetric }$
in a common metric space.

\end{theorem}

\paraid{p-c3f214f2a9fe}
The next lemma realizes a prescribed correspondence in a common metric space
with a bound on the distances between corresponding points.
We use \Cref{lem:correspondence-embedding} in the proof of
\Cref{lem:common-gh-embedding}
to embed a convergent sequence of compact metric spaces into one compact metric space.

\begin{lemma}\label{lem:correspondence-embedding}
\paraid{p-77a82a051441}
Let
$\MetricSpace{\BaseCarrier}{\MetricSymbol}$
and
$\MetricSpace{\ComparisonCarrier}{\ComparisonMetric}$
be nonempty compact metric spaces,
let
$\Correspondence\subset\BaseCarrier\times\ComparisonCarrier$
be a correspondence,
and let
$\SmallError>0$
satisfy
$\dis(\Correspondence)\leq2\SmallError$.
Then there is a metric
$\AmbientMetric$
on the disjoint union
$\BaseCarrier\sqcup\ComparisonCarrier$
extending both metrics such that for every
$(\BasePoint,\ComparisonPoint)\in\Correspondence$
we have
\[
 \AmbientMetric(\BasePoint,\ComparisonPoint)\leq\SmallError.
\]
In particular,
$\HausdorffDistance{\AmbientMetric}(\BaseCarrier,\ComparisonCarrier)\leq\SmallError$.
\end{lemma}

\begin{proof}
\paraid{p-1dc4840c170b}
This follows from
\cite[proof of Theorem~7.3.25, pp.~257--258]{BuragoBuragoIvanov2001}.
\end{proof}

\begin{theorem}[{Consequence of \cite[arXiv v1, Theorem~7.18]{Tuzhilin2020}}]\label{thm:uniform-nets}
\paraid{p-b1b76db6a357}
If
$\MetricSpace{\BaseCarrier_\SequenceIndex}{\MetricSymbol_\SequenceIndex}
 \to\MetricSpace{\BaseCarrier}{\MetricSymbol}$
in
$\GHSpace$,
then for every
$\ErrorTolerance>0$
there is a positive integer
$\NetSize$
such that each
$\BaseCarrier_\SequenceIndex$
admits an
$\ErrorTolerance$-net
of cardinality at most
$\NetSize$.
\end{theorem}

\paraid{p-e041a46858f0}
Let
$\MetricSpace{\BaseCarrier}{\MetricSymbol}$
be a metric space, and let
$\SeparationScale>0$.
A subset
$S\subset\BaseCarrier$
is called
\emph{$\SeparationScale$-separated}
if
\[
 \forall x,y\in S,\qquad
 x\ne y\Longrightarrow\MetricSymbol(x,y)\ge\SeparationScale.
\]

\begin{corollary}\label{cor:uniform-separated-sets}
\paraid{p-f9bb5a038d1d}
Let
$\MetricSpace{\BaseCarrier_\SequenceIndex}{\MetricSymbol_\SequenceIndex}
 \to\MetricSpace{\BaseCarrier}{\MetricSymbol}$
in
$\GHSpace$,
and let
$\SeparationScale>0$.
Then there is a positive integer
$\NetSize$
such that every
$\SeparationScale$-separated
subset of every
$\BaseCarrier_\SequenceIndex$
has cardinality at most
$\NetSize$.
\end{corollary}

\begin{proof}
\paraid{p-0e6ff4e87f94}
By \Cref{thm:uniform-nets},
choose
$\SeparationScale/3$-nets
$\mathcal{N}_\SequenceIndex\subset\BaseCarrier_\SequenceIndex$
in all
$\BaseCarrier_\SequenceIndex$
with
$\Card(\mathcal{N}_\SequenceIndex)\leq\NetSize$.
For an arbitrary
$\SeparationScale$-separated subset
$S_\SequenceIndex\subset\BaseCarrier_\SequenceIndex$,
choose for each point of
$S_\SequenceIndex$
a point of
$\mathcal{N}_\SequenceIndex$
within distance
$\SeparationScale/3$.
This assignment is injective,
because two points with the same image would have distance less than
$2\SeparationScale/3$.
Consequently,
$\Card(S_\SequenceIndex)
 \leq \Card(\mathcal{N}_\SequenceIndex)\leq \NetSize$.
This proves the uniform cardinality bound.
\end{proof}

\paraid{p-47d23e54815b}
We use completeness and separability of the Gromov--Hausdorff space
$\GHSpace$.

\begin{theorem}[{\cite[arXiv v1, Corollary~7.24]{Tuzhilin2020}}]\label{lem:gh-basics}
\paraid{p-fa3b1c89e980}
The metric space
$\MetricSpace{\GHSpace}{\GHDistance}$
is complete and separable.
\end{theorem}

\paraid{p-2215668b1fa1}
Gromov constructs a common compact metric space admitting isometric embeddings
of every member of a family of compact metric spaces with uniformly bounded
diameters and uniformly bounded covering numbers at each positive radius
\cite[Section~6, pp.~64--65]{Gromov1981}.
See also \cite[arXiv v1, Theorem~7.19]{Tuzhilin2020}.
In the next theorem,
we also arrange Hausdorff convergence of the embedded spaces.

\begin{theorem}\label{lem:common-gh-embedding}
\paraid{p-fc625e3c11fb}
Every Cauchy sequence
$\{\MetricSpace{\BaseCarrier_\SequenceIndex}{\MetricSymbol_\SequenceIndex}\}_{\SequenceIndex\in\NonnegativeIntegers}$
in
$\GHSpace$
admits isometric embeddings into a compact metric space
$\MetricSpace{\AmbientCarrier}{\AmbientMetric}$
whose images converge in Hausdorff distance to a nonempty compact subset.
If its Gromov--Hausdorff limit is
$\MetricSpace{\BaseCarrier}{\MetricSymbol}$,
then we may also embed
$\MetricSpace{\BaseCarrier}{\MetricSymbol}$
into the same ambient space.
Identifying these spaces with their images, we obtain
\[
 \HausdorffDistance{\AmbientMetric}(\BaseCarrier_\SequenceIndex,\BaseCarrier)\longrightarrow0.
\]
\end{theorem}

\begin{proof}
\paraid{p-0e319197acba}
By \Cref{lem:gh-basics},
let
$\MetricSpace{\BaseCarrier}{\MetricSymbol}$
be the limit and put
\[
 \ErrorTolerance_\SequenceIndex
 =\GHDistance(\MetricSpace{\BaseCarrier_\SequenceIndex}{\MetricSymbol_\SequenceIndex},
              \MetricSpace{\BaseCarrier}{\MetricSymbol})+2^{-\SequenceIndex}.
\]
Using \Cref{thm:gh-correspondence,lem:correspondence-embedding}
and gluing along the common copy of
$\BaseCarrier$
\cite[author's manuscript, Theorem~27.2]{Villani2009},
we obtain a metric
$\AmbientMetric$
on the union of
$\BaseCarrier$
and all
$\BaseCarrier_\SequenceIndex$
extending their metrics and satisfying
\[
 \HausdorffDistance{\AmbientMetric}
 (\BaseCarrier_\SequenceIndex,\BaseCarrier)
 \leq\ErrorTolerance_\SequenceIndex\longrightarrow0.
\]

\paraid{p-82732ff2ca3f}
For every
$\SmallError>0$,
the tail lies in the
$\SmallError/2$-neighborhood of
$\BaseCarrier$.
A finite
$\SmallError/2$-net in
$\BaseCarrier$,
together with finite
$\SmallError$-nets in the finitely many remaining spaces,
forms a finite
$\SmallError$-net for the union.
The completion of this union is therefore compact and retains the displayed
Hausdorff bounds.
\end{proof}

\paraid{p-159a1141fda7}
We also recall the compactness criterion used for families of functions.

\begin{theorem}[{Arzel\`a--Ascoli, \cite[Theorem~6.26]{Muscat2024}}]\label{thm:arzela-ascoli}
\paraid{p-017c33e53ae1}
Let
$\MetricSpace{\BaseCarrier}{\MetricSymbol}$
be nonempty and compact,
and let
$\MetricSpace{\ComparisonCarrier}{\ComparisonMetric}$
be complete.
Equip the set
$\ContinuousFunctions(\BaseCarrier,\ComparisonCarrier)$
of continuous maps with the uniform metric
\[
 \UniformMetric(\Mapping,\SecondMapping)
 =\sup_{\BasePoint\in\BaseCarrier}
   \ComparisonMetric(\Mapping(\BasePoint),\SecondMapping(\BasePoint)).
\]
Then a family
$\FunctionFamily\subset\ContinuousFunctions(\BaseCarrier,\ComparisonCarrier)$
has compact closure if and only if its combined image
$\bigcup_{\Mapping\in\FunctionFamily}\Mapping(\BaseCarrier)$
is totally bounded and the family
$\FunctionFamily$
is equicontinuous.
Here equicontinuity means that for every
$\ErrorTolerance>0$
there is
$\SmallError>0$
such that for every
$\Mapping\in\FunctionFamily$
and every
$\BasePoint,\ComparisonPoint\in\BaseCarrier$
with
$\MetricSymbol(\BasePoint,\ComparisonPoint)<\SmallError$,
we have
\[
 \ComparisonMetric(\Mapping(\BasePoint),\Mapping(\ComparisonPoint))<\ErrorTolerance.
\]
\end{theorem}

\paraid{p-f1481eeb5719}
The cited theorem characterizes total boundedness in the uniform metric.
Completeness of the target
$\MetricSpace{\ComparisonCarrier}{\ComparisonMetric}$
implies completeness of the function space
$\ContinuousFunctions(\BaseCarrier,\ComparisonCarrier)$
with the uniform metric
$\UniformMetric$,
which yields the formulation in terms of compact closure above.

\subsection{Pseudometrics and metric quotients}\label{subsec:preliminaries-quotients}
\paraid{p-dcc3379498b8}
We prove estimates that also apply when distinct points are identified.

\paraid{p-f455c68e0d8f}
The diameter is continuous on
$\GHSpace$,
by \cite[arXiv v1, Example~6.29]{Tuzhilin2020},
which states that for all nonempty compact metric spaces
$\MetricSpace{\BaseCarrier}{\MetricSymbol}$
and
$\MetricSpace{\ComparisonCarrier}{\ComparisonMetric}$
we have
\[
 \abs{\diam\MetricSpace{\BaseCarrier }{\MetricSymbol }-\diam\MetricSpace{\ComparisonCarrier }{\ComparisonMetric }}\le 2\GHDistance(\MetricSpace{\BaseCarrier }{\MetricSymbol },\MetricSpace{\ComparisonCarrier }{\ComparisonMetric }).
\]
A continuous pseudometric
$\Pseudometric $
on a compact space
$\BaseCarrier $
defines an equivalence relation as follows.
For
$\BasePoint,\ComparisonPoint\in\BaseCarrier$,
write
$\BasePoint\sim_\Pseudometric\ComparisonPoint$
if
$\Pseudometric(\BasePoint,\ComparisonPoint)=0$.
Its metric quotient is
\[
 \MetricQuotient{\BaseCarrier }{\Pseudometric }
 =\MetricSpace{\BaseCarrier /{\sim_\Pseudometric }}{\overline{\Pseudometric }},
\]
where for representatives
$\BasePoint,\ComparisonPoint\in\BaseCarrier$
the quotient metric satisfies
\[
 \overline{\Pseudometric }([\BasePoint ],[\ComparisonPoint ])=\Pseudometric (\BasePoint ,\ComparisonPoint ).
\]
This quotient is compact.
When
$\Pseudometric $
is a metric,
it is isometric to
$\MetricSpace{\BaseCarrier }{\Pseudometric }$.
For the comparison of a pseudometric quotient with a metric on the same set,
see the author's
\cite[Proposition~2.12]{Ishiki2023FractalDimensions}.
We use the following correspondence form,
which also compares two metric quotients.

\begin{proposition}\label{lem:pseudometric-comparison}
\paraid{p-03dd7915a29a}
Let
$\BaseCarrier$
and
$\ComparisonCarrier$
be nonempty compact spaces.
Let
$\Correspondence\subset\BaseCarrier\times\ComparisonCarrier$
be a correspondence.
For continuous pseudometrics
$\Pseudometric$
on
$\BaseCarrier$
and
$\ComparisonPseudometric$
on
$\ComparisonCarrier$,
put
\begin{equation}\label{eq:correspondence-pseudometric-error}
 \PseudometricError_\Correspondence(\Pseudometric,\ComparisonPseudometric)
 =\sup_{(\BasePoint_0,\ComparisonPoint_0),(\BasePoint_1,\ComparisonPoint_1)\in\Correspondence}
 |\Pseudometric(\BasePoint_0,\BasePoint_1)
       -\ComparisonPseudometric(\ComparisonPoint_0,\ComparisonPoint_1)|.
\end{equation}
\begin{enumerate}[label=\textup{(Q\arabic*)},ref=\textup{(Q\arabic*)}]
\item\label{item:quotient-1}
\paraid{p-31226097263c}
The quotient spaces satisfy
\begin{equation}\label{eq:pseudometric-correspondence-bound}
 \GHDistance(\MetricQuotient{\BaseCarrier}{\Pseudometric},
              \MetricQuotient{\ComparisonCarrier}{\ComparisonPseudometric})
 \leq\tfrac12\PseudometricError_\Correspondence(\Pseudometric,\ComparisonPseudometric).
\end{equation}
\item\label{item:quotient-2}
\paraid{p-ef3a293af6ca}
If
$\Pseudometric_0,\Pseudometric_1$
are continuous pseudometrics on
$\BaseCarrier$
and
$\ComparisonPseudometric_0,\ComparisonPseudometric_1$
are continuous pseudometrics on
$\ComparisonCarrier$,
then
\begin{equation}\label{eq:pseudometric-uniform-comparison}
 \bigl|\norm{\Pseudometric_0-\Pseudometric_1}_\infty
       -\norm{\ComparisonPseudometric_0-\ComparisonPseudometric_1}_\infty\bigr|
 \leq\PseudometricError_\Correspondence(\Pseudometric_0,\ComparisonPseudometric_0)
       +\PseudometricError_\Correspondence(\Pseudometric_1,\ComparisonPseudometric_1).
\end{equation}
\item\label{item:quotient-3}
\paraid{p-b2b1a301c639}
Let the pseudometrics satisfy the hypotheses in \ref{item:quotient-2}
and let
$\HomotopyTime\in[0,1]$.
Then
\begin{equation}\label{eq:pseudometric-interpolation}
 \PseudometricError_\Correspondence((1-\HomotopyTime)\Pseudometric_0+\HomotopyTime\Pseudometric_1,
                   (1-\HomotopyTime)\ComparisonPseudometric_0+\HomotopyTime\ComparisonPseudometric_1)
 \leq(1-\HomotopyTime)\PseudometricError_\Correspondence(\Pseudometric_0,\ComparisonPseudometric_0)
       +\HomotopyTime\PseudometricError_\Correspondence(\Pseudometric_1,\ComparisonPseudometric_1).
\end{equation}
\item\label{item:quotient-4}
\paraid{p-45401a311836}
For continuous pseudometrics
$\Pseudometric_0,\Pseudometric_1$
on a nonempty compact space
$\BaseCarrier$,
\begin{equation}\label{eq:pseudometric-bound}
 \GHDistance(\MetricQuotient{\BaseCarrier}{\Pseudometric_0},\MetricQuotient{\BaseCarrier}{\Pseudometric_1})
 \leq\tfrac12\norm{\Pseudometric_0-\Pseudometric_1}_\infty.
\end{equation}
\end{enumerate}
\end{proposition}

\begin{proof}
\textbf{Proof of \ref{item:quotient-1}. Correspondence estimate.}
\paraid{p-04c1d56a107a}
Let
$\QuotientMap\colon\BaseCarrier\to\BaseCarrier/{\sim_\Pseudometric}$
and
$\SecondQuotientMap\colon\ComparisonCarrier\to\ComparisonCarrier/{\sim_\ComparisonPseudometric}$
be the quotient maps.
The set
\begin{equation}\label{eq:quotient-correspondence}
 \{(\QuotientMap(\BasePoint),\SecondQuotientMap(\ComparisonPoint))
       \mid(\BasePoint,\ComparisonPoint)\in\Correspondence\}
\end{equation}
is a correspondence because both quotient maps and both coordinate
projections of
$\Correspondence$
are surjective.
For two elements of the correspondence in
\eqref{eq:quotient-correspondence},
definition \eqref{eq:correspondence-pseudometric-error} implies
\begin{equation}\label{eq:quotient-correspondence-error}
 |\Pseudometric(\BasePoint_0,\BasePoint_1)
   -\ComparisonPseudometric(\ComparisonPoint_0,\ComparisonPoint_1)|
 \leq\PseudometricError_\Correspondence(\Pseudometric,\ComparisonPseudometric).
\end{equation}
Applying \Cref{thm:gh-correspondence} to \eqref{eq:quotient-correspondence-error}
proves \eqref{eq:pseudometric-correspondence-bound}.

\textbf{Proof of \ref{item:quotient-2}. Uniform comparison.}
\paraid{p-1e5ed86ba147}
Fix
$(\BasePoint_0,\ComparisonPoint_0),(\BasePoint_1,\ComparisonPoint_1)\in\Correspondence$.
The triangle inequality on the real line implies
\begin{equation}\label{eq:pseudometric-pointwise-comparison}
 |\Pseudometric_0(\BasePoint_0,\BasePoint_1)
    -\Pseudometric_1(\BasePoint_0,\BasePoint_1)|
 \leq|\ComparisonPseudometric_0(\ComparisonPoint_0,\ComparisonPoint_1)
         -\ComparisonPseudometric_1(\ComparisonPoint_0,\ComparisonPoint_1)|
 +\PseudometricError_\Correspondence(\Pseudometric_0,\ComparisonPseudometric_0)
       +\PseudometricError_\Correspondence(\Pseudometric_1,\ComparisonPseudometric_1).
\end{equation}
Every pair in
$\BaseCarrier\times\BaseCarrier$
has a corresponding pair in
$\ComparisonCarrier\times\ComparisonCarrier$.
Taking the supremum in \eqref{eq:pseudometric-pointwise-comparison},
we obtain
\begin{equation}\label{eq:pseudometric-sup-comparison}
 \norm{\Pseudometric_0-\Pseudometric_1}_\infty
 \leq\norm{\ComparisonPseudometric_0-\ComparisonPseudometric_1}_\infty
      +\PseudometricError_\Correspondence(\Pseudometric_0,\ComparisonPseudometric_0)
      +\PseudometricError_\Correspondence(\Pseudometric_1,\ComparisonPseudometric_1).
\end{equation}
Applying \eqref{eq:pseudometric-sup-comparison} after interchanging the carriers
proves \eqref{eq:pseudometric-uniform-comparison}.

\textbf{Proof of \ref{item:quotient-3}. Interpolation.}
\paraid{p-94f3dbc66479}
Fix
$(\BasePoint_0,\ComparisonPoint_0),(\BasePoint_1,\ComparisonPoint_1)\in\Correspondence$.
The difference of the interpolated pseudometrics at these pairs is
\begin{equation}\label{eq:pseudometric-pointwise-interpolation}
 (1-\HomotopyTime)
   \bigl(\Pseudometric_0(\BasePoint_0,\BasePoint_1)
       -\ComparisonPseudometric_0(\ComparisonPoint_0,\ComparisonPoint_1)\bigr)
 +\HomotopyTime
   \bigl(\Pseudometric_1(\BasePoint_0,\BasePoint_1)
       -\ComparisonPseudometric_1(\ComparisonPoint_0,\ComparisonPoint_1)\bigr).
\end{equation}
The absolute value of \eqref{eq:pseudometric-pointwise-interpolation} is bounded by the weighted sum of the two errors,
since both weights are nonnegative.
Taking the supremum and using \eqref{eq:correspondence-pseudometric-error}
proves \eqref{eq:pseudometric-interpolation}.

\textbf{Proof of \ref{item:quotient-4}. A common carrier.}
\paraid{p-95c5873e8f8e}
Use
$\Correspondence=\{(\BasePoint,\BasePoint)\mid\BasePoint\in\BaseCarrier\}$.
By \eqref{eq:correspondence-pseudometric-error},
\begin{equation}\label{eq:pseudometric-diagonal-error}
 \PseudometricError_\Correspondence(\Pseudometric_0,\Pseudometric_1)
 =\norm{\Pseudometric_0-\Pseudometric_1}_\infty.
\end{equation}
Substituting \eqref{eq:pseudometric-diagonal-error} into
\eqref{eq:pseudometric-correspondence-bound} proves \eqref{eq:pseudometric-bound}.
\end{proof}

% Retired paragraph IDs.
% \paraid{p-1248fc16354c}
% \paraid{p-83a90d959336}

\subsection{Probability measures and weak convergence}\label{subsec:preliminaries-probability}
\paraid{p-553202822150}
We collect the measure-theoretic tools used for averaging probabilities
and comparing distance operators on varying compact spaces.

\paraid{p-a5436302dabd}
For a metrizable space
$\ProbabilityCarrier $,
let
$\BorelSets(\ProbabilityCarrier )$
be its Borel
$\SigmaAlgebraPrefix$-algebra, namely, the
$\SigmaAlgebraPrefix$-algebra generated by the open subsets of
$\ProbabilityCarrier$.
A finite Borel measure
$\FirstMeasure$
on
$\ProbabilityCarrier$
is called \emph{Radon} if it is inner regular by compact sets,
meaning that for every
$\BorelSubset\in\BorelSets(\ProbabilityCarrier)$
we have
\[
 \FirstMeasure(\BorelSubset)
 =\sup\{\FirstMeasure(\CompactSubset)\mid
          \CompactSubset\subset\BorelSubset\text{ is compact}\}.
\]
We denote by
$\ProbabilityMeasures(\ProbabilityCarrier)$
the set of Radon measures with total mass one.
A \emph{probability law} on
$\ProbabilityCarrier$
is an element of
$\ProbabilityMeasures(\ProbabilityCarrier)$.
For Radon probabilities
$\FirstMeasure\in\ProbabilityMeasures(\ProbabilityCarrier)$
and
$\SecondMeasure\in\ProbabilityMeasures(\SecondProbabilityCarrier)$,
the notation
$\FirstMeasure\otimes\SecondMeasure$
denotes their product probability on
$\ProbabilityCarrier\times\SecondProbabilityCarrier$.
For every
$\BorelSubset\in\BorelSets(\ProbabilityCarrier)$
and
$\ClosedTestSet\in\BorelSets(\SecondProbabilityCarrier)$,
it satisfies
\[
 (\FirstMeasure\otimes\SecondMeasure)(\BorelSubset\times\ClosedTestSet)
 =\FirstMeasure(\BorelSubset)\SecondMeasure(\ClosedTestSet).
\]
These identities characterize the product probability.
For a Radon probability
$\FirstMeasure\in\ProbabilityMeasures(\ProbabilityCarrier)$,
its \emph{support} is
\[
 \supp(\FirstMeasure)
 =\{\BasePoint\in\ProbabilityCarrier\mid
    \FirstMeasure(\OpenSubset)>0\text{ for every open }\OpenSubset\subset\ProbabilityCarrier
    \text{ with }\BasePoint\in\OpenSubset\}.
\]
We say that
$\FirstMeasure$
has \emph{full support} if
$\supp(\FirstMeasure)=\ProbabilityCarrier$.
Equivalently,
\[
 \forall\OpenSubset\subset\ProbabilityCarrier\text{ open and nonempty}\quad
 \FirstMeasure(\OpenSubset)>0.
\]

\paraid{p-2c1ffb2b3240}
If
$\Mapping \colon \ProbabilityCarrier \to \SecondProbabilityCarrier $
is continuous and
$\FirstMeasure \in\ProbabilityMeasures(\ProbabilityCarrier )$,
its pushforward is the probability
$\Mapping \Pushforward\FirstMeasure $
defined for every
$\BorelSubset\in\BorelSets(\SecondProbabilityCarrier)$
by
\[
 (\Mapping\Pushforward\FirstMeasure)(\BorelSubset)
 =\FirstMeasure(\Mapping^{-1}(\BorelSubset)).
\]
For a nonempty metrizable space
$\ProbabilityCarrier$,
we write
\[
 \BoundedContinuousFunctions(\ProbabilityCarrier)
 =\{\TestFunction\colon\ProbabilityCarrier\to\RealNumbers\mid
       \TestFunction\text{ is continuous and bounded}\},
\]
and
\[
 \norm{\TestFunction}_\infty=\sup_{\BasePoint\in\ProbabilityCarrier}|\TestFunction(\BasePoint)|.
\]
For any compact metric space
$\BaseCarrier$,
we define
\[
 \ContinuousFunctions(\BaseCarrier)=\ContinuousFunctions(\BaseCarrier,\RealNumbers)
 =\{\TestFunction\colon\BaseCarrier\to\RealNumbers\mid\TestFunction\text{ is continuous}\},
\]
and
\[
 \norm{\TestFunction}_\infty=\max_{\BasePoint\in\BaseCarrier}|\TestFunction(\BasePoint)|.
\]
This is a real Banach space.
With continuous functions taking values in
$\ComplexNumbers$,
we obtain the complex Banach space
$\ContinuousFunctions(\BaseCarrier,\ComplexNumbers)$.

\paraid{p-bca1c91c25a9}
We equip
$\ProbabilityMeasures(\ProbabilityCarrier)$
with the weak topology,
the coarsest topology for which,
for every
$\TestFunction\in\BoundedContinuousFunctions(\ProbabilityCarrier)$,
the map
\[
 \ProbabilityMeasures(\ProbabilityCarrier)\to\RealNumbers
 \]
 by
 \[\FirstMeasure\longmapsto\int_\ProbabilityCarrier\TestFunction\,
 \IntegrationDifferential\FirstMeasure
\]
is continuous.
We write
$\FirstMeasure_\SequenceIndex\to\FirstMeasure$
weakly when for every
$\TestFunction\in\BoundedContinuousFunctions(\ProbabilityCarrier)$
we have
\[
 \int_\ProbabilityCarrier\TestFunction\,\IntegrationDifferential\FirstMeasure_\SequenceIndex
 \longrightarrow
 \int_\ProbabilityCarrier\TestFunction\,\IntegrationDifferential\FirstMeasure.
\]
On a compact space all continuous functions are bounded.
For bounded complex-valued continuous test functions the same convergence follows by
applying the definition to their real and imaginary parts.

\paraid{p-f85199b8afe5}
We construct averages of measures by specifying their integrals of continuous functions.
The Riesz--Markov--Kakutani theorem states that every positive normalized
linear functional on the real continuous functions on a nonempty compact
metric space is integration against a unique Radon probability measure.

\begin{theorem}[{Riesz--Markov--Kakutani, \cite[Theorem~7.10.4 and Corollary~7.10.5]{Bogachev2007II}}]\label{thm:riesz-markov-kakutani}
\paraid{p-37b86f9f43f7}
Let
$\MetricSpace{\BaseCarrier}{\MetricSymbol}$
be a nonempty compact metric space.
Let
\[
 \RepresentingFunctional\colon\ContinuousFunctions(\BaseCarrier,\RealNumbers)\to\RealNumbers
\]
be a linear functional.
Assume that for every
$\TestFunction\in\ContinuousFunctions(\BaseCarrier,\RealNumbers)$
with
$\TestFunction\geq0$
we have
\[
 \RepresentingFunctional(\TestFunction)\geq0.
\]
Let
$1$
denote the constant function with value one and assume that
\[
 \RepresentingFunctional(1)=1.
\]
Then there is a unique Radon probability
$\FirstMeasure\in\ProbabilityMeasures(\BaseCarrier)$
such that for every
$\TestFunction\in\ContinuousFunctions(\BaseCarrier,\RealNumbers)$,
we have
\[
 \RepresentingFunctional(\TestFunction)
 =\int_\BaseCarrier\TestFunction\,\IntegrationDifferential\FirstMeasure.
\]
\end{theorem}

\paraid{p-d1b3a84b65e1}
Positivity and normalization imply
$|\RepresentingFunctional(\TestFunction)|\leq\norm{\TestFunction}_\infty$,
so the functional
$\RepresentingFunctional$
is automatically continuous.
The cited representation theorem applies,
and evaluating at the constant function
$1$
proves that the representing measure
$\FirstMeasure$
has total mass one.

\begin{theorem}[{Portmanteau, \cite[Theorem~2.1]{Billingsley1999}, \cite[Theorem~8.2.3, Corollary~8.2.4(a)]{Bogachev2007II}}]\label{thm:portmanteau}
\paraid{p-e85c034ef572}
Let
$\ProbabilityCarrier$
be a metrizable space and let
$\FirstMeasure,\FirstMeasure_\SequenceIndex\in\ProbabilityMeasures(\ProbabilityCarrier)$
for
$\SequenceIndex\in\NonnegativeIntegers$.
Then the following conditions are equivalent.
\begin{enumerate}[label=\textup{(P\arabic*)},ref=\textup{(P\arabic*)}]
\item\label{item:portmanteau-weak}
\paraid{p-7c0d10bcd5a8}
For every
$\TestFunction\in\BoundedContinuousFunctions(\ProbabilityCarrier)$,
\[
 \lim_{\SequenceIndex\to\infty}\int_{\ProbabilityCarrier}\TestFunction\,\IntegrationDifferential\FirstMeasure_\SequenceIndex
 =\int_{\ProbabilityCarrier}\TestFunction\,\IntegrationDifferential\FirstMeasure.
\]
\item
\paraid{p-d9674e60f731}
For every closed subset
$\ClosedTestSet\subset\ProbabilityCarrier$,
\[
 \limsup_{\SequenceIndex\to\infty}\FirstMeasure_\SequenceIndex(\ClosedTestSet)
 \leq\FirstMeasure(\ClosedTestSet).
\]
\item
\paraid{p-533dc3ffdbf5}
For every open subset
$\OpenSubset\subset\ProbabilityCarrier$,
\[
 \FirstMeasure(\OpenSubset)
 \leq\liminf_{\SequenceIndex\to\infty}\FirstMeasure_\SequenceIndex(\OpenSubset).
\]
\end{enumerate}
\end{theorem}

\paraid{p-02818be37bab}
We also use the following consequence for measures on varying compact sets.

\begin{lemma}\label{lem:limit-support}
\paraid{p-a127ca5a2f91}
Let
$\FirstCompactSet_\SequenceIndex,\FirstCompactSet$
be nonempty compact subsets of a compact metric space
$\MetricSpace{\AmbientCarrier}{\AmbientMetric}$
with
$\HausdorffDistance{\AmbientMetric}(\FirstCompactSet_\SequenceIndex,\FirstCompactSet)\to0$.
If
$\FirstMeasure_\SequenceIndex\to\FirstMeasure$
weakly in
$\ProbabilityMeasures(\AmbientCarrier)$
and
$\FirstMeasure_\SequenceIndex(\FirstCompactSet_\SequenceIndex)=1$,
then
$\FirstMeasure(\FirstCompactSet)=1$.
\end{lemma}

\begin{proof}
\paraid{p-323518997d9a}
Define the function
$\TestFunction\colon\AmbientCarrier\to\RealNumbers$
by
\[
 \TestFunction(\IntegrationPoint)
 =\dist_{\AmbientMetric}(\IntegrationPoint,\FirstCompactSet).
\]
Since
$\AmbientCarrier$
is compact,
$\TestFunction$
is bounded and continuous.
Weak convergence and
$\FirstMeasure_\SequenceIndex(\FirstCompactSet_\SequenceIndex)=1$
imply
\[
 \begin{aligned}
 0\leq\int_\AmbientCarrier\TestFunction(\IntegrationPoint)
       \,\IntegrationDifferential\FirstMeasure(\IntegrationPoint)
 &=\lim_\SequenceIndex\int_\AmbientCarrier\TestFunction(\IntegrationPoint)
       \,\IntegrationDifferential\FirstMeasure_\SequenceIndex(\IntegrationPoint)\\
 &\leq\lim_\SequenceIndex
       \HausdorffDistance{\AmbientMetric}(\FirstCompactSet_\SequenceIndex,\FirstCompactSet)
 =0.
 \end{aligned}
\]
Since
$\TestFunction^{-1}(\{0\})=\FirstCompactSet$,
we obtain
$\FirstMeasure(\FirstCompactSet)=1$.
\end{proof}

\paraid{p-241d95ad9f1f}
Let
$\MetricSpace{\AmbientCarrier}{\AmbientMetric}$
be a compact metric space.
For
$\BorelSubset\in\BorelSets(\AmbientCarrier)$
and
$\SmallError>0$,
put
\[
 \BorelSubset^{\SmallError}
 =\{\BasePoint\in\AmbientCarrier\mid
       \dist_{\AmbientMetric}(\BasePoint,\BorelSubset)<\SmallError\}.
\]
For probabilities
$\FirstMeasure,\SecondMeasure\in\ProbabilityMeasures(\AmbientCarrier)$,
we define the \emph{L\'evy--Prokhorov distance} by
\[
 \ProkhorovDistance{\AmbientMetric}(\FirstMeasure,\SecondMeasure)
 =\inf\left\{\SmallError>0\;\middle|\;
 \begin{gathered}
 \text{for every }\BorelSubset\in\BorelSets(\AmbientCarrier),\\
 \FirstMeasure(\BorelSubset)\leq\SecondMeasure(\BorelSubset^{\SmallError})+\SmallError,\\
 \SecondMeasure(\BorelSubset)\leq\FirstMeasure(\BorelSubset^{\SmallError})+\SmallError
 \end{gathered}\right\}.
\]
This is a metric on
$\ProbabilityMeasures(\AmbientCarrier)$
and induces weak convergence
\cite[Section~6, pp.~72--73, Theorem~6.8]{Billingsley1999},
\cite[author's manuscript, Definition~1.12, Proposition~1.13 and Lemma~1.15]{ShioyaMetricMeasure}.

\begin{lemma}[{\cite[arXiv v5, Example~2.1(iii)]{Khezeli2023Framework}}]\label{lem:prokhorov-isometric-embedding}
\paraid{p-8ef64efd5d7f}
Let
$\MetricSpace{\BaseCarrier}{\MetricSymbol}$
and
$\MetricSpace{\AmbientCarrier}{\AmbientMetric}$
be compact metric spaces, and let
$\FirstEmbedding\colon\MetricSpace{\BaseCarrier}{\MetricSymbol}\to\MetricSpace{\AmbientCarrier}{\AmbientMetric}$
be an isometric embedding.
For all
$\FirstMeasure,\SecondMeasure\in\ProbabilityMeasures(\BaseCarrier)$,
we have
\[
 \ProkhorovDistance{\AmbientMetric}
 (\FirstEmbedding\Pushforward\FirstMeasure,\FirstEmbedding\Pushforward\SecondMeasure)
 =\ProkhorovDistance{\MetricSymbol}(\FirstMeasure,\SecondMeasure).
\]
\end{lemma}

% \paraid{p-d40306374206}

\paraid{p-b04d75056e40}
Every Borel probability on a Polish space is Radon
\cite[Theorem~7.1.7]{Bogachev2007II}.
We use the following theorem for Radon probability measures.

% \paraid{p-c3d38955d956}

\begin{theorem}[{\cite[Theorems~8.9.3 and~8.9.4]{Bogachev2007II}}]\label{thm:probability-topology}
\paraid{p-250872571fa5}
Let
$\ProbabilityCarrier$
be a separable metrizable space.
Then
$\ProbabilityMeasures(\ProbabilityCarrier)$
is separable and metrizable in the weak topology.
If
$\ProbabilityCarrier$
is Polish,
then
$\ProbabilityMeasures(\ProbabilityCarrier)$
is Polish.
If
$\ProbabilityCarrier$
is compact,
then
$\ProbabilityMeasures(\ProbabilityCarrier)$
is compact.
\end{theorem}

% \paraid{p-4be0d8c030c3}
% \paraid{p-489f0dc304a2}



\paraid{p-fd786a47d1fa}
Weak convergence is uniform over compact families of continuous test functions.

\begin{lemma}\label{lem:uniform-weak-integrals}
\paraid{p-97f9eb06751a}
Let
$\AmbientCarrier $
be a nonempty compact metric space,
let
$\FirstMeasure _\SequenceIndex \to\FirstMeasure $
weakly in
$\ProbabilityMeasures(\AmbientCarrier )$,
and let
$\CompactTestFamily \subset\ContinuousFunctions(\AmbientCarrier )$
be nonempty and compact in the uniform norm.
Then
\begin{equation}\label{eq:uniform-test-integrals}
 \sup_{\Mapping \in\CompactTestFamily }\left|\int_\AmbientCarrier  \Mapping \,\IntegrationDifferential \FirstMeasure _\SequenceIndex -\int_\AmbientCarrier  \Mapping \,\IntegrationDifferential \FirstMeasure \right|\to0.
\end{equation}
The same conclusion holds for complex continuous functions.
\end{lemma}

\begin{proof}
\paraid{p-3a52e680c1c2}
For each
$\SequenceIndex$,
choose
$\Mapping_\SequenceIndex\in\CompactTestFamily$
attaining the supremum in \eqref{eq:uniform-test-integrals}.
By compactness of
$\CompactTestFamily$,
every subsequence has a further subsequence
$\{\Mapping_{n_k}\}_{k\in\NonnegativeIntegers}$
converging uniformly to some
$\Mapping\in\CompactTestFamily$.
For these indices,
\[
 \left|\int_\AmbientCarrier\Mapping_{n_k}\,
           \IntegrationDifferential(\FirstMeasure_{n_k}-\FirstMeasure)\right|
 \leq2\norm{\Mapping_{n_k}-\Mapping}_\infty
      +\left|\int_\AmbientCarrier\Mapping\,
           \IntegrationDifferential(\FirstMeasure_{n_k}-\FirstMeasure)\right|
 \longrightarrow0.
\]
Thus the whole sequence of suprema tends to
$0$.
The same argument applies to complex functions,
since weak convergence controls both their real and imaginary parts.
\end{proof}

\paraid{p-9b10954ce9e8}
We use \Cref{lem:uniform-weak-integrals} to prove uniform convergence
of integrals whose integrands vary uniformly with compact parameters.

\begin{lemma}\label{lem:parameter-integrals}
\paraid{p-dc7ad8c030f6}
Let
$\CompactParameterSpace$
and
$\AmbientCarrier$
be nonempty compact metric spaces.
Let
$\FirstMeasure_\SequenceIndex\to\FirstMeasure$
weakly in
$\ProbabilityMeasures(\AmbientCarrier)$.
Let
$\Integrand_\SequenceIndex,\Integrand\colon\CompactParameterSpace\times\AmbientCarrier\to\ComplexNumbers$
be continuous real-valued or complex-valued functions.
Assume that
$\norm{\Integrand_\SequenceIndex-\Integrand}_\infty\to0$.
Then
\[
 \sup_{\ParameterPoint\in\CompactParameterSpace}
 \left|\int_\AmbientCarrier\Integrand_\SequenceIndex(\ParameterPoint,\IntegrationPoint)\,\IntegrationDifferential\FirstMeasure_\SequenceIndex(\IntegrationPoint)
 -\int_\AmbientCarrier\Integrand(\ParameterPoint,\IntegrationPoint)\,\IntegrationDifferential\FirstMeasure(\IntegrationPoint)\right|\to0.
\]
\end{lemma}

\begin{proof}
\paraid{p-03fa713f9e12}
The triangle inequality yields
\begin{equation}\label{eq:parameter-integral-error}
\begin{aligned}
&\sup_{\ParameterPoint\in\CompactParameterSpace}
 \left|
 \int_\AmbientCarrier
 \Integrand_\SequenceIndex(\ParameterPoint,\IntegrationPoint)
 \,\IntegrationDifferential\FirstMeasure_\SequenceIndex(\IntegrationPoint)
 -\int_\AmbientCarrier
 \Integrand(\ParameterPoint,\IntegrationPoint)
 \,\IntegrationDifferential\FirstMeasure(\IntegrationPoint)
 \right|\\
&\leq
 \norm{\Integrand_\SequenceIndex-\Integrand}_\infty
 +\sup_{\ParameterPoint\in\CompactParameterSpace}
 \left|
 \int_\AmbientCarrier
 \Integrand(\ParameterPoint,\IntegrationPoint)
 \,\IntegrationDifferential\FirstMeasure_\SequenceIndex(\IntegrationPoint)
 -\int_\AmbientCarrier
 \Integrand(\ParameterPoint,\IntegrationPoint)
 \,\IntegrationDifferential\FirstMeasure(\IntegrationPoint)
 \right|.
\end{aligned}
\end{equation}
The uniform error
$\norm{\Integrand_\SequenceIndex-\Integrand}_\infty$
in \eqref{eq:parameter-integral-error} tends to
$0$
by assumption.
The family
$\{\Integrand(\ParameterPoint,\cdot)\mid\ParameterPoint\in\CompactParameterSpace\}$
is compact in the uniform norm,
so \Cref{lem:uniform-weak-integrals} proves that the supremum of the integral differences
on the right-hand side of \eqref{eq:parameter-integral-error} tends to
$0$.
By \eqref{eq:parameter-integral-error},
the parameter integrals converge uniformly on
$\CompactParameterSpace$.
\end{proof}



\paraid{p-c858cdb38f3d}
We use the following continuous functions to detect the support of a measure
and approximate a distance profile by an integral.

\begin{lemma}\label{lem:metric-bump}
\paraid{p-e99ec25288a5}
Let
$\MetricSpace{\BaseCarrier}{\MetricSymbol}$
be a nonempty compact metric space,
let
$\FirstMeasure\in\ProbabilityMeasures(\BaseCarrier)$,
and let
$\Radius>0$.
For
$\BasePoint\in\BaseCarrier$,
define
\[
 \BumpFunction{\BaseCarrier}{\MetricSymbol}{\BasePoint}{\Radius}
 \colon\BaseCarrier\to[0,1]
\]
by
\begin{equation}\label{eq:metric-bump-definition}
 \BumpFunction{\BaseCarrier}{\MetricSymbol}{\BasePoint}{\Radius}(\ComparisonPoint)
 =\max\{1-\MetricSymbol(\BasePoint,\ComparisonPoint)/\Radius,0\}.
\end{equation}
Then for all
$\BasePoint,\BasePoint_0,\ComparisonPoint,\ComparisonPoint_0\in\BaseCarrier$,
\begin{equation}\label{eq:metric-bump-lipschitz}
 \left|\BumpFunction{\BaseCarrier}{\MetricSymbol}{\BasePoint}{\Radius}(\ComparisonPoint)
 -\BumpFunction{\BaseCarrier}{\MetricSymbol}{\BasePoint_0}{\Radius}(\ComparisonPoint_0)\right|
 \leq\frac{\MetricSymbol(\BasePoint,\BasePoint_0)
              +\MetricSymbol(\ComparisonPoint,\ComparisonPoint_0)}{\Radius}.
\end{equation}
Moreover,
\begin{equation}\label{eq:metric-bump-integral}
 \frac12\FirstMeasure\bigl(\MetricBall(\BasePoint,\Radius/2;\MetricSymbol)\bigr)
 \leq\int_\BaseCarrier
       \BumpFunction{\BaseCarrier}{\MetricSymbol}{\BasePoint}{\Radius}
       \,\IntegrationDifferential\FirstMeasure
 \leq\FirstMeasure\bigl(\MetricBall(\BasePoint,\Radius;\MetricSymbol)\bigr).
\end{equation}
\end{lemma}

\begin{proof}
\paraid{p-886f82531c59}
The function
$\SecondParameter\mapsto\max\{1-\SecondParameter/\Radius,0\}$
is
$1/\Radius$-Lipschitz.
The triangle inequality for
$\MetricSymbol$
therefore proves \eqref{eq:metric-bump-lipschitz}.
The pointwise bounds
\[
 \tfrac12\Indicator_{\MetricBall(\BasePoint,\Radius/2;\MetricSymbol)}
 \leq\BumpFunction{\BaseCarrier}{\MetricSymbol}{\BasePoint}{\Radius}
 \leq\Indicator_{\MetricBall(\BasePoint,\Radius;\MetricSymbol)}
\]
imply \eqref{eq:metric-bump-integral} by integration.
\end{proof}

\begin{lemma}\label{lem:normalized-metric-bump}
\paraid{p-ac4da2b9d373}
Let
$\MetricSpace{\BaseCarrier}{\MetricSymbol}$
be a nonempty compact metric space,
let
$\FirstMeasure$
be a probability measure on
$\BaseCarrier$
with full support,
and let
$\Radius>0$.
For every
$\BasePoint\in\BaseCarrier$,
the function
\begin{equation}\label{eq:normalized-bump-definition}
 \NormalizedBump{\BaseCarrier}{\MetricSymbol}{\BasePoint}{\Radius}{\FirstMeasure}
 =\frac{\BumpFunction{\BaseCarrier}{\MetricSymbol}{\BasePoint}{\Radius/2}}
 {\displaystyle\int_\BaseCarrier
      \BumpFunction{\BaseCarrier}{\MetricSymbol}{\BasePoint}{\Radius/2}
      \,\IntegrationDifferential\FirstMeasure}
\end{equation}
is well-defined and continuous,
and satisfies
\begin{equation}\label{eq:normalized-bump-properties-1}
 \NormalizedBump{\BaseCarrier}{\MetricSymbol}{\BasePoint}{\Radius}{\FirstMeasure}\geq0,
\end{equation}
and
\begin{equation}\label{eq:normalized-bump-properties-2}
 \int_\BaseCarrier
 \NormalizedBump{\BaseCarrier}{\MetricSymbol}{\BasePoint}{\Radius}{\FirstMeasure}
 \,\IntegrationDifferential\FirstMeasure=1,
\end{equation}
and
\begin{equation}\label{eq:normalized-bump-properties-3}
 \supp\left(\NormalizedBump{\BaseCarrier}{\MetricSymbol}{\BasePoint}{\Radius}{\FirstMeasure}\right)
 \subset\MetricBall(\BasePoint,\Radius;\MetricSymbol).
\end{equation}
\end{lemma}

\begin{proof}
\paraid{p-84a66f6444a3}
Fix
$\BasePoint\in\BaseCarrier$.
Since
$\supp(\FirstMeasure)=\BaseCarrier$,
\eqref{eq:metric-bump-integral} with radius
$\Radius/2$
proves that the denominator in \eqref{eq:normalized-bump-definition} is positive.
Continuity and \eqref{eq:normalized-bump-properties-1}--\eqref{eq:normalized-bump-properties-2}
follow from \eqref{eq:metric-bump-lipschitz} and \eqref{eq:normalized-bump-definition}.
By \eqref{eq:metric-bump-definition},
\[
 \supp\left(\NormalizedBump{\BaseCarrier}{\MetricSymbol}{\BasePoint}{\Radius}{\FirstMeasure}\right)
 \subset\{\ComparisonPoint\in\BaseCarrier\mid
              \MetricSymbol(\BasePoint,\ComparisonPoint)\leq\Radius/2\}
 \subset\MetricBall(\BasePoint,\Radius;\MetricSymbol).
\]
This proves \eqref{eq:normalized-bump-properties-3}.
\end{proof}

\paraid{p-0c1ccd024c12}
The next lemma compares product measures when one factor is a point mass.
It will be used to average laws over a convergent family of compact spaces.

\begin{lemma}[{\cite[Chapter~III, Lemma~1.1]{Parthasarathy2005}}]\label{lem:dirac-product-convergence}
\paraid{p-45f1a6509071}
Let
$\CompactParameterSpace$
be a compact metric space and let
$\ProbabilityCarrier$
be a Polish space.
Let
$\ParameterPoint_\SequenceIndex\to\ParameterPoint$
in
$\CompactParameterSpace$
and
$\FirstMeasure_\SequenceIndex\to\FirstMeasure$
weakly in
$\ProbabilityMeasures(\ProbabilityCarrier)$.
For a point
$\ParameterPoint\in\CompactParameterSpace$,
write
$\DiracProbability_\ParameterPoint$
for the probability concentrated at
$\ParameterPoint$.
Then
\begin{equation}\label{eq:dirac-product-convergence}
 \DiracProbability_{\ParameterPoint_\SequenceIndex}\otimes\FirstMeasure_\SequenceIndex
 \longrightarrow\DiracProbability_\ParameterPoint\otimes\FirstMeasure
 \quad\text{weakly in }\ProbabilityMeasures(\CompactParameterSpace\times\ProbabilityCarrier).
\end{equation}
\end{lemma}

% \paraid{p-c62c3005a39f}



\paraid{p-2cd0585c4b58}
For comparisons with metric measure geometry,
we also recall the box distance.

\begin{definition}\label{def:box-distance}
\paraid{p-b4f453559431}
An \emph{mm-space} is a complete separable metric space
$(X,d_X)$
with a Borel probability measure
$\mu_X$.
Two mm-spaces are \emph{mm-isomorphic} if there is a measure-preserving
isometry between their supports.
Let
$\MMClassSpace$
denote the set of these classes.
Put
$\BoxParameterInterval=[0,1)$
and let
$\BoxLebesgueMeasure$
be Lebesgue measure
on this interval.
A \emph{parameter} of
$X$
is a Borel map
$\BoxFirstParameter\colon\BoxParameterInterval\to X$
with
$(\BoxFirstParameter)_*\BoxLebesgueMeasure=\mu_X$.
Every mm-space has a parameter.
The \emph{box distance}
$\BoxDistance(X,Y)$
is the infimum of all
$\ErrorTolerance>0$
for which there are parameters
$\BoxFirstParameter$
of
$X$
and
$\BoxSecondParameter$
of
$Y$
and a Borel set
$\BoxLargeSubset\subset\BoxParameterInterval$
satisfying
\[
 \BoxLebesgueMeasure(\BoxLargeSubset)\geq1-\ErrorTolerance,
\]
and such that for every
$s,t\in\BoxLargeSubset$
we have
\[
 \left|d_X(\BoxFirstParameter(s),\BoxFirstParameter(t))
       -d_Y(\BoxSecondParameter(s),\BoxSecondParameter(t))\right|
 \leq\ErrorTolerance.
\]
It is a metric on
$\MMClassSpace$.
The topology it induces is called the \emph{box topology}
\cite[arXiv v1, Definitions~2.1, 2.6, 2.7 and Theorem~2.8]{KazukawaNakajimaShioya2024}.
\end{definition}

\subsection{Partitions of unity, contractions and absolute retracts}\label{subsec:preliminaries-topology}
\paraid{p-33a116cd3f4a}
We recall partitions of unity,
an extension theorem for continuous real functions,
and the terminology for contractions and absolute retracts.

\paraid{p-4f4ac6ddbc21}
A family of subsets of a topological space is \emph{locally finite} if every point
has a neighborhood meeting only finitely many members of the family.
A cover \emph{refines} another cover if each member of the first is contained
in a member of the second.
A Hausdorff space is \emph{paracompact} if every open cover admits a locally finite
open refinement.

\paraid{p-bcbfe50f69ab}
Let
$X$
be a topological space and let
$I$
be an index set.
A family of continuous functions
$\lambda_i\colon X\to[0,1]$,
indexed by
$i\in I$,
is a \emph{locally finite partition of unity} if the supports
\[
 \supp\lambda_i=\overline{\{x\in X\mid\lambda_i(x)\ne0\}}
\]
form a locally finite family and for every
$x\in X$
we have
\[
 \sum_{i\in I}\lambda_i(x)=1.
\]
It is \emph{subordinate} to an open cover if each support is contained in a member
of that cover.

\begin{theorem}[{\cite[Theorems~1.3.33 and~1.3.38]{CobzasMiculescuNicolae2019}}]\label{thm:metric-partition-unity}
\paraid{p-e8f70dfc07eb}
Every metrizable space is paracompact.
Every open cover of a metrizable space admits a locally finite partition of unity
subordinate to it.
\end{theorem}

\begin{theorem}[{Tietze--Urysohn, \cite[Theorem~2.1.8]{Engelking1989}}]\label{thm:tietze}
\paraid{p-fb8cfff8e576}
Let
$\ExtensionDomain$
be a normal space,
let
$\ClosedDomain\subset\ExtensionDomain$
be closed,
and let
$\FunctionBound\geq0$.
Then every continuous function
$\Mapping\colon\ClosedDomain\to[-\FunctionBound,\FunctionBound]$
has a continuous extension
$\TietzeExtension\colon\ExtensionDomain\to[-\FunctionBound,\FunctionBound]$.
\end{theorem}

\paraid{p-c07e9f5bb5bf}
A continuous map
$\Mapping \colon \ClosedDomain \to \ExtensionTarget $
is \emph{null-homotopic} if there exist a point
$\ConstantPoint \in \ExtensionTarget $
and a continuous map
$\ContractionHomotopy \colon \ClosedDomain \times\UnitInterval\to \ExtensionTarget $
such that for every
$\FirstScale\in\ClosedDomain$
we have
\[
 \ContractionHomotopy (\FirstScale ,0)=\Mapping (\FirstScale ),
\]
and
\[
 \ContractionHomotopy (\FirstScale ,1)=\ConstantPoint.
\]
A nonempty space is \emph{contractible} if its identity map is null-homotopic.

\paraid{p-457cc7aecb3b}
The following proposition shows that scaling distances defines a contraction.
Write
$\SingletonSpace=\MetricSpace{\{0\}}{0}$
for the one-point metric space.

\begin{proposition}\label{lem:scaling-contraction}
\paraid{p-3f7987feb185}
For compact metric spaces
$\MetricSpace{\BaseCarrier }{\MetricSymbol }$
and
$\MetricSpace{\ComparisonCarrier }{\ComparisonMetric }$
and real numbers
$\FirstScale ,\SecondScale \geq0$,
we have
\begin{equation}\label{eq:scaling}
 \begin{aligned}
 \GHDistance(\MetricQuotient{\BaseCarrier }{\FirstScale  \MetricSymbol },\MetricQuotient{\ComparisonCarrier }{\SecondScale  \ComparisonMetric })
 &\leq \FirstScale \GHDistance(\MetricSpace{\BaseCarrier }{\MetricSymbol },\MetricSpace{\ComparisonCarrier }{\ComparisonMetric })
       +\frac{|\FirstScale -\SecondScale |}{2}\diam\MetricSpace{\ComparisonCarrier }{\ComparisonMetric }.
 \end{aligned}
\end{equation}
Consequently the map
\[
 \ContractionHomotopy \colon\GHSpace\times\UnitInterval\to\GHSpace
\]
defined by
\[
 \ContractionHomotopy (\MetricSpace{\BaseCarrier }{\MetricSymbol },\HomotopyTime )=\MetricQuotient{\BaseCarrier }{(1-\HomotopyTime )\MetricSymbol }
\]
is continuous and satisfies
\[
 \ContractionHomotopy (\MetricSpace{\BaseCarrier }{\MetricSymbol },0)=\MetricSpace{\BaseCarrier }{\MetricSymbol },
\]
and
\[
 \ContractionHomotopy (\MetricSpace{\BaseCarrier }{\MetricSymbol },1)=\SingletonSpace.
\]
The contraction fixes the one-point space at every time,
so
\[
 \ContractionHomotopy (\SingletonSpace,\HomotopyTime )=\SingletonSpace.
\]
In particular,
$\GHSpace$
is contractible.
\end{proposition}

\begin{proof}
\paraid{p-98b9e8081dac}
By the triangle inequality and the scaling identities in
\cite[arXiv v1, Examples~6.31 and~6.32]{Tuzhilin2020},
we obtain \eqref{eq:scaling},
also when one of the scaling factors is zero.
Since the diameter is continuous on
$\GHSpace$,
the estimate proves continuity of the map
\[
 \GHSpace\times[0,\infty)\to\GHSpace,
 \qquad
 (\MetricSpace{\BaseCarrier}{\MetricSymbol},\FirstScale)
       \longmapsto\MetricQuotient{\BaseCarrier}{\FirstScale\MetricSymbol},
\]
including at
$\FirstScale=0$.
The three identities for
$\ContractionHomotopy $
follow directly from its definition.
\end{proof}

\paraid{p-3e5d9433b03f}
For balls in Gromov--Hausdorff space we use
a different  symbol,
\[
 \GHball(\MetricSpace{\BaseCarrier}{\MetricSymbol},\Radius)
 =\{\MetricSpace{\ComparisonCarrier}{\ComparisonMetric}\in\GHSpace\mid
 \GHDistance(\MetricSpace{\BaseCarrier}{\MetricSymbol},\MetricSpace{\ComparisonCarrier}{\ComparisonMetric})<\Radius\}.
\]
The correspondence formula also implies
$\GHDistance(\MetricSpace{\BaseCarrier }{\MetricSymbol },\SingletonSpace)=\diam\MetricSpace{\BaseCarrier }{\MetricSymbol }/2$.

\paraid{p-263af0d7c090}
We state extension and retraction properties relative to a class of
topological spaces.

\begin{definition}\label{def:extension}
\paraid{p-3d55c72c99aa}
Let
$\TopologicalSpaceClass$
be a class of topological spaces.
A topological space
$\ExtensionTarget$
is an \emph{absolute extensor for $\TopologicalSpaceClass$},
or AE for
$\TopologicalSpaceClass$,
if for every space
$\BaseCarrier\in\TopologicalSpaceClass$,
every closed subset
$\ClosedDomain\subset\BaseCarrier$,
and every continuous map
$\Mapping\colon\ClosedDomain\to\ExtensionTarget$,
there exists a continuous map
$\ContinuousExtension\colon\BaseCarrier\to\ExtensionTarget$
such that
$\ContinuousExtension|_\ClosedDomain=\Mapping$.
It is an \emph{absolute neighborhood extensor for $\TopologicalSpaceClass$},
or ANE for
$\TopologicalSpaceClass$,
if for every space
$\BaseCarrier\in\TopologicalSpaceClass$,
every closed subset
$\ClosedDomain\subset\BaseCarrier$,
and every continuous map
$\Mapping\colon\ClosedDomain\to\ExtensionTarget$,
there exist an open neighborhood
$\OpenNeighborhood$
of
$\ClosedDomain$
in
$\BaseCarrier$
and a continuous map
$\ContinuousExtension\colon\OpenNeighborhood\to\ExtensionTarget$
such that
$\ContinuousExtension|_\ClosedDomain=\Mapping$.
A space
$\ExtensionTarget\in\TopologicalSpaceClass$
is an \emph{absolute retract for $\TopologicalSpaceClass$},
or AR for
$\TopologicalSpaceClass$,
if for every space
$\BaseCarrier\in\TopologicalSpaceClass$
and every embedding
$\FirstEmbedding\colon\ExtensionTarget\to\BaseCarrier$
with closed image,
there exists a continuous map
$\RetractionMap\colon\BaseCarrier\to\ExtensionTarget$
such that
$\RetractionMap\circ\FirstEmbedding=\IdentityMap_\ExtensionTarget$.
A space
$\ExtensionTarget\in\TopologicalSpaceClass$
is an \emph{absolute neighborhood retract for $\TopologicalSpaceClass$},
or ANR for
$\TopologicalSpaceClass$,
if for every space
$\BaseCarrier\in\TopologicalSpaceClass$
and every embedding
$\FirstEmbedding\colon\ExtensionTarget\to\BaseCarrier$
with closed image,
there exist an open neighborhood
$\OpenNeighborhood$
of
$\FirstEmbedding(\ExtensionTarget)$
in
$\BaseCarrier$
and a continuous map
$\RetractionMap\colon\OpenNeighborhood\to\ExtensionTarget$
such that
$\RetractionMap\circ\FirstEmbedding=\IdentityMap_\ExtensionTarget$.
\end{definition}

\paraid{p-cb58f0b73b6e}
Unless otherwise stated,
ANR means ANR for all metrizable spaces,
and the same convention applies to AE,
ANE,
and AR.
Spaces described by these unqualified terms are assumed to be metrizable.

\begin{theorem}[{\cite[Theorem~12.1]{Dugundji1958}}]\label{thm:retract-extensor}
\paraid{p-a393bc2f3507}
For a metrizable target in the category of metrizable spaces,
the AE and AR properties are equivalent,
and the ANE and ANR properties are equivalent.
\end{theorem}

\begin{definition}\label{def:local-contractibility}
\paraid{p-e8a842a80650}
A space
$\ExtensionTarget $
is \emph{locally contractible at}
$\BasePoint \in \ExtensionTarget $
if for every open neighborhood
$\LargerNeighborhood $
of
$\BasePoint $
there is an open neighborhood
$\SmallerNeighborhood \subset \LargerNeighborhood $
of
$\BasePoint $
such that the inclusion
$\SmallerNeighborhood \to \LargerNeighborhood $
is null-homotopic.
\end{definition}

\paraid{p-6fea115b8d45}
The homotopy may be chosen to end at the constant map with value
$\BasePoint $.
Indeed,
evaluating the homotopy at
$\BasePoint $
defines a path from
$\BasePoint $
to the constant value,
and its reversed path can be appended to the homotopy path of every point of
$\SmallerNeighborhood$.
For
$\GHSpace$,
this definition is equivalent to the following condition.
For every
$\MetricSpace{\BaseCarrier }{\MetricSymbol }\in\GHSpace$
and
$\ErrorTolerance >0$,
there exist
$0<\ContractionRadius <\ErrorTolerance $
and a continuous map
\[
 \ContractionHomotopy \colon \GHball(\MetricSpace{\BaseCarrier }{\MetricSymbol },\ContractionRadius )\times\UnitInterval
          \to \GHball(\MetricSpace{\BaseCarrier }{\MetricSymbol },\ErrorTolerance )
\]
such that
$\ContractionHomotopy (\MetricSpace{\ComparisonCarrier }{\ComparisonMetric },0)=\MetricSpace{\ComparisonCarrier }{\ComparisonMetric }$
and
$\ContractionHomotopy (\MetricSpace{\ComparisonCarrier }{\ComparisonMetric },1)=\MetricSpace{\BaseCarrier }{\MetricSymbol }$.


\clearpage
\part{The topology of Gromov--Hausdorff space I: Continuous invariant measures}\label{part:measures}
\begin{quote}
\small
\noindent\textbf{Abstract.}
\paraid{p-4e86714a292d}
We assign a full-support Borel probability measure to every nonempty compact metric space. The assignment is equivariant under isometries. When the spaces converge in Hausdorff distance in a common compact metric space, the selected measures converge weakly. We prove that the subspace of measured compact spaces whose probability measures are invariant and have full support is Polish for the Gromov--Hausdorff--Prokhorov distance and that its projection onto the Gromov--Hausdorff space is an open surjection. We obtain the required assignment by continuously selecting probability laws on the fibers and averaging with respect to these laws.

\par\smallskip
\noindent\textbf{Keywords.} Gromov--Hausdorff space,\newline Gromov--Hausdorff--Prokhorov distance, invariant probability measure, full support, continuous selection.

\par\smallskip
\noindent\textbf{2020 Mathematics Subject Classification.} Primary 54E35; Secondary 54C65, 60B05, 60B10.

\end{quote}

\section{Introduction to Part I}\label{sec:intro-measures}
\paraid{p-e72a35de4049}
Let
$\GHSpace$
denote the space of isometry classes of nonempty compact metric spaces
with the Gromov--Hausdorff distance
$\GHDistance$.
In this paper,
we construct a measure for each compact metric space.
The assignment is continuous in the Gromov--Hausdorff topology.
Specifically,
if the spaces are isometrically embedded in a common compact metric space
and converge in Hausdorff distance,
then the pushed-forward measures converge weakly.

\paraid{p-96038747b89a}
Larrieu
\cite[arXiv v4, Theorem~3]{Larrieu2014InvariantProbabilities}
describes a construction of Borel probability measures on nonempty compact
metric spaces that are invariant under isometries between open subsets.
The construction uses a fixed generalized limit.
The cited result does not assert full support or continuity as the underlying
compact metric space varies.
Leinster and Roff
\cite[Theorem~1.1 and Definition~9.2]{LeinsterRoff2021Entropy}
study probability measures that maximize their family of metric entropies.
They define a uniform measure by taking a weak limit of maximizing measures
as the metric is multiplied by a factor tending to infinity,
under the assumptions that the maximizing measure is unique for all sufficiently
large factors and that the limit exists.
The  problem here concerns the ordinary Gromov--Hausdorff topology
on the underlying compact spaces and weak convergence of measures in
common ambient spaces.

\paraid{p-6e73278858c1}
In our main result,
\Cref{thm:invariant-measures},
we assign to every nonempty compact metric space
$\MetricSpace{\BaseCarrier}{\MetricSymbol}$
a Borel probability measure
$\SelectedMeasure{\BaseCarrier}{\MetricSymbol}$
with full support,
so that
\[
 \supp(\SelectedMeasure{\BaseCarrier}{\MetricSymbol})=\BaseCarrier.
\]
For every isometry
$\IdentifyingIsometry\colon\MetricSpace{\BaseCarrier}{\MetricSymbol}
\to\MetricSpace{\ComparisonCarrier}{\ComparisonMetric}$,
we have
\[
 \IdentifyingIsometry\Pushforward\SelectedMeasure{\BaseCarrier}{\MetricSymbol}
 =\SelectedMeasure{\ComparisonCarrier}{\ComparisonMetric}.
\]
This means well-definedness.
Let compact subspaces
$\BaseCarrier_\SequenceIndex$
and
$\BaseCarrier$
of a compact metric space
$\MetricSpace{\AmbientCarrier}{\AmbientMetric}$
carry the restricted metrics
$\MetricSymbol_\SequenceIndex$
and
$\MetricSymbol$.
Then
\[
 \HausdorffDistance{\AmbientMetric}(\BaseCarrier_\SequenceIndex,\BaseCarrier)\to0
 \quad\Longrightarrow\quad
 \SelectedMeasure{\BaseCarrier_\SequenceIndex}{\MetricSymbol_\SequenceIndex}
 \to\SelectedMeasure{\BaseCarrier}{\MetricSymbol}
 \quad\text{weakly in }\ProbabilityMeasures(\AmbientCarrier).
\]
We regard the selected measures as probabilities on the common ambient space.
Since
$\supp(\SelectedMeasure{\BaseCarrier}{\MetricSymbol})=\BaseCarrier$,
for every
$1\leq\IntegrabilityExponent<\infty$
and every
$\HilbertFunction\in\ContinuousFunctions(\BaseCarrier)\setminus\{0\}$,
we have
\[
 \norm{\HilbertFunction}_{\LebesgueSpace^\IntegrabilityExponent
   (\BaseCarrier,\SelectedMeasure{\BaseCarrier}{\MetricSymbol})}^{\IntegrabilityExponent}
 =\int_\BaseCarrier|\HilbertFunction(\BasePoint)|^{\IntegrabilityExponent}
   \,\IntegrationDifferential\SelectedMeasure{\BaseCarrier}{\MetricSymbol}(\BasePoint)
 >0.
\]
Isometry compatibility and weak continuity permit the corresponding
integrals to be compared as the compact spaces vary.

\paraid{p-a9e34dfe60e6}
To obtain an invariant full-support probability,
for a fixed compact metric space,
we average a full-support probability over its compact isometry group.
Rouyer's results
\cite[arXiv v1, Theorems~2 and~4]{Rouyer2011}
imply that compact metric spaces with trivial isometry group are dense in
$\GHSpace$.
Thus the map assigning to each compact metric space its isometry group,
equipped with the uniform metric, is discontinuous with respect to the
Gromov--Hausdorff distance.
We therefore construct the invariant probabilities simultaneously and continuously.
We work in the Gromov--Hausdorff--Prokhorov space of compact spaces
with probability measures.
In this space,
we retain the whole underlying compact space even when
the measure has smaller support.
Let
$\InvariantMeasuredSpaces$
denote the subspace of invariant full-support measured spaces.
We prove that
$\InvariantMeasuredSpaces$
is Polish and that the forgetful map,
\[
 \MeasureProjection\colon\InvariantMeasuredSpaces\to\GHSpace,
 \qquad
 \MeasureProjection\MeasuredSpace{\BaseCarrier}{\MetricSymbol}{\FirstMeasure}
 =\MetricSpace{\BaseCarrier}{\MetricSymbol},
\]
is a continuous open surjection.
By the theorems of Banakh and Radul
\cite[Theorem~1.2, p.~20]{BanakhRadul1999}
and Valov (\Cref{thm:valov}),
there is a continuous map
\[
 \LawSection\colon\GHSpace\to\ProbabilityMeasures(\InvariantMeasuredSpaces),
\]
satisfying
\[
 \MeasureProjection\Pushforward\LawSection\MetricSpace{\BaseCarrier}{\MetricSymbol}
 =\DiracProbability_{\MetricSpace{\BaseCarrier}{\MetricSymbol}}.
\]
Fix a compact metric space
$\MetricSpace{\BaseCarrier}{\MetricSymbol}$.
Each point of the fiber
$\MeasureProjection^{-1}(\MetricSpace{\BaseCarrier}{\MetricSymbol})$
can be represented by a compact metric space isometric to
$\MetricSpace{\BaseCarrier}{\MetricSymbol}$,
together with an invariant full-support probability measure.
We choose an isometry from that space to
$\MetricSpace{\BaseCarrier}{\MetricSymbol}$
and push the measure forward to
$\BaseCarrier$.
The resulting probability does not depend on the chosen isometry,
since the measure is invariant under isometries.
It also does not depend on the representative,
since two representatives of the same class are related by a measure-preserving isometry.
We average these probabilities with respect to
$\LawSection\MetricSpace{\BaseCarrier}{\MetricSymbol}$,
using the Riesz--Markov--Kakutani theorem
(\Cref{thm:riesz-markov-kakutani}).
We prove continuity of the selected measures by comparing integrals in a
common compact ambient space.

\paraid{p-b5920f7184fb}
The law and its average have different domains.
For the fixed representative
$\MetricSpace{\BaseCarrier}{\MetricSymbol}$,
put
$\MeasureFiber{\BaseCarrier}
 =\MeasureProjection^{-1}(\MetricSpace{\BaseCarrier}{\MetricSymbol})$.
For every
$\FiberPoint\in\MeasureFiber{\BaseCarrier}$,
let
$\FiberMeasure{\BaseCarrier}{\FiberPoint}$
be the transported probability on
$\BaseCarrier$
described above.
For every continuous real function
$\TestFunction$
on
$\BaseCarrier$,
the selected measure satisfies
\begin{equation}\label{eq:part-i-average-interface}
 \int_\BaseCarrier\TestFunction(\BasePoint)
 \,\IntegrationDifferential\SelectedMeasure{\BaseCarrier}{\MetricSymbol}(\BasePoint)
 =\int_{\MeasureFiber{\BaseCarrier}}
 \left(\int_\BaseCarrier\TestFunction(\BasePoint)
 \,\IntegrationDifferential\FiberMeasure{\BaseCarrier}{\FiberPoint}(\BasePoint)\right)
 \,\IntegrationDifferential\LawSection\MetricSpace{\BaseCarrier}{\MetricSymbol}(\FiberPoint).
\end{equation}
The inner integral uses an invariant full-support probability on the
fixed carrier,
and the outer integral uses a law on measured isomorphism classes.
\Cref{lem:continuous-barycenter} shows that this averaging preserves
invariance and full support.

\paraid{p-3423d1d0cc93}
As an application,
we construct a topological embedding
$s$
of
$\GHSpace$
into the Gromov--Hausdorff--Prokhorov space by attaching the selected measure.
The forgetful map is a continuous left inverse of
$s$.
We prove that the image
$s(\GHSpace)$
is a retract of the Gromov--Hausdorff--Prokhorov space and is contained in
the subspace of invariant full-support probabilities
(Corollary~\ref{cor:ghp-retract}).

\paraid{p-7734020dafc8}
For metric measure spaces,
Kazukawa,
Nakajima,
and Shioya
\cite[arXiv v1, Theorems~1.4--1.6]{KazukawaNakajimaShioya2024}
prove contractibility and local path connectedness for both the box and
concentration topologies.
The section and retraction constructed here relate the ordinary
Gromov--Hausdorff space to the Gromov--Hausdorff--Prokhorov space.

\medskip
\noindent\textbf{Dependence on the other parts.}
\paraid{p-f205aa0a2b93}
We prove the selection theorem and the retraction using the common
preliminaries and the results recalled in Section~\ref{sec:prelim-measures}.
In Part~\ref{part:spectral},
we use \Cref{thm:invariant-measures}
to construct local models.
The later absolute retract argument combines these local models.
The selected measures enter that argument in the construction of the local models.

\medskip
\noindent\textbf{Organization.}
\paraid{p-0dd0534490f6}
Section~\ref{sec:prelim-measures} contains the preliminaries for this part.
In Subsection~\ref{subsec:measures-ghp},
we use the probability background from Subsection~\ref{subsec:preliminaries-probability}
to introduce the Gromov--Hausdorff--Prokhorov space and prove common ambient
realization results.
Subsection~\ref{subsec:measures-preliminaries-laws} recalls Haar probability
and Valov's theorem.
In Subsection~\ref{subsec:measures-open-projection},
we prove that the space of metric measured spaces with invariant full-support
measures is Polish
and that its projection onto the unmeasured Gromov--Hausdorff space is open.
In Subsection~\ref{subsec:measures-laws-and-averages},
we construct continuous probability laws on the fibers
and average them to prove \Cref{thm:invariant-measures}.
In Subsection~\ref{subsec:measures-retract},
we identify the selected measured spaces as a retract
of the Gromov--Hausdorff--Prokhorov space.

\medskip
\noindent\textbf{Conventions and notation.}
\paraid{p-33d1f4dd98d2}
All compact metric spaces are nonempty.
When a metric space represents a point of
$\GHSpace$,
we identify it with its isometry class.
An isometry is a surjective isometric embedding.
We denote by
$\ProbabilityMeasures(\DomainSpace)$
the set of Radon probability measures on a metrizable space
$\DomainSpace$,
equipped with the topology of weak convergence against bounded continuous real functions.
Subscripts in
$\SelectedMeasure{\BaseCarrier}{\MetricSymbol}$
refer to the specified metric space
$\MetricSpace{\BaseCarrier}{\MetricSymbol}$.
For isometrically embedded subsets of a common ambient space,
we use the same symbols for the subsets and their restricted metrics.
Sequences are indexed by
$\NonnegativeIntegers=\{0,1,2,\ldots\}$.


\section{Preliminaries}\label{sec:prelim-measures}
\paraid{p-6a3715dd0289}
We first introduce the Gromov--Hausdorff--Prokhorov space used in the
selection construction.
We then recall Haar probability and Valov's theorem on continuous laws.
The notation for probabilities and weak convergence is fixed in
Subsection~\ref{subsec:preliminaries-probability}.

\subsection{The Gromov--Hausdorff--Prokhorov space}\label{subsec:measures-ghp}
\paraid{p-ca05461b41c0}
Using the probability measures and weak topology recalled in
Subsection~\ref{subsec:preliminaries-probability},
we define a metric on the space of measured compact spaces.
At this stage a measure need not have full support.

\begin{definition}\label{def:measured-space}
\paraid{p-d2ece563f1da}
A \emph{measured compact metric space} is a triple
$\MeasuredSpace{\BaseCarrier}{\MetricSymbol}{\FirstMeasure}$
consisting of a nonempty compact metric space
$\MetricSpace{\BaseCarrier}{\MetricSymbol}$
and a probability
$\FirstMeasure\in\ProbabilityMeasures(\BaseCarrier)$.
For two such triples
$\MeasuredSpace{\BaseCarrier}{\MetricSymbol}{\FirstMeasure}$
and
$\MeasuredSpace{\ComparisonCarrier}{\ComparisonMetric}{\SecondMeasure}$,
an \emph{isomorphism} from
$\MeasuredSpace{\BaseCarrier }{\MetricSymbol }{\FirstMeasure }$
to
$\MeasuredSpace{\ComparisonCarrier }{\ComparisonMetric }{\SecondMeasure }$
is an isometry
$\IdentifyingIsometry \colon\MetricSpace{\BaseCarrier }{\MetricSymbol }\to\MetricSpace{\ComparisonCarrier }{\ComparisonMetric }$
such that
$\IdentifyingIsometry \Pushforward\FirstMeasure =\SecondMeasure $.
\end{definition}

\paraid{p-38ea40c9599f}
We denote by
$\MeasuredSpaces$
the set of isomorphism classes of measured compact metric spaces.

\begin{definition}\label{def:ghp-distance}
\paraid{p-d28a49b96a0a}
Let
$\MeasuredSpace{\BaseCarrier}{\MetricSymbol}{\FirstMeasure}$
and
$\MeasuredSpace{\ComparisonCarrier}{\ComparisonMetric}{\SecondMeasure}$
be representatives of two elements of
$\MeasuredSpaces$.
Let
$\MetricSpace{\AmbientCarrier}{\AmbientMetric}$
range over compact metric spaces and let
$\FirstEmbedding\colon\MetricSpace{\BaseCarrier}{\MetricSymbol}\to\MetricSpace{\AmbientCarrier}{\AmbientMetric}$
and
$\SecondEmbedding\colon\MetricSpace{\ComparisonCarrier}{\ComparisonMetric}\to\MetricSpace{\AmbientCarrier}{\AmbientMetric}$
range over isometric embeddings.
Let
$\mathcal S$
be the collection of quadruples
$(\AmbientCarrier,\AmbientMetric,\FirstEmbedding,\SecondEmbedding)$
consisting of a compact metric space and two isometric embeddings with these domains.
Taking the infimum over
$\mathcal S$,
we define the \emph{Gromov--Hausdorff--Prokhorov distance} by
\begin{equation}\label{eq:ghp-distance}
 \begin{split}
 &\GHPDistance(\MeasuredSpace{\BaseCarrier }{\MetricSymbol }{\FirstMeasure },\MeasuredSpace{\ComparisonCarrier }{\ComparisonMetric }{\SecondMeasure })\\
 &\qquad=\inf_{(\AmbientCarrier,\AmbientMetric,\FirstEmbedding,\SecondEmbedding)\in\mathcal S}
 \max\{\HausdorffDistance{\AmbientMetric }(\FirstEmbedding (\BaseCarrier ),\SecondEmbedding (\ComparisonCarrier )),
          \ProkhorovDistance{\AmbientMetric }(\FirstEmbedding \Pushforward\FirstMeasure ,\SecondEmbedding \Pushforward\SecondMeasure )\}.
 \end{split}
\end{equation}
We use the maximum convention for this unpointed distance.
These conventions agree with
\cite[arXiv v5, Example~2.1(iii), Definition~2.4 and equation~(8)]{Khezeli2023Framework},
with probability measures as the additional structure.
\end{definition}

\paraid{p-bc49bdb8e003}
For the measure selection,
we use the subspace of invariant probabilities with full support.

\begin{definition}\label{def:invariant-measured-space}
\paraid{p-ad3282ce0247}
Let
$\InvariantMeasuredSpaces\subset\MeasuredSpaces$
consist of the classes represented by triples
$\MeasuredSpace{\BaseCarrier}{\MetricSymbol}{\FirstMeasure}$
satisfying
\begin{equation}\label{eq:invariant-full-support-1}
 \supp(\FirstMeasure )=\BaseCarrier ,
\end{equation}
and for every
$\GroupElement\in\Isom\MetricSpace{\BaseCarrier}{\MetricSymbol}$,
\begin{equation}\label{eq:invariant-full-support-2}
 \GroupElement \Pushforward\FirstMeasure =\FirstMeasure.
\end{equation}
\end{definition}

\paraid{p-76c8d59831e8}
For full-support measures,
the isomorphism in \Cref{def:measured-space} agrees with the usual mm-isomorphism in
\cite[arXiv v1, Definition~2.1]{KazukawaNakajimaShioya2024}.
The cited definition compares the supports of the measures.
Our definition compares the entire compact spaces even when the measures
have smaller supports.
Thus
the Hausdorff distance term in
\eqref{eq:ghp-distance}
retains points outside the support.


\begin{theorem}[{\cite[arXiv v5, Theorems~2.6 and~2.12 and Example~2.1(iii)]{Khezeli2023Framework}}]\label{thm:ghp-polish}
\paraid{p-f754c0ec1769}
The metric
$\GHPDistance$
makes
$\MeasuredSpaces$
a complete separable metric space.
\end{theorem}

\paraid{p-b89da4d94310}
We next realize a convergent sequence in a common compact ambient space.

\begin{lemma}\label{lem:common-measured-embedding}
\paraid{p-3b87b5a0068d}
If
\[
 \GHPDistance(\MeasuredSpace{\BaseCarrier _\SequenceIndex }{\MetricSymbol _\SequenceIndex }{\FirstMeasure _\SequenceIndex },
                   \MeasuredSpace{\BaseCarrier }{\MetricSymbol }{\FirstMeasure })\longrightarrow0,
\]
then there exist a compact metric space
$\MetricSpace{\AmbientCarrier }{\AmbientMetric }$
and isometric embeddings
\[
 \FirstEmbedding _\SequenceIndex \colon\MetricSpace{\BaseCarrier _\SequenceIndex }{\MetricSymbol _\SequenceIndex }\to\MetricSpace{\AmbientCarrier }{\AmbientMetric },
\]
and
\[
 \FirstEmbedding \colon\MetricSpace{\BaseCarrier }{\MetricSymbol }\to\MetricSpace{\AmbientCarrier }{\AmbientMetric }
\]
such that
\begin{equation}\label{eq:common-measured-embedding-1}
 \HausdorffDistance{\AmbientMetric }(\FirstEmbedding _\SequenceIndex (\BaseCarrier _\SequenceIndex ),\FirstEmbedding (\BaseCarrier ))\longrightarrow0,
\end{equation}
and
\begin{equation}\label{eq:common-measured-embedding-2}
 (\FirstEmbedding _\SequenceIndex )\Pushforward\FirstMeasure _\SequenceIndex \longrightarrow\FirstEmbedding \Pushforward\FirstMeasure
 \quad\text{weakly in }\ProbabilityMeasures(\AmbientCarrier ).
\end{equation}
\end{lemma}

\begin{proof}
\paraid{p-a7da142773e7}
By \cite[arXiv v5, Lemma~2.5 and Example~2.1(iii)]{Khezeli2023Framework},
we obtain a common compact ambient space in which the carriers
converge in Hausdorff distance and the pushed-forward probabilities
converge in the L\'evy--Prokhorov distance.
On a compact metric space,
convergence in the L\'evy--Prokhorov distance is equivalent to weak
convergence
\cite[author's manuscript, Lemma~1.15]{ShioyaMetricMeasure}.
These are exactly the two conclusions in
\eqref{eq:common-measured-embedding-1} and \eqref{eq:common-measured-embedding-2}.
\end{proof}

\subsection{Haar probability and continuous laws}\label{subsec:measures-preliminaries-laws}
\paraid{p-c03d5071df48}
We use Haar probability to average over the isometry group.

\paraid{p-d2f74a62001a}
By the next lemma,
the full isometry group is compact,
so we can apply
the Haar probability theorem.

\begin{lemma}\label{lem:compact-isometry-group}
\paraid{p-bd8a5b8d04a8}
Let
$\MetricSpace{\BaseCarrier}{\MetricSymbol}$
be a nonempty compact metric space.
Then
$\Isom\MetricSpace{\BaseCarrier}{\MetricSymbol}$,
equipped with the uniform topology,
is a compact metrizable topological group.
\end{lemma}

\begin{proof}
\paraid{p-e174430bf1e9}
Let
$\IsometryGroup=\Isom\MetricSpace{\BaseCarrier}{\MetricSymbol}$.
For
$\Mapping,\SecondMapping\in\ContinuousFunctions(\BaseCarrier,\BaseCarrier)$,
define the uniform metric by
\[
 \UniformMetric(\Mapping,\SecondMapping)
 =\sup_{\BasePoint\in\BaseCarrier}
    \MetricSymbol(\Mapping(\BasePoint),\SecondMapping(\BasePoint)).
\]
Since every element of
$\IsometryGroup$
is an isometry and
$\BaseCarrier$
is compact,
the Arzel\`a--Ascoli theorem (\Cref{thm:arzela-ascoli}) implies that
$\IsometryGroup$
has compact closure in
$\ContinuousFunctions(\BaseCarrier,\BaseCarrier)$.

\paraid{p-6888b18e99f0}
A uniform limit
$\Mapping\in\ContinuousFunctions(\BaseCarrier,\BaseCarrier)$
of elements of
$\IsometryGroup$
preserves distances and is surjective by
\cite[Theorem~1.6.14]{BuragoBuragoIvanov2001}.
Thus
$\IsometryGroup$
is closed and compact.
Let
$\GroupElement_\SequenceIndex\to\GroupElement$
and
$\SecondGroupElement_\SequenceIndex\to\SecondGroupElement$
be two convergent sequences in
$\IsometryGroup$.
From the estimates
\[
 \UniformMetric(\GroupElement_\SequenceIndex\circ\SecondGroupElement_\SequenceIndex,
               \GroupElement\circ\SecondGroupElement)
 \leq\UniformMetric(\GroupElement_\SequenceIndex,\GroupElement)
       +\UniformMetric(\SecondGroupElement_\SequenceIndex,\SecondGroupElement),
\]
and
\[
 \UniformMetric(\GroupElement_\SequenceIndex^{-1},\GroupElement^{-1})
 =\UniformMetric(\GroupElement_\SequenceIndex,\GroupElement),
\]
it follows  that composition and inversion are continuous.
Hence
$\IsometryGroup$
is a compact metrizable topological group.
\end{proof}

\begin{theorem}[{Haar probability, \cite[Theorem~5.14 and Chapter~5, Section~3]{DiestelSpalsbury2014}}]\label{thm:haar}
\paraid{p-c37dcbf3e97d}
Every compact metrizable topological group
$\ActingGroup$
admits a unique left invariant Radon probability
$\HaarProbability\in\ProbabilityMeasures(\ActingGroup)$.
Thus
\[
 \HaarProbability(\ActingGroup)=1,
\]
and for every
$\GroupElement\in\ActingGroup$
and every
$\BorelSubset\in\BorelSets(\ActingGroup)$,
we have
\[
 \HaarProbability(\GroupElement\BorelSubset)=\HaarProbability(\BorelSubset).
\]
This probability is also right invariant and has full support.
\end{theorem}

\paraid{p-09eb6dcfaa8e}
We apply the following consequence of Valov's theorem
\cite[arXiv v2, Theorem~1.1]{Valov2009}.
The probabilities in this formulation are Radon probabilities.
The theorem imposes no compact-support restriction.
For Polish spaces, the same conclusion follows from the earlier theorem
of Banakh and Radul \cite[Theorem~1.2, p.~20]{BanakhRadul1999}.

\begin{theorem}[{Valov, \cite[arXiv v2, Theorem~1.1]{Valov2009}}]\label{thm:valov}
\paraid{p-ea9930a1ae3f}
Let
$\OpenSurjection \colon \LawTotalSpace \to \LawBaseSpace $
be an open continuous surjection between completely metrizable spaces.
Then there is a continuous map
$\LawSection\colon \LawBaseSpace \to\ProbabilityMeasures(\LawTotalSpace )$
such that
$\OpenSurjection \Pushforward\LawSection(\LawBasePoint )=\DiracProbability _\LawBasePoint $
for every
$\LawBasePoint \in \LawBaseSpace $.
\end{theorem}

\paraid{p-07d85ea78312}
Here
$\DiracProbability _\LawBasePoint $
is the Dirac probability at
$\LawBasePoint $.
Valov's theorem yields a continuous right inverse of the induced
map on probabilities.
Composing this right inverse with
$\LawBasePoint \mapsto\DiracProbability _\LawBasePoint $
proves this formulation.

\begin{remark}\label{rem:laws-averaging-operator}
\paraid{p-2db988c38721}
The law in \Cref{thm:valov} also defines a linear averaging operator.
Assume that
$\LawBaseSpace\ne\emptyset$,
and define the averaging operator
\[
 \AveragingOperator\colon
 \BoundedContinuousFunctions(\LawTotalSpace,\RealNumbers)
 \to\BoundedContinuousFunctions(\LawBaseSpace,\RealNumbers).
\]
For every
$\TestFunction\in\BoundedContinuousFunctions(\LawTotalSpace,\RealNumbers)$
and
$\LawBasePoint\in\LawBaseSpace$,
put
\[
 (\AveragingOperator\TestFunction)(\LawBasePoint)
 =\int_{\LawTotalSpace}\TestFunction(\FiberPoint)
   \,\IntegrationDifferential\LawSection(\LawBasePoint)(\FiberPoint).
\]
Since
$\LawSection$
is a weakly continuous family of probabilities,
$\AveragingOperator$
is a positive linear operator satisfying
\begin{equation}\label{eq:law-averaging-identities-1}
 \AveragingOperator1=1.
\end{equation}
For every
$\TestFunction\in\BoundedContinuousFunctions(\LawTotalSpace,\RealNumbers)$,
\begin{equation}\label{eq:law-averaging-identities-2}
 \|\AveragingOperator\TestFunction\|_\infty\leq\|\TestFunction\|_\infty.
\end{equation}
The fiber support condition implies that,
for every
$f\in\BoundedContinuousFunctions(\LawBaseSpace,\RealNumbers)$,
\begin{equation}\label{eq:law-averaging-identities-3}
 \AveragingOperator(f\circ\OpenSurjection)=f.
\end{equation}
Define
$\OpenSurjection^*\colon
 \BoundedContinuousFunctions(\LawBaseSpace,\RealNumbers)
 \to\BoundedContinuousFunctions(\LawTotalSpace,\RealNumbers)$
by
$\OpenSurjection^*f=f\circ\OpenSurjection$.
By \eqref{eq:law-averaging-identities-1}--\eqref{eq:law-averaging-identities-3},
$\OpenSurjection^*\AveragingOperator$
is a projection of norm one onto
$\OpenSurjection^*\BoundedContinuousFunctions(\LawBaseSpace,\RealNumbers)$.
\end{remark}


\section{Continuous invariant probability measures}\label{sec:measures}
\paraid{p-4a506ce908d0}
Its construction therefore requires a probability measure on that space.
We choose these measures so that isometries preserve them,
every nonempty open set has positive measure,
and convergence of the metric spaces implies weak convergence of the measures.
The last requirement concerns measures on different spaces.
We compare them after isometric embeddings into one compact metric space.

\paraid{p-bc9128d8b5de}
The construction has three steps.
We first prove that the invariant full-support subspace
$\InvariantMeasuredSpaces$
is Polish.
We then prove that forgetting the measure is an open map.
Finally,
we use Valov's theorem (\Cref{thm:valov}) to obtain a continuous probability law on each fiber
and average the measures represented by that law to construct the required measure.

\subsection{The open projection to unmeasured spaces}\label{subsec:measures-open-projection}
\paraid{p-fef9f75c2978}
We next restrict the measured space
$\MeasuredSpaces$
to invariant full-support measures
and prove that forgetting the measure is an open map.
We use uniform control of integrals to prove that the full-support condition
is a countable intersection of open conditions.
By taking limits of isometries,
we prove that the invariance condition is also
a countable intersection of open conditions and that the averaged measures
used to lift a convergent sequence of compact spaces converge weakly.
The required uniform estimates are in
\Cref{lem:uniform-weak-integrals,lem:parameter-integrals}.

\paraid{p-573911841fb0}
In the next proposition,
we extract a convergent subsequence of graphs of isometries
when the compact spaces vary.

\begin{proposition}\label{lem:isometry-limits}
\paraid{p-7c4e086d9e70}
Let
$\BaseCarrier _\SequenceIndex ,\BaseCarrier \subset \AmbientCarrier $
be nonempty compact subsets of a compact metric space
$\MetricSpace{\AmbientCarrier }{\AmbientMetric }$
with
$\HausdorffDistance{\AmbientMetric }(\BaseCarrier _\SequenceIndex ,\BaseCarrier )\to0$.
Equip each subset of
$\AmbientCarrier$
with the restricted ambient metric, and equip the product
$\AmbientCarrier\times\AmbientCarrier$
with the maximum product metric.
For any isometries
$\Isometry _\SequenceIndex \colon \BaseCarrier _\SequenceIndex \to \BaseCarrier _\SequenceIndex $,
a subsequence of their graphs converges in
$\CompactHyperspace(\AmbientCarrier \times \AmbientCarrier )$
to the graph of an isometry
$\Isometry \colon \BaseCarrier \to \BaseCarrier $.
Along that subsequence,
let
$\GraphProbability_\SequenceIndex\in\ProbabilityMeasures(\BaseCarrier_\SequenceIndex)$
and
$\GraphProbability\in\ProbabilityMeasures(\BaseCarrier)$
satisfy
$\GraphProbability_\SequenceIndex\to\GraphProbability$
weakly as probabilities on
$\AmbientCarrier$.
Then
\[
 (\Isometry_\SequenceIndex)\Pushforward\GraphProbability_\SequenceIndex
 \longrightarrow\Isometry\Pushforward\GraphProbability
 \quad\text{weakly in }\ProbabilityMeasures(\AmbientCarrier).
\]
\end{proposition}

\begin{proof}
\paraid{p-b6a94ad3a354}
We first identify the limit of the graphs and then the limit of the
probabilities supported on them.
For all
$\BasePoint_0,\ComparisonPoint_0,\BasePoint_1,\ComparisonPoint_1\in\AmbientCarrier$,
write
\[
 \AmbientMetric_\times((\BasePoint_0,\ComparisonPoint_0),(\BasePoint_1,\ComparisonPoint_1))
 =\max\{\AmbientMetric(\BasePoint_0,\BasePoint_1),
          \AmbientMetric(\ComparisonPoint_0,\ComparisonPoint_1)\}.
\]
For each
$\SequenceIndex$,
put
\[
 \Correspondence_\SequenceIndex
 =\{(\BasePoint,\Isometry_\SequenceIndex(\BasePoint))\mid
       \BasePoint\in\BaseCarrier_\SequenceIndex\}.
\]
By compactness of
$\CompactHyperspace(\AmbientCarrier\times\AmbientCarrier)$,
choose a subsequence whose graphs converge to a nonempty compact set
$\Correspondence\subset\AmbientCarrier\times\AmbientCarrier$.
For the remainder of this proof,
reindex the selected spaces,
isometries,
probabilities,
and graphs by
$\SequenceIndex\in\NonnegativeIntegers$,
retaining the notation
$\BaseCarrier_\SequenceIndex$,
$\Isometry_\SequenceIndex$,
$\GraphProbability_\SequenceIndex$,
and
$\Correspondence_\SequenceIndex$.
Then
\begin{equation}\label{eq:isometry-graph-convergence}
 \HausdorffDistance{\AmbientMetric_\times}
 (\Correspondence_\SequenceIndex,\Correspondence)\longrightarrow0.
\end{equation}

\textbf{Step 1.
The limiting graph.}
\paraid{p-7e16f5aad565}
Let
$\CoordinateProjection_1,\CoordinateProjection_2\colon
 \AmbientCarrier\times\AmbientCarrier\to\AmbientCarrier$
be the coordinate projections.
They are
$1$-Lipschitz
for
$\AmbientMetric_\times$,
and
$\CoordinateProjection_i(\Correspondence_\SequenceIndex)=\BaseCarrier_\SequenceIndex$
for
$i=1,2$,
because each
$\Isometry_\SequenceIndex$
is surjective.
For each
$i=1,2$,
we apply \Cref{lem:hausdorff-lipschitz-images} to
$\CoordinateProjection_i$
with
$L=1$.
Together with the triangle inequality,
this yields
\[
 \begin{aligned}
 \HausdorffDistance{\AmbientMetric}(\CoordinateProjection_i(\Correspondence),\BaseCarrier)
 &\leq\HausdorffDistance{\AmbientMetric}
       (\CoordinateProjection_i(\Correspondence),\CoordinateProjection_i(\Correspondence_\SequenceIndex))
       +\HausdorffDistance{\AmbientMetric}(\BaseCarrier_\SequenceIndex,\BaseCarrier)\\
 &\leq\HausdorffDistance{\AmbientMetric_\times}(\Correspondence,\Correspondence_\SequenceIndex)
       +\HausdorffDistance{\AmbientMetric}(\BaseCarrier_\SequenceIndex,\BaseCarrier)
       \longrightarrow0.
 \end{aligned}
\]
Thus
$\CoordinateProjection_1(\Correspondence)=\CoordinateProjection_2(\Correspondence)=\BaseCarrier$.
In particular,
$\Correspondence\subset\BaseCarrier\times\BaseCarrier$.

\paraid{p-49be4c2dec12}
Fix
$(\BasePoint_0,\ComparisonPoint_0),(\BasePoint_1,\ComparisonPoint_1)\in\Correspondence$.
For each
$i\in\{0,1\}$
and each
$\SequenceIndex\in\NonnegativeIntegers$,
choose a pair
$(\BasePoint_{i,\SequenceIndex},\ComparisonPoint_{i,\SequenceIndex})\in\Correspondence_\SequenceIndex$
whose distance from
$(\BasePoint_i,\ComparisonPoint_i)$
in the metric
$\AmbientMetric_\times$
is minimal.
Such a pair exists because
$\Correspondence_\SequenceIndex$
is nonempty and compact.
By \eqref{eq:isometry-graph-convergence},
\[
 \begin{aligned}
 \max\{\AmbientMetric(\BasePoint_{i,\SequenceIndex},\BasePoint_i),
        \AmbientMetric(\ComparisonPoint_{i,\SequenceIndex},\ComparisonPoint_i)\}
 &=\dist_{\AmbientMetric_\times}
   ((\BasePoint_i,\ComparisonPoint_i),\Correspondence_\SequenceIndex)\\
 &\leq\HausdorffDistance{\AmbientMetric_\times}
   (\Correspondence_\SequenceIndex,\Correspondence)\longrightarrow0.
 \end{aligned}
\]
Thus
$\BasePoint_{i,\SequenceIndex}\to\BasePoint_i$
and
$\ComparisonPoint_{i,\SequenceIndex}\to\ComparisonPoint_i$.
Since
$\ComparisonPoint_{i,\SequenceIndex}=\Isometry_\SequenceIndex(\BasePoint_{i,\SequenceIndex})$,
we obtain
\begin{equation}\label{eq:limit-graph-distance}
 \begin{aligned}
 \AmbientMetric(\BasePoint_0,\BasePoint_1)
 &=\lim_\SequenceIndex\AmbientMetric(\BasePoint_{0,\SequenceIndex},\BasePoint_{1,\SequenceIndex})\\
 &=\lim_\SequenceIndex\AmbientMetric(
      \Isometry_\SequenceIndex(\BasePoint_{0,\SequenceIndex}),
      \Isometry_\SequenceIndex(\BasePoint_{1,\SequenceIndex}))\\
 &=\lim_\SequenceIndex\AmbientMetric(\ComparisonPoint_{0,\SequenceIndex},\ComparisonPoint_{1,\SequenceIndex})
 =\AmbientMetric(\ComparisonPoint_0,\ComparisonPoint_1).
 \end{aligned}
\end{equation}
In particular,
if
 $(\BasePoint,\ComparisonPoint_0),(\BasePoint,\ComparisonPoint_1)\in\Correspondence$,
 then
 $\AmbientMetric(\ComparisonPoint_0,\ComparisonPoint_1)=0$.
 Hence
 $\ComparisonPoint_0=\ComparisonPoint_1$.
Since the first projection is surjective and each first coordinate has a unique
second coordinate,
we define the map
$\Isometry\colon\BaseCarrier\to\BaseCarrier$
by
\[
 \Correspondence=\{(\BasePoint,\Isometry(\BasePoint))\mid\BasePoint\in\BaseCarrier\}.
\]
Since the second projection is surjective,
we have
$\Isometry(\BaseCarrier)=\BaseCarrier$,
and for every
$\BasePoint_0,\BasePoint_1\in\BaseCarrier$,
\eqref{eq:limit-graph-distance} implies
\[
 \AmbientMetric(\Isometry(\BasePoint_0),\Isometry(\BasePoint_1))
 =\AmbientMetric(\BasePoint_0,\BasePoint_1).
\]
Hence
$\Isometry$
is an isometry.

\textbf{Step 2.
Probabilities on the graphs.}
\paraid{p-338e44ba371c}
For every
$\SequenceIndex\in\NonnegativeIntegers$,
define the map
\[
 \GraphMap_\SequenceIndex\colon\BaseCarrier_\SequenceIndex\to\AmbientCarrier\times\AmbientCarrier
\]
by
\[
 \GraphMap_\SequenceIndex(\BasePoint)=(\BasePoint,\Isometry_\SequenceIndex(\BasePoint)).
\]
We also define the map
\[
 \GraphMap\colon\BaseCarrier\to\AmbientCarrier\times\AmbientCarrier
\]
by
\[
 \GraphMap(\BasePoint)=(\BasePoint,\Isometry(\BasePoint)).
\]
Set
\[
 \GraphLaw_\SequenceIndex=(\GraphMap_\SequenceIndex)\Pushforward\GraphProbability_\SequenceIndex.
\]
By compactness of
$\ProbabilityMeasures(\AmbientCarrier\times\AmbientCarrier)$,
every subsequence of
$\GraphLaw_\SequenceIndex$
has a further weakly convergent subsequence.
Write the indices of any such subsequence as
$n_k$
and denote its limit by
$\GraphLaw\in\ProbabilityMeasures(\AmbientCarrier\times\AmbientCarrier)$.
Then
\[
 \GraphLaw_{n_k}\longrightarrow\GraphLaw
 \quad\text{weakly in }\ProbabilityMeasures(\AmbientCarrier\times\AmbientCarrier).
\]
Since
$\GraphLaw_{n_k}(\Correspondence_{n_k})=1$,
Equation \eqref{eq:isometry-graph-convergence},
\Cref{lem:limit-support},
and continuity of the first projection
$\CoordinateProjection_1$
imply
\begin{equation}\label{eq:limit-graph-law-support}
 \GraphLaw(\Correspondence)=1,
\end{equation}
and
\[
 (\CoordinateProjection_1)\Pushforward\GraphLaw
 =\lim_{k\to\infty}(\CoordinateProjection_1)\Pushforward\GraphLaw_{n_k}
 =\lim_{k\to\infty}\GraphProbability_{n_k}
 =\GraphProbability.
\]
On
$\Correspondence$
we have
\begin{equation}\label{eq:graph-projection-inverse}
 \GraphMap\circ\CoordinateProjection_1=\operatorname{id}_\Correspondence.
\end{equation}
Regard
$\GraphLaw$
as a probability on
$\Correspondence$.
By \eqref{eq:graph-projection-inverse},
we obtain
\[
 \GraphLaw
 =(\GraphMap\circ\CoordinateProjection_1)\Pushforward\GraphLaw
 =\GraphMap\Pushforward\bigl((\CoordinateProjection_1)\Pushforward\GraphLaw\bigr)
 =\GraphMap\Pushforward\GraphProbability.
\]
Every weakly convergent subsequence has this same limit.
Compactness of the probability space
$\ProbabilityMeasures(\AmbientCarrier\times\AmbientCarrier)$
consequently implies
\[
 \GraphLaw_\SequenceIndex\longrightarrow\GraphMap\Pushforward\GraphProbability
 \quad\text{weakly in }\ProbabilityMeasures(\AmbientCarrier\times\AmbientCarrier).
\]
Applying the continuous second projection
$\CoordinateProjection_2$
yields
\[
 (\Isometry_\SequenceIndex)\Pushforward\GraphProbability_\SequenceIndex
 =(\CoordinateProjection_2)\Pushforward\GraphLaw_\SequenceIndex
 \longrightarrow(\CoordinateProjection_2)\Pushforward(\GraphMap\Pushforward\GraphProbability)
 =\Isometry\Pushforward\GraphProbability.
\]
This is the required weak convergence on
$\AmbientCarrier$.
\end{proof}

\begin{proposition}\label{lem:full-support-gdelta}
\paraid{p-380dfd13f917}
The subset
\[
 \mathscr S_{\mathrm{full}}
 =\{\MeasuredSpace{\BaseCarrier}{\MetricSymbol}{\FirstMeasure}\in\MeasuredSpaces
       \mid\supp(\FirstMeasure)=\BaseCarrier\}
\]
is a
$\CountableIntersectionType$
subset of
$\MeasuredSpaces$.
\end{proposition}

\begin{proof}
\paraid{p-0aa6b81d8d5c}
For
$\MetricSpace{\BaseCarrier}{\MetricSymbol}\in\GHSpace$,
$\FirstMeasure\in\ProbabilityMeasures(\BaseCarrier)$,
and
$\Radius>0$,
define the ball mass function at each
$\BasePoint\in\BaseCarrier$
\[
 \BallMassFunction_\Radius\colon\BaseCarrier\to[0,1]
\]
by
\begin{equation}\label{eq:ball-mass-functions-1}
 \BallMassFunction_\Radius(\BasePoint)
 =\int_\BaseCarrier\BumpFunction{\BaseCarrier}{\MetricSymbol}{\BasePoint}{\Radius}
       \,\IntegrationDifferential\FirstMeasure,
\end{equation}
and define
\[
 \MinimumBallMass_\Radius\colon\MeasuredSpaces\to[0,1]
\]
by
\begin{equation}\label{eq:ball-mass-functions-2}
 \MinimumBallMass_\Radius\MeasuredSpace{\BaseCarrier}{\MetricSymbol}{\FirstMeasure}
 =\min_{\BasePoint\in\BaseCarrier}\BallMassFunction_\Radius(\BasePoint).
\end{equation}
By \eqref{eq:metric-bump-lipschitz},
the function
$\BallMassFunction_\Radius$
is
$1/\Radius$-Lipschitz,
so its minimum exists.
To confirm the continuity of
$\MinimumBallMass_\Radius$,
let
$\MeasuredSpace{\BaseCarrier_\SequenceIndex}{\MetricSymbol_\SequenceIndex}{\FirstMeasure_\SequenceIndex}$
converge to
$\MeasuredSpace{\BaseCarrier}{\MetricSymbol}{\FirstMeasure}$
in
$\MeasuredSpaces$.
Use \Cref{lem:common-measured-embedding} to realize the carriers in
$\MetricSpace{\AmbientCarrier}{\AmbientMetric}$
with Hausdorff convergence and
$\FirstMeasure_\SequenceIndex\to\FirstMeasure$
weakly on
$\AmbientCarrier$.
Define ambient extensions of the ball mass functions
by integrating the ambient bump functions against the measures supported on the respective carriers.
For every
$\IntegrationPoint\in\AmbientCarrier$,
define
\begin{equation}\label{eq:ambient-ball-mass-functions-1}
 \BallMassFunction_{\Radius,\SequenceIndex}(\IntegrationPoint)
 =\int_\AmbientCarrier
    \BumpFunction{\AmbientCarrier}{\AmbientMetric}{\IntegrationPoint}{\Radius}
    \,\IntegrationDifferential\FirstMeasure_\SequenceIndex,
\end{equation}
and
\begin{equation}\label{eq:ambient-ball-mass-functions-2}
 \tilde{\BallMassFunction}_\Radius(\IntegrationPoint)
 =\int_\AmbientCarrier
    \BumpFunction{\AmbientCarrier}{\AmbientMetric}{\IntegrationPoint}{\Radius}
    \,\IntegrationDifferential\FirstMeasure.
\end{equation}
We have
$\tilde{\BallMassFunction}_\Radius|_\BaseCarrier=\BallMassFunction_\Radius$,
and
$\BallMassFunction_{\Radius,\SequenceIndex}|_{\BaseCarrier_\SequenceIndex}$
is the ball mass function in \eqref{eq:ball-mass-functions-1} for
$\MeasuredSpace{\BaseCarrier_\SequenceIndex}{\MetricSymbol_\SequenceIndex}{\FirstMeasure_\SequenceIndex}$.
By \Cref{lem:parameter-integrals} and \eqref{eq:metric-bump-lipschitz},
\begin{equation}\label{eq:uniform-ball-mass-convergence}
 \norm{\BallMassFunction_{\Radius,\SequenceIndex}-\tilde{\BallMassFunction}_\Radius}_\infty
 \longrightarrow0.
\end{equation}
Using \eqref{eq:uniform-ball-mass-convergence} and the
$1/\Radius$-Lipschitz bound,
we obtain
\begin{equation}\label{eq:minimum-ball-mass-convergence}
 |\min_{\BasePoint\in\BaseCarrier_\SequenceIndex}\BallMassFunction_{\Radius,\SequenceIndex}(\BasePoint)
      -\min_{\BasePoint\in\BaseCarrier}\tilde{\BallMassFunction}_\Radius(\BasePoint)|
 \leq\norm{\BallMassFunction_{\Radius,\SequenceIndex}-\tilde{\BallMassFunction}_\Radius}_\infty
      +\frac{\HausdorffDistance{\AmbientMetric}(\BaseCarrier_\SequenceIndex,\BaseCarrier)}{\Radius}
 \longrightarrow0.
\end{equation}
Thus
$\MinimumBallMass _\Radius $
is continuous on
$\MeasuredSpaces$.

\paraid{p-871abef1344e}
If
$\supp(\FirstMeasure )=\BaseCarrier $,
then \eqref{eq:metric-bump-integral} implies that for every
$\BasePoint \in \BaseCarrier $,
\[
 \BallMassFunction _\Radius (\BasePoint )\geq\tfrac12\FirstMeasure (\MetricBall(\BasePoint ,\Radius /2;\MetricSymbol))>0.
\]
Continuity and compactness imply
$\MinimumBallMass _\Radius >0$.
Conversely,
if
$\BasePoint \notin\supp(\FirstMeasure )$,
choose $r>0$ such that
$\FirstMeasure(\MetricBall(\BasePoint,r;\MetricSymbol))=0$.
Choose $\CoordinateIndex\in\NonnegativeIntegers$ with $2^{-\CoordinateIndex}<r$.
Then $\FirstMeasure(\MetricBall(\BasePoint,2^{-\CoordinateIndex};\MetricSymbol))=0$.
Then
$\BallMassFunction _{2^{-\CoordinateIndex} }(\BasePoint )=0$.
We have proved that
$\supp(\FirstMeasure)=\BaseCarrier$
if and only if for every
$\CoordinateIndex\in\NonnegativeIntegers$
we have
\[
 \MinimumBallMass_{2^{-\CoordinateIndex}}
 \MeasuredSpace{\BaseCarrier}{\MetricSymbol}{\FirstMeasure}>0.
\]
For each
$\CoordinateIndex\in\NonnegativeIntegers$,
define the subset
\[
 \PositiveBallMassSet _\CoordinateIndex =\{\FiberPoint \in\MeasuredSpaces\mid \MinimumBallMass _{2^{-\CoordinateIndex} }(\FiberPoint )>0\}.
\]
Each
$\PositiveBallMassSet_\CoordinateIndex$
is open by continuity of
$\MinimumBallMass_{2^{-\CoordinateIndex}}$.

\paraid{p-39a712fdc965}
Consequently,
\[
 \mathscr S_{\mathrm{full}}
 =\bigcap_{\CoordinateIndex\in\NonnegativeIntegers}\PositiveBallMassSet_\CoordinateIndex.
\]
Thus
$\mathscr S_{\mathrm{full}}$
is a countable intersection of open subsets of
$\MeasuredSpaces$,
as required.
\end{proof}

\begin{proposition}\label{lem:invariant-measures-gdelta}
\paraid{p-fc58c5a6a993}
Let
$\mathscr S_{\mathrm{inv}}\subset\MeasuredSpaces$
consist of the classes
$\MeasuredSpace{\BaseCarrier}{\MetricSymbol}{\FirstMeasure}$
whose measures are invariant under every isometry of
$\MetricSpace{\BaseCarrier}{\MetricSymbol}$.
Then
$\mathscr S_{\mathrm{inv}}$
is a
$\CountableIntersectionType$
subset of
$\MeasuredSpaces$.
\end{proposition}

\begin{proof}
\paraid{p-9a02df523b4a}
Fix
$\CoordinateIndex\in\NonnegativeIntegers$.
Define
\[
 \begin{aligned}
 \NonInvariantSet _\CoordinateIndex &=\bigl\{\FiberPoint=\MeasuredSpace{\BaseCarrier }{\MetricSymbol }{\FirstMeasure }\in\MeasuredSpaces\bigm|
       \exists\GroupElement \in\Isom\MetricSpace{\BaseCarrier }{\MetricSymbol }\\
    &\qquad\qquad\ProkhorovDistance{\MetricSymbol }(\FirstMeasure ,\GroupElement \Pushforward\FirstMeasure )\geq2^{-\CoordinateIndex} \bigr\}.
 \end{aligned}
\]
We prove that
$\NonInvariantSet_\CoordinateIndex$
is closed.
Take a sequence
\[
 \FiberPoint_\SequenceIndex
 =\MeasuredSpace{\BaseCarrier_\SequenceIndex}{\MetricSymbol_\SequenceIndex}{\FirstMeasure_\SequenceIndex}
 \in\NonInvariantSet_\CoordinateIndex
\]
converging to
$\FiberPoint=\MeasuredSpace{\BaseCarrier}{\MetricSymbol}{\FirstMeasure}\in\MeasuredSpaces$.
For each
$\SequenceIndex$,
choose
$\GroupElement_\SequenceIndex\in\Isom\MetricSpace{\BaseCarrier_\SequenceIndex}{\MetricSymbol_\SequenceIndex}$
such that
\[
 \ProkhorovDistance{\MetricSymbol_\SequenceIndex}
 (\FirstMeasure_\SequenceIndex,
  (\GroupElement_\SequenceIndex)\Pushforward\FirstMeasure_\SequenceIndex)
 \geq2^{-\CoordinateIndex}.
\]
By \Cref{lem:common-measured-embedding},
there exist a compact metric space
$\MetricSpace{\AmbientCarrier}{\AmbientMetric}$
and isometric embeddings
$\FirstEmbedding_\SequenceIndex\colon\MetricSpace{\BaseCarrier_\SequenceIndex}{\MetricSymbol_\SequenceIndex}\to\MetricSpace{\AmbientCarrier}{\AmbientMetric}$
and
$\FirstEmbedding\colon\MetricSpace{\BaseCarrier}{\MetricSymbol}\to\MetricSpace{\AmbientCarrier}{\AmbientMetric}$
such that
\[
 \HausdorffDistance{\AmbientMetric}
 (\FirstEmbedding_\SequenceIndex(\BaseCarrier_\SequenceIndex),\FirstEmbedding(\BaseCarrier))\to0,
\]
and
\[
 (\FirstEmbedding_\SequenceIndex)\Pushforward\FirstMeasure_\SequenceIndex
 \to\FirstEmbedding\Pushforward\FirstMeasure
 \quad\text{weakly in }\ProbabilityMeasures(\AmbientCarrier).
\]
Identify the carriers with these images and the measures with these pushforwards.
For each
$\SequenceIndex$,
on the embedded carrier we use
$\GroupElement_\SequenceIndex$
to denote the conjugate
$\FirstEmbedding_\SequenceIndex\circ\GroupElement_\SequenceIndex\circ\FirstEmbedding_\SequenceIndex^{-1}$.
By \Cref{lem:isometry-limits},
we obtain a strictly increasing sequence of indices
$n_\ell$
and an isometry
$\GroupElement\in\Isom\MetricSpace{\BaseCarrier}{\MetricSymbol}$
for which
\[
 (\GroupElement_{n_\ell})\Pushforward\FirstMeasure_{n_\ell}
 \to\GroupElement\Pushforward\FirstMeasure
 \quad\text{weakly}.
\]
By \Cref{lem:prokhorov-isometric-embedding},
isometric embeddings preserve the L\'evy--Prokhorov distance.
Thus the triangle inequality on
$\ProbabilityMeasures(\AmbientCarrier)$
yields
\[
 \begin{aligned}
 &\left|
 \ProkhorovDistance{\MetricSymbol_{n_\ell}}
 (\FirstMeasure_{n_\ell},(\GroupElement_{n_\ell})\Pushforward\FirstMeasure_{n_\ell})
 -\ProkhorovDistance{\MetricSymbol}(\FirstMeasure,\GroupElement\Pushforward\FirstMeasure)
 \right|\\
 &\quad=\left|
 \ProkhorovDistance{\AmbientMetric}
 (\FirstMeasure_{n_\ell},(\GroupElement_{n_\ell})\Pushforward\FirstMeasure_{n_\ell})
 -\ProkhorovDistance{\AmbientMetric}(\FirstMeasure,\GroupElement\Pushforward\FirstMeasure)
 \right|\\
 &\quad\leq
 \ProkhorovDistance{\AmbientMetric}(\FirstMeasure_{n_\ell},\FirstMeasure)
 +\ProkhorovDistance{\AmbientMetric}
 ((\GroupElement_{n_\ell})\Pushforward\FirstMeasure_{n_\ell},\GroupElement\Pushforward\FirstMeasure)
 \longrightarrow0.
 \end{aligned}
\]
Both terms on the right tend to
$0$
because
$\FirstMeasure_{n_\ell}\to\FirstMeasure$
and
$(\GroupElement_{n_\ell})\Pushforward\FirstMeasure_{n_\ell}\to\GroupElement\Pushforward\FirstMeasure$
weakly in
$\ProbabilityMeasures(\AmbientCarrier)$,
and
$\ProkhorovDistance{\AmbientMetric}$
induces weak convergence.
Consequently,
\[
 \ProkhorovDistance{\MetricSymbol }(\FirstMeasure ,\GroupElement \Pushforward\FirstMeasure )
 =\lim_{\ell\to\infty} \ProkhorovDistance{\MetricSymbol _{n_\ell}}(\FirstMeasure _{n_\ell} ,(\GroupElement _{n_\ell} )\Pushforward\FirstMeasure _{n_\ell} )\geq2^{-\CoordinateIndex} .
\]
Thus
$\FiberPoint\in\NonInvariantSet_\CoordinateIndex$,
and hence
$\NonInvariantSet_\CoordinateIndex$
is closed.
Since the L\'evy--Prokhorov distance is a metric,
we see that
\[
 \mathscr S_{\mathrm{inv}}
 =\bigcap_{\CoordinateIndex\in\NonnegativeIntegers}
       (\MeasuredSpaces\setminus\NonInvariantSet_\CoordinateIndex).
\]
Each set in this intersection is open.
\end{proof}

\begin{theorem}\label{prop:invariant-measured-polish}
\paraid{p-a10d59349312}
The space
$\InvariantMeasuredSpaces$
is Polish.
\end{theorem}

\begin{proof}
\paraid{p-463534d4ce84}
By \Cref{lem:full-support-gdelta,lem:invariant-measures-gdelta},
\[
 \InvariantMeasuredSpaces
 =\mathscr S_{\mathrm{full}}\cap\mathscr S_{\mathrm{inv}}
\]
is a
$\CountableIntersectionType$
subset of
$\MeasuredSpaces$.
\Cref{thm:ghp-polish} states that
$\MeasuredSpaces$
is Polish,
so this subspace is Polish.
\end{proof}

\paraid{p-47cbaeca7f5a}
Define
\[
 \MeasureProjection\colon\InvariantMeasuredSpaces\to\GHSpace
\]
by
\[
 \MeasureProjection\MeasuredSpace{\BaseCarrier}{\MetricSymbol}{\FirstMeasure}
 =\MetricSpace{\BaseCarrier}{\MetricSymbol}.
\]
Namely, this map forgets the measure.
\begin{lemma}\label{lem:invariant-full-support}
\paraid{p-ab6dab25fbcd}
Every nonempty compact metric space
$\MetricSpace{\BaseCarrier}{\MetricSymbol}$
admits an invariant probability with full support.
\end{lemma}

\begin{proof}
\paraid{p-fee9db58d6fe}
Fix a dense sequence
$\{\BasePoint_\CoordinateIndex\}_{\CoordinateIndex\in\NonnegativeIntegers}$
in
$\BaseCarrier$
and put
\[
 \ApproximatingProbability
 =\sum_{\CoordinateIndex\in\NonnegativeIntegers}
   2^{-(\CoordinateIndex+1)}\DiracProbability_{\BasePoint_\CoordinateIndex}.
\]
Since
$\sum_{\CoordinateIndex\in\NonnegativeIntegers}2^{-(\CoordinateIndex+1)}=1$,
this is a probability on
$\BaseCarrier$.
By the density of
$\{\BasePoint_\CoordinateIndex\}_{\CoordinateIndex\in\NonnegativeIntegers}$,
we see that
$\supp(\ApproximatingProbability)=\BaseCarrier$.

\paraid{p-b8095b9a1aa1}
Put
$\IsometryGroup=\Isom\MetricSpace{\BaseCarrier}{\MetricSymbol}$.
By \Cref{lem:compact-isometry-group},
the group
$\IsometryGroup$
is compact and metrizable.
By the Haar theorem (\Cref{thm:haar}),
we obtain a normalized left invariant probability
$\HaarProbability _\IsometryGroup $
on
$\IsometryGroup $.
In the iterated integral below,
the outer variable
$\GroupElement\in\IsometryGroup$
is integrated against
$\HaarProbability_\IsometryGroup$,
and the inner variable
$\BasePoint\in\BaseCarrier$
is integrated against
$\ApproximatingProbability$.
Define the average
$\overline\ApproximatingProbability$
by requiring that for every
$\TestFunction\in\ContinuousFunctions(\BaseCarrier)$,
\[
 \int_\BaseCarrier \TestFunction \,\IntegrationDifferential \overline\ApproximatingProbability
 =\int_\IsometryGroup \int_\BaseCarrier \TestFunction (\GroupElement \BasePoint )\,\IntegrationDifferential \ApproximatingProbability (\BasePoint )\,\IntegrationDifferential \HaarProbability _\IsometryGroup (\GroupElement ).
\]
By the Riesz--Markov--Kakutani theorem (\Cref{thm:riesz-markov-kakutani}),
this positive normalized linear functional represents a unique probability on
$\BaseCarrier $.
Left invariance of
$\HaarProbability _\IsometryGroup $
implies
$\TranslatingIsometry \Pushforward\overline\ApproximatingProbability =\overline\ApproximatingProbability $
for every
$\TranslatingIsometry \in \IsometryGroup $.
For every nonempty open set
$\OpenSubset\subset\BaseCarrier$,
choose a nonzero continuous function
$\TestFunction\colon\BaseCarrier\to[0,1]$
that vanishes outside
$\OpenSubset$.
For each
$\GroupElement\in\IsometryGroup$,
the function
$\TestFunction\circ\GroupElement$
is nonzero and nonnegative,
so its integral against the full-support measure
$\ApproximatingProbability$
is positive.
The defining identity for the average therefore implies
\[
 \overline\ApproximatingProbability(\OpenSubset)
 \geq\int_\IsometryGroup\int_\BaseCarrier
       \TestFunction(\GroupElement\BasePoint)
       \,\IntegrationDifferential\ApproximatingProbability(\BasePoint)
       \,\IntegrationDifferential\HaarProbability_\IsometryGroup(\GroupElement)>0.
\]
Hence
$\MeasuredSpace{\BaseCarrier }{\MetricSymbol }{\overline\ApproximatingProbability }\in\InvariantMeasuredSpaces$,
as claimed.
\end{proof}

\begin{proposition}\label{lem:convergent-invariant-lifts}
\paraid{p-a0c55c41bf20}
Let
$\FiberPoint=\MeasuredSpace{\BaseCarrier}{\MetricSymbol}{\FirstMeasure}\in\InvariantMeasuredSpaces$
and assume that
$\MetricSpace{\BaseCarrier_\SequenceIndex}{\MetricSymbol_\SequenceIndex}
 \to\MetricSpace{\BaseCarrier}{\MetricSymbol}$
in
$\GHSpace$.
Then there are invariant probabilities
$\SecondMeasure_\SequenceIndex$
with full support on
$\BaseCarrier_\SequenceIndex$
such that
\[
 \MeasuredSpace{\BaseCarrier_\SequenceIndex}{\MetricSymbol_\SequenceIndex}{\SecondMeasure_\SequenceIndex}
 \longrightarrow\FiberPoint
 \quad\text{in }\InvariantMeasuredSpaces.
\]
\end{proposition}

\begin{proof}
\paraid{p-c601850135cd}
By \Cref{lem:common-gh-embedding},
choose isometric embeddings into a compact
$\MetricSpace{\AmbientCarrier }{\AmbientMetric }$
and identify the spaces with their images.
Let $\iota_\SequenceIndex\colon\BaseCarrier_\SequenceIndex\hookrightarrow\AmbientCarrier$
and $\iota\colon\BaseCarrier\hookrightarrow\AmbientCarrier$ be the inclusions.
By the Hausdorff continuity of the probability-measure functor
\cite[arXiv v5, Lemma~2.13 and Example~2.1(iii)]{Khezeli2023Framework},
there are probabilities
$\FirstMeasure_\SequenceIndex\in\ProbabilityMeasures(\BaseCarrier_\SequenceIndex)$
such that
$(\iota_\SequenceIndex)_*\FirstMeasure_\SequenceIndex\to\iota_*\FirstMeasure$
weakly on $\AmbientCarrier$.
Define
$\bar\mu_\SequenceIndex=(\iota_\SequenceIndex)_*\FirstMeasure_\SequenceIndex$
and $\bar\mu=\iota_*\FirstMeasure$ in $\ProbabilityMeasures(\AmbientCarrier)$.
Then
\[
 \ProkhorovDistance{\AmbientMetric}(\bar\mu_\SequenceIndex,\bar\mu)\longrightarrow0.
\]
Thus $\bar\mu_\SequenceIndex\to\bar\mu$ weakly in
$\ProbabilityMeasures(\AmbientCarrier)$.
Choose a probability
$\FullSupportProbability_\SequenceIndex$
with full support on each
$\BaseCarrier_\SequenceIndex$,
for example by assigning positive weights summing to one to a dense sequence.
For
$0<\MixingWeight_\SequenceIndex<1$
with
$\MixingWeight_\SequenceIndex\to0$,
put
\[
 \ApproximatingProbability_\SequenceIndex
 =(1-\MixingWeight_\SequenceIndex)\FirstMeasure_\SequenceIndex
   +\MixingWeight_\SequenceIndex\FullSupportProbability_\SequenceIndex.
\]
This probability has full support.
For every real continuous function
$\TestFunction$
on
$\AmbientCarrier$,
\[
 \left|\int_{\BaseCarrier_\SequenceIndex}(\TestFunction\circ\iota_\SequenceIndex)\,\IntegrationDifferential\ApproximatingProbability_\SequenceIndex
       -\int_{\BaseCarrier_\SequenceIndex}(\TestFunction\circ\iota_\SequenceIndex)\,\IntegrationDifferential\FirstMeasure_\SequenceIndex\right|
 \leq2\MixingWeight_\SequenceIndex\norm{\TestFunction}_\infty\longrightarrow0.
\]
Hence
$(\iota_\SequenceIndex)_*\ApproximatingProbability_\SequenceIndex\to\bar\mu$
weakly.

\paraid{p-0cb8c42117db}
Average
$\ApproximatingProbability _\SequenceIndex $
with respect to the normalized Haar probability on
$\IsometryGroup _\SequenceIndex =\Isom\MetricSpace{\BaseCarrier _\SequenceIndex }{\MetricSymbol _\SequenceIndex}$
and denote the resulting probability by
$\SecondMeasure _\SequenceIndex $.
The averaging construction in \Cref{lem:invariant-full-support}
shows that these probabilities are invariant and have full support.
For each
$\TestFunction \in\ContinuousFunctions(\AmbientCarrier )$,
we claim that
\begin{equation}\label{eq:uniform-orbit-measures}
 \sup_{\GroupElement \in \IsometryGroup _\SequenceIndex }
 \left|\int_{\BaseCarrier_\SequenceIndex}(\TestFunction\circ\iota_\SequenceIndex) \,\IntegrationDifferential (\GroupElement \Pushforward\ApproximatingProbability _\SequenceIndex )-\int_\BaseCarrier(\TestFunction\circ\iota) \,\IntegrationDifferential \FirstMeasure \right|
 \longrightarrow0.
\end{equation}
For the sake of contradiction,
suppose that \eqref{eq:uniform-orbit-measures} fails.
Choose
$\ErrorTolerance>0$
and a strictly increasing sequence of indices
$\{n_k\}_{k\in\NonnegativeIntegers}$
and for each
$k\in\NonnegativeIntegers$
choose
$\GroupElement_{n_k}\in\IsometryGroup_{n_k}$
satisfying
\begin{equation}\label{eq:orbit-integral-lower-bound}
 \left|\int_{\BaseCarrier_{n_k}}(\TestFunction\circ\iota_{n_k})\,
       \IntegrationDifferential((\GroupElement_{n_k})\Pushforward\ApproximatingProbability_{n_k})
       -\int_\BaseCarrier(\TestFunction\circ\iota)\,\IntegrationDifferential\FirstMeasure\right|
 \geq\ErrorTolerance.
\end{equation}
By \Cref{lem:isometry-limits},
there exist a strictly increasing sequence of indices
$\{k(j)\}_{j\in\NonnegativeIntegers}$
and an isometry
$\GroupElement\in\Isom\MetricSpace{\BaseCarrier}{\MetricSymbol}$
such that
\[
 (\iota_{n_{k(j)}}\circ\GroupElement_{n_{k(j)}})_*\ApproximatingProbability_{n_{k(j)}}
 \to(\iota\circ\GroupElement)_*\FirstMeasure=\bar\mu.
\]
Apply this weak convergence to the fixed function
$\TestFunction\in\ContinuousFunctions(\AmbientCarrier)$.
By the pushforward integral identity,
\[
 \left|\int_{\BaseCarrier_{n_{k(j)}}}(\TestFunction\circ\iota_{n_{k(j)}})
 \,\IntegrationDifferential((\GroupElement_{n_{k(j)}})_*\ApproximatingProbability_{n_{k(j)}})
 -\int_\BaseCarrier(\TestFunction\circ\iota)\,\IntegrationDifferential\FirstMeasure\right|
 \longrightarrow0.
\]
This contradicts the lower bound $\ErrorTolerance>0$ in
\eqref{eq:orbit-integral-lower-bound}.
Therefore,
averaging \eqref{eq:uniform-orbit-measures} over
$\IsometryGroup _\SequenceIndex $
yields
$(\iota_\SequenceIndex)_*\SecondMeasure_\SequenceIndex\to\bar\mu$
weakly.
Define
\[
 \FiberPoint_\SequenceIndex
 :=\MeasuredSpace{\BaseCarrier_\SequenceIndex}{\MetricSymbol_\SequenceIndex}{\SecondMeasure_\SequenceIndex}.
\]
Consequently,
\[
 \FiberPoint_\SequenceIndex\longrightarrow\FiberPoint,
\]
and
\[
 \MeasureProjection(\FiberPoint _\SequenceIndex )=\MetricSpace{\BaseCarrier _\SequenceIndex }{\MetricSymbol _\SequenceIndex}.
\]
Thus
$\{\FiberPoint_\SequenceIndex\}_{\SequenceIndex\in\NonnegativeIntegers}$
is the required sequence of invariant lifts.
This finishes the proof.
\end{proof}

\begin{theorem}\label{prop:open-invariant-projection}
\paraid{p-909c842b7494}
The map
\[
 \MeasureProjection\colon\InvariantMeasuredSpaces\to\GHSpace
\]
defined by
\[
 \MeasureProjection\MeasuredSpace{\BaseCarrier}{\MetricSymbol}{\FirstMeasure}
 =\MetricSpace{\BaseCarrier}{\MetricSymbol},
\]
is a continuous open surjection.
\end{theorem}

\begin{proof}
\paraid{p-34f1f298a7ea}
The inequality
\[
 \GHDistance(\MetricSpace{\BaseCarrier }{\MetricSymbol },\MetricSpace{\ComparisonCarrier }{\ComparisonMetric })
 \leq\GHPDistance(\MeasuredSpace{\BaseCarrier }{\MetricSymbol }{\FirstMeasure },\MeasuredSpace{\ComparisonCarrier }{\ComparisonMetric }{\SecondMeasure })
\]
shows that
$\MeasureProjection$
is continuous.
By \Cref{lem:invariant-full-support},
the map
$\MeasureProjection$
is surjective.
To prove openness,
let
$\OpenSubset \subset\InvariantMeasuredSpaces$
be open,
and let
$\FiberPoint \in \OpenSubset $.
For the sake of contradiction,
suppose that
$\MeasureProjection(\OpenSubset )$
contains no neighborhood of
$\MeasureProjection(\FiberPoint )$.
For each
$\SequenceIndex\in\NonnegativeIntegers$,
choose
\[
 \MetricSpace{\BaseCarrier _\SequenceIndex }{\MetricSymbol _\SequenceIndex}\in
 \GHball(\MeasureProjection(\FiberPoint ),2^{-\SequenceIndex} )\setminus\MeasureProjection(\OpenSubset ).
\]
By \Cref{lem:convergent-invariant-lifts},
there exist
$\FiberPoint_\SequenceIndex\in\InvariantMeasuredSpaces$
such that
\[
 \MeasureProjection(\FiberPoint_\SequenceIndex)
 =\MetricSpace{\BaseCarrier_\SequenceIndex}{\MetricSymbol_\SequenceIndex},
\]
and
\[
 \FiberPoint_\SequenceIndex\to\FiberPoint.
\]
Since
$\OpenSubset$
is an open neighborhood of
$\FiberPoint$,
we obtain
$\FiberPoint _\SequenceIndex \in \OpenSubset $
for all sufficiently large
$\SequenceIndex $.
This implies
$\MetricSpace{\BaseCarrier _\SequenceIndex }{\MetricSymbol _\SequenceIndex}\in\MeasureProjection(\OpenSubset )$,
a contradiction.
Thus
$\MeasureProjection(\OpenSubset )$
is open.
\end{proof}

\subsection{Continuous laws and their averages}\label{subsec:measures-laws-and-averages}
\paraid{p-e99b54a5841a}
We now select probability laws on the fibers of the forgetful map
and average the invariant measures with respect to these laws.

\paraid{p-f921de5bf828}
The averaging map in \Cref{lem:continuous-barycenter} is the standard
barycenter map on the space of probability measures.
Its construction and weak continuity appear in
\cite[proof of Theorem~8.10.5, p.~218]{Bogachev2007II}.
We include a proof using the Riesz--Markov--Kakutani theorem
(\Cref{thm:riesz-markov-kakutani}),
together with the full-support and invariance properties needed below.
We then prove continuity as the compact metric space varies by comparing
integrals in a common compact ambient space.

\begin{proposition}\label{lem:continuous-barycenter}
\paraid{p-ee76895903f0}
Let
$\MetricSpace{\BaseCarrier}{\MetricSymbol}$
be a nonempty compact metric space.
There exists a unique map
\[
 \operatorname{bar}_\BaseCarrier\colon
 \ProbabilityMeasures(\ProbabilityMeasures(\BaseCarrier))
 \to\ProbabilityMeasures(\BaseCarrier)
\]
such that for every
$\FirstProbabilityLaw\in\ProbabilityMeasures(\ProbabilityMeasures(\BaseCarrier))$
and every
$\TestFunction\in\ContinuousFunctions(\BaseCarrier)$,
\begin{equation}\label{eq:continuous-barycenter}
 \int_\BaseCarrier\TestFunction\,
       \IntegrationDifferential\operatorname{bar}_\BaseCarrier(\FirstProbabilityLaw)
 =\int_{\ProbabilityMeasures(\BaseCarrier)}
       \left(\int_\BaseCarrier\TestFunction\,\IntegrationDifferential\FirstMeasure\right)
       \,\IntegrationDifferential\FirstProbabilityLaw(\FirstMeasure).
\end{equation}
This map is continuous for the weak topologies.
If
$\FirstProbabilityLaw$
is concentrated on a Borel set of full-support measures,
then
$\operatorname{bar}_\BaseCarrier(\FirstProbabilityLaw)$
has full support.
For each
$\GroupElement\in\Isom\MetricSpace{\BaseCarrier}{\MetricSymbol}$,
assume that
\begin{equation}\label{eq:barycenter-invariant-law}
 \FirstProbabilityLaw\left(
 \left\{\FirstMeasure\in\ProbabilityMeasures(\BaseCarrier)
 \,\middle|\,
 \GroupElement\Pushforward\FirstMeasure=\FirstMeasure\right\}
 \right)=1.
\end{equation}
Then
\[
 \GroupElement\Pushforward\operatorname{bar}_\BaseCarrier(\FirstProbabilityLaw)
 =\operatorname{bar}_\BaseCarrier(\FirstProbabilityLaw).
\]
\end{proposition}

\begin{proof}
\paraid{p-11561faad095}
For
$\TestFunction\in\ContinuousFunctions(\BaseCarrier)$,
define
$\widehat\TestFunction\colon\ProbabilityMeasures(\BaseCarrier)\to\RealNumbers$
as follows.
For every
$\FirstMeasure\in\ProbabilityMeasures(\BaseCarrier)$,
set
\[
 \widehat\TestFunction(\FirstMeasure)
 =\int_\BaseCarrier\TestFunction\,\IntegrationDifferential\FirstMeasure.
\]
This function is continuous and bounded in absolute value by
$\norm{\TestFunction}_\infty$.
Fix
$\FirstProbabilityLaw\in\ProbabilityMeasures(\ProbabilityMeasures(\BaseCarrier))$.
Define the functional
$\RepresentingFunctional\colon\ContinuousFunctions(\BaseCarrier)\to\RealNumbers$
by integrating these functions,
\[
 \RepresentingFunctional(\TestFunction)
 =\int_{\ProbabilityMeasures(\BaseCarrier)}\widehat\TestFunction
       \,\IntegrationDifferential\FirstProbabilityLaw.
\]
This functional is positive and linear,
and
$\RepresentingFunctional(1)=1$.
\Cref{thm:riesz-markov-kakutani} therefore yields a unique representing
probability satisfying \eqref{eq:continuous-barycenter}.
If
$\FirstProbabilityLaw_\SequenceIndex\to\FirstProbabilityLaw$
weakly in
$\ProbabilityMeasures(\ProbabilityMeasures(\BaseCarrier))$,
then for every
$\TestFunction\in\ContinuousFunctions(\BaseCarrier)$,
\[
 \int_{\ProbabilityMeasures(\BaseCarrier)}\widehat\TestFunction
       \,\IntegrationDifferential\FirstProbabilityLaw_\SequenceIndex
 \longrightarrow
 \int_{\ProbabilityMeasures(\BaseCarrier)}\widehat\TestFunction
       \,\IntegrationDifferential\FirstProbabilityLaw.
\]
Thus the barycenters converge weakly.
Since both probability spaces are metrizable in their weak topologies,
this proves continuity.

\paraid{p-055551975ef9}
Assume that
$\FirstProbabilityLaw$
is concentrated on a Borel set of full-support measures.
For a nonempty open subset
$\OpenSubset\subset\BaseCarrier$,
choose a continuous function
$\TestFunction\colon\BaseCarrier\to[0,1]$
that is nonzero and vanishes outside
$\OpenSubset$.
For every
$\FirstMeasure\in\ProbabilityMeasures(\BaseCarrier)$
with
$\supp(\FirstMeasure)=\BaseCarrier$,
we have
$\widehat\TestFunction(\FirstMeasure)>0$.
By the hypothesis and \eqref{eq:continuous-barycenter},
\[
 \operatorname{bar}_\BaseCarrier(\FirstProbabilityLaw)(\OpenSubset)
 \geq\int_{\ProbabilityMeasures(\BaseCarrier)}
       \widehat\TestFunction(\FirstMeasure)
       \,\IntegrationDifferential\FirstProbabilityLaw(\FirstMeasure)>0.
\]
This proves full support.

\paraid{p-5cc5a2ea18dd}
Fix an isometry
$\GroupElement\in\Isom\MetricSpace{\BaseCarrier}{\MetricSymbol}$
for which \eqref{eq:barycenter-invariant-law} holds.
For every
$\TestFunction\in\ContinuousFunctions(\BaseCarrier)$
and every
$\FirstMeasure\in\ProbabilityMeasures(\BaseCarrier)$
with
$\GroupElement\Pushforward\FirstMeasure=\FirstMeasure$,
we have
\begin{equation}\label{eq:fiber-invariant-integrals}
 \int_\BaseCarrier\TestFunction\circ\GroupElement\,
       \IntegrationDifferential\FirstMeasure
 =\int_\BaseCarrier\TestFunction\,
       \IntegrationDifferential\FirstMeasure.
\end{equation}
By \eqref{eq:barycenter-invariant-law},
integrating \eqref{eq:fiber-invariant-integrals} against
$\FirstProbabilityLaw$
and using \eqref{eq:continuous-barycenter} proves invariance of
$\operatorname{bar}_\BaseCarrier(\FirstProbabilityLaw)$
under
$\GroupElement$.
\end{proof}

\begin{theorem}\label{thm:invariant-measures}
\paraid{p-b7a56605e33d}
For every nonempty compact metric space
$\MetricSpace{\BaseCarrier}{\MetricSymbol}$,
there exists a probability measure
$\SelectedMeasure{\BaseCarrier}{\MetricSymbol}\in\ProbabilityMeasures(\BaseCarrier)$.
These measures can be chosen simultaneously to satisfy the three properties below.
\begin{enumerate}[label=\textup{(M\arabic*)},ref=\textup{(M\arabic*)}]
\item
\paraid{p-09630bcb2594}
For every
$\MetricSpace{\BaseCarrier}{\MetricSymbol}$,
we have
$\supp(\SelectedMeasure{\BaseCarrier}{\MetricSymbol})=\BaseCarrier$.
\item
\paraid{p-1eec8fcfc366}
For every pair of nonempty compact metric spaces
$\MetricSpace{\BaseCarrier}{\MetricSymbol}$
and
$\MetricSpace{\ComparisonCarrier }{\ComparisonMetric }$
and every isometry
$\IdentifyingIsometry \colon\MetricSpace{\BaseCarrier }{\MetricSymbol }\to\MetricSpace{\ComparisonCarrier }{\ComparisonMetric }$,
we have
\[
 \IdentifyingIsometry \Pushforward\SelectedMeasure{\BaseCarrier}{\MetricSymbol}=\SelectedMeasure{\ComparisonCarrier}{\ComparisonMetric}.
\]
\item\label{item:invariant-measures-continuity}
\paraid{p-4f2c8f1817cc}
Let
$\MetricSpace{\BaseCarrier _\SequenceIndex }{\MetricSymbol _\SequenceIndex }$
and
$\MetricSpace{\BaseCarrier }{\MetricSymbol }$
be nonempty compact metric spaces.
Let
$\MetricSpace{\AmbientCarrier}{\AmbientMetric}$
be any compact metric space and let
$\FirstEmbedding _\SequenceIndex \colon\MetricSpace{\BaseCarrier _\SequenceIndex }{\MetricSymbol _\SequenceIndex }\to\MetricSpace{\AmbientCarrier }{\AmbientMetric }$
and
$\FirstEmbedding \colon\MetricSpace{\BaseCarrier }{\MetricSymbol }\to\MetricSpace{\AmbientCarrier }{\AmbientMetric }$
be isometric embeddings.
Then
\[
 \HausdorffDistance{\AmbientMetric }(\FirstEmbedding _\SequenceIndex (\BaseCarrier _\SequenceIndex ),\FirstEmbedding (\BaseCarrier ))\to0
 \quad\Longrightarrow\quad
 (\FirstEmbedding _\SequenceIndex )\Pushforward\SelectedMeasure{\BaseCarrier _\SequenceIndex}{\MetricSymbol _\SequenceIndex}\to\FirstEmbedding \Pushforward\SelectedMeasure{\BaseCarrier}{\MetricSymbol}
 \quad\text{weakly in }\ProbabilityMeasures(\AmbientCarrier ).
\]
\end{enumerate}
\end{theorem}

\begin{proof}
\paraid{p-941b13a492c1}
By \Cref{prop:invariant-measured-polish,prop:open-invariant-projection},
the map
$\MeasureProjection\colon\InvariantMeasuredSpaces\to\GHSpace$
is an open continuous surjection between Polish spaces.
By Valov's theorem (\Cref{thm:valov}),
we obtain a continuous map
$\LawSection\colon\GHSpace\to\ProbabilityMeasures(\InvariantMeasuredSpaces)$
such that for every
$\MetricSpace{\BaseCarrier}{\MetricSymbol}\in\GHSpace$
we have
$\MeasureProjection\Pushforward\LawSection\MetricSpace{\BaseCarrier }{\MetricSymbol }=\DiracProbability _{\MetricSpace{\BaseCarrier }{\MetricSymbol }}$.
Fix a representative
$\MetricSpace{\BaseCarrier }{\MetricSymbol }$
and set
\[
 \MeasureFiber{\BaseCarrier }=\MeasureProjection^{-1}(\MetricSpace{\BaseCarrier }{\MetricSymbol }).
\]
This is a closed Polish subspace of
$\InvariantMeasuredSpaces$,
and
$\LawSection\MetricSpace{\BaseCarrier }{\MetricSymbol }(\MeasureFiber{\BaseCarrier })=1$.
We construct a map
$\FiberMeasureMap{\BaseCarrier}\colon\MeasureFiber{\BaseCarrier}\to\ProbabilityMeasures(\BaseCarrier)$
on the fixed carrier as follows.
For each class
$\FiberPoint\in\MeasureFiber{\BaseCarrier}$,
choose a representative
$\MeasuredSpace{\ComparisonCarrier}{\ComparisonMetric}{\SecondMeasure}$
of $\FiberPoint$ and an isometry
$\IdentifyingIsometry \colon\MetricSpace{\ComparisonCarrier }{\ComparisonMetric }\to\MetricSpace{\BaseCarrier }{\MetricSymbol }$.
We denote the pushforward measure on this fixed representative by
\[
 \FiberMeasure{\BaseCarrier}{\FiberPoint}=\IdentifyingIsometry\Pushforward\SecondMeasure.
\]
For every
$\BorelSubset\in\BorelSets(\BaseCarrier)$,
this means that
\[
 \FiberMeasure{\BaseCarrier}{\FiberPoint}(\BorelSubset)
 =\SecondMeasure(\IdentifyingIsometry^{-1}(\BorelSubset)).
\]
Set
$\FiberMeasureMap{\BaseCarrier}(\FiberPoint)=\FiberMeasure{\BaseCarrier}{\FiberPoint}$.
If
$\SecondIdentifyingIsometry \colon\MetricSpace{\ComparisonCarrier }{\ComparisonMetric }\to\MetricSpace{\BaseCarrier }{\MetricSymbol }$
is another isometry,
then
$\SecondIdentifyingIsometry ^{-1}\IdentifyingIsometry \in\Isom\MetricSpace{\ComparisonCarrier }{\ComparisonMetric }$
and
\[
 \IdentifyingIsometry \Pushforward\SecondMeasure
 =\SecondIdentifyingIsometry \Pushforward\bigl((\SecondIdentifyingIsometry ^{-1}\IdentifyingIsometry )\Pushforward\SecondMeasure \bigr)
 =\SecondIdentifyingIsometry \Pushforward\SecondMeasure .
\]
Thus the measure is independent of the chosen isometry from this representative.
Now let $(Y_0,e_0,\nu_0)$ be another representative of $\FiberPoint$.
Choose a measure-preserving isometry
$h\colon\MetricSpace{\ComparisonCarrier}{\ComparisonMetric}\to(Y_0,e_0)$,
so that $h_*\SecondMeasure=\nu_0$,
and an isometry $k\colon(Y_0,e_0)\to\MetricSpace{\BaseCarrier}{\MetricSymbol}$.
Independence of the transport isometry on $\ComparisonCarrier$ implies
\[
 k_*\nu_0=(k\circ h)_*\SecondMeasure
 =\IdentifyingIsometry_*\SecondMeasure.
\]
Hence the measure is also independent of the representative of $\FiberPoint$.

\paraid{p-218d0d8ad567}
We first prove continuity of this assignment of measures under the hypotheses of \ref{item:invariant-measures-continuity}.
Fix a sequence
$\{\MetricSpace{\BaseCarrier_\SequenceIndex}{\MetricSymbol_\SequenceIndex}\}_{\SequenceIndex\in\NonnegativeIntegers}$,
a nonempty compact metric space
$\MetricSpace{\BaseCarrier}{\MetricSymbol}$,
a compact metric space
$\MetricSpace{\AmbientCarrier}{\AmbientMetric}$,
and prescribed isometric embeddings
$\FirstEmbedding_\SequenceIndex$
and
$\FirstEmbedding$
as in \ref{item:invariant-measures-continuity}, with
\[
 \HausdorffDistance{\AmbientMetric}(\FirstEmbedding_\SequenceIndex(\BaseCarrier_\SequenceIndex),\FirstEmbedding(\BaseCarrier))\to0.
\]
Identify
$\BaseCarrier_\SequenceIndex$
and
$\BaseCarrier$
with their images under the isometric embeddings into
$\AmbientCarrier$.
If
$\FiberPoint _\SequenceIndex \in\MeasureFiber{\BaseCarrier _\SequenceIndex }$
and
$\FiberPoint \in\MeasureFiber{\BaseCarrier }$
satisfy
$\FiberPoint _\SequenceIndex \to\FiberPoint $
in
$\InvariantMeasuredSpaces$,
then
\begin{equation}\label{eq:fiber-pullbacks}
 \FiberMeasure{\BaseCarrier _\SequenceIndex}{\FiberPoint _\SequenceIndex}\longrightarrow\FiberMeasure{\BaseCarrier}{\FiberPoint}
 \quad\text{weakly in }\ProbabilityMeasures(\AmbientCarrier ).
\end{equation}
Indeed,
let
$\SecondMeasure$
be the limit of any weakly convergent subsequence
\[
 \{\FiberMeasure{\BaseCarrier_{n_k}}{\FiberPoint_{n_k}}\}_{k\in\NonnegativeIntegers}.
\]
Such subsequences exist by compactness of
$\ProbabilityMeasures(\AmbientCarrier)$.
The limit is concentrated on
$\BaseCarrier $,
by \Cref{lem:limit-support}.
By the definition of
$\GHPDistance$
and weak convergence on
$\AmbientCarrier $,
we obtain
\[
 \FiberPoint _{n_k} \longrightarrow\MeasuredSpace{\BaseCarrier }{\MetricSymbol }{\SecondMeasure }.
\]
By uniqueness of limits,
$\MeasuredSpace{\BaseCarrier }{\MetricSymbol }{\SecondMeasure }=\FiberPoint $.
Thus there exists an isometry
$\GroupElement\colon\BaseCarrier\to\BaseCarrier$
such that
$\SecondMeasure=\GroupElement\Pushforward\FiberMeasure{\BaseCarrier}{\FiberPoint}$.
The measure
$\FiberMeasure{\BaseCarrier}{\FiberPoint}$
is invariant,
so
\begin{equation}\label{eq:fiber-measure-limit}
 \SecondMeasure=\FiberMeasure{\BaseCarrier}{\FiberPoint}.
\end{equation}
Compactness of
$\ProbabilityMeasures(\AmbientCarrier)$
and uniqueness of the subsequential limit in
\eqref{eq:fiber-measure-limit} imply \eqref{eq:fiber-pullbacks}.
With the space
$\BaseCarrier $
fixed,
this also proves continuity of
$\FiberMeasureMap{\BaseCarrier }$.

\paraid{p-39f42cc41be0}
Push the law on the fiber forward by
$\FiberMeasureMap{\BaseCarrier}$
and define
\[
 \FirstProbabilityLaw_\BaseCarrier
 =(\FiberMeasureMap{\BaseCarrier})\Pushforward
   \left(\LawSection\MetricSpace{\BaseCarrier}{\MetricSymbol}
                 |_{\MeasureFiber{\BaseCarrier}}\right).
\]
Using the barycenter in \Cref{lem:continuous-barycenter},
we define
\[
 \SelectedMeasure{\BaseCarrier}{\MetricSymbol}
 =\operatorname{bar}_\BaseCarrier(\FirstProbabilityLaw_\BaseCarrier).
\]
Define
$\RepresentingFunctional_\BaseCarrier\colon\ContinuousFunctions(\BaseCarrier)\to\RealNumbers$
as follows.
For every
$\TestFunction\in\ContinuousFunctions(\BaseCarrier)$,
set
\begin{equation}\label{eq:average-functional}
 \RepresentingFunctional_\BaseCarrier(\TestFunction)
 =\int_{\MeasureFiber{\BaseCarrier}}
       \left(\int_\BaseCarrier\TestFunction\,
       \IntegrationDifferential\FiberMeasure{\BaseCarrier}{\FiberPoint}\right)
       \IntegrationDifferential\LawSection
       \MetricSpace{\BaseCarrier}{\MetricSymbol}(\FiberPoint).
\end{equation}
Substituting
$\FirstProbabilityLaw_\BaseCarrier$
for
$\FirstProbabilityLaw$
in the identity \eqref{eq:continuous-barycenter} gives
\begin{equation}\label{eq:measure-barycenter}
 \int_\BaseCarrier\TestFunction\,
       \IntegrationDifferential\SelectedMeasure{\BaseCarrier}{\MetricSymbol}
 =\RepresentingFunctional_\BaseCarrier(\TestFunction).
\end{equation}
Figure~\ref{fig:measure-average} shows the two levels of probability measures.
\begin{figure}[H]
\centering
\begin{tikzpicture}[>=Stealth,node distance=9mm,font=\small]
\node (law) {$\LawSection\MetricSpace{\BaseCarrier}{\MetricSymbol}
 \in\ProbabilityMeasures(\MeasureFiber{\BaseCarrier})$};
\node[below=of law] (transport) {$
 (\FiberMeasureMap{\BaseCarrier})_*\LawSection\MetricSpace{\BaseCarrier}{\MetricSymbol}
 \in\ProbabilityMeasures(\ProbabilityMeasures(\BaseCarrier))$};
\node[below=of transport] (average) {$
 \SelectedMeasure{\BaseCarrier}{\MetricSymbol}\in\ProbabilityMeasures(\BaseCarrier)$};
\draw[->] (law) -- node[right]{transport by $\FiberMeasureMap{\BaseCarrier}$} (transport);
\draw[->] (transport) -- node[right]{average} (average);
\end{tikzpicture}
\caption{For a fixed representative
$\MetricSpace{\BaseCarrier}{\MetricSymbol}$,
the law
$\LawSection\MetricSpace{\BaseCarrier}{\MetricSymbol}$
on the fiber
$\MeasureFiber{\BaseCarrier}$
is transported to a law on
$\ProbabilityMeasures(\BaseCarrier)$.
The averaged measure
$\SelectedMeasure{\BaseCarrier}{\MetricSymbol}$
is defined by \eqref{eq:average-functional} and \eqref{eq:measure-barycenter}.}
\label{fig:measure-average}
\end{figure}


\paraid{p-4bf68c702f9c}
For a compact metric space
$\MetricSpace{\BaseCarrier}{\MetricSymbol}$,
we denote by
$\FullSupportProbabilities{\BaseCarrier}$
the set of full-support probabilities on
$\BaseCarrier$,
defined by
\[
 \FullSupportProbabilities{\BaseCarrier}
 =\{\FirstMeasure\in\ProbabilityMeasures(\BaseCarrier)
       \mid\supp\FirstMeasure=\BaseCarrier\}.
\]
For each
$\GroupElement\in\Isom\MetricSpace{\BaseCarrier}{\MetricSymbol}$,
we denote by
$\InvariantProbabilities{\BaseCarrier}{\GroupElement}$
the set of probabilities invariant under
$\GroupElement$,
defined by
\[
 \InvariantProbabilities{\BaseCarrier}{\GroupElement}
 =\{\FirstMeasure\in\ProbabilityMeasures(\BaseCarrier)
       \mid\GroupElement\Pushforward\FirstMeasure=\FirstMeasure\}.
\]
The law
$\FirstProbabilityLaw_\BaseCarrier$
is concentrated on
\[
 \FullSupportProbabilities{\BaseCarrier}
 \cap\bigcap_{\GroupElement\in\Isom\MetricSpace{\BaseCarrier}{\MetricSymbol}}
       \InvariantProbabilities{\BaseCarrier}{\GroupElement}.
\]
Choose a countable base $\CountableBasis$ for
$\BaseCarrier$.
A probability
$\FirstMeasure\in\ProbabilityMeasures(\BaseCarrier)$
has full support if and only if
$\FirstMeasure(\OpenSubset)>0$
for every nonempty
$\OpenSubset\in\CountableBasis$.
The weak topology on
$\ProbabilityMeasures(\BaseCarrier)$
is metrizable.
For each nonempty
$\OpenSubset\in\CountableBasis$
and each sequence
$\{\FirstMeasure_\SequenceIndex\}_{\SequenceIndex\in\NonnegativeIntegers}$
that converges weakly to
$\FirstMeasure\in\ProbabilityMeasures(\BaseCarrier)$,
\Cref{thm:portmanteau}
applied to the open set
$\OpenSubset$
gives
\[
 \FirstMeasure(\OpenSubset)
 \leq\liminf_{\SequenceIndex\to\infty}
          \FirstMeasure_\SequenceIndex(\OpenSubset).
\]
The sequential criterion for lower semicontinuity in a metrizable space now
shows that the map
\[
 \ProbabilityMeasures(\BaseCarrier)\to[0,1],
 \qquad\FirstMeasure\mapsto\FirstMeasure(\OpenSubset)
\]
is lower semicontinuous for the weak topology.
The strict superlevel set at $0$ is open, so
\[
 \{\FirstMeasure\in\ProbabilityMeasures(\BaseCarrier)
      \mid \FirstMeasure(\OpenSubset)>0\}
\]
is open in
$\ProbabilityMeasures(\BaseCarrier)$.
The basis criterion yields
\[
 \FullSupportProbabilities{\BaseCarrier}
 =\bigcap_{\substack{\OpenSubset\in\CountableBasis\\
                      \OpenSubset\ne\emptyset}}
       \{\FirstMeasure\in\ProbabilityMeasures(\BaseCarrier)
           \mid\FirstMeasure(\OpenSubset)>0\}.
\]
Since the basis is countable and each set in this intersection is open,
$\FullSupportProbabilities{\BaseCarrier}$
is a countable intersection of open sets, and in particular is Borel.
For each
$\GroupElement\in\Isom\MetricSpace{\BaseCarrier}{\MetricSymbol}$,
the set
$\InvariantProbabilities{\BaseCarrier}{\GroupElement}$
is closed in
$\ProbabilityMeasures(\BaseCarrier)$,
as we now verify.
For every
$\FirstMeasure\in\ProbabilityMeasures(\BaseCarrier)$
and every
$\TestFunction\in\ContinuousFunctions(\BaseCarrier)$,
the pushforward satisfies
\[
 \int_\BaseCarrier\TestFunction\,
       \IntegrationDifferential(\GroupElement\Pushforward\FirstMeasure)
 =\int_\BaseCarrier(\TestFunction\circ\GroupElement)\,
       \IntegrationDifferential\FirstMeasure.
\]
Since
$\GroupElement$
is continuous,
$\TestFunction\circ\GroupElement$
is continuous.
Let
$\{\FirstMeasure_\SequenceIndex\}_{\SequenceIndex\in\NonnegativeIntegers}$
be a sequence in
$\ProbabilityMeasures(\BaseCarrier)$
that converges weakly to
$\FirstMeasure$.
For every
$\TestFunction\in\ContinuousFunctions(\BaseCarrier)$,
the integral identity gives
\[
 \int_\BaseCarrier\TestFunction\,
       \IntegrationDifferential(\GroupElement\Pushforward\FirstMeasure_\SequenceIndex)
 \longrightarrow
 \int_\BaseCarrier\TestFunction\,
       \IntegrationDifferential(\GroupElement\Pushforward\FirstMeasure).
\]
Hence the pushforward map
\[
 \GroupElement\Pushforward\colon
 \ProbabilityMeasures(\BaseCarrier)\to\ProbabilityMeasures(\BaseCarrier)
\]
is continuous for the weak topology.
Now let
$\{\FirstMeasure_\SequenceIndex\}_{\SequenceIndex\in\NonnegativeIntegers}$
be a sequence in
$\InvariantProbabilities{\BaseCarrier}{\GroupElement}$
that converges weakly to
$\FirstMeasure\in\ProbabilityMeasures(\BaseCarrier)$.
Continuity and invariance give
\[
 \GroupElement\Pushforward\FirstMeasure_\SequenceIndex
 =\FirstMeasure_\SequenceIndex\longrightarrow\FirstMeasure,
 \qquad
 \GroupElement\Pushforward\FirstMeasure_\SequenceIndex
 \longrightarrow\GroupElement\Pushforward\FirstMeasure.
\]
Uniqueness of weak limits gives
$\GroupElement\Pushforward\FirstMeasure=\FirstMeasure$.
Thus
$\InvariantProbabilities{\BaseCarrier}{\GroupElement}$
is sequentially closed.
Since the weak topology is metrizable,
$\InvariantProbabilities{\BaseCarrier}{\GroupElement}$
is closed.
\Cref{lem:continuous-barycenter} therefore implies that
$\SelectedMeasure{\BaseCarrier}{\MetricSymbol}$
is invariant and has full support.
Finally,
for an isometry
$\IdentifyingIsometry \colon\MetricSpace{\BaseCarrier }{\MetricSymbol }\to\MetricSpace{\ComparisonCarrier }{\ComparisonMetric }$,
we have the equality of isometry classes
\[
 \MetricSpace{\BaseCarrier}{\MetricSymbol}
 =\MetricSpace{\ComparisonCarrier}{\ComparisonMetric}
 \quad\text{in }\GHSpace.
\]
Thus the fibers satisfy
\[
 \MeasureFiber{\BaseCarrier}=\MeasureFiber{\ComparisonCarrier}
 \quad\text{as subsets of }\InvariantMeasuredSpaces,
\]
and,
since the map
$\LawSection$
is defined on GH isometry classes,
\[
 \LawSection\MetricSpace{\BaseCarrier}{\MetricSymbol}
 =\LawSection\MetricSpace{\ComparisonCarrier}{\ComparisonMetric}.
\]
Fix the representatives
$\MetricSpace{\BaseCarrier}{\MetricSymbol}$
and
$\MetricSpace{\ComparisonCarrier}{\ComparisonMetric}$
and the isometry
$\IdentifyingIsometry$
between them.
For every
$\FiberPoint\in\MeasureFiber{\BaseCarrier}$,
the definition by transport implies
\[
 \FiberMeasure{\ComparisonCarrier}{\FiberPoint}
 =\IdentifyingIsometry\Pushforward\FiberMeasure{\BaseCarrier}{\FiberPoint}.
\]
For every
$\TestFunction\in\ContinuousFunctions(\ComparisonCarrier)$,
apply \eqref{eq:average-functional} and \eqref{eq:measure-barycenter} on
$\BaseCarrier$
to
$\TestFunction\circ\IdentifyingIsometry$.
The change-of-variables identity for each fiber measure then implies
\[
 \begin{aligned}
 \int_\ComparisonCarrier\TestFunction\,
   \IntegrationDifferential(\IdentifyingIsometry\Pushforward\SelectedMeasure{\BaseCarrier}{\MetricSymbol})
 &=\int_{\MeasureFiber{\BaseCarrier}}
     \left(\int_\BaseCarrier\TestFunction\circ\IdentifyingIsometry\,
       \IntegrationDifferential\FiberMeasure{\BaseCarrier}{\FiberPoint}\right)
       \IntegrationDifferential\LawSection\MetricSpace{\BaseCarrier}{\MetricSymbol}(\FiberPoint)\\
 &=\int_{\MeasureFiber{\ComparisonCarrier}}
     \left(\int_\ComparisonCarrier\TestFunction\,
       \IntegrationDifferential\FiberMeasure{\ComparisonCarrier}{\FiberPoint}\right)
       \IntegrationDifferential\LawSection\MetricSpace{\ComparisonCarrier}{\ComparisonMetric}(\FiberPoint)\\
 &=\int_\ComparisonCarrier\TestFunction\,
       \IntegrationDifferential\SelectedMeasure{\ComparisonCarrier}{\ComparisonMetric}.
 \end{aligned}
\]
Uniqueness of the representing measure proves the required compatibility
with isometries.

\paraid{p-7ae020536d8d}
It remains to prove continuity as the underlying spaces vary.
We use the prescribed embeddings into
$\AmbientCarrier$
and identify each carrier with its image.
Fix a real continuous function
$\TestFunction\in\ContinuousFunctions(\AmbientCarrier)$.
We express the integral of
$\TestFunction$
against each selected measure using one bounded continuous extension
on a fixed space.

\textbf{Step 1.
A closed domain for the fiber integrals.}
\paraid{p-72ddde2e12f8}
Put
\[
 \SequenceParameterSpace=\{0\}\cup\{2^{-\SequenceIndex}\mid\SequenceIndex\in\NonnegativeIntegers\}.
\]
Set
\[
 \ParameterCarrier_0=\BaseCarrier.
\]
For each
$\SequenceIndex\in\NonnegativeIntegers$,
set
\[
 \ParameterCarrier_{2^{-\SequenceIndex}}=\BaseCarrier_\SequenceIndex.
\]
All carriers have the restricted ambient metric.
The map
$\ParameterTime\mapsto[\ParameterCarrier_\ParameterTime]$
from
$\SequenceParameterSpace$
to
$\GHSpace$
is continuous.
Here the brackets denote the isometry class with the restricted metric.
Consequently,
\[
 \FiberIntegralDomain
 =\{(\ParameterTime,\FiberPoint)\in\SequenceParameterSpace\times\InvariantMeasuredSpaces
       \mid\MeasureProjection(\FiberPoint)=[\ParameterCarrier_\ParameterTime]\}
\]
is closed,
since it is the equality set of two continuous maps to the metric space
$\GHSpace$.
Define a function
$\FiberIntegral\colon\FiberIntegralDomain\to\RealNumbers$
by
\[
 \FiberIntegral(\ParameterTime,\FiberPoint)
 =\int_{\ParameterCarrier_\ParameterTime}\TestFunction\,
       \IntegrationDifferential\FiberMeasure{\ParameterCarrier_\ParameterTime}{\FiberPoint}.
\]
Continuity at parameter
$0$
follows from \eqref{eq:fiber-pullbacks}.
At an isolated parameter,
continuity follows from \eqref{eq:fiber-pullbacks}
with the carrier fixed.
Moreover,
$|\FiberIntegral|\leq\norm{\TestFunction}_\infty$.

\textbf{Step 2.
Extension and weak convergence of the laws.}
\paraid{p-b2e0811e32ee}
The product
$\SequenceParameterSpace\times\InvariantMeasuredSpaces$
is metrizable and hence normal.
By the Tietze--Urysohn theorem (\Cref{thm:tietze}),
$\FiberIntegral$
extends to a continuous function
$\ExtendedFiberIntegral\colon\SequenceParameterSpace\times\InvariantMeasuredSpaces\to\RealNumbers$
with
$|\ExtendedFiberIntegral|\leq\norm{\TestFunction}_\infty$.
Write
\[
 \ParameterLaw_\ParameterTime=\LawSection([\ParameterCarrier_\ParameterTime]).
\]
Continuity of
$\LawSection$
implies
$\ParameterLaw_{2^{-\SequenceIndex}}\to\ParameterLaw_0$
weakly.
\Cref{lem:dirac-product-convergence},
with
$\CompactParameterSpace=\SequenceParameterSpace$
and
$\ProbabilityCarrier=\InvariantMeasuredSpaces$,
implies
\begin{equation}\label{eq:parameter-law-products}
 \DiracProbability_{2^{-\SequenceIndex}}\otimes\ParameterLaw_{2^{-\SequenceIndex}}
 \longrightarrow\DiracProbability_0\otimes\ParameterLaw_0
 \quad\text{weakly on }\SequenceParameterSpace\times\InvariantMeasuredSpaces.
\end{equation}

\textbf{Step 3.
Integrals of the averaged measures.}
\paraid{p-b288c7e8154f}
Each product measure in \eqref{eq:parameter-law-products} is concentrated on
$\FiberIntegralDomain$.
By \eqref{eq:average-functional} and \eqref{eq:measure-barycenter},
\[
 \begin{aligned}
 \int_\AmbientCarrier\TestFunction\,\IntegrationDifferential\SelectedMeasure{\BaseCarrier_\SequenceIndex}{\MetricSymbol_\SequenceIndex}
 &=\int\ExtendedFiberIntegral\,
   \IntegrationDifferential(\DiracProbability_{2^{-\SequenceIndex}}\otimes\ParameterLaw_{2^{-\SequenceIndex}})\\
 &\longrightarrow\int\ExtendedFiberIntegral\,
   \IntegrationDifferential(\DiracProbability_0\otimes\ParameterLaw_0)
 =\int_\AmbientCarrier\TestFunction\,\IntegrationDifferential\SelectedMeasure{\BaseCarrier}{\MetricSymbol}.
 \end{aligned}
\]
Since this holds for every real continuous
$\TestFunction$
on the compact ambient space
$\AmbientCarrier$,
the selected measures
$\SelectedMeasure{\BaseCarrier_\SequenceIndex}{\MetricSymbol_\SequenceIndex}$
converge weakly to
$\SelectedMeasure{\BaseCarrier}{\MetricSymbol}$
as asserted.
\end{proof}

\subsection{A retract of the measured space}\label{subsec:measures-retract}
\paraid{p-e9cc85069378}
We now realize
$\GHSpace$
as a retract of the measured space.
By \Cref{thm:invariant-measures},
we obtain a continuous section of the forgetful map from
$\MeasuredSpaces$
to
$\GHSpace$.

\begin{corollary}\label{cor:ghp-retract}
\paraid{p-49d998de0004}
The map
\[
 s\colon\GHSpace\to\MeasuredSpaces
\]
defined by
\[
 s(\MetricSpace{\BaseCarrier}{\MetricSymbol})
 =\MeasuredSpace{\BaseCarrier}{\MetricSymbol}
   {\SelectedMeasure{\BaseCarrier}{\MetricSymbol}}
\]
is a topological embedding.
Its image is contained in
$\InvariantMeasuredSpaces$
and is a retract of
$\MeasuredSpaces$.
Consequently,
$\GHSpace$
is homeomorphic to a retract of the Gromov--Hausdorff--Prokhorov space
$\MeasuredSpaces$,
also when the latter is restricted to full-support probabilities or to
invariant full-support probabilities.
\end{corollary}

\begin{proof}
\paraid{p-170912b3d7e5}
Define the forgetful map on the whole measured space
\[
 p\colon\MeasuredSpaces\to\GHSpace
\]
by
\[
 p(\MeasuredSpace{\BaseCarrier}{\MetricSymbol}{\FirstMeasure})
 =\MetricSpace{\BaseCarrier}{\MetricSymbol}.
\]
Isomorphic measured spaces have isometric underlying spaces,
so this map is well defined.
By \eqref{eq:ghp-distance},
\[
 \GHDistance\bigl(p(\MeasuredSpace{\BaseCarrier}{\MetricSymbol}{\FirstMeasure}),
                   p(\MeasuredSpace{\ComparisonCarrier}{\ComparisonMetric}{\SecondMeasure})\bigr)
 \leq\GHPDistance(\MeasuredSpace{\BaseCarrier}{\MetricSymbol}{\FirstMeasure},
                   \MeasuredSpace{\ComparisonCarrier}{\ComparisonMetric}{\SecondMeasure}).
\]
Thus
$p$
is continuous.
By \Cref{thm:invariant-measures},
$s$
is well defined and takes its values in
$\InvariantMeasuredSpaces$.
Let
$\{\MetricSpace{\BaseCarrier_\SequenceIndex}{\MetricSymbol_\SequenceIndex}\}_{\SequenceIndex\in\NonnegativeIntegers}$
converge to
$\MetricSpace{\BaseCarrier}{\MetricSymbol}$
in
$\GHSpace$.
By \Cref{lem:common-gh-embedding},
choose compact representatives in a common compact metric space
$\MetricSpace{\AmbientCarrier}{\AmbientMetric}$
whose Hausdorff distances to
$\BaseCarrier$
tend to
$0$.
By \Cref{thm:invariant-measures},
\[
 \SelectedMeasure{\BaseCarrier_\SequenceIndex}{\MetricSymbol_\SequenceIndex}
 \to\SelectedMeasure{\BaseCarrier}{\MetricSymbol}
 \quad\text{weakly in }\ProbabilityMeasures(\AmbientCarrier).
\]
The L\'evy--Prokhorov distance metrizes weak convergence,
so \eqref{eq:ghp-distance} implies
\[
 \GHPDistance\bigl(
 s\MetricSpace{\BaseCarrier_\SequenceIndex}{\MetricSymbol_\SequenceIndex},
 s\MetricSpace{\BaseCarrier}{\MetricSymbol}\bigr)\to0.
\]
This proves continuity of
$s$.

\paraid{p-f1d4cfcbf0d7}
Since
$p\circ s=\IdentityMap_{\GHSpace}$,
$s$
is a topological embedding and
\[
 r=s\circ p\colon\MeasuredSpaces\to s(\GHSpace)
\]
is a retraction.
Its restrictions to the full-support subspace and to
$\InvariantMeasuredSpaces$
are retractions onto the same image.
\end{proof}

\paraid{p-fee6b65c4104}
The retract conclusion of \Cref{cor:ghp-retract} is a special case of
the absolute retract property of
$\GHSpace$,
which we prove in \Cref{thm:absolute-extensor}.


\clearpage
\part{The topology of Gromov--Hausdorff space II: Local coordinates of Gromov--Hausdorff space}\label{part:spectral}
\begin{quote}
\small
\noindent\textbf{Abstract.}
\paraid{p-3eaab1fe064b}
We construct local finite-dimensional approximations of compact metric spaces that vary continuously in the Gromov--Hausdorff topology. Near every nonempty compact metric space and for every positive error tolerance, we obtain a fixed finite dimension, a varying norm, and continuous point maps whose induced pseudometrics approximate the original distances uniformly. We choose the point maps to be equivariant under orthogonal representations of the isometry groups. The maps and norms are well defined up to orthogonal changes of coordinates. Using continuous invariant full-support probabilities, we approximate distance functions in finite spectral subspaces of the distance operator by unique best approximations in finite $L^p$ norms. We prove continuity of the local models using continuity of the spectral subspaces and the minimizers.

\par\smallskip
\noindent\textbf{Keywords.} Gromov--Hausdorff space, finite-dimensional approximation, distance operator, spectral subspace, best approximation.

\par\smallskip
\noindent\textbf{2020 Mathematics Subject Classification.} Primary 54E35; Secondary 47A55, 47A58.

\end{quote}

\section{Introduction to Part II}\label{sec:intro-spectral}
\paraid{p-c05f83eda3c9}
We represent a compact metric space by its distance functions in
an infinite-dimensional function space.
Using the continuous invariant probabilities of Part~\ref{part:measures},
we construct finite-dimensional approximations of this representation
that are equivariant under isometries and vary continuously with the compact space.
Let
$\GHSpace$
denote the space of isometry classes of nonempty compact metric spaces.
Our main result is \Cref{thm:local-models}.

\paraid{p-79876a468ed5}
The symmetry of the input space is part of the approximation problem.
Rouyer's results
\cite[arXiv v1, Theorems~2 and~4]{Rouyer2011}
imply that compact metric spaces with trivial isometry group are dense in
$\GHSpace$.
The approximation must also remain continuous at spaces with nontrivial
isometries.
We retain whole finite spectral subspaces and allow changes of their
orthonormal bases.
The resulting orthogonal actions express both the symmetry of the input
and the coordinate changes needed when the input varies.


\paraid{p-a0f40ab05c70}
For every nonempty compact metric space
$\MetricSpace{\BaseCarrier}{\MetricSymbol}$
and every
$\LocalTolerance>0$,
we construct a neighborhood
$\ModelNeighborhood\subset\GHSpace$
and a fixed integer
$\ModelDimension\geq1$.
For each compact metric space
$\MetricSpace{\ComparisonCarrier}{\ComparisonMetric}$
whose isometry class belongs to
$\ModelNeighborhood$,
there is a norm
$\NormSymbol_\ComparisonCarrier$
on
$\RealNumbers^\ModelDimension$
and a continuous map
\[
 \FeatureCoordinates_\ComparisonCarrier\colon\ComparisonCarrier
 \to\RealNumbers^\ModelDimension.
\]
For
$\ComparisonPoint,\IntegrationPoint\in\ComparisonCarrier$,
define the induced pseudometric by
\[
 \Pseudometric_\ComparisonCarrier(\ComparisonPoint,\IntegrationPoint)
 =\NormSymbol_\ComparisonCarrier
 (\FeatureCoordinates_\ComparisonCarrier(\ComparisonPoint)
 -\FeatureCoordinates_\ComparisonCarrier(\IntegrationPoint)).
\]
Its uniform error is
\[
 \LocalError\MetricSpace{\ComparisonCarrier}{\ComparisonMetric}
 =\sup_{\ComparisonPoint,\IntegrationPoint\in\ComparisonCarrier}
 |\Pseudometric_\ComparisonCarrier(\ComparisonPoint,\IntegrationPoint)
 -\ComparisonMetric(\ComparisonPoint,\IntegrationPoint)|.
\]
The function
$\LocalError$
is continuous on
$\ModelNeighborhood$
and satisfies
$\LocalError\MetricSpace{\BaseCarrier}{\MetricSymbol}<\LocalTolerance$.
We can
shrink
$\ModelNeighborhood$
so that
$\LocalError\MetricSpace{\ComparisonCarrier}{\ComparisonMetric}<\LocalTolerance$
for every
$\MetricSpace{\ComparisonCarrier}{\ComparisonMetric}\in\ModelNeighborhood$.
The dimension and the neighborhood may depend on the center
$\MetricSpace{\BaseCarrier}{\MetricSymbol}$
and the tolerance
$\LocalTolerance$.
Within that neighborhood,
the image
$\FeatureCoordinates_\ComparisonCarrier(\ComparisonCarrier)$
is a compact subset of one fixed vector space equipped with the varying norm
$\NormSymbol_\ComparisonCarrier$.
We also use
$\NormSymbol_\ComparisonCarrier$
to denote the metric
$(\FirstVector,\SecondVector)\mapsto\NormSymbol_\ComparisonCarrier(\FirstVector-\SecondVector)$
on
$\RealNumbers^\ModelDimension$.
The correspondence
$\{(\ComparisonPoint,\FeatureCoordinates_\ComparisonCarrier(\ComparisonPoint))
 \mid\ComparisonPoint\in\ComparisonCarrier\}$
has distortion at most
$\LocalError\MetricSpace{\ComparisonCarrier}{\ComparisonMetric}$.
Then \Cref{thm:gh-correspondence} implies
\[
 \GHDistance\!\left(
 \MetricSpace{\ComparisonCarrier}{\ComparisonMetric},
 \MetricSpace{\FeatureCoordinates_\ComparisonCarrier(\ComparisonCarrier)}{\NormSymbol_\ComparisonCarrier}
 \right)<\frac{\LocalTolerance}{2}.
\]
We summarize the continuity assertion of \Cref{thm:local-models}\ref{item:model-continuity} for the norms and point maps.
Let
$\{\MetricSpace{\ComparisonCarrier_\SequenceIndex}{\ComparisonMetric_\SequenceIndex}\}_{\SequenceIndex\in\NonnegativeIntegers}$
be a sequence in
$\ModelNeighborhood$,
and let
$\MetricSpace{\ComparisonCarrier}{\ComparisonMetric}\in\ModelNeighborhood$.
Assume that these spaces are embedded isometrically in a common compact ambient space
and that their carriers converge in Hausdorff distance to
$\ComparisonCarrier$.
After fixing these spaces and embeddings, choose a coordinate pair
$(\FeatureCoordinates_\ComparisonCarrier,\NormSymbol_\ComparisonCarrier)$
for the limit space and coordinate pairs
$\{(\FeatureCoordinates_{\ComparisonCarrier_\SequenceIndex},\NormSymbol_{\ComparisonCarrier_\SequenceIndex})\}_{\SequenceIndex\in\NonnegativeIntegers}$
for the spaces in the sequence so that
\[
 \sup_{\substack{\CoefficientVector\in\RealNumbers^\ModelDimension\\\norm{\CoefficientVector}_2=1}}
 \bigl|\NormSymbol_{\ComparisonCarrier_\SequenceIndex}(\CoefficientVector)
       -\NormSymbol_\ComparisonCarrier(\CoefficientVector)\bigr|\to0,
\]
and for every sequence
$\ComparisonPoint_\SequenceIndex\in\ComparisonCarrier_\SequenceIndex$
and every
$\ComparisonPoint\in\ComparisonCarrier$
with
$\ComparisonPoint_\SequenceIndex\to\ComparisonPoint$
in the common ambient space,
the coordinate values satisfy
\[
 \FeatureCoordinates_{\ComparisonCarrier_\SequenceIndex}(\ComparisonPoint_\SequenceIndex)
 \to\FeatureCoordinates_\ComparisonCarrier(\ComparisonPoint).
\]
For each
$\OrthogonalChange\in\OrthogonalGroup(\ModelDimension)$,
the corresponding change of orthonormal basis acts on the pair by
\[
 (\FeatureCoordinates_\ComparisonCarrier,\NormSymbol_\ComparisonCarrier)
 \longmapsto
 (\OrthogonalChange\circ\FeatureCoordinates_\ComparisonCarrier,
  \NormSymbol_\ComparisonCarrier\circ\OrthogonalChange^{-1}).
\]
This action preserves the induced pseudometric.

\paraid{p-210d146c099f}
The feature map is equivariant with respect to an orthogonal representation
of the isometry group of each input space.
In Remark~\ref{rem:local-isometry-representation},
we construct a continuous
homomorphism
\[
 \IsometryRepresentation{\ComparisonCarrier}{\ComparisonMetric}
 \colon\Isom\MetricSpace{\ComparisonCarrier}{\ComparisonMetric}
 \to\OrthogonalGroup(\ModelDimension)
\]
that satisfies equivariance and preserves the chosen norm.
For every
$\FirstIsometry\in\Isom\MetricSpace{\ComparisonCarrier}{\ComparisonMetric}$
and every
$\ComparisonPoint\in\ComparisonCarrier$,
the equivariance identity is
\[
 \FeatureCoordinates_\ComparisonCarrier(\FirstIsometry\ComparisonPoint)
 =\IsometryRepresentation{\ComparisonCarrier}{\ComparisonMetric}(\FirstIsometry)
   \FeatureCoordinates_\ComparisonCarrier(\ComparisonPoint).
\]
For every
$\FirstIsometry\in\Isom\MetricSpace{\ComparisonCarrier}{\ComparisonMetric}$
and every
$\CoefficientVector\in\RealNumbers^\ModelDimension$,
the norm-preservation identity is
\[
 \NormSymbol_\ComparisonCarrier
   \bigl(\IsometryRepresentation{\ComparisonCarrier}{\ComparisonMetric}(\FirstIsometry)
          \CoefficientVector\bigr)
 =\NormSymbol_\ComparisonCarrier(\CoefficientVector).
\]
Thus the varying isometry groups act through one fixed orthogonal group
on the model neighborhood.
Orthogonal changes of basis remain relevant even when the input space
has trivial isometry group.

\paraid{p-852d0b9cfe26}
We first recall two examples of spectral methods in geometry that provide context for our construction.
Bates
\cite[arXiv v1, Theorem~2]{Bates2014Eigenfunctions}
proves that closed connected Riemannian manifolds of fixed dimension at least two
and unit volume,
with a common lower bound on Ricci curvature and a common positive lower bound
on injectivity radius,
admit smooth embeddings by a uniformly bounded number of Laplacian eigenfunctions.
Kuwae and Shioya study spectral structures associated with nonnegative
self-adjoint operators on varying Hilbert spaces.
Under compact convergence of these structures,
they prove eventual equality of the dimensions of spectral subspaces
for bounded intervals whose endpoints lie outside the limiting spectrum
\cite[Theorem~2.6]{KuwaeShioya2003Spectral}.
In the present  paper,
 we use the distance kernel on arbitrary compact metric spaces
and approximate the distance uniformly.
We use the continuous invariant full-support assignment of probabilities
constructed in \Cref{thm:invariant-measures} of Part~\ref{part:measures}.
For
$\MetricSpace{\BaseCarrier}{\MetricSymbol}$,
the associated distance operator
$\DistanceOperator{\BaseCarrier}{\MetricSymbol}
 \colon\LebesgueSpace^2(\BaseCarrier,\SelectedMeasure{\BaseCarrier}{\MetricSymbol})
 \to\LebesgueSpace^2(\BaseCarrier,\SelectedMeasure{\BaseCarrier}{\MetricSymbol})$
is a self-adjoint Hilbert--Schmidt integral operator.
For every
$f\in\LebesgueSpace^2(\BaseCarrier,\SelectedMeasure{\BaseCarrier}{\MetricSymbol})$
and
$\BasePoint\in\BaseCarrier$,
its value is
\[
 (\DistanceOperator{\BaseCarrier}{\MetricSymbol}f)(\BasePoint)
 =\int_{\BaseCarrier}\MetricSymbol(\BasePoint,\IntegrationPoint)
 f(\IntegrationPoint)\,d\SelectedMeasure{\BaseCarrier}{\MetricSymbol}(\IntegrationPoint).
\]
Maria,
Oudot,
and Solomon study this operator and its spectral coordinates
\cite[Section~3]{MariaOudotSolomon2020}.
We approximate the distance functions by unique best approximations in
finite spectral subspaces,
using
$\LebesgueSpace^\IntegrabilityExponent$
norms.
Here
$2\leq\IntegrabilityExponent<\infty$,
and the norm is taken with respect to the selected probability measure.
We choose the exponent so that the distances between the approximants
approximate the original distances uniformly.
We retain whole eigenspaces and account for eigenvalue multiplicities
by orthogonal changes of basis.
Using continuity of the spectral subspaces and of the best approximations,
we prove the asserted continuity of the norms and point maps.

% Retired paragraph ID from the introduction to former Subsections 8.5 and 8.6.
% \paraid{p-fc6db22afc8b}

\paraid{p-cfb0b7538b10}
By the interpolation result in
\cite[Theorem~1.1]{Ishiki2026Interpolation},
we obtain a uniform error bound when extending prescribed metrics from a
discrete family of closed subsets of a fixed metrizable space.
Here we use \Cref{thm:local-models} to construct maps and norms on varying
compact spaces for the subsequent absolute retract argument.

\medskip
\noindent\textbf{Dependence on the other parts.}
\paraid{p-b4ac06e8df4d}
\Cref{thm:invariant-measures} is the input from Part~\ref{part:measures}.
We prove the spectral approximation and its applications using this
assignment and the preliminaries of Section~\ref{sec:prelim-spectral}.
Part~\ref{part:topology} uses \Cref{thm:local-models} to represent the varying
norms in fixed Banach spaces and combine the local point maps.
The orthogonal actions determine the compact group used in that construction.

\medskip
\noindent\textbf{Organization.}
\paraid{p-6ad9276d62d7}
In Section~\ref{sec:prelim-spectral},
we recall best approximation in
$\LebesgueSpace^\IntegrabilityExponent$,
spectral decomposition and perturbation,
and square roots of positive definite matrices.
In Subsection~\ref{subsec:spectral-distance-operator},
we define the distance operator and study its finite
spectral subspaces.
In Subsection~\ref{subsec:spectral-distance-approximation},
we prove approximation of distance functions using finite
spectral subspaces and unique best approximations in finite
$\LebesgueSpace^\IntegrabilityExponent$
norms.
In Subsection~\ref{subsec:spectral-subspace-continuity},
we prove continuity of the finite spectral subspaces
and construct orthonormal bases whose continuous extensions converge uniformly
on a common compact ambient space,
as asserted in \ref{item:spectral-basis-uniform-convergence}
of \Cref{lem:spectral-continuity}.
In Subsection~\ref{subsec:spectral-local-models},
we combine these results to construct the local maps and
varying norms in \Cref{thm:local-models}.

\medskip
\noindent\textbf{Conventions and notation.}
\paraid{p-21fe15dbd1cc}
All compact metric spaces are nonempty,
and finite-dimensional coordinate spaces are real.
Weak convergence of probabilities and the uniform integral estimates are
recalled in Subsection~\ref{subsec:preliminaries-probability}.
Function spaces are real unless complex scalars are explicitly indicated
for spectral perturbation.
An isometry is a surjective isometric embedding.
The pair
$\MetricSpace{\BaseCarrier}{\MetricSymbol}$
specifies both the underlying set and its metric.
In particular,
$\SelectedMeasure{\BaseCarrier}{\MetricSymbol}$
and
$\DistanceOperator{\BaseCarrier}{\MetricSymbol}$
depend on that pair.
When compact spaces are embedded isometrically into a common ambient
space,
we use the same symbols for the embedded subsets and the restricted
metrics.
The notation
$\IdentityMap_\BaseCarrier$
denotes the identity map on
$\BaseCarrier$.
We omit the subscript when its domain is specified.
Sequences are indexed by
$\NonnegativeIntegers=\{0,1,2,\ldots\}$.


\section{Preliminaries}\label{sec:prelim-spectral}
\paraid{p-920415618816}
We recall best approximation in
$\LebesgueSpace^\IntegrabilityExponent$,
the spectral theorem for compact self-adjoint operators,
and the perturbation results used to compare finite spectral subspaces.
The last subsection recalls a matrix theorem used to normalize bases of
finite spectral subspaces.

\subsection{Best approximation in Lebesgue spaces}\label{subsec:spectral-preliminaries-lp}
\paraid{p-57b625728f09}
Fix a nonempty compact metric space
$\MetricSpace{\BaseCarrier}{\MetricSymbol}$.
For a probability
$\FirstMeasure $
on
$\BaseCarrier $
and an exponent
$1\leq\IntegrabilityExponent<\infty$,
we identify measurable functions
$\HilbertFunction\colon\BaseCarrier\to\RealNumbers$
that are equal
$\FirstMeasure$-almost everywhere.
We denote by
$\LebesgueSpace^\IntegrabilityExponent (\BaseCarrier ,\FirstMeasure )$
the space of equivalence classes with finite norm
\[
 \norm{\HilbertFunction }_\IntegrabilityExponent =\left(\int_\BaseCarrier |\HilbertFunction |^\IntegrabilityExponent \,\IntegrationDifferential \FirstMeasure \right)^{1/\IntegrabilityExponent }.
\]

\begin{theorem}[{Best approximation, a consequence of \cite[Lemma~3.2.36(b), Theorem~5.6.1]{BreitGmeineder2026}}]\label{thm:best-approximation}
\paraid{p-b7f7a8e7fdb8}
Let
$1<\IntegrabilityExponent<\infty$,
let
$\FirstMeasure$
be a probability measure on
$\BaseCarrier$,
and let
$\ClosedDomain\subset\LebesgueSpace^{\IntegrabilityExponent}(\BaseCarrier,\FirstMeasure)$
be nonempty,
closed,
and convex.
Then for every
$\HilbertFunction\in\LebesgueSpace^{\IntegrabilityExponent}(\BaseCarrier,\FirstMeasure)$,
there is a unique
$\BestApproximation\in\ClosedDomain$
such that
\[
 \norm{\HilbertFunction-\BestApproximation}_{\IntegrabilityExponent}
 =\inf_{\Vector\in\ClosedDomain}\norm{\HilbertFunction-\Vector}_{\IntegrabilityExponent}.
\]
\end{theorem}

\begin{proof}
\paraid{p-6dd85192f3b4}
By the cited Theorem~5.6.1,
$\LebesgueSpace^{\IntegrabilityExponent}(\BaseCarrier,\FirstMeasure)$
is uniformly convex.
The translated set
$\{\Vector-\HilbertFunction\mid\Vector\in\ClosedDomain\}$
is nonempty,
closed,
and convex.
By Lemma~3.2.36(b) of the cited source,
we obtain a unique element of minimum norm in this set.
Translating that element by
$\HilbertFunction$
proves the assertion.
\end{proof}

\subsection{Compact operators and spectral decomposition}\label{subsec:spectral-preliminaries-operators}
\paraid{p-17e565d7cc93}
For a Hilbert space
$\HilbertDomain $,
we call a bounded linear operator
$\CompactOperator \colon \HilbertDomain \to \HilbertDomain $
\emph{compact} if the image of the closed unit ball of
$\HilbertDomain$
has compact closure.
We call
$\CompactOperator$
\emph{self-adjoint} if for every
$\HilbertFunction,\HilbertTestFunction\in\HilbertDomain$
we have
\[
 \langle \CompactOperator \HilbertFunction ,\HilbertTestFunction \rangle_\HilbertDomain =\langle \HilbertFunction ,\CompactOperator \HilbertTestFunction \rangle_\HilbertDomain.
\]
A nonzero vector
$\HilbertFunction $
is an \emph{eigenvector} for the eigenvalue
$\Eigenvalue $
if
$\CompactOperator \HilbertFunction =\Eigenvalue  \HilbertFunction $.
For an operator on a function space we also call it an eigenfunction.

\paraid{p-e197c5fac09a}
For a real Hilbert space
$\HilbertDomain$,
write
$\HilbertDomain_{\ComplexNumbers}=\HilbertDomain\oplus\ImaginaryUnit\HilbertDomain$
for its complexification.
We denote the complex-linear extension of
$\CompactOperator$
by the same symbol.
For
$\HilbertFunction,\HilbertTestFunction\in\HilbertDomain$,
set
\[
 \CompactOperator(\HilbertFunction+\ImaginaryUnit\HilbertTestFunction)
 =\CompactOperator\HilbertFunction+\ImaginaryUnit\CompactOperator\HilbertTestFunction.
\]
For a complex Banach space,
write
$\IdentityMap$
for its identity operator.
For a bounded operator
$\CompactOperator$
on that space,
its \emph{spectrum}
$\spec(\CompactOperator)$
is the set of complex numbers
$\Eigenvalue$
for which
$\Eigenvalue\IdentityMap-\CompactOperator$
has no bounded inverse.
For an operator on a real Hilbert space,
we use the spectrum of its
complex-linear extension to
$\HilbertDomain_{\ComplexNumbers}$.
For a self-adjoint operator this set is real.
In finite dimension it is exactly the set of eigenvalues.
In infinite dimension,
the number  $0$ can belong to the spectrum without being an
eigenvalue.

\begin{theorem}[{Riesz--Schauder and the self-adjoint spectral theorem, \cite[Theorems~14.19 and~15.22]{Muscat2024}}]\label{prop:spectral-basics}
\paraid{p-0edefea64bb3}
Let
$\CompactOperator $
be a compact self-adjoint operator on a separable real Hilbert space
$\HilbertDomain $.
Then every nonzero spectral value is a real eigenvalue with finite-dimensional
eigenspace.
For each
$\SpectralCutoff >0$,
there are only finitely many eigenvalues,
counted with multiplicity,
with absolute value greater than
$\SpectralCutoff $.
There is a finite or countable orthonormal family
$\{\Eigenfunction _\CoordinateIndex \}_{\CoordinateIndex \in \EigenIndexSet }$
indexed by a subset
$\EigenIndexSet\subset\NonnegativeIntegers$,
with nonzero eigenvalues
$\{\Eigenvalue _\CoordinateIndex \}_{\CoordinateIndex \in \EigenIndexSet }$
such that
\[
 \HilbertDomain =\ker \CompactOperator \ \mathbin{\oplus}\
       \overline{\LinearSpan }\{\Eigenfunction _\CoordinateIndex \mid \CoordinateIndex \in \EigenIndexSet \}.
\]
Every
$\HilbertFunction \in \HilbertDomain $
admits the decomposition
\[
 \HilbertFunction =\HilbertFunction _0+\sum_{\CoordinateIndex \in \EigenIndexSet }\langle \HilbertFunction ,\Eigenfunction _\CoordinateIndex \rangle \Eigenfunction _\CoordinateIndex .
\]
The operator acts on this decomposition by
\[
 \CompactOperator \HilbertFunction =\sum_{\CoordinateIndex \in \EigenIndexSet }\Eigenvalue _\CoordinateIndex \langle \HilbertFunction ,\Eigenfunction _\CoordinateIndex \rangle \Eigenfunction _\CoordinateIndex .
\]
The remaining component satisfies
\[
 \HilbertFunction _0\in\ker \CompactOperator .
\]
Both sums converge in
$\HilbertDomain $.
If
$\EigenIndexSet $
is infinite,
its eigenvalues tend to
$0$
in any enumeration.
\end{theorem}

\paraid{p-c81130e9d92f}
This is the self-adjoint case of the spectral theorem for compact normal
operators
\cite[Theorems~14.19 and~15.22]{Muscat2024}.
In particular,
a nonzero compact self-adjoint operator has a nonzero eigenvalue.

\subsection{Spectral perturbation and generalized eigenspaces}\label{subsec:spectral-preliminaries-perturbation}
\paraid{p-f87a26fde631}
Let
$\BanachAmbientSpace$
be a complex Banach space.
Write
\[
 \ClosedUnitBall=\{\FirstVector\in\BanachAmbientSpace\mid\norm{\FirstVector}\leq1\}
\]
for its closed unit ball.
A family of bounded linear operators
$\{\AmbientOperator_\SequenceIndex\colon\BanachAmbientSpace\to\BanachAmbientSpace\}_{\SequenceIndex\in\NonnegativeIntegers}$
is \emph{collectively compact} if
\[
 \overline{\bigcup_{\SequenceIndex\in\NonnegativeIntegers}
 \AmbientOperator_\SequenceIndex(\ClosedUnitBall)}
 \text{ is compact}.
\]
Strong convergence to a bounded linear operator
$\AmbientOperator\colon\BanachAmbientSpace\to\BanachAmbientSpace$
means that for every
$\FirstVector\in\BanachAmbientSpace$
we have
$\norm{\AmbientOperator_\SequenceIndex\FirstVector-\AmbientOperator\FirstVector}\to0$.

\begin{theorem}[{Riesz--Schauder, \cite[Theorem~14.19]{Muscat2024}}]\label{thm:compact-spectrum}
\paraid{p-1b1391ae5587}
Let
$\CompactOperator\colon\BanachAmbientSpace\to\BanachAmbientSpace$
be a compact linear operator on a complex Banach space.
Then every
$\Eigenvalue\in\spec(\CompactOperator)\setminus\{0\}$
is an eigenvalue,
and
\[
 0<\dim\ker(\CompactOperator-\Eigenvalue\IdentityMap)<\infty.
\]
The nonzero spectrum is at most countable and has no accumulation point
other than possibly zero.
\end{theorem}

\begin{theorem}[{Spectral inclusion, \cite[Theorem~5.3]{AnselonePalmer1968}}]\label{thm:spectral-inclusion}
\paraid{p-73a5ad7d4740}
Let
$\AmbientOperator_\SequenceIndex,\AmbientOperator\colon\BanachAmbientSpace\to\BanachAmbientSpace$
be bounded linear operators on a complex Banach space.
Assume that
$\AmbientOperator_\SequenceIndex\to\AmbientOperator$
strongly and that
$\{\AmbientOperator_\SequenceIndex-\AmbientOperator\}_{\SequenceIndex\in\NonnegativeIntegers}$
is collectively compact.
Then for every open set
$\SpectralNeighborhood\subset\ComplexNumbers$
containing
$\spec(\AmbientOperator)$,
there is
$\SequenceIndex_0\in\NonnegativeIntegers$
such that for every
$\SequenceIndex\geq\SequenceIndex_0$
we have
\[
 \spec(\AmbientOperator_\SequenceIndex)\subset\SpectralNeighborhood.
\]
\end{theorem}

\begin{definition}[{Riesz projection, \cite[Definition~14.23 and Proposition~14.28]{Muscat2024}}]\label{def:riesz-projection}
\paraid{p-52d4e9265055}
Let
$\CompactOperator\colon\BanachAmbientSpace\to\BanachAmbientSpace$
be a bounded linear operator on a complex Banach space.
Let
$\SpectralContour$
be a positively oriented finite union of rectifiable Jordan curves
avoiding
$\spec(\CompactOperator)$
and having winding number either
$0$
or
$1$
at every point of
$\spec(\CompactOperator)$.
The \emph{Riesz projection} of
$\CompactOperator$
is the bounded linear operator
$\RieszProjection{\CompactOperator}{\SpectralContour}\colon\BanachAmbientSpace\to\BanachAmbientSpace$
defined by
\[
 \RieszProjection{\CompactOperator}{\SpectralContour}
 =\frac{1}{2\CircleConstant\ImaginaryUnit}
 \int_{\SpectralContour}(\ResolventParameter\IdentityMap-\CompactOperator)^{-1}
 \,\IntegrationDifferential\ResolventParameter,
\]
where the line integral is taken in the operator norm.
If
$\Eigenvalue\in\spec(\CompactOperator)$
and
$\SpectralContour$
has winding number
$1$
at
$\Eigenvalue$
and
$0$
at every point of
$\spec(\CompactOperator)\setminus\{\Eigenvalue\}$,
then we call
$\RieszProjection{\CompactOperator}{\SpectralContour}$
the \emph{Riesz projection associated with
$\Eigenvalue$}.
\end{definition}

\begin{lemma}[{\cite[Examples~14.24(1) and Proposition~14.28]{Muscat2024}}]\label{lem:riesz-projection-properties}
\paraid{p-b805e4feeaaf}
Let
$\CompactOperator$
and
$\SpectralContour$
be as in Definition~\ref{def:riesz-projection}.
Then we have
\[
 \RieszProjection{\CompactOperator}{\SpectralContour}^2
 =\RieszProjection{\CompactOperator}{\SpectralContour},
\]
and
\[
 \CompactOperator\RieszProjection{\CompactOperator}{\SpectralContour}
 =\RieszProjection{\CompactOperator}{\SpectralContour}\CompactOperator.
\]
Moreover,
$\RieszProjection{\CompactOperator}{\SpectralContour}$
depends only on the subset of
$\spec(\CompactOperator)$
at which
$\SpectralContour$
has winding number
$1$.
\end{lemma}

\begin{theorem}[{Convergence of spectral projections, \cite[Proposition~6.3]{AnselonePalmer1968}}]\label{thm:riesz-convergence}
\paraid{p-ab1c2cec6e08}
Let
$\AmbientOperator_\SequenceIndex,\AmbientOperator\colon\BanachAmbientSpace\to\BanachAmbientSpace$
satisfy the operator hypotheses of
\Cref{thm:spectral-inclusion}.
Let
$\SpectralContour$
be as in Definition~\ref{def:riesz-projection} for
$\AmbientOperator$,
and put
\[
 \AmbientProjection=\RieszProjection{\AmbientOperator}{\SpectralContour}.
\]
Then for all sufficiently large
$\SequenceIndex$,
$\SpectralContour$
satisfies the conditions of Definition~\ref{def:riesz-projection} for
$\AmbientOperator_\SequenceIndex$,
and the projections
$\AmbientProjection_\SequenceIndex\colon\BanachAmbientSpace\to\BanachAmbientSpace$
defined by
\[
 \AmbientProjection_\SequenceIndex
 =\RieszProjection{\AmbientOperator_\SequenceIndex}{\SpectralContour}
\]
satisfy
\[
 \dim\range\AmbientProjection_\SequenceIndex=\dim\range\AmbientProjection.
\]
For every
$\FirstVector\in\BanachAmbientSpace$,
we also have
\[
 \norm{\AmbientProjection_\SequenceIndex\FirstVector-\AmbientProjection\FirstVector}\to0.
\]
If
$\AmbientOperator_\SequenceIndex$
and
$\AmbientOperator$
are compact and
$\SpectralContour$
has winding number
$0$
at
$0$,
then
$\dim\range\AmbientProjection_\SequenceIndex$
and
$\dim\range\AmbientProjection$
are finite.
\end{theorem}

\paraid{p-90dc1dadcd0a}
For a bounded linear operator
$\CompactOperator\colon\BanachAmbientSpace\to\BanachAmbientSpace$
on a complex Banach space and
$\Eigenvalue\in\ComplexNumbers$,
we write
\[
 \GeneralizedEigenspace{\CompactOperator}{\Eigenvalue}
 =\bigcup_{\RootOrder\geq1}
  \ker\bigl((\CompactOperator-\Eigenvalue\IdentityMap)^\RootOrder\bigr)
\]
for the generalized eigenspace at
$\Eigenvalue$.

\begin{theorem}[{\cite[Section~7, p.~429]{AnselonePalmer1968}, \cite[Proposition~14.28 and Examples~14.29]{Muscat2024}}]\label{lem:riesz-range-generalized-eigenspace}
\paraid{p-99a6a24b0cfc}
Let
$\CompactOperator\colon\BanachAmbientSpace\to\BanachAmbientSpace$
be a compact linear operator on a complex Banach space.
For every
$\Eigenvalue\in\spec(\CompactOperator)\setminus\{0\}$,
write
$\EigenvalueProjection{\CompactOperator}{\Eigenvalue}$
for the Riesz projection associated with
$\Eigenvalue$.
Then
\[
 \range\EigenvalueProjection{\CompactOperator}{\Eigenvalue}
 =\GeneralizedEigenspace{\CompactOperator}{\Eigenvalue}.
\]
More generally, let
$\SelectedEigenvalues\subset\spec(\CompactOperator)\setminus\{0\}$
be a finite set and let
$\SpectralContour$
be as in Definition~\ref{def:riesz-projection} for
$\CompactOperator$.
Assume that
$\SpectralContour$
has winding number
$1$
at every point of
$\SelectedEigenvalues$
and
$0$
at every point of
$\spec(\CompactOperator)\setminus\SelectedEigenvalues$.
Then we have
\[
 \RieszProjection{\CompactOperator}{\SpectralContour}
 =\sum_{\Eigenvalue\in\SelectedEigenvalues}
   \EigenvalueProjection{\CompactOperator}{\Eigenvalue},
\]
and
\[
 \range\RieszProjection{\CompactOperator}{\SpectralContour}
 =\bigoplus_{\Eigenvalue\in\SelectedEigenvalues}
   \GeneralizedEigenspace{\CompactOperator}{\Eigenvalue}.
\]
\end{theorem}

\subsection{Square roots of positive definite matrices}\label{subsec:spectral-preliminaries-matrices}
\paraid{p-e1b4c2b41113}
We use the following result to construct orthonormal bases from Gram matrices.

\begin{theorem}[{\cite[Theorems~7.29 and~7.39]{Axler2024}}]\label{thm:matrix-square-root}
\paraid{p-836db9990f19}
Every real symmetric matrix has an orthonormal basis of eigenvectors.
If it is positive definite,
then it has a unique symmetric positive definite square root.
In an orthonormal eigenbasis,
this square root is the diagonal matrix whose entries are the positive
square roots of the original eigenvalues.
\end{theorem}


\section{Finite-dimensional approximation of distances}\label{sec:spectral}
\paraid{p-d1bcda9ea36a}
We now approximate the distance functions on a compact metric space by
functions in a finite-dimensional vector space.
The selected probability measure makes this vector space invariant under
isometries.
Using the spectral results from Section~\ref{sec:prelim-spectral},
we first prove approximation at one space.
We then compare the resulting functions and coordinates on nearby spaces.

\paraid{p-a4c5f187e680}
A simple construction explains why finite-dimensional coordinates can
approximate a metric.
Let
$\MetricSpace{\BaseCarrier}{\MetricSymbol}$
be nonempty and compact,
fix
$\ErrorTolerance>0$,
and choose a finite
$\ErrorTolerance$-net
$\{\NetPoint_0,\ldots,\NetPoint_{\NetSize-1}\}$.
Define
\[
 \NetCoordinateMap\colon\BaseCarrier\to\RealNumbers^{\NetSize}
\]
by
\[
 \NetCoordinateMap(\BasePoint)
 =\bigl(\MetricSymbol(\BasePoint,\NetPoint_0),\ldots,
        \MetricSymbol(\BasePoint,\NetPoint_{\NetSize-1})\bigr).
\]
For all
$\BasePoint,\ComparisonPoint\in\BaseCarrier$,
\begin{equation}\label{eq:finite-net-distance-approximation}
 \MetricSymbol(\BasePoint,\ComparisonPoint)-2\ErrorTolerance
 \leq\norm{\NetCoordinateMap(\BasePoint)-\NetCoordinateMap(\ComparisonPoint)}_\infty
 \leq\MetricSymbol(\BasePoint,\ComparisonPoint).
\end{equation}

\paraid{p-37280d9fcd1a}
Finite families of functions also underlie Shioya's approximation of metric
measure spaces
\cite[author's manuscript, Section~4.4, Definition~4.37, Theorem~4.43 and Corollary~4.44]{ShioyaMetricMeasure}.
More precisely,
for every mm-space
$(X,d_X,\mu_X)$
and every
$\ErrorTolerance>0$,
there is an integer
$N\geq1$
for which we can choose real
$1$-Lipschitz
functions
$f_1,\ldots,f_N$
on
$X$
such that the map
$\Phi=(f_1,\ldots,f_N)$
satisfies
\[
 \BoxDistance\bigl(X,(\RealNumbers^N,\norm{\cdot}_\infty,
                         \Phi_*\mu_X)\bigr)<\ErrorTolerance.
\]
Here
$\BoxDistance$
is the distance in \Cref{def:box-distance}.
Estimate \eqref{eq:finite-net-distance-approximation}
concerns uniform error over all pairs of points
of one compact space.
Our construction also needs continuity throughout a Gromov--Hausdorff neighborhood
and compatibility with isometries.
A finite net in one space does not determine such continuous choices.
Finite spectral subspaces provide the coordinate classes of \Cref{thm:local-models}.
For a convergent sequence of spaces,
\ref{item:model-continuity} guarantees coordinate pairs with the stated
convergence of norms and coordinate maps.
Condition \ref{item:model-equivariance} specifies their transformation under isometries.
Their coordinates are the coefficients of unique best approximations.
For each compact metric space
$\MetricSpace{\ComparisonCarrier}{\ComparisonMetric}$,
we define a norm
$\NormSymbol_\ComparisonCarrier$
on the coordinate space.
Under a change of orthonormal basis,
each coordinate vector $v$ is replaced by
$\OrthogonalChange v$.
The norm in the new coordinates is therefore
$\NormSymbol_\ComparisonCarrier\circ\OrthogonalChange^{-1}$,
since a vector $v$ in the new coordinates represents
$\OrthogonalChange^{-1}v$
in the old coordinates.
Thus the distances between coordinate vectors do not depend on the basis.

\paraid{p-1bee81df97ab}
Throughout this section,
$\SelectedMeasure{\BaseCarrier}{\MetricSymbol}$
denotes the assignment of \Cref{thm:invariant-measures}.
The two subscripts denote the carrier
$\BaseCarrier$
and its metric
$\MetricSymbol$.

\subsection{The distance operator and its spectral subspaces}\label{subsec:spectral-distance-operator}
\paraid{p-834e7a7b7a65}
Fix a nonempty compact metric space
$\MetricSpace{\BaseCarrier}{\MetricSymbol}$.
Equip the space
\[
 \HilbertDomain _\BaseCarrier =\LebesgueSpace^2(\BaseCarrier ,\SelectedMeasure{\BaseCarrier}{\MetricSymbol})
\]
with the inner product
\[
 \langle \HilbertFunction ,\HilbertTestFunction \rangle_{\HilbertDomain _\BaseCarrier }=\int_\BaseCarrier  \HilbertFunction  \HilbertTestFunction \,\IntegrationDifferential \SelectedMeasure{\BaseCarrier}{\MetricSymbol}.
\]
This is a separable real Hilbert space.
Here separability follows from density of continuous functions in
$\LebesgueSpace^2$
and the separability of
$\ContinuousFunctions(\BaseCarrier )$
for compact metric spaces
\cite[Section~6.4]{Muscat2024}.
We define the \emph{distance operator}
$\DistanceOperator{\BaseCarrier}{\MetricSymbol}\colon\HilbertDomain_\BaseCarrier\to\HilbertDomain_\BaseCarrier$
as follows.
For every
$\HilbertFunction\in\HilbertDomain_\BaseCarrier$
and
$\BasePoint\in\BaseCarrier$,
set
\[
 \DistanceOperator{\BaseCarrier}{\MetricSymbol}\HilbertFunction (\BasePoint )
 =\int_\BaseCarrier  \MetricSymbol (\BasePoint ,\IntegrationPoint )\HilbertFunction (\IntegrationPoint )\,\IntegrationDifferential \SelectedMeasure{\BaseCarrier}{\MetricSymbol}(\IntegrationPoint ).
\]
The integral is independent of the representative of
$\HilbertFunction $.
This is the distance kernel operator of
\cite[Definition~11]{MariaOudotSolomon2020}.

\begin{lemma}\label{lem:distance-operator}
\paraid{p-fd7cf639024a}
The operator
$\DistanceOperator{\BaseCarrier}{\MetricSymbol}\colon \HilbertDomain _\BaseCarrier \to \HilbertDomain _\BaseCarrier $
is compact and self-adjoint.
It maps
$\HilbertDomain _\BaseCarrier $
into
$\ContinuousFunctions(\BaseCarrier )$,
and for every
$\HilbertFunction \in \HilbertDomain _\BaseCarrier $
and
$\BasePoint ,\ComparisonPoint \in \BaseCarrier $,
\begin{equation}\label{eq:operator-uniform-bound-1}
 \norm{\DistanceOperator{\BaseCarrier}{\MetricSymbol}\HilbertFunction }_\infty
 \leq\diam\MetricSpace{\BaseCarrier }{\MetricSymbol }\norm{\HilbertFunction }_2,
\end{equation}
and
\begin{equation}\label{eq:operator-uniform-bound-2}
 |\DistanceOperator{\BaseCarrier}{\MetricSymbol}\HilbertFunction (\BasePoint )-\DistanceOperator{\BaseCarrier}{\MetricSymbol}\HilbertFunction (\ComparisonPoint )|
 \leq \MetricSymbol (\BasePoint ,\ComparisonPoint )\norm{\HilbertFunction }_2.
\end{equation}
Every eigenfunction with nonzero eigenvalue has a unique continuous
representative.
\end{lemma}

\paraid{p-e83149336582}
Compactness and self-adjointness of the distance operator are also proved in
\cite[Propositions~12 and~13]{MariaOudotSolomon2020}.

\begin{proof}
\paraid{p-3eab3569fc53}
Let
$\HilbertFunction\in\HilbertDomain_\BaseCarrier$
and
$\BasePoint,\ComparisonPoint\in\BaseCarrier$.
By the Cauchy--Schwarz inequality
\cite[Lemma~3.2.27]{BreitGmeineder2026}
and
$\SelectedMeasure{\BaseCarrier}{\MetricSymbol}(\BaseCarrier)=1$,
we obtain
\[
 \int_\BaseCarrier |\HilbertFunction |\,\IntegrationDifferential \SelectedMeasure{\BaseCarrier}{\MetricSymbol}\leq\norm{\HilbertFunction }_2.
\]
Consequently,
\begin{equation}\label{eq:distance-operator-pointwise-bounds-1}
|\DistanceOperator{\BaseCarrier}{\MetricSymbol}\HilbertFunction (\BasePoint )|
 \leq\int_\BaseCarrier  \MetricSymbol (\BasePoint ,\IntegrationPoint )|\HilbertFunction (\IntegrationPoint )|\,\IntegrationDifferential \SelectedMeasure{\BaseCarrier}{\MetricSymbol}(\IntegrationPoint )
 \leq\diam\MetricSpace{\BaseCarrier }{\MetricSymbol }\norm{\HilbertFunction }_2,
\end{equation}
and
\begin{equation}\label{eq:distance-operator-pointwise-bounds-2}
\begin{aligned}
|\DistanceOperator{\BaseCarrier}{\MetricSymbol}\HilbertFunction (\BasePoint )-\DistanceOperator{\BaseCarrier}{\MetricSymbol}\HilbertFunction (\ComparisonPoint )|
 &\leq\int_\BaseCarrier |\MetricSymbol (\BasePoint ,\IntegrationPoint )-\MetricSymbol (\ComparisonPoint ,\IntegrationPoint )|\,|\HilbertFunction (\IntegrationPoint )|\,\IntegrationDifferential \SelectedMeasure{\BaseCarrier}{\MetricSymbol}(\IntegrationPoint )\\
 &\leq \MetricSymbol (\BasePoint ,\ComparisonPoint )\norm{\HilbertFunction }_2.
\end{aligned}
\end{equation}
The estimates \eqref{eq:distance-operator-pointwise-bounds-1} and
\eqref{eq:distance-operator-pointwise-bounds-2},
together with the Arzel\`a--Ascoli theorem (\Cref{thm:arzela-ascoli}),
show that
$\DistanceOperator{\BaseCarrier}{\MetricSymbol}\colon
\HilbertDomain_\BaseCarrier\to\ContinuousFunctions(\BaseCarrier)$
is compact.
Since the inclusion
$\ContinuousFunctions(\BaseCarrier)\to\HilbertDomain_\BaseCarrier$
is bounded,
the operator on
$\HilbertDomain_\BaseCarrier$
is compact as well.
Symmetry of the kernel and Fubini's theorem
\cite[Theorem~2.8.20(a)]{BreitGmeineder2026}
imply self-adjointness.

\paraid{p-9cf941c90fec}
If
$\DistanceOperator{\BaseCarrier}{\MetricSymbol}\HilbertFunction =\Eigenvalue  \HilbertFunction $
in
$\HilbertDomain _\BaseCarrier $
and
$\Eigenvalue \ne0$,
then
$\Eigenvalue ^{-1}\DistanceOperator{\BaseCarrier}{\MetricSymbol}\HilbertFunction $
is a continuous representative of
$\HilbertFunction $.
The representative is unique because
$\SelectedMeasure{\BaseCarrier}{\MetricSymbol}$
has full support.
\end{proof}

\paraid{p-f6f6a5a89ba1}
For every eigenfunction with nonzero eigenvalue, we use its unique continuous representative.
Apply \Cref{prop:spectral-basics} to
$\DistanceOperator{\BaseCarrier}{\MetricSymbol}$
on
$\HilbertDomain_\BaseCarrier$,
and choose an orthonormal family
$\{\Eigenfunction_\CoordinateIndex\}_{\CoordinateIndex\in\EigenIndexSet}$
with corresponding nonzero eigenvalues
$\{\Eigenvalue_\CoordinateIndex\}_{\CoordinateIndex\in\EigenIndexSet}$
as in that theorem.
For
$\SpectralCutoff >0$,
put
\[
 \EigenIndexSet _{\BaseCarrier ,\SpectralCutoff }=\{\CoordinateIndex \in \EigenIndexSet \mid|\Eigenvalue _\CoordinateIndex |>\SpectralCutoff \},
\]
and
\[
 \SpectralSubspace{\BaseCarrier }{\SpectralCutoff }
 =\LinearSpan \{\Eigenfunction _\CoordinateIndex \mid \CoordinateIndex \in \EigenIndexSet _{\BaseCarrier ,\SpectralCutoff }\}.
\]
This is a finite-dimensional space of continuous functions.
Define the orthogonal projection
$\SpectralProjection{\BaseCarrier}{\SpectralCutoff}
 \colon\HilbertDomain_\BaseCarrier\to\SpectralSubspace{\BaseCarrier}{\SpectralCutoff}$
by
\begin{equation}\label{eq:spectral-projection-definition}
 \SpectralProjection{\BaseCarrier }{\SpectralCutoff }\HilbertFunction
 =\sum_{\CoordinateIndex \in \EigenIndexSet _{\BaseCarrier ,\SpectralCutoff }}\langle \HilbertFunction ,\Eigenfunction _\CoordinateIndex \rangle_{\HilbertDomain _\BaseCarrier }\Eigenfunction _\CoordinateIndex .
\end{equation}
Equivalently,
\[
 \begin{aligned}
 \SpectralSubspace{\BaseCarrier }{\SpectralCutoff }
 =\bigoplus_{\substack{\Eigenvalue \in\spec(\DistanceOperator{\BaseCarrier}{\MetricSymbol})\\|\Eigenvalue |>\SpectralCutoff }}
       \ker(\DistanceOperator{\BaseCarrier}{\MetricSymbol}-\Eigenvalue \IdentityMap),
 \end{aligned}
\]
and
\[
 \range\SpectralProjection{\BaseCarrier }{\SpectralCutoff }=\SpectralSubspace{\BaseCarrier }{\SpectralCutoff }.
\]
These definitions do not depend on the orthonormal basis
$\{\Eigenfunction_\CoordinateIndex\}_{\CoordinateIndex\in\EigenIndexSet}$.
For every
$\HilbertFunction\in\HilbertDomain_\BaseCarrier$,
the vector
$\SpectralProjection{\BaseCarrier}{\SpectralCutoff}\HilbertFunction$
converges in the norm of
$\HilbertDomain_\BaseCarrier$
to the orthogonal projection of
$\HilbertFunction$
onto
$(\ker\DistanceOperator{\BaseCarrier}{\MetricSymbol})^\perp$
as
$\SpectralCutoff\downarrow0$.

\subsection{Approximation of distance functions}\label{subsec:spectral-distance-approximation}
\paraid{p-a700e9644abe}
We first approximate distance profiles in the uniform norm by functions in
finite-dimensional spectral subspaces.
For
$\BasePoint\in\BaseCarrier$,
we call the function
\[
 \MetricSymbol_\BasePoint=\MetricSymbol(\BasePoint,\cdot)
 \in\ContinuousFunctions(\BaseCarrier)
\]
the \emph{distance profile} of
$\BasePoint$.
Its values are the distances from
$\BasePoint$
to all points of
$\BaseCarrier$.
The map
\[
 \BaseCarrier\to\ContinuousFunctions(\BaseCarrier),
 \qquad\BasePoint\longmapsto\MetricSymbol_\BasePoint
\]
is an isometric embedding for the uniform norm,
because
\[
 \norm{\MetricSymbol_\BasePoint-\MetricSymbol_\ComparisonPoint}_\infty
 =\MetricSymbol(\BasePoint,\ComparisonPoint).
\]
This is the distance-function construction of the Kuratowski embedding.
Here the functions
$\MetricSymbol_\BasePoint$
are bounded because
$\MetricSpace{\BaseCarrier}{\MetricSymbol}$
is compact.
The following theorem approximates all distance profiles in one finite
spectral subspace.

\begin{theorem}\label{lem:spectral-density}
\paraid{p-c4a771b9c692}
For every nonempty compact metric space
$\MetricSpace{\BaseCarrier }{\MetricSymbol }$
and every
$\SmallError >0$,
there is
$\SpectralCutoff >0$
such that
$\SpectralCutoff ,-\SpectralCutoff \notin\spec(\DistanceOperator{\BaseCarrier}{\MetricSymbol})$
and
\begin{equation}\label{eq:spectral-uniform-approximation}
 \sup_{\BasePoint \in \BaseCarrier }\inf_{\Vector \in\SpectralSubspace{\BaseCarrier }{\SpectralCutoff }}
          \norm{\MetricSymbol _\BasePoint -\Vector }_\infty<\SmallError .
\end{equation}
If $\Card(\BaseCarrier)\geq2$,
then we may also require
$\dim\SpectralSubspace{\BaseCarrier }{\SpectralCutoff }\geq1$.
\end{theorem}

\begin{proof}
\paraid{p-53612876bfd1}
Let
$\UniformSpectralSpan $
be the closed linear span in
$\ContinuousFunctions(\BaseCarrier )$
of the continuous representatives of all nonzero eigenfunctions.
For
$\HilbertFunction \in \HilbertDomain _\BaseCarrier $,
write
$\HilbertFunction =\HilbertFunction _0+\HilbertFunction _\perp$
with
$\HilbertFunction _0\in\ker\DistanceOperator{\BaseCarrier}{\MetricSymbol}$
and
$\HilbertFunction _\perp\perp\ker\DistanceOperator{\BaseCarrier}{\MetricSymbol}$.
Choose a decreasing sequence of positive real numbers
$\{\SpectralCutoff_m\}_{m\in\NonnegativeIntegers}$
converging to
$0$
such that, for every $m\in\NonnegativeIntegers$,
$\SpectralCutoff_m,-\SpectralCutoff_m\notin\spec(\DistanceOperator{\BaseCarrier}{\MetricSymbol})$.
Such a sequence exists because the nonzero spectrum is countable.
By the Riesz--Schauder and self-adjoint spectral theorem (\Cref{prop:spectral-basics}),
we have
$\norm{\SpectralProjection{\BaseCarrier }{\SpectralCutoff_m}\HilbertFunction -\HilbertFunction _\perp}_2\to0$.
Since
$\HilbertFunction_0\in\ker\DistanceOperator{\BaseCarrier}{\MetricSymbol}$,
the continuous function
$\DistanceOperator{\BaseCarrier}{\MetricSymbol}\HilbertFunction_0$
vanishes almost everywhere.
By full support of
$\SelectedMeasure{\BaseCarrier}{\MetricSymbol}$,
for every
$\BasePoint\in\BaseCarrier$
we have
\begin{equation}\label{eq:continuous-kernel-vanishing}
 (\DistanceOperator{\BaseCarrier}{\MetricSymbol}\HilbertFunction_0)(\BasePoint)=0.
\end{equation}
Using \eqref{eq:operator-uniform-bound-1},
\[
 \norm{\DistanceOperator{\BaseCarrier}{\MetricSymbol}\SpectralProjection{\BaseCarrier }{\SpectralCutoff_m}\HilbertFunction
                 -\DistanceOperator{\BaseCarrier}{\MetricSymbol}\HilbertFunction }_\infty
 \leq\diam\MetricSpace{\BaseCarrier }{\MetricSymbol }
       \norm{\SpectralProjection{\BaseCarrier }{\SpectralCutoff_m}\HilbertFunction -\HilbertFunction _\perp}_2\to0.
\]
For every
$\SpectralCutoff>0$,
we have
\[
 \DistanceOperator{\BaseCarrier}{\MetricSymbol}
 \SpectralProjection{\BaseCarrier}{\SpectralCutoff}\HilbertFunction
 \in\SpectralSubspace{\BaseCarrier}{\SpectralCutoff}
 \subset\UniformSpectralSpan.
\]
Since
$\UniformSpectralSpan$
is closed in the uniform norm,
it follows that
$\DistanceOperator{\BaseCarrier}{\MetricSymbol}\HilbertFunction \in \UniformSpectralSpan $.

\paraid{p-4c9880842102}
Fix
$\BasePoint \in \BaseCarrier $
and
$\Radius >0$.
Use the normalized function
$\NormalizedBump{\BaseCarrier}{\MetricSymbol}{\BasePoint}{\Radius}{\SelectedMeasure{\BaseCarrier}{\MetricSymbol}}$
from \eqref{eq:normalized-bump-definition}.
By \eqref{eq:normalized-bump-properties-1},
\eqref{eq:normalized-bump-properties-2} and \eqref{eq:normalized-bump-properties-3},
the normalized function is nonnegative,
has integral one,
and is supported in
$\MetricBall(\BasePoint,\Radius;\MetricSymbol)$.
For every
$\ComparisonPoint \in \BaseCarrier $,
\begin{equation}\label{eq:bump-distance-approximation}
 \begin{aligned}
 &|\DistanceOperator{\BaseCarrier}{\MetricSymbol}\NormalizedBump{\BaseCarrier}{\MetricSymbol}{\BasePoint}{\Radius}{\SelectedMeasure{\BaseCarrier}{\MetricSymbol}}(\ComparisonPoint )-\MetricSymbol _\BasePoint (\ComparisonPoint )|\\
 &=\left|\int_\BaseCarrier (\MetricSymbol (\ComparisonPoint ,\IntegrationPoint )-\MetricSymbol (\ComparisonPoint ,\BasePoint ))\NormalizedBump{\BaseCarrier}{\MetricSymbol}{\BasePoint}{\Radius}{\SelectedMeasure{\BaseCarrier}{\MetricSymbol}}(\IntegrationPoint )
                                      \,\IntegrationDifferential \SelectedMeasure{\BaseCarrier}{\MetricSymbol}(\IntegrationPoint )\right|\\
 &\leq\int_\BaseCarrier  \MetricSymbol (\BasePoint ,\IntegrationPoint )\NormalizedBump{\BaseCarrier}{\MetricSymbol}{\BasePoint}{\Radius}{\SelectedMeasure{\BaseCarrier}{\MetricSymbol}}(\IntegrationPoint )\,\IntegrationDifferential \SelectedMeasure{\BaseCarrier}{\MetricSymbol}(\IntegrationPoint )
 \leq \Radius .
 \end{aligned}
\end{equation}
Since \eqref{eq:bump-distance-approximation} holds for every
$\Radius>0$
and
$\UniformSpectralSpan$
is closed in the uniform norm,
we obtain
$\MetricSymbol_\BasePoint\in\UniformSpectralSpan$.

\paraid{p-0c97576ebc0b}
Recall that
\[
 \norm{\MetricSymbol _\BasePoint -\MetricSymbol _\ComparisonPoint }_\infty=\MetricSymbol (\BasePoint ,\ComparisonPoint ),
\]
as established before the theorem.
Choose finitely many points
$\BasePoint_0,\ldots,\BasePoint_{\NetSize-1}$
whose
$\SmallError /3$-balls
cover
$\BaseCarrier $.
For each
$\NetIndex $,
choose a finite sum
$\Vector _\NetIndex $
of nonzero eigenfunctions with
$\norm{\MetricSymbol _{\BasePoint _\NetIndex }-\Vector _\NetIndex }_\infty<\SmallError /3$.
Let
$\ApproximationEigenvalues$
be the finite set of eigenvalues occurring in these sums.
Choose
$\SpectralCutoff>0$
such that
\[
 \SpectralCutoff,-\SpectralCutoff\notin\spec(\DistanceOperator{\BaseCarrier}{\MetricSymbol}),
\]
and such that for every
$\Eigenvalue\in\ApproximationEigenvalues$
we have
\[
 \SpectralCutoff<|\Eigenvalue|.
\]
Note that even if
$\ApproximationEigenvalues=\emptyset$,
it suffices to choose a positive real number
$\SpectralCutoff$
satisfying
$\SpectralCutoff,-\SpectralCutoff\notin\spec(\DistanceOperator{\BaseCarrier}{\MetricSymbol})$.
Such a real number exists because the nonzero spectrum is countable.
Then every
$\Vector _\NetIndex $
belongs to
$\SpectralSubspace{\BaseCarrier }{\SpectralCutoff }$.
For
$\MetricSymbol (\BasePoint ,\BasePoint _\NetIndex )<\SmallError /3$,
\[
 \inf_{\Vector \in\SpectralSubspace{\BaseCarrier }{\SpectralCutoff }}\norm{\MetricSymbol _\BasePoint -\Vector }_\infty
 \leq\norm{\MetricSymbol _\BasePoint -\MetricSymbol _{\BasePoint _\NetIndex }}_\infty+\norm{\MetricSymbol _{\BasePoint _\NetIndex }-\Vector _\NetIndex }_\infty
 <2\SmallError /3.
\]
This proves \eqref{eq:spectral-uniform-approximation}.

\paraid{p-f6773180e565}
If
$\BaseCarrier $
is not a singleton,
choose
$\BasePoint ,\IntegrationPoint \in \BaseCarrier $
with
$\MetricSymbol (\BasePoint ,\IntegrationPoint )>0$.
The integral
$\DistanceOperator{\BaseCarrier}{\MetricSymbol}1(\BasePoint )$
is positive because the distance from
$\BasePoint $
is positive on a neighborhood of
$\IntegrationPoint $
of positive measure.
Thus
$\DistanceOperator{\BaseCarrier}{\MetricSymbol}\ne0$.
By the Riesz--Schauder and self-adjoint spectral theorem (\Cref{prop:spectral-basics}),
it has a nonzero eigenvalue.
If necessary,
choose a smaller positive
$\SpectralCutoff$
below the absolute value of this eigenvalue and still satisfying
$\pm\SpectralCutoff\notin\spec(\DistanceOperator{\BaseCarrier}{\MetricSymbol})$.
This retains the approximation in \eqref{eq:spectral-uniform-approximation} and ensures
$\dim\SpectralSubspace{\BaseCarrier}{\SpectralCutoff}\geq1$.
\end{proof}

\paraid{p-9e170f5d7db0}
The uniform distance between the distance profiles is the original distance,
but a best approximation in a finite-dimensional subspace need not be
unique for the uniform norm.
We minimize a finite
$\LebesgueSpace^\IntegrabilityExponent $-norm
instead.
For
$1<\IntegrabilityExponent <\infty$,
the best approximation exists and is unique by \Cref{thm:best-approximation}.
\Cref{lem:lp-distance} shows that the
$\LebesgueSpace^\IntegrabilityExponent$-distances
between distance profiles
approximate the original distances uniformly when
$\IntegrabilityExponent$
is sufficiently large.
We then combine \eqref{eq:lp-distance-error} with finite spectral approximation in
\eqref{eq:finite-spectral-error}.

\paraid{p-c8d27acdc339}
To compare the best approximations,
we use the following continuity statement for their coefficients.

\begin{lemma}\label{lem:minimizer-continuity}
\paraid{p-df64ce5d76d4}
Let
$\ModelDimension\geq1$,
and let continuous functions
$\ApproximationObjective_\SequenceIndex,\ApproximationObjective\colon
 \RealNumbers^\ModelDimension\to\RealNumbers$
converge uniformly on every compact subset.
Assume that each
$\CoefficientVector_\SequenceIndex$
minimizes
$\ApproximationObjective_\SequenceIndex$,
that these vectors form a bounded sequence,
and that
$\ApproximationObjective$
has a unique minimizer
$\CoefficientVector$.
Then
$\CoefficientVector_\SequenceIndex\to\CoefficientVector$.
\end{lemma}

\begin{proof}
\paraid{p-586499da67ea}
By boundedness of $\CoefficientVector_j$,
choose
$\CoefficientBound>0$
such that
\[
 \{\CoefficientVector_\SequenceIndex\mid\SequenceIndex\in\NonnegativeIntegers\}
 \subset\{\CompetitorCoefficients\in\RealNumbers^\ModelDimension
       \mid\norm{\CompetitorCoefficients}_2\leq\CoefficientBound\}.
\]
This closed ball is compact,
so every subsequence has a further convergent subsequence.
Let
$\ComparisonCoefficients$
denote the limit of such a subsequence
$\{\CoefficientVector_{n_k}\}_{k\in\NonnegativeIntegers}$.
The limit satisfies
$\norm{\ComparisonCoefficients}_2\leq\CoefficientBound$,
and the assumed uniform convergence implies
\begin{equation}\label{eq:minimizer-uniform-objective-error}
 \sup_{\substack{\CompetitorCoefficients\in\RealNumbers^\ModelDimension\\\norm{\CompetitorCoefficients}_2\leq\CoefficientBound}}
 |\ApproximationObjective_{n_k}(\CompetitorCoefficients)
       -\ApproximationObjective(\CompetitorCoefficients)|\longrightarrow0.
\end{equation}
Since $\ApproximationObjective$ is continuous and
$\CoefficientVector_{n_k}\to\ComparisonCoefficients$,
the difference between its values at these two coefficient vectors tends to $0$.
The objective error
$|\ApproximationObjective_{n_k}(\CoefficientVector_{n_k})
 -\ApproximationObjective(\CoefficientVector_{n_k})|$
tends to $0$ by \eqref{eq:minimizer-uniform-objective-error}.
Consequently,
\begin{equation}\label{eq:minimizer-objective-limit}
 \begin{aligned}
 &|\ApproximationObjective_{n_k}(\CoefficientVector_{n_k})
        -\ApproximationObjective(\ComparisonCoefficients)|
 \leq\\
 &|\ApproximationObjective_{n_k}(\CoefficientVector_{n_k})
        -\ApproximationObjective(\CoefficientVector_{n_k})|+|\ApproximationObjective(\CoefficientVector_{n_k})
        -\ApproximationObjective(\ComparisonCoefficients)|\longrightarrow0.
 \end{aligned}
\end{equation}
For a fixed competitor
$\CompetitorCoefficients\in\RealNumbers^\ModelDimension$,
minimality at each index and convergence at that competitor now imply
\begin{equation}\label{eq:limit-minimizer-comparison}
 \ApproximationObjective(\ComparisonCoefficients)
 =\lim_{k\to\infty}\ApproximationObjective_{n_k}(\CoefficientVector_{n_k})
 \leq\lim_{k\to\infty}\ApproximationObjective_{n_k}(\CompetitorCoefficients)
 =\ApproximationObjective(\CompetitorCoefficients).
\end{equation}
Thus
$\ComparisonCoefficients=\CoefficientVector$
by uniqueness.
Relative compactness and uniqueness of the subsequential limit imply
$\CoefficientVector_\SequenceIndex\to\CoefficientVector$.
\end{proof}

\begin{lemma}\label{lem:lp-distance}
\paraid{p-c7556a5ad9a5}
Let
$\MetricSpace{\BaseCarrier }{\MetricSymbol }$
be a nonempty compact metric space and set
$\DiameterBound =\diam\MetricSpace{\BaseCarrier }{\MetricSymbol }$.
Then for every
$\Radius >0$,
the number
\[
 \BallMassBound _\Radius =\inf_{\BasePoint \in \BaseCarrier }\SelectedMeasure{\BaseCarrier}{\MetricSymbol}(\MetricBall(\BasePoint ,\Radius;\MetricSymbol))
\]
is positive.
Moreover,
for every
$\Radius>0$,
every
$2\leq \IntegrabilityExponent <\infty$,
and every pair
$\BasePoint ,\ComparisonPoint \in \BaseCarrier $,
\begin{equation}\label{eq:lp-distance-error}
 0\leq \MetricSymbol (\BasePoint ,\ComparisonPoint )-\norm{\MetricSymbol _\BasePoint -\MetricSymbol _\ComparisonPoint }_{\LebesgueSpace^\IntegrabilityExponent (\SelectedMeasure{\BaseCarrier}{\MetricSymbol})}
 \leq \DiameterBound (1-\BallMassBound _\Radius ^{1/\IntegrabilityExponent })+2\Radius .
\end{equation}
\end{lemma}

\begin{proof}
\paraid{p-ce0329a6de36}
In this proof, $\norm{\cdot}_\IntegrabilityExponent$ denotes the norm on
$\LebesgueSpace^\IntegrabilityExponent(\BaseCarrier,\SelectedMeasure{\BaseCarrier}{\MetricSymbol})$.
Choose a finite
$\Radius /4$-net
$\{\IntegrationPoint_0,\ldots,\IntegrationPoint_{\NetSize-1}\}$.
For each
$\BasePoint \in \BaseCarrier $,
some
$\NetIndex $
satisfies
$\MetricSymbol (\BasePoint ,\IntegrationPoint _\NetIndex )<\Radius /4$,
and then
$\MetricBall(\IntegrationPoint _\NetIndex ,\Radius /4;\MetricSymbol)\subset\MetricBall(\BasePoint ,\Radius;\MetricSymbol)$.
Therefore,
\[
 \BallMassBound _\Radius \geq\min_{0\leq \NetIndex<\NetSize }
          \SelectedMeasure{\BaseCarrier}{\MetricSymbol}(\MetricBall(\IntegrationPoint _\NetIndex ,\Radius /4;\MetricSymbol))>0.
\]
For all
$\IntegrationPoint \in \BaseCarrier $,
$|\MetricSymbol _\BasePoint (\IntegrationPoint )-\MetricSymbol _\ComparisonPoint (\IntegrationPoint )|\leq \MetricSymbol (\BasePoint ,\ComparisonPoint )$.
For
$\IntegrationPoint \in\MetricBall(\BasePoint ,\Radius;\MetricSymbol)$,
\[
 \MetricSymbol _\ComparisonPoint (\IntegrationPoint )-\MetricSymbol _\BasePoint (\IntegrationPoint )\geq \MetricSymbol (\BasePoint ,\ComparisonPoint )-2\MetricSymbol (\BasePoint ,\IntegrationPoint )>\MetricSymbol (\BasePoint ,\ComparisonPoint )-2\Radius .
\]
Consequently,
\begin{equation}\label{eq:lp-distance-two-sided-bound}
 \BallMassBound _\Radius ^{1/\IntegrabilityExponent }\max\{\MetricSymbol (\BasePoint ,\ComparisonPoint )-2\Radius ,0\}
 \leq\norm{\MetricSymbol _\BasePoint -\MetricSymbol _\ComparisonPoint }_\IntegrabilityExponent\leq \MetricSymbol (\BasePoint ,\ComparisonPoint ).
\end{equation}
Since
$\SelectedMeasure{\BaseCarrier}{\MetricSymbol}$
is a probability and the ball masses are positive,
$0<\BallMassBound_\Radius\leq1$.
Consequently,
$0<\BallMassBound_\Radius^{1/\IntegrabilityExponent}\leq1$.
From the lower bound in \eqref{eq:lp-distance-two-sided-bound} and
\[
\max\{\MetricSymbol(\BasePoint,\ComparisonPoint)-2\Radius,0\}
 \geq\MetricSymbol(\BasePoint,\ComparisonPoint)-2\Radius
 \]
it follows that
\[
 \begin{aligned}
 \MetricSymbol(\BasePoint,\ComparisonPoint)
   -\norm{\MetricSymbol_\BasePoint-\MetricSymbol_\ComparisonPoint}_\IntegrabilityExponent
 &\leq \MetricSymbol(\BasePoint,\ComparisonPoint)
 -\BallMassBound_\Radius^{1/\IntegrabilityExponent}
 \max\{\MetricSymbol(\BasePoint,\ComparisonPoint)-2\Radius,0\}\\
 &\leq \MetricSymbol(\BasePoint,\ComparisonPoint)
 -\BallMassBound_\Radius^{1/\IntegrabilityExponent}
 \bigl(\MetricSymbol(\BasePoint,\ComparisonPoint)-2\Radius\bigr)\\
 &= (1-\BallMassBound_\Radius^{1/\IntegrabilityExponent})
      \MetricSymbol(\BasePoint,\ComparisonPoint)
       +2\Radius\BallMassBound_\Radius^{1/\IntegrabilityExponent}\\
 &\leq
 (1-\BallMassBound_\Radius^{1/\IntegrabilityExponent})
      \DiameterBound+2\Radius.
 \end{aligned}
\]
The last inequality uses
$\MetricSymbol(\BasePoint,\ComparisonPoint)\leq\DiameterBound$,
$1-\BallMassBound_\Radius^{1/\IntegrabilityExponent}\geq0$,
and
$\BallMassBound_\Radius^{1/\IntegrabilityExponent}\leq1$.
This proves \eqref{eq:lp-distance-error}.
\end{proof}

\paraid{p-533de80974c4}
Fix
$\SpectralSubspace{\BaseCarrier }{\SpectralCutoff }$
and
$2\leq \IntegrabilityExponent <\infty$.
Write $\norm{\cdot}_\IntegrabilityExponent$ for the norm on
$\LebesgueSpace^\IntegrabilityExponent(\BaseCarrier,\SelectedMeasure{\BaseCarrier}{\MetricSymbol})$.
For each
$\BasePoint \in \BaseCarrier $,
let
$\Vector _\BasePoint $
be the unique minimizer of
$\Vector \mapsto\norm{\MetricSymbol _\BasePoint -\Vector }_\IntegrabilityExponent ^\IntegrabilityExponent $
over
$\SpectralSubspace{\BaseCarrier }{\SpectralCutoff }$.
The spectral subspace
$\SpectralSubspace{\BaseCarrier}{\SpectralCutoff}$
is finite-dimensional,
and hence closed and convex in
$\LebesgueSpace^\IntegrabilityExponent(\BaseCarrier,\SelectedMeasure{\BaseCarrier}{\MetricSymbol})$.
By \Cref{thm:best-approximation},
this minimizer exists and is unique.
Since
$\SelectedMeasure{\BaseCarrier}{\MetricSymbol}$
has full support,
the natural map
$\ContinuousFunctions(\BaseCarrier)\to
 \LebesgueSpace^\IntegrabilityExponent(\BaseCarrier,\SelectedMeasure{\BaseCarrier}{\MetricSymbol})$
is injective.

\paraid{p-f10c0269a0dd}
We call the function
$\Vector_\BasePoint\in\SpectralSubspace{\BaseCarrier}{\SpectralCutoff}$
the \emph{finite-dimensional approximation of the distance profile}.
The map
$\BaseCarrier\to\SpectralSubspace{\BaseCarrier}{\SpectralCutoff}$,
$\BasePoint\mapsto\Vector_\BasePoint$,
need not be injective.
We use these approximating functions to approximate distances with a controlled error.

\paraid{p-774eca26292a}
If \eqref{eq:spectral-uniform-approximation} holds,
then
$\norm{\MetricSymbol _\BasePoint -\Vector _\BasePoint }_\IntegrabilityExponent <\SmallError $.
Define
$\Pseudometric_\BaseCarrier\colon\BaseCarrier\times\BaseCarrier\to[0,\infty)$
by
$\Pseudometric_\BaseCarrier(\BasePoint,\ComparisonPoint)
 =\norm{\Vector_\BasePoint-\Vector_\ComparisonPoint}_\IntegrabilityExponent$.
The norm triangle inequality implies
\[
 \left|\Pseudometric _\BaseCarrier (\BasePoint ,\ComparisonPoint )-\norm{\MetricSymbol _\BasePoint -\MetricSymbol _\ComparisonPoint }_\IntegrabilityExponent \right|
 \leq\norm{\Vector _\BasePoint -\MetricSymbol _\BasePoint }_\IntegrabilityExponent +\norm{\Vector _\ComparisonPoint -\MetricSymbol _\ComparisonPoint }_\IntegrabilityExponent <2\SmallError .
\]
Together with \Cref{lem:lp-distance},
this yields
\begin{equation}\label{eq:finite-spectral-error}
 \norm{\Pseudometric _\BaseCarrier -\MetricSymbol }_\infty
 \leq2\SmallError +\DiameterBound (1-\BallMassBound _\Radius ^{1/\IntegrabilityExponent })+2\Radius .
\end{equation}
For a prescribed error tolerance,
choose
$\Radius $
small,
then choose a sufficiently large finite
$\IntegrabilityExponent $,
and finally choose
$\SmallError $
small and a corresponding spectral subspace.

\subsection{Continuity of finite spectral subspaces}\label{subsec:spectral-subspace-continuity}
\paraid{p-f52fe507fe4d}
The distance operators for different compact spaces act on different Hilbert spaces.
We extend the distance kernels to one compact metric space
$\MetricSpace{\AmbientCarrier}{\AmbientMetric}$
to obtain operators on the fixed Banach space
$\ContinuousFunctions(\AmbientCarrier ,\ComplexNumbers)$.
We apply the spectral inclusion and spectral projection convergence theorems
of Anselone and Palmer,
\Cref{thm:spectral-inclusion,thm:riesz-convergence}
in Subsection~\ref{subsec:spectral-preliminaries-perturbation},
to these ambient operators.
By spectral inclusion,
the gaps at the chosen real numbers
$\pm\SpectralCutoff$
persist for all sufficiently large indices.
By convergence of the spectral projections,
the retained spectral subspaces
have the same dimension for all sufficiently large indices.

\paraid{p-040fa54a5428}
By \Cref{lem:parameter-integrals},
the parameter-dependent integrals used below converge uniformly.

\paraid{p-f33149a827cc}
The next proposition compares the generalized eigenspaces of the ambient
kernel operator with those of the self-adjoint distance operator.

\begin{proposition}\label{lem:ambient-eigenspaces}
\paraid{p-c756883ac07c}
Let
$\ComparisonCarrier$
be a nonempty compact subset of a compact metric space
$\MetricSpace{\AmbientCarrier}{\AmbientMetric}$,
and put
\[
 \ComparisonMetric=\AmbientMetric|_{\ComparisonCarrier\times\ComparisonCarrier},
\]
and
\[
 \DiameterBound=\diam\MetricSpace{\AmbientCarrier}{\AmbientMetric}.
\]
Regard
$\DistanceOperator{\ComparisonCarrier}{\ComparisonMetric}$
as its complex-linear extension
\[
 \DistanceOperator{\ComparisonCarrier}{\ComparisonMetric}\colon
 \LebesgueSpace^2(\ComparisonCarrier,\SelectedMeasure{\ComparisonCarrier}{\ComparisonMetric};\ComplexNumbers)
 \to
 \LebesgueSpace^2(\ComparisonCarrier,\SelectedMeasure{\ComparisonCarrier}{\ComparisonMetric};\ComplexNumbers).
\]
Define
$\RestrictionOperator _\ComparisonCarrier \colon\ContinuousFunctions(\AmbientCarrier ,\ComplexNumbers)\to
                       \LebesgueSpace^2(\ComparisonCarrier ,\SelectedMeasure{\ComparisonCarrier}{\ComparisonMetric};\ComplexNumbers)$
as follows.
For every
$\HilbertFunction\in\ContinuousFunctions(\AmbientCarrier,\ComplexNumbers)$,
let
$\RestrictionOperator_\ComparisonCarrier\HilbertFunction$
be the equivalence class of
$\HilbertFunction|_\ComparisonCarrier$.
Define
$\KernelExtensionOperator _\ComparisonCarrier \colon\LebesgueSpace^2(\ComparisonCarrier ,\SelectedMeasure{\ComparisonCarrier}{\ComparisonMetric};\ComplexNumbers)
                       \to\ContinuousFunctions(\AmbientCarrier ,\ComplexNumbers)$
as follows.
For every
$\HilbertTestFunction\in\LebesgueSpace^2(\ComparisonCarrier,\SelectedMeasure{\ComparisonCarrier}{\ComparisonMetric};\ComplexNumbers)$,
using
$\SecondVector\in\ComparisonCarrier$
as the integration variable,
define for every
$\IntegrationPoint\in\AmbientCarrier$
\[
 (\KernelExtensionOperator_\ComparisonCarrier\HilbertTestFunction)(\IntegrationPoint)
 =\int_\ComparisonCarrier
 \AmbientMetric(\IntegrationPoint,\SecondVector)\HilbertTestFunction(\SecondVector)
 \,\IntegrationDifferential\SelectedMeasure{\ComparisonCarrier}{\ComparisonMetric}(\SecondVector).
\]
The integral depends only on the equivalence class of
$\HilbertTestFunction$.
The bounded continuous kernel defines a continuous function on
$\AmbientCarrier$.
Set
\[
 \CarrierAmbientOperator_\ComparisonCarrier
 =\KernelExtensionOperator_\ComparisonCarrier\RestrictionOperator_\ComparisonCarrier
 \colon\ContinuousFunctions(\AmbientCarrier,\ComplexNumbers)
 \to\ContinuousFunctions(\AmbientCarrier,\ComplexNumbers).
\]
Then
$\CarrierAmbientOperator_\ComparisonCarrier$
is compact.
For every
$\Eigenvalue\in\ComplexNumbers\setminus\{0\}$,
the generalized eigenspaces satisfy
\[
 \GeneralizedEigenspace{\CarrierAmbientOperator_\ComparisonCarrier}{\Eigenvalue}
 =\ker(\CarrierAmbientOperator_\ComparisonCarrier-\Eigenvalue\IdentityMap),
\]
and
\[
 \GeneralizedEigenspace{\DistanceOperator{\ComparisonCarrier}{\ComparisonMetric}}{\Eigenvalue}
 =\ker(\DistanceOperator{\ComparisonCarrier}{\ComparisonMetric}-\Eigenvalue\IdentityMap).
\]
Moreover, for every
$\Eigenvalue\in\ComplexNumbers\setminus\{0\}$
and every integer
$\RootOrder\geq1$,
\[
 \ker\bigl((\CarrierAmbientOperator_\ComparisonCarrier-\Eigenvalue\IdentityMap)^\RootOrder\bigr)
 =\ker(\CarrierAmbientOperator_\ComparisonCarrier-\Eigenvalue\IdentityMap),
\]
and
\[
 \ker\bigl((\DistanceOperator{\ComparisonCarrier}{\ComparisonMetric}-\Eigenvalue\IdentityMap)^\RootOrder\bigr)
 =\ker(\DistanceOperator{\ComparisonCarrier}{\ComparisonMetric}-\Eigenvalue\IdentityMap).
\]
Restriction induces the linear isomorphism
\[
 \RestrictionOperator_\ComparisonCarrier\colon
 \GeneralizedEigenspace{\CarrierAmbientOperator_\ComparisonCarrier}{\Eigenvalue}
 \to
 \GeneralizedEigenspace{\DistanceOperator{\ComparisonCarrier}{\ComparisonMetric}}{\Eigenvalue},
\]
Its inverse is the restricted map
\[
 \Eigenvalue^{-1}\KernelExtensionOperator_\ComparisonCarrier\colon
 \GeneralizedEigenspace{\DistanceOperator{\ComparisonCarrier}{\ComparisonMetric}}{\Eigenvalue}
 \to\GeneralizedEigenspace{\CarrierAmbientOperator_\ComparisonCarrier}{\Eigenvalue}.
\]
Moreover,
\[
 \spec(\CarrierAmbientOperator_\ComparisonCarrier)\setminus\{0\}
 =\spec(\DistanceOperator{\ComparisonCarrier}{\ComparisonMetric})\setminus\{0\}
 \subset[-\DiameterBound,\DiameterBound].
\]
\end{proposition}

\begin{proof}
\textbf{Step 1. Compactness and factorization.}
\paraid{p-03d00ab877e8}
Since
$\SelectedMeasure{\ComparisonCarrier}{\ComparisonMetric}$
has total mass one,
for every
$\HilbertFunction\in\ContinuousFunctions(\AmbientCarrier,\ComplexNumbers)$
we have
\[
 \norm{\RestrictionOperator_\ComparisonCarrier\HilbertFunction}_2
 \leq\norm{\HilbertFunction}_\infty.
\]
For every
$\HilbertTestFunction\in\LebesgueSpace^2(\ComparisonCarrier,\SelectedMeasure{\ComparisonCarrier}{\ComparisonMetric};\ComplexNumbers)$
and
$\IntegrationPoint,\IntegrationPoint_0\in\AmbientCarrier$,
the Cauchy--Schwarz inequality yields
\[
 \norm{\KernelExtensionOperator_\ComparisonCarrier\HilbertTestFunction}_\infty
 \leq\DiameterBound\norm{\HilbertTestFunction}_2,
\]
and
\[
 |\KernelExtensionOperator_\ComparisonCarrier\HilbertTestFunction(\IntegrationPoint)
  -\KernelExtensionOperator_\ComparisonCarrier\HilbertTestFunction(\IntegrationPoint_0)|
 \leq\AmbientMetric(\IntegrationPoint,\IntegrationPoint_0)
       \norm{\HilbertTestFunction}_2.
\]
Thus
$\RestrictionOperator_\ComparisonCarrier$
and
$\KernelExtensionOperator_\ComparisonCarrier$
are bounded,
and the Arzel\`a--Ascoli theorem (\Cref{thm:arzela-ascoli}) implies that
$\KernelExtensionOperator_\ComparisonCarrier$
is compact.
Consequently,
$\CarrierAmbientOperator_\ComparisonCarrier$
is compact.
By their definitions,
\begin{equation}\label{eq:operator-factorization-1}
 \CarrierAmbientOperator _\ComparisonCarrier =\KernelExtensionOperator _\ComparisonCarrier \RestrictionOperator _\ComparisonCarrier ,
\end{equation}
and
\begin{equation}\label{eq:operator-factorization-2}
 \DistanceOperator{\ComparisonCarrier}{\ComparisonMetric}=\RestrictionOperator _\ComparisonCarrier \KernelExtensionOperator _\ComparisonCarrier ,
\end{equation}
and
\begin{equation}\label{eq:operator-factorization-3}
 \RestrictionOperator _\ComparisonCarrier  \CarrierAmbientOperator _\ComparisonCarrier =\DistanceOperator{\ComparisonCarrier}{\ComparisonMetric}\RestrictionOperator _\ComparisonCarrier .
\end{equation}

\paraid{p-55a3ce893fd0}
The intertwining identity is displayed in Figure~\ref{fig:operator-comparison}.
\begin{figure}[H]
\centering
\begin{tikzcd}[column sep=large,row sep=large]
 \ContinuousFunctions(\AmbientCarrier,\ComplexNumbers)
 \arrow[r,"\CarrierAmbientOperator_\ComparisonCarrier"]
 \arrow[d,"\RestrictionOperator_\ComparisonCarrier"'] &
 \ContinuousFunctions(\AmbientCarrier,\ComplexNumbers)
 \arrow[d,"\RestrictionOperator_\ComparisonCarrier"] \\
 \LebesgueSpace^2(\ComparisonCarrier,\SelectedMeasure{\ComparisonCarrier}{\ComparisonMetric};\ComplexNumbers)
 \arrow[r,"\DistanceOperator{\ComparisonCarrier}{\ComparisonMetric}"']
 \arrow[ur,"\KernelExtensionOperator_\ComparisonCarrier" description] &
 \LebesgueSpace^2(\ComparisonCarrier,\SelectedMeasure{\ComparisonCarrier}{\ComparisonMetric};\ComplexNumbers)
\end{tikzcd}
\caption{Restriction intertwines the ambient operator
$\CarrierAmbientOperator_\ComparisonCarrier$
and the distance operator
$\DistanceOperator{\ComparisonCarrier}{\ComparisonMetric}$.
The factorizations
$\CarrierAmbientOperator_\ComparisonCarrier=\KernelExtensionOperator_\ComparisonCarrier\RestrictionOperator_\ComparisonCarrier$
and
$\DistanceOperator{\ComparisonCarrier}{\ComparisonMetric}=\RestrictionOperator_\ComparisonCarrier\KernelExtensionOperator_\ComparisonCarrier$
explain the comparison of their nonzero spectral subspaces.}
\label{fig:operator-comparison}
\end{figure}


\textbf{Step 2. Generalized eigenspaces and restriction.}
\paraid{p-3b9fc77c0141}
Fix
$\Eigenvalue\in\ComplexNumbers\setminus\{0\}$
and an integer
$\RootOrder\geq1$.
By the self-adjoint spectral theorem (\Cref{prop:spectral-basics}),
\begin{equation}\label{eq:self-adjoint-power-kernel}
 \ker\bigl((\DistanceOperator{\ComparisonCarrier}{\ComparisonMetric}
               -\Eigenvalue\IdentityMap)^\RootOrder\bigr)
 =\ker(\DistanceOperator{\ComparisonCarrier}{\ComparisonMetric}
               -\Eigenvalue\IdentityMap).
\end{equation}
We prove the equality
\begin{equation}\label{eq:restricted-power-kernel}
 \RestrictionOperator_\ComparisonCarrier
 \left(\ker\bigl((\CarrierAmbientOperator_\ComparisonCarrier
                    -\Eigenvalue\IdentityMap)^\RootOrder\bigr)\right)
 =\ker(\DistanceOperator{\ComparisonCarrier}{\ComparisonMetric}
                    -\Eigenvalue\IdentityMap)
\end{equation}
by both inclusions.
Take
$\HilbertTestFunction\in
 \RestrictionOperator_\ComparisonCarrier
 \left(\ker\bigl((\CarrierAmbientOperator_\ComparisonCarrier
                 -\Eigenvalue\IdentityMap)^\RootOrder\bigr)\right)$.
Then there is
$\HilbertFunction\in\ker\bigl((\CarrierAmbientOperator_\ComparisonCarrier
 -\Eigenvalue\IdentityMap)^\RootOrder\bigr)$
with
$\HilbertTestFunction=\RestrictionOperator_\ComparisonCarrier\HilbertFunction$.
By \eqref{eq:operator-factorization-3},
\[
 (\DistanceOperator{\ComparisonCarrier}{\ComparisonMetric}-\Eigenvalue\IdentityMap)^\RootOrder
 \HilbertTestFunction
 =\RestrictionOperator_\ComparisonCarrier
  (\CarrierAmbientOperator_\ComparisonCarrier-\Eigenvalue\IdentityMap)^\RootOrder
  \HilbertFunction=0.
\]
Equation \eqref{eq:self-adjoint-power-kernel} therefore implies
$\HilbertTestFunction\in\ker(\DistanceOperator{\ComparisonCarrier}{\ComparisonMetric}
 -\Eigenvalue\IdentityMap)$.

\paraid{p-8a2730244eab}
Conversely,
take
$\HilbertTestFunction\in\ker(\DistanceOperator{\ComparisonCarrier}{\ComparisonMetric}
 -\Eigenvalue\IdentityMap)$
and put
$\HilbertFunction=\Eigenvalue^{-1}\KernelExtensionOperator_\ComparisonCarrier\HilbertTestFunction$.
Using \eqref{eq:operator-factorization-1} and \eqref{eq:operator-factorization-2},
we obtain
\begin{equation}\label{eq:eigenfunction-extension-1}
 \RestrictionOperator_\ComparisonCarrier\HilbertFunction
 =\Eigenvalue^{-1}\DistanceOperator{\ComparisonCarrier}{\ComparisonMetric}
    \HilbertTestFunction=\HilbertTestFunction,
\end{equation}
and
\begin{equation}\label{eq:eigenfunction-extension-2}
 \CarrierAmbientOperator_\ComparisonCarrier\HilbertFunction
 =\KernelExtensionOperator_\ComparisonCarrier\HilbertTestFunction
 =\Eigenvalue\HilbertFunction.
\end{equation}
Thus
$\HilbertFunction\in\ker\bigl((\CarrierAmbientOperator_\ComparisonCarrier
 -\Eigenvalue\IdentityMap)^\RootOrder\bigr)$,
and \eqref{eq:eigenfunction-extension-1} implies
$\HilbertTestFunction\in
 \RestrictionOperator_\ComparisonCarrier
 \left(\ker\bigl((\CarrierAmbientOperator_\ComparisonCarrier
                 -\Eigenvalue\IdentityMap)^\RootOrder\bigr)\right)$.
This proves \eqref{eq:restricted-power-kernel}.

\paraid{p-7ddba0188af2}
The restriction operator
$\RestrictionOperator_\ComparisonCarrier$
is injective on
$\ker\bigl((\CarrierAmbientOperator_\ComparisonCarrier-\Eigenvalue\IdentityMap)^\RootOrder\bigr)$.
To see this,
take
$\HilbertFunction$
in this kernel with
$\RestrictionOperator_\ComparisonCarrier\HilbertFunction=0$.
Then \eqref{eq:operator-factorization-1} implies
$\CarrierAmbientOperator_\ComparisonCarrier\HilbertFunction=0$,
so
\[
 0=(\CarrierAmbientOperator_\ComparisonCarrier-\Eigenvalue\IdentityMap)^\RootOrder
     \HilbertFunction=(-\Eigenvalue)^\RootOrder\HilbertFunction.
\]
Since
$\Eigenvalue\ne0$,
we have
$\HilbertFunction=0$.

\paraid{p-60e424cf6d26}
We next prove
\begin{equation}\label{eq:ambient-power-kernel}
 \ker\bigl((\CarrierAmbientOperator_\ComparisonCarrier-\Eigenvalue\IdentityMap)^\RootOrder\bigr)
 =\ker(\CarrierAmbientOperator_\ComparisonCarrier-\Eigenvalue\IdentityMap).
\end{equation}
Take
$\HilbertFunction\in\ker\bigl((\CarrierAmbientOperator_\ComparisonCarrier
 -\Eigenvalue\IdentityMap)^\RootOrder\bigr)$
and set
\[
 \HilbertTestFunction=\RestrictionOperator_\ComparisonCarrier\HilbertFunction,
\]
and
\[
 \LiftedEigenfunction=\Eigenvalue^{-1}\KernelExtensionOperator_\ComparisonCarrier\HilbertTestFunction.
\]
Equations \eqref{eq:restricted-power-kernel},
\eqref{eq:eigenfunction-extension-1} and \eqref{eq:eigenfunction-extension-2}
imply
\[
 \LiftedEigenfunction\in\ker(\CarrierAmbientOperator_\ComparisonCarrier-\Eigenvalue\IdentityMap),
\]
and
\[
 \RestrictionOperator_\ComparisonCarrier\LiftedEigenfunction
 =\RestrictionOperator_\ComparisonCarrier\HilbertFunction.
\]
Both
$\HilbertFunction$
and
$\LiftedEigenfunction$
belong to
\[
 \ker\bigl((\CarrierAmbientOperator_\ComparisonCarrier
 -\Eigenvalue\IdentityMap)^\RootOrder\bigr),
\]
the kernel on the left of \eqref{eq:ambient-power-kernel}.
Injectivity of restriction on this kernel implies
$\HilbertFunction=\LiftedEigenfunction$,
so
$\HilbertFunction\in\ker(\CarrierAmbientOperator_\ComparisonCarrier-\Eigenvalue\IdentityMap)$.
Conversely,
if
$\HilbertFunction\in\ker(\CarrierAmbientOperator_\ComparisonCarrier-\Eigenvalue\IdentityMap)$,
then
\[
 (\CarrierAmbientOperator_\ComparisonCarrier-\Eigenvalue\IdentityMap)^\RootOrder\HilbertFunction
 =(\CarrierAmbientOperator_\ComparisonCarrier-\Eigenvalue\IdentityMap)^{\RootOrder-1}0=0.
\]
Hence
$\HilbertFunction\in\ker\bigl((\CarrierAmbientOperator_\ComparisonCarrier
 -\Eigenvalue\IdentityMap)^\RootOrder\bigr)$.
This proves \eqref{eq:ambient-power-kernel}.
Taking the union over
$\RootOrder\geq1$
in \eqref{eq:self-adjoint-power-kernel} and \eqref{eq:ambient-power-kernel}
identifies the corresponding generalized eigenspaces.
In particular,
restriction induces the linear isomorphism
\[
 \RestrictionOperator_\ComparisonCarrier\colon
 \ker(\CarrierAmbientOperator_\ComparisonCarrier-\Eigenvalue\IdentityMap)
 \to\ker(\DistanceOperator{\ComparisonCarrier}{\ComparisonMetric}-\Eigenvalue\IdentityMap),
\]
Its inverse is the restricted map
\[
 \Eigenvalue^{-1}
 \left.\KernelExtensionOperator_\ComparisonCarrier\right|
 _{\ker(\DistanceOperator{\ComparisonCarrier}{\ComparisonMetric}-\Eigenvalue\IdentityMap)}\colon
 \ker(\DistanceOperator{\ComparisonCarrier}{\ComparisonMetric}-\Eigenvalue\IdentityMap)
 \to\ker(\CarrierAmbientOperator_\ComparisonCarrier-\Eigenvalue\IdentityMap).
\]

\textbf{Step 3. Equality of the nonzero spectra.}
\paraid{p-7bd77f5fa755}
Both operators are compact,
so the Riesz--Schauder theorem (\Cref{thm:compact-spectrum}) ensures
that every nonzero spectral value is an eigenvalue.
By \Cref{lem:distance-operator},
$\norm{\DistanceOperator{\ComparisonCarrier}{\ComparisonMetric}}
 \leq\diam\MetricSpace{\ComparisonCarrier}{\ComparisonMetric}
 \leq\DiameterBound$.
The standard spectral bound for bounded operators
\cite[Theorem~14.3]{Muscat2024}
implies
\[
 \spec(\DistanceOperator{\ComparisonCarrier}{\ComparisonMetric})
 \subset
 \{\Eigenvalue\in\ComplexNumbers\mid
   |\Eigenvalue|\leq
   \norm{\DistanceOperator{\ComparisonCarrier}{\ComparisonMetric}}\}.
\]
Since
$\DistanceOperator{\ComparisonCarrier}{\ComparisonMetric}$
is self-adjoint, its spectrum is real. Thus
\[
 \spec(\CarrierAmbientOperator _\ComparisonCarrier )\setminus\{0\}
 =\spec(\DistanceOperator{\ComparisonCarrier}{\ComparisonMetric})\setminus\{0\}
 \subset[-\DiameterBound,\DiameterBound].
\]
This proves the equality and the bound for the nonzero spectra,
and finishes the proof.
\end{proof}

\paraid{p-6a4c7f8bfb4d}
We first compare the ambient Riesz projection with the orthogonal projection
for a fixed compact carrier.
We then prove convergence of the ambient projections as the carrier varies.

\begin{proposition}\label{lem:projection-restriction}
\paraid{p-7625e731f8c2}
Let
$\ComparisonCarrier$
be a nonempty compact subset of a compact metric space
$\MetricSpace{\AmbientCarrier}{\AmbientMetric}$,
and put
$\ComparisonMetric=\AmbientMetric|_{\ComparisonCarrier\times\ComparisonCarrier}$.
Use the operators of \Cref{lem:ambient-eigenspaces}.
Fix
$\SpectralCutoff>0$
with
$\pm\SpectralCutoff\notin\spec(\DistanceOperator{\ComparisonCarrier}{\ComparisonMetric})$.
Let
$\AmbientProjection_\ComparisonCarrier$
be the Riesz projection of
$\CarrierAmbientOperator_\ComparisonCarrier$
associated with the eigenvalues of absolute value greater than
$\SpectralCutoff$,
as in Definition~\ref{def:riesz-projection}.
Then this projection preserves real functions,
and restriction induces the complex-linear isomorphism
\[
 \RestrictionOperator_\ComparisonCarrier\colon
 \range\AmbientProjection_\ComparisonCarrier
 \to\SpectralSubspace{\ComparisonCarrier}{\SpectralCutoff}
   +\ImaginaryUnit\SpectralSubspace{\ComparisonCarrier}{\SpectralCutoff}.
\]
Use
$\SpectralProjection{\ComparisonCarrier}{\SpectralCutoff}$
also for the complex-linear extension of the orthogonal projection.
For every
$\HilbertFunction\in\ContinuousFunctions(\AmbientCarrier,\ComplexNumbers)$
we have
\[
 (\AmbientProjection_\ComparisonCarrier\HilbertFunction)|_\ComparisonCarrier
 =\SpectralProjection{\ComparisonCarrier}{\SpectralCutoff}
   (\HilbertFunction|_\ComparisonCarrier).
\]
\end{proposition}

\begin{proof}
\paraid{p-8812db27373e}
By \Cref{lem:ambient-eigenspaces},
\begin{equation}\label{eq:fixed-carrier-nonzero-spectra}
 \spec(\CarrierAmbientOperator_\ComparisonCarrier)\setminus\{0\}
 =\spec(\DistanceOperator{\ComparisonCarrier}{\ComparisonMetric})\setminus\{0\}
 \subset\RealNumbers.
\end{equation}
Choose a bounded open set
$\SpectralRegion\subset\ComplexNumbers$
that is a union of two rectangles symmetric about the real axis and satisfies
$0\notin\overline{\SpectralRegion}$.
By \eqref{eq:fixed-carrier-nonzero-spectra}
and the hypothesis
$\pm\SpectralCutoff\notin\spec(\DistanceOperator{\ComparisonCarrier}{\ComparisonMetric})$,
we can choose
$\SpectralRegion$
so that for each operator
$\CompactOperator\in\{\CarrierAmbientOperator_\ComparisonCarrier,
\DistanceOperator{\ComparisonCarrier}{\ComparisonMetric}\}$,
\begin{equation}\label{eq:fixed-carrier-selected-spectrum}
 \SpectralRegion\cap\spec(\CompactOperator)
 =\{\Eigenvalue\in\spec(\CompactOperator)\mid|\Eigenvalue|>\SpectralCutoff\},
\end{equation}
and
\[
 \partial\SpectralRegion\cap\spec(\CompactOperator)=\emptyset.
\]
Orient
$\SpectralContour=\partial\SpectralRegion$
positively.
By Definition~\ref{def:riesz-projection}
and \Cref{lem:riesz-projection-properties},
we have
\[
 \AmbientProjection_\ComparisonCarrier
 =\RieszProjection{\CarrierAmbientOperator_\ComparisonCarrier}{\SpectralContour}.
\]
By \Cref{lem:riesz-range-generalized-eigenspace},
\begin{equation}\label{eq:fixed-carrier-riesz-range}
 \range\AmbientProjection_\ComparisonCarrier
 =\bigoplus_{\Eigenvalue\in\SpectralRegion\cap\spec(\CarrierAmbientOperator_\ComparisonCarrier)}
 \GeneralizedEigenspace{\CarrierAmbientOperator_\ComparisonCarrier}{\Eigenvalue}.
\end{equation}
For each
$\Eigenvalue\in\SpectralRegion\cap\spec(\CarrierAmbientOperator_\ComparisonCarrier)$,
\Cref{lem:ambient-eigenspaces} shows that the restriction operator
$\RestrictionOperator_\ComparisonCarrier$
maps
$\GeneralizedEigenspace{\CarrierAmbientOperator_\ComparisonCarrier}{\Eigenvalue}$
isomorphically onto
$\ker(\DistanceOperator{\ComparisonCarrier}{\ComparisonMetric}-\Eigenvalue\IdentityMap)$.
Taking the direct sum in \eqref{eq:fixed-carrier-riesz-range}
and using \eqref{eq:fixed-carrier-selected-spectrum},
we obtain the isomorphism
\[
 \RestrictionOperator_\ComparisonCarrier\colon
 \range\AmbientProjection_\ComparisonCarrier
 \to\bigoplus_{\substack{\Eigenvalue\in\spec(\DistanceOperator{\ComparisonCarrier}{\ComparisonMetric})\\
 |\Eigenvalue|>\SpectralCutoff}}
 \ker(\DistanceOperator{\ComparisonCarrier}{\ComparisonMetric}-\Eigenvalue\IdentityMap).
\]
By the definition of
$\SpectralSubspace{\ComparisonCarrier}{\SpectralCutoff}$,
this direct sum is
$\SpectralSubspace{\ComparisonCarrier}{\SpectralCutoff}
+\ImaginaryUnit\SpectralSubspace{\ComparisonCarrier}{\SpectralCutoff}$.
We prove the identity relating projection and restriction for
an arbitrary continuous function.
For every
$\ResolventParameter\in\SpectralContour$,
\eqref{eq:operator-factorization-3} implies
\[
 (\ResolventParameter\IdentityMap-\DistanceOperator{\ComparisonCarrier}{\ComparisonMetric})
 \RestrictionOperator_\ComparisonCarrier
 =\RestrictionOperator_\ComparisonCarrier
 (\ResolventParameter\IdentityMap-\CarrierAmbientOperator_\ComparisonCarrier).
\]
For
$\ResolventParameter\in\SpectralContour$,
both operators in parentheses are invertible.
Multiplying by their inverses on the left and on the right,
respectively,
we obtain
\begin{equation}\label{eq:resolvent-restriction}
 \RestrictionOperator_\ComparisonCarrier
 (\ResolventParameter\IdentityMap-\CarrierAmbientOperator_\ComparisonCarrier)^{-1}
 =(\ResolventParameter\IdentityMap-\DistanceOperator{\ComparisonCarrier}{\ComparisonMetric})^{-1}
 \RestrictionOperator_\ComparisonCarrier.
\end{equation}
Both sides are bounded maps from
$\ContinuousFunctions(\AmbientCarrier,\ComplexNumbers)$
to
$\LebesgueSpace^2(\ComparisonCarrier,\SelectedMeasure{\ComparisonCarrier}{\ComparisonMetric};\ComplexNumbers)$.
Both resolvents in \eqref{eq:resolvent-restriction} depend continuously on
$\ResolventParameter\in\SpectralContour$
in the operator norm.
The line integrals can therefore be computed as limits of Riemann sums.
Since
$\RestrictionOperator_\ComparisonCarrier$
is bounded and linear,
it commutes with these sums and their limits.
For every
$\HilbertFunction\in\ContinuousFunctions(\AmbientCarrier,\ComplexNumbers)$,
we consequently obtain
\begin{equation}\label{eq:restricted-contour-integral}
 \begin{aligned}
 \RestrictionOperator_\ComparisonCarrier\AmbientProjection_\ComparisonCarrier\HilbertFunction
 &=\RestrictionOperator_\ComparisonCarrier
   \left(\frac{1}{2\CircleConstant\ImaginaryUnit}
   \int_\SpectralContour
   (\ResolventParameter\IdentityMap-\CarrierAmbientOperator_\ComparisonCarrier)^{-1}
   \HilbertFunction\,\IntegrationDifferential\ResolventParameter\right)\\
 &=\frac{1}{2\CircleConstant\ImaginaryUnit}
   \int_\SpectralContour\RestrictionOperator_\ComparisonCarrier
   (\ResolventParameter\IdentityMap-\CarrierAmbientOperator_\ComparisonCarrier)^{-1}
   \HilbertFunction\,\IntegrationDifferential\ResolventParameter\\
 &=\frac{1}{2\CircleConstant\ImaginaryUnit}
   \int_\SpectralContour
   (\ResolventParameter\IdentityMap-\DistanceOperator{\ComparisonCarrier}{\ComparisonMetric})^{-1}
   \RestrictionOperator_\ComparisonCarrier\HilbertFunction
   \,\IntegrationDifferential\ResolventParameter.
 \end{aligned}
\end{equation}
In \eqref{eq:restricted-contour-integral},
the first equality is the definition of
$\AmbientProjection_\ComparisonCarrier$,
the second uses the bounded linearity of
$\RestrictionOperator_\ComparisonCarrier$,
and the third uses \eqref{eq:resolvent-restriction}.

\paraid{p-b3df5075d3de}
We next prove that the last line integral in
\eqref{eq:restricted-contour-integral} equals
$\SpectralProjection{\ComparisonCarrier}{\SpectralCutoff}
 (\RestrictionOperator_\ComparisonCarrier\HilbertFunction)$.
Here
$\SpectralProjection{\ComparisonCarrier}{\SpectralCutoff}$
also denotes the complex-linear extension of the orthogonal projection
defined in \eqref{eq:spectral-projection-definition}.
Let
$\HilbertTestFunction$
be an eigenvector of
$\DistanceOperator{\ComparisonCarrier}{\ComparisonMetric}$
with eigenvalue
$\Eigenvalue$.
For every
$\ResolventParameter\in\SpectralContour$,
\begin{equation}\label{eq:eigenvector-resolvent}
 (\ResolventParameter\IdentityMap-\DistanceOperator{\ComparisonCarrier}{\ComparisonMetric})^{-1}
 \HilbertTestFunction
 =\frac{1}{\ResolventParameter-\Eigenvalue}\HilbertTestFunction.
\end{equation}
Substituting \eqref{eq:eigenvector-resolvent} into the line integral
and using the winding numbers of
$\SpectralContour$,
we obtain
\begin{equation}\label{eq:eigenvector-contour-integral}
 \begin{aligned}
 &\frac{1}{2\CircleConstant\ImaginaryUnit}
 \int_\SpectralContour
 (\ResolventParameter\IdentityMap-\DistanceOperator{\ComparisonCarrier}{\ComparisonMetric})^{-1}
 \HilbertTestFunction\,\IntegrationDifferential\ResolventParameter\\
 &\qquad=
 \left(\frac{1}{2\CircleConstant\ImaginaryUnit}
 \int_\SpectralContour
 \frac{\IntegrationDifferential\ResolventParameter}{\ResolventParameter-\Eigenvalue}\right)
 \HilbertTestFunction
 =
 \begin{cases}
 \HilbertTestFunction,&|\Eigenvalue|>\SpectralCutoff,\\
 0,&|\Eigenvalue|<\SpectralCutoff.
 \end{cases}
 \end{aligned}
\end{equation}
The scalar line integral equals one when
$\Eigenvalue$
is inside
$\SpectralContour$
and zero when it is outside.
The second case includes
$\Eigenvalue=0$,
so the line integral vanishes on
$\ker\DistanceOperator{\ComparisonCarrier}{\ComparisonMetric}$.
Take the closure in the
$\LebesgueSpace^2$-norm.
By the self-adjoint spectral theorem (\Cref{prop:spectral-basics}),
\begin{equation}\label{eq:hilbert-eigenspace-density}
 \begin{aligned}
 &\overline{
   \ker\DistanceOperator{\ComparisonCarrier}{\ComparisonMetric}
   \mathbin{\oplus}
   \bigoplus_{\substack{
     \Eigenvalue\in\spec(\DistanceOperator{\ComparisonCarrier}{\ComparisonMetric})\\
     \Eigenvalue\ne0}}
   \ker(\DistanceOperator{\ComparisonCarrier}{\ComparisonMetric}-\Eigenvalue\IdentityMap)
   }\\
 &\qquad=
 \LebesgueSpace^2(\ComparisonCarrier,\SelectedMeasure{\ComparisonCarrier}{\ComparisonMetric};\ComplexNumbers).
 \end{aligned}
\end{equation}
By \eqref{eq:eigenvector-contour-integral},
\eqref{eq:spectral-projection-definition},
and orthogonality of eigenspaces for distinct eigenvalues,
the operator defined by the line integral and
$\SpectralProjection{\ComparisonCarrier}{\SpectralCutoff}$
agree on each summand in \eqref{eq:hilbert-eigenspace-density}.
By linearity and boundedness,
the density in \eqref{eq:hilbert-eigenspace-density} proves
\begin{equation}\label{eq:hilbert-riesz-projection}
 \frac{1}{2\CircleConstant\ImaginaryUnit}
 \int_\SpectralContour
 (\ResolventParameter\IdentityMap-\DistanceOperator{\ComparisonCarrier}{\ComparisonMetric})^{-1}
 \,\IntegrationDifferential\ResolventParameter
 =\SpectralProjection{\ComparisonCarrier}{\SpectralCutoff}
\end{equation}
as operators on
$\LebesgueSpace^2(\ComparisonCarrier,\SelectedMeasure{\ComparisonCarrier}{\ComparisonMetric};\ComplexNumbers)$.
Substituting \eqref{eq:hilbert-riesz-projection} into
\eqref{eq:restricted-contour-integral},
we obtain
\[
 \RestrictionOperator_\ComparisonCarrier\AmbientProjection_\ComparisonCarrier\HilbertFunction
 =\SpectralProjection{\ComparisonCarrier}{\SpectralCutoff}
    \RestrictionOperator_\ComparisonCarrier\HilbertFunction.
\]
Since
$\RestrictionOperator_\ComparisonCarrier\HilbertFunction=\HilbertFunction|_\ComparisonCarrier$,
we conclude that
\begin{equation}\label{eq:projection-restriction}
 (\AmbientProjection_\ComparisonCarrier\HilbertFunction)|_\ComparisonCarrier
 =\SpectralProjection{\ComparisonCarrier}{\SpectralCutoff}
   (\HilbertFunction|_\ComparisonCarrier).
\end{equation}
On the left we project in
$\ContinuousFunctions(\AmbientCarrier,\ComplexNumbers)$
and then restrict to
$\ComparisonCarrier$.
On the right we first restrict and then apply the projection on
$\LebesgueSpace^2(\ComparisonCarrier,\SelectedMeasure{\ComparisonCarrier}{\ComparisonMetric};\ComplexNumbers)$.
Both sides of \eqref{eq:projection-restriction} are continuous on
$\ComparisonCarrier$.
Since
$\SelectedMeasure{\ComparisonCarrier}{\ComparisonMetric}$
has full support,
their equality in
$\LebesgueSpace^2$
also implies pointwise equality on
$\ComparisonCarrier$.

\paraid{p-4e9939cad1e5}
Finally,
the distance kernels are real and
$\SpectralRegion$
is invariant under complex conjugation.
Conjugating the line integrals and reversing the orientation of
$\SpectralContour$
therefore yields
\[
 \overline{\AmbientProjection_\ComparisonCarrier\HilbertFunction}
 =\AmbientProjection_\ComparisonCarrier\overline{\HilbertFunction}.
\]
Thus
$\AmbientProjection_\ComparisonCarrier$
preserves real functions.
\end{proof}

\begin{proposition}\label{lem:spectral-projections}
\paraid{p-b10cee5c3ca1}
Let
$\BaseCarrier _\SequenceIndex ,\BaseCarrier $
be nonempty compact subsets of a compact metric space
$\MetricSpace{\AmbientCarrier }{\AmbientMetric }$,
with restricted metrics
\[
 \MetricSymbol_\SequenceIndex=\AmbientMetric|_{\BaseCarrier_\SequenceIndex\times\BaseCarrier_\SequenceIndex},
\]
and
\[
 \MetricSymbol=\AmbientMetric|_{\BaseCarrier\times\BaseCarrier},
\]
such that
\[
 \HausdorffDistance{\AmbientMetric }(\BaseCarrier _\SequenceIndex ,\BaseCarrier )\to0,
\]
and let $\iota_\SequenceIndex\colon\BaseCarrier_\SequenceIndex\hookrightarrow\AmbientCarrier$
and $\iota\colon\BaseCarrier\hookrightarrow\AmbientCarrier$ be the inclusions.
Define probabilities on $\AmbientCarrier$ by
$\bar\mu_\SequenceIndex=(\iota_\SequenceIndex)_*\SelectedMeasure{\BaseCarrier_\SequenceIndex}{\MetricSymbol_\SequenceIndex}$
and $\bar\mu=\iota_*\SelectedMeasure{\BaseCarrier}{\MetricSymbol}$.
Fix
$\SpectralCutoff >0$
with
$\SpectralCutoff ,-\SpectralCutoff \notin\spec(\DistanceOperator{\BaseCarrier}{\MetricSymbol})$.
Use the operators of \Cref{lem:ambient-eigenspaces} and set
\[
 \AmbientOperator_\SequenceIndex=\CarrierAmbientOperator_{\BaseCarrier_\SequenceIndex},
\]
and
\[
 \AmbientOperator=\CarrierAmbientOperator_\BaseCarrier.
\]
Then for all sufficiently large
$\SequenceIndex$,
\[
 \pm\SpectralCutoff
 \notin\spec(\DistanceOperator{\BaseCarrier_\SequenceIndex}{\MetricSymbol_\SequenceIndex}),
\]
and
\[
 \dim\SpectralSubspace{\BaseCarrier_\SequenceIndex}{\SpectralCutoff}
 =\dim\SpectralSubspace{\BaseCarrier}{\SpectralCutoff}
 =\SpectralRank.
\]
Moreover,
there are finite-rank Riesz projections
\[
 \AmbientProjection_\SequenceIndex,\AmbientProjection
 \colon\ContinuousFunctions(\AmbientCarrier,\ComplexNumbers)
 \to\ContinuousFunctions(\AmbientCarrier,\ComplexNumbers)
\]
associated with the eigenvalues of absolute value greater than
$\SpectralCutoff$
of
$\AmbientOperator_\SequenceIndex$
and
$\AmbientOperator$,
respectively.
Their ranges are the finite-dimensional sums
\[
 \range\AmbientProjection
 =\bigoplus_{\substack{\Eigenvalue\in\spec(\AmbientOperator)\\|\Eigenvalue|>\SpectralCutoff}}
 \GeneralizedEigenspace{\AmbientOperator}{\Eigenvalue},
\]
and, for every sufficiently large $\SequenceIndex$,
\[
 \range\AmbientProjection_\SequenceIndex
 =\bigoplus_{\substack{\Eigenvalue\in\spec(\AmbientOperator_\SequenceIndex)\\|\Eigenvalue|>\SpectralCutoff}}
 \GeneralizedEigenspace{\AmbientOperator_\SequenceIndex}{\Eigenvalue}.
\]
Both projections preserve real functions and satisfy,
for every
$\HilbertFunction\in\ContinuousFunctions(\AmbientCarrier,\ComplexNumbers)$,
\[
 \norm{\AmbientProjection_\SequenceIndex\HilbertFunction
       -\AmbientProjection\HilbertFunction}_\infty\to0.
\]
\end{proposition}

\begin{proof}
\textbf{Step 1. Strong convergence and collective compactness.}
\paraid{p-424dfa7891fa}
By \Cref{thm:invariant-measures}\ref{item:invariant-measures-continuity} applied to the inclusions
$\iota_\SequenceIndex$
and
$\iota$,
we have
$\bar\mu_\SequenceIndex\to\bar\mu$
weakly.
For every
$\SequenceIndex\in\NonnegativeIntegers$,
$\HilbertFunction\in\ContinuousFunctions(\AmbientCarrier,\ComplexNumbers)$,
and
$\IntegrationPoint\in\AmbientCarrier$,
the definitions imply
\[
\AmbientOperator _\SequenceIndex \HilbertFunction (\IntegrationPoint )=\int_\AmbientCarrier  \AmbientMetric (\IntegrationPoint ,\SecondVector )\HilbertFunction (\SecondVector )\,\IntegrationDifferential \bar\mu_\SequenceIndex(\SecondVector ),
\]
and
\[
\AmbientOperator \HilbertFunction (\IntegrationPoint )=\int_\AmbientCarrier  \AmbientMetric (\IntegrationPoint ,\SecondVector )\HilbertFunction (\SecondVector )\,\IntegrationDifferential \bar\mu(\SecondVector ).
\]
For each fixed
$\HilbertFunction $,
the family
$\{\AmbientMetric (\IntegrationPoint ,\cdot)\HilbertFunction (\cdot)\mid \IntegrationPoint \in \AmbientCarrier \}$
is compact in
$\ContinuousFunctions(\AmbientCarrier ,\ComplexNumbers)$.
By \Cref{lem:parameter-integrals} with parameter
$\IntegrationPoint\in\AmbientCarrier$,
\begin{equation}\label{eq:ambient-strong-convergence}
 \norm{\AmbientOperator _\SequenceIndex \HilbertFunction -\AmbientOperator \HilbertFunction }_\infty\to0.
\end{equation}
Thus
$\AmbientOperator_\SequenceIndex\to\AmbientOperator$
strongly.
Writing
$\DiameterBound =\diam\MetricSpace{\AmbientCarrier }{\AmbientMetric }$,
for every
$\SequenceIndex\in\NonnegativeIntegers$,
$\HilbertFunction\in\ContinuousFunctions(\AmbientCarrier,\ComplexNumbers)$
with
$\norm{\HilbertFunction}_\infty\leq1$,
and
$\IntegrationPoint,\IntegrationPoint_0\in\AmbientCarrier$,
we obtain
\begin{equation}\label{eq:ambient-operator-bounds-1}
 \norm{\AmbientOperator _\SequenceIndex \HilbertFunction }_\infty\leq \DiameterBound ,
\end{equation}
and
\begin{equation}\label{eq:ambient-operator-bounds-2}
 |\AmbientOperator _\SequenceIndex \HilbertFunction (\IntegrationPoint )-\AmbientOperator _\SequenceIndex \HilbertFunction (\IntegrationPoint _0)|\leq \AmbientMetric (\IntegrationPoint ,\IntegrationPoint _0).
\end{equation}
Hence
$\bigcup_\SequenceIndex  \AmbientOperator _\SequenceIndex \{\HilbertFunction \mid\norm{\HilbertFunction }_\infty\leq1\}$
has compact closure by the Arzel\`a--Ascoli theorem (\Cref{thm:arzela-ascoli}).
Thus
$\{\AmbientOperator_\SequenceIndex\}_{\SequenceIndex\in\NonnegativeIntegers}$
is collectively compact.
Since
$\AmbientOperator$
is compact by \Cref{lem:ambient-eigenspaces},
the family
$\{\AmbientOperator_\SequenceIndex-\AmbientOperator\}_{\SequenceIndex\in\NonnegativeIntegers}$
is also collectively compact.

\textbf{Step 2. A common boundary and spectral projections.}
\paraid{p-8737a7c3af89}
By \Cref{lem:ambient-eigenspaces} and the hypothesis on
$\SpectralCutoff$,
we have
\[
 \spec(\AmbientOperator)
 \subset\ComplexNumbers\setminus\{-\SpectralCutoff,\SpectralCutoff\}.
\]
The set
$\ComplexNumbers\setminus\{-\SpectralCutoff,\SpectralCutoff\}$
is open.
The strong convergence in \eqref{eq:ambient-strong-convergence}
and the collective compactness of
$\{\AmbientOperator_\SequenceIndex-\AmbientOperator\}_{\SequenceIndex\in\NonnegativeIntegers}$
verify the hypotheses on operators in the statement of
\Cref{thm:spectral-inclusion}.
Applying that theorem with this open set
$\ComplexNumbers\setminus\{-\SpectralCutoff,\SpectralCutoff\}$,
we obtain an index
$\SequenceIndex_0\in\NonnegativeIntegers$
such that for every
$\SequenceIndex\geq\SequenceIndex_0$
we have
\begin{equation}\label{eq:ambient-cutoff-exclusion}
 \{-\SpectralCutoff,\SpectralCutoff\}
 \cap\bigl(\spec(\AmbientOperator)\cup\spec(\AmbientOperator_\SequenceIndex)\bigr)
 =\emptyset.
\end{equation}
\Cref{lem:ambient-eigenspaces} then implies
\[
 \pm\SpectralCutoff
 \notin\spec(\DistanceOperator{\BaseCarrier_\SequenceIndex}{\MetricSymbol_\SequenceIndex}).
\]
This proves the claimed exclusion of
$\pm\SpectralCutoff$.

\paraid{p-4c4bdee6941a}
We next choose a bounded open set whose boundary encloses the eigenvalues of absolute
value greater than
$\SpectralCutoff$
for both
$\AmbientOperator$
and
$\AmbientOperator_\SequenceIndex$.
Choose
$\SpectralOuterBound >\max\{\DiameterBound ,\SpectralCutoff \}+1$
and let
$\SpectralRegion $
be the union of the two rectangles
\[
 \begin{gathered}
 \{\ResolventParameter \in\ComplexNumbers\mid \SpectralCutoff <\RealPart \ResolventParameter <\SpectralOuterBound ,\ |\ImaginaryPart \ResolventParameter |<1\},\\
 \{\ResolventParameter \in\ComplexNumbers\mid-\SpectralOuterBound <\RealPart \ResolventParameter <-\SpectralCutoff ,\ |\ImaginaryPart \ResolventParameter |<1\}.
 \end{gathered}
\]
Write
$\SpectralContour =\partial\SpectralRegion $
with positive orientation.
For every
$\SequenceIndex\in\NonnegativeIntegers$,
\Cref{lem:ambient-eigenspaces} implies
\begin{equation}\label{eq:ambient-spectral-bound}
 \spec(\AmbientOperator)\cup\spec(\AmbientOperator_\SequenceIndex)
 \subset[-\DiameterBound,\DiameterBound].
\end{equation}
The choice of
$\SpectralOuterBound$
implies
\begin{equation}\label{eq:contour-real-intersection}
 \SpectralContour\cap[-\DiameterBound,\DiameterBound]
 \subset\{-\SpectralCutoff,\SpectralCutoff\}.
\end{equation}
Combining \eqref{eq:contour-real-intersection} with \eqref{eq:ambient-cutoff-exclusion}
and \eqref{eq:ambient-spectral-bound},
we obtain for every
$\SequenceIndex\geq\SequenceIndex_0$
\begin{equation}\label{eq:common-spectral-contour}
 \SpectralContour\cap
 \bigl(\spec(\AmbientOperator)\cup\spec(\AmbientOperator_\SequenceIndex)\bigr)
 =\emptyset.
\end{equation}
Moreover,
\begin{equation}\label{eq:spectral-region-selection-1}
 \spec(\AmbientOperator)\cap\SpectralRegion
 =\{\Eigenvalue\in\spec(\AmbientOperator)\mid|\Eigenvalue|>\SpectralCutoff\},
\end{equation}
and
\begin{equation}\label{eq:spectral-region-selection-2}
 \spec(\AmbientOperator_\SequenceIndex)\cap\SpectralRegion
 =\{\Eigenvalue\in\spec(\AmbientOperator_\SequenceIndex)\mid|\Eigenvalue|>\SpectralCutoff\},
\end{equation}
and
\begin{equation}\label{eq:spectral-region-selection-3}
 0\notin\overline\SpectralRegion.
\end{equation}
By \eqref{eq:common-spectral-contour},
we use the same boundary for both operators
for all sufficiently large indices.
Equations \eqref{eq:spectral-region-selection-1} and \eqref{eq:spectral-region-selection-2}
show that this boundary encloses precisely the eigenvalues with absolute value
greater than
$\SpectralCutoff$.
Figure~\ref{fig:spectral-contour} shows the two components of this boundary.
\begin{figure}[H]
\centering
\begin{tikzpicture}[>=Stealth,x=1cm,y=1cm,font=\small]
\fill[black!6] (-4,-.8) rectangle (-1,.8);
\fill[black!6] (1,-.8) rectangle (4,.8);
\draw[->] (-4.5,0)--(4.7,0) node[right] {$\RealPart\ResolventParameter$};
\draw[->] (0,-1.05)--(0,1.3) node[above] {$\ImaginaryPart\ResolventParameter$};
\draw[thick] (-4,-.8) rectangle (-1,.8);
\draw[thick] (1,-.8) rectangle (4,.8);
\draw[->,thick] (-3.2,-.8)--(-2.5,-.8);
\draw[->,thick] (2,-.8)--(2.7,-.8);
\foreach \x in {-3.5,-2.8,-1.7,1.6,2.5,3.4} \fill (\x,0) circle (1.5pt);
\foreach \x in {-.55,-.3,.25,.5} \fill[black!45] (\x,0) circle (1.2pt);
\node[below left] at (0,0) {$0$};
\node[below] at (-4,-.85) {$-\SpectralOuterBound$};
\node[below] at (-1,-.85) {$-\SpectralCutoff$};
\node[below] at (1,-.85) {$\SpectralCutoff$};
\node[below] at (4,-.85) {$\SpectralOuterBound$};
\node[above] at (-2.5,.85) {$\SpectralContour$};
\node[above] at (2.5,.85) {$\SpectralContour$};
\end{tikzpicture}
\caption{The positively oriented boundary
$\SpectralContour=\partial\SpectralRegion$
encloses the retained positive and negative spectral values.
The dots are schematic.
The gaps at
$\pm\SpectralCutoff$
allow the same boundary to be used for
$\AmbientOperator_\SequenceIndex$
for all sufficiently large
$\SequenceIndex$.}
\label{fig:spectral-contour}
\end{figure}


\paraid{p-341c679243de}
Using Definition~\ref{def:riesz-projection},
we define the projections
$\AmbientProjection,\AmbientProjection_\SequenceIndex
 \colon\ContinuousFunctions(\AmbientCarrier,\ComplexNumbers)
 \to\ContinuousFunctions(\AmbientCarrier,\ComplexNumbers)$
by
\[
 \AmbientProjection=\RieszProjection{\AmbientOperator}{\SpectralContour},
\]
and
\[
 \AmbientProjection_\SequenceIndex
 =\RieszProjection{\AmbientOperator_\SequenceIndex}{\SpectralContour}.
\]
For
$\HilbertFunction\in\ContinuousFunctions(\AmbientCarrier,\ComplexNumbers)$
and each fixed
$\IntegrationPoint\in\AmbientCarrier$,
we have
\begin{equation}\label{eq:riesz-pointwise-integral}
 (\AmbientProjection\HilbertFunction)(\IntegrationPoint)
 =\frac{1}{2\CircleConstant\ImaginaryUnit}
 \int_\SpectralContour
 \bigl[(\ResolventParameter\IdentityMap-\AmbientOperator)^{-1}
       \HilbertFunction\bigr](\IntegrationPoint)
 \,\IntegrationDifferential\ResolventParameter.
\end{equation}
The integration variable
$\ResolventParameter$
ranges over the boundary
$\SpectralContour$,
while the point
$\IntegrationPoint\in\AmbientCarrier$
is fixed.
Evaluation at
$\IntegrationPoint$
is a bounded linear functional on
$\ContinuousFunctions(\AmbientCarrier,\ComplexNumbers)$,
so it commutes with the line integral.
Equation \eqref{eq:riesz-pointwise-integral} also holds for
$\AmbientProjection_\SequenceIndex$
with
$\AmbientOperator_\SequenceIndex$
in place of
$\AmbientOperator$.

\textbf{Step 3. Convergence and dimensions of the ranges.}
\paraid{p-a982b8ced6ab}
By \Cref{lem:riesz-range-generalized-eigenspace},
we obtain
\begin{equation}\label{eq:ambient-projection-range}
 \range\AmbientProjection
 =\bigoplus_{\Eigenvalue\in\spec(\AmbientOperator)\cap\SpectralRegion}
   \GeneralizedEigenspace{\AmbientOperator}{\Eigenvalue}.
\end{equation}
For all sufficiently large
$\SequenceIndex$,
we also have
\begin{equation}\label{eq:approximating-projection-range}
 \range\AmbientProjection_\SequenceIndex
 =\bigoplus_{\Eigenvalue\in\spec(\AmbientOperator_\SequenceIndex)\cap\SpectralRegion}
   \GeneralizedEigenspace{\AmbientOperator_\SequenceIndex}{\Eigenvalue}.
\end{equation}
The sums are finite by \Cref{thm:compact-spectrum},
since
$0\notin\overline\SpectralRegion$.
By \Cref{lem:ambient-eigenspaces},
for every
$\Eigenvalue\in\spec(\AmbientOperator)\cap\SpectralRegion$,
\begin{equation}\label{eq:projection-range-eigenspace}
 \GeneralizedEigenspace{\AmbientOperator}{\Eigenvalue}
 =\ker(\AmbientOperator-\Eigenvalue\IdentityMap).
\end{equation}
For every sufficiently large
$\SequenceIndex$
and every
$\Eigenvalue\in\spec(\AmbientOperator_\SequenceIndex)\cap\SpectralRegion$,
the same proposition implies
\begin{equation}\label{eq:approximating-range-eigenspace}
 \GeneralizedEigenspace{\AmbientOperator_\SequenceIndex}{\Eigenvalue}
 =\ker(\AmbientOperator_\SequenceIndex-\Eigenvalue\IdentityMap).
\end{equation}
We can now apply \Cref{thm:riesz-convergence}
on the fixed Banach space
$\ContinuousFunctions(\AmbientCarrier,\ComplexNumbers)$.
The hypotheses on operators in the statement of
\Cref{thm:riesz-convergence}
follow from \eqref{eq:ambient-strong-convergence}
and the collective compactness of
$\{\AmbientOperator_\SequenceIndex-\AmbientOperator\}_{\SequenceIndex\in\NonnegativeIntegers}$.
Equations \eqref{eq:common-spectral-contour},
\eqref{eq:spectral-region-selection-1},
\eqref{eq:spectral-region-selection-2} and \eqref{eq:spectral-region-selection-3}
verify its hypotheses on
$\SpectralContour$.
For all sufficiently large
$\SequenceIndex$,
we obtain
\begin{equation}\label{eq:spectral-projection-convergence-1}
 \dim\range \AmbientProjection _\SequenceIndex =\dim\range \AmbientProjection .
\end{equation}
For every
$\HilbertFunction\in\ContinuousFunctions(\AmbientCarrier,\ComplexNumbers)$,
we also obtain
\begin{equation}\label{eq:spectral-projection-convergence-2}
 \norm{\AmbientProjection _\SequenceIndex \HilbertFunction -\AmbientProjection \HilbertFunction }_\infty\to0.
\end{equation}
The dimensions are finite because
$0\notin\overline\SpectralRegion $.
The convergence in \eqref{eq:spectral-projection-convergence-2} means
uniform convergence on
$\AmbientCarrier$
for each fixed function
$\HilbertFunction\in\ContinuousFunctions(\AmbientCarrier,\ComplexNumbers)$.

\paraid{p-60dd00c4dbee}
We now compare the projections on
$\ContinuousFunctions(\AmbientCarrier,\ComplexNumbers)$
with the spectral subspaces of the distance operators.
For
$\ComparisonCarrier=\BaseCarrier$
or
$\ComparisonCarrier=\BaseCarrier_\SequenceIndex$
with
$\SequenceIndex$
sufficiently large,
write
$\AmbientProjection_\BaseCarrier=\AmbientProjection$
and
$\AmbientProjection_{\BaseCarrier_\SequenceIndex}=\AmbientProjection_\SequenceIndex$.
By \eqref{eq:spectral-region-selection-1} and
\eqref{eq:spectral-region-selection-2},
$\SpectralRegion$
contains exactly those $\Eigenvalue\in\spec(\AmbientOperator)$ with
$|\Eigenvalue|>\SpectralCutoff$,
and exactly those $\Eigenvalue\in\spec(\AmbientOperator_\SequenceIndex)$ with
$|\Eigenvalue|>\SpectralCutoff$.
\Cref{lem:projection-restriction} induces the complex-linear isomorphism
\[
 \RestrictionOperator_\ComparisonCarrier\colon
 \range\AmbientProjection_\ComparisonCarrier
 \to
 \SpectralSubspace{\ComparisonCarrier}{\SpectralCutoff}
 +\ImaginaryUnit\SpectralSubspace{\ComparisonCarrier}{\SpectralCutoff}.
\]
Here the space on the right is the complexification of the real space
$\SpectralSubspace{\ComparisonCarrier}{\SpectralCutoff}$.
Its complex dimension equals
$\dim_\RealNumbers\SpectralSubspace{\ComparisonCarrier}{\SpectralCutoff}$.
The dimension equality in \eqref{eq:spectral-projection-convergence-1}
therefore yields
\begin{equation}\label{eq:spectral-dimension-identification}
 \begin{aligned}
 \dim_\RealNumbers\SpectralSubspace{\BaseCarrier_\SequenceIndex}{\SpectralCutoff}
 &=\dim_\ComplexNumbers\range\AmbientProjection_\SequenceIndex
 =\dim_\ComplexNumbers\range\AmbientProjection\\
 &=\dim_\RealNumbers\SpectralSubspace{\BaseCarrier}{\SpectralCutoff}
 =\SpectralRank.
 \end{aligned}
\end{equation}
\paraid{p-7a7a27ecec95}
Finally, applying \Cref{lem:projection-restriction} to the limit carrier
and to each sufficiently late carrier $\BaseCarrier_\SequenceIndex$ shows that
$\AmbientProjection$ and $\AmbientProjection_\SequenceIndex$ preserve real
functions.
\end{proof}

\begin{theorem}\label{lem:spectral-continuity}
\paraid{p-56449b80711e}
Let
$\BaseCarrier _\SequenceIndex ,\BaseCarrier $
be compact subsets of a compact metric space
$\MetricSpace{\AmbientCarrier }{\AmbientMetric }$,
with restricted metrics
\[
 \MetricSymbol_\SequenceIndex=\AmbientMetric|_{\BaseCarrier_\SequenceIndex\times\BaseCarrier_\SequenceIndex},
\]
and
\[
 \MetricSymbol=\AmbientMetric|_{\BaseCarrier\times\BaseCarrier},
\]
such that
\[
 \HausdorffDistance{\AmbientMetric }(\BaseCarrier _\SequenceIndex ,\BaseCarrier )\to0,
\]
and let $\iota_\SequenceIndex\colon\BaseCarrier_\SequenceIndex\hookrightarrow\AmbientCarrier$
and $\iota\colon\BaseCarrier\hookrightarrow\AmbientCarrier$ be the inclusions.
Define probabilities on $\AmbientCarrier$ by
$\bar\mu_\SequenceIndex=(\iota_\SequenceIndex)_*\SelectedMeasure{\BaseCarrier_\SequenceIndex}{\MetricSymbol_\SequenceIndex}$
and $\bar\mu=\iota_*\SelectedMeasure{\BaseCarrier}{\MetricSymbol}$.
Fix
$\SpectralCutoff >0$
with
$\SpectralCutoff ,-\SpectralCutoff \notin\spec(\DistanceOperator{\BaseCarrier}{\MetricSymbol})$.
Put
$\SpectralRank=\dim\SpectralSubspace{\BaseCarrier}{\SpectralCutoff}$.
There exists $\SequenceIndex_0\in\NonnegativeIntegers$ such that, for every
$\SequenceIndex\geq\SequenceIndex_0$,
\begin{equation}\label{eq:stable-spectral-endpoints}
 \SpectralCutoff,-\SpectralCutoff
 \notin\spec(\DistanceOperator{\BaseCarrier_\SequenceIndex}{\MetricSymbol_\SequenceIndex}),
\end{equation}
and
\begin{equation}\label{eq:stable-spectral-rank}
 \dim\SpectralSubspace{\BaseCarrier_\SequenceIndex}{\SpectralCutoff}
 =\SpectralRank.
\end{equation}
If
$\SpectralRank>0$,
fix any real orthonormal basis
$\{\Eigenfunction_\BasisIndex\}_{0\leq\BasisIndex<\SpectralRank}$
of
$\SpectralSubspace{\BaseCarrier}{\SpectralCutoff}$.
For the fixed basis
$\{\Eigenfunction_\BasisIndex\}_{0\leq\BasisIndex<\SpectralRank}$
and the sequence
$\{(\BaseCarrier_\SequenceIndex,\bar\mu_\SequenceIndex)\}_{\SequenceIndex\in\NonnegativeIntegers}$,
increase
$\SequenceIndex_0$
if necessary.
Then there exist families of functions
\[
 \{\Eigenfunction_{\SequenceIndex,\BasisIndex}
   \colon\BaseCarrier_\SequenceIndex\to\RealNumbers\}_{\SequenceIndex\geq\SequenceIndex_0,\,0\leq\BasisIndex<\SpectralRank},
\]
\[
 \{\BasisExtension_{\SequenceIndex,\BasisIndex}
   \colon\AmbientCarrier\to\RealNumbers\}_{\SequenceIndex\geq\SequenceIndex_0,\,0\leq\BasisIndex<\SpectralRank},
\]
and
\[
 \{\BasisExtension_\BasisIndex
   \colon\AmbientCarrier\to\RealNumbers\}_{0\leq\BasisIndex<\SpectralRank}
\]
with the following properties.
\begin{enumerate}[label=\textup{(B\arabic*)},ref=\textup{(B\arabic*)}]
\item\label{item:spectral-basis-orthonormality}
For every
$\SequenceIndex\geq\SequenceIndex_0$,
the family
$\{\Eigenfunction_{\SequenceIndex,\BasisIndex}\}_{0\leq\BasisIndex<\SpectralRank}$
is a real orthonormal basis of
$\SpectralSubspace{\BaseCarrier_\SequenceIndex}{\SpectralCutoff}$.
\item\label{item:spectral-basis-extensions}
For every
$\SequenceIndex\geq\SequenceIndex_0$
and
$0\leq\BasisIndex<\SpectralRank$,
the functions
$\BasisExtension_{\SequenceIndex,\BasisIndex},\BasisExtension_\BasisIndex
 \colon\AmbientCarrier\to\RealNumbers$
are continuous and satisfy
\[
 \BasisExtension_{\SequenceIndex,\BasisIndex}|_{\BaseCarrier_\SequenceIndex}
 =\Eigenfunction_{\SequenceIndex,\BasisIndex},
\]
and
\[
 \BasisExtension_\BasisIndex|_\BaseCarrier=\Eigenfunction_\BasisIndex.
\]
\item\label{item:spectral-basis-uniform-convergence}
The extensions also satisfy
\begin{equation}\label{eq:aligned-basis-extensions}
 \max_{0\leq\BasisIndex<\SpectralRank}
 \norm{\BasisExtension_{\SequenceIndex,\BasisIndex}-\BasisExtension_\BasisIndex}_\infty
 \longrightarrow0.
\end{equation}
\item\label{item:spectral-basis-correspondence-convergence}
For every sequence of correspondences
$\{\Correspondence_\SequenceIndex\subset\BaseCarrier_\SequenceIndex\times\BaseCarrier\}_{\SequenceIndex\geq\SequenceIndex_0}$
satisfying
\[
 \sup_{(\BasePoint_\SequenceIndex,\BasePoint)\in\Correspondence_\SequenceIndex}
 \AmbientMetric(\BasePoint_\SequenceIndex,\BasePoint)\longrightarrow0,
\]
the bases
$\{\Eigenfunction_{\SequenceIndex,\BasisIndex}\}_{0\leq\BasisIndex<\SpectralRank}$
and
$\{\Eigenfunction_\BasisIndex\}_{0\leq\BasisIndex<\SpectralRank}$
satisfy
\begin{equation}\label{eq:aligned-bases}
 \max_{0\leq\BasisIndex<\SpectralRank}
 \sup_{(\BasePoint_\SequenceIndex,\BasePoint)\in\Correspondence_\SequenceIndex}
 |\Eigenfunction_{\SequenceIndex,\BasisIndex}(\BasePoint_\SequenceIndex)
       -\Eigenfunction_\BasisIndex(\BasePoint)|\longrightarrow0.
\end{equation}
\end{enumerate}
\end{theorem}

\begin{proof}
\paraid{p-dc3fbce63b61}
By \Cref{thm:invariant-measures}\ref{item:invariant-measures-continuity} applied to the inclusions
$\iota_\SequenceIndex$
and
$\iota$,
we have
$\bar\mu_\SequenceIndex\to\bar\mu$
weakly.
By \Cref{lem:spectral-projections},
we obtain \eqref{eq:stable-spectral-endpoints}
and \eqref{eq:stable-spectral-rank}.
Let
$\AmbientProjection_\SequenceIndex,\AmbientProjection$
be the projections constructed in that proposition.
If
$\SpectralRank=0$,
the stable spectral subspaces are zero and no basis or extension data are required.
Assume that
$\SpectralRank>0$.
To construct the real orthonormal bases
$\{\Eigenfunction_{\SequenceIndex,\BasisIndex}\}_{0\leq\BasisIndex<\SpectralRank}$
of
$\SpectralSubspace{\BaseCarrier_\SequenceIndex}{\SpectralCutoff}$
in \ref{item:spectral-basis-orthonormality},
we first project continuous extensions of the prescribed basis
$\{\Eigenfunction_\BasisIndex\}_{0\leq\BasisIndex<\SpectralRank}$.
The resulting functions need not be orthonormal for the varying measures
$\SelectedMeasure{\BaseCarrier_\SequenceIndex}{\MetricSymbol_\SequenceIndex}$.
We use the Gram matrix to normalize these functions without changing their limit.
We then verify \ref{item:spectral-basis-extensions}
and \ref{item:spectral-basis-uniform-convergence}
for the normalized extensions and deduce
\ref{item:spectral-basis-correspondence-convergence}
from uniform continuity.

\ProofStep{Projected extensions and convergence of the Gram matrices.}\label{step:spectral-basis-projection}
\paraid{p-b8f758726161}
For every integer
$0\leq\BasisIndex<\SpectralRank$,
use the Tietze--Urysohn theorem (\Cref{thm:tietze}) to choose a continuous real extension
$\widetilde\Eigenfunction_\BasisIndex\colon\AmbientCarrier\to\RealNumbers$
of the prescribed function
$\Eigenfunction_\BasisIndex\colon\BaseCarrier\to\RealNumbers$.
For every sufficiently large
$\SequenceIndex\in\NonnegativeIntegers$
and every integer
$0\leq\BasisIndex<\SpectralRank$,
put
\[
 \ProjectedBasisFunction _{\SequenceIndex ,\BasisIndex }=\AmbientProjection _\SequenceIndex \widetilde \Eigenfunction _\BasisIndex .
\]
For every integer
$0\leq\BasisIndex<\SpectralRank$,
put
\[
 \ProjectedBasisFunction _\BasisIndex =\AmbientProjection \widetilde \Eigenfunction _\BasisIndex .
\]
Applying \eqref{eq:spectral-projection-convergence-2} to each fixed function
$\widetilde\Eigenfunction_\BasisIndex$,
we obtain
$\norm{\ProjectedBasisFunction _{\SequenceIndex ,\BasisIndex }-\ProjectedBasisFunction _\BasisIndex }_\infty\to0$
for every
$0\leq\BasisIndex<\SpectralRank$.
Moreover,
\eqref{eq:projection-restriction} implies
\[
 \RestrictionOperator_\BaseCarrier\ProjectedBasisFunction_\BasisIndex
 =\SpectralProjection{\BaseCarrier}{\SpectralCutoff}
   \RestrictionOperator_\BaseCarrier\widetilde\Eigenfunction_\BasisIndex
 =\SpectralProjection{\BaseCarrier}{\SpectralCutoff}\Eigenfunction_\BasisIndex
 =\Eigenfunction_\BasisIndex.
\]
Thus
$\ProjectedBasisFunction_\BasisIndex|_\BaseCarrier=\Eigenfunction_\BasisIndex$.
Equation \eqref{eq:projection-restriction} on
$\BaseCarrier_\SequenceIndex$
shows that
\[
 \ProjectedBasisFunction_{\SequenceIndex,\BasisIndex}|_{\BaseCarrier_\SequenceIndex}
 =\SpectralProjection{\BaseCarrier_\SequenceIndex}{\SpectralCutoff}
   (\widetilde\Eigenfunction_\BasisIndex|_{\BaseCarrier_\SequenceIndex})
 \in\SpectralSubspace{\BaseCarrier_\SequenceIndex}{\SpectralCutoff}.
\]
Using
$\IntegrationPoint\in\AmbientCarrier$
as the integration variable,
define the real Gram matrix
\[
 \GramMatrix _\SequenceIndex =\left(\int_\AmbientCarrier  \ProjectedBasisFunction _{\SequenceIndex ,\BasisIndex }(\IntegrationPoint)\ProjectedBasisFunction _{\SequenceIndex ,\SecondBasisIndex }(\IntegrationPoint)\,\IntegrationDifferential \bar\mu_\SequenceIndex(\IntegrationPoint)
                                     \right)_{0\leq \BasisIndex ,\SecondBasisIndex<\SpectralRank }.
\]
For every pair of integers
$0\leq\BasisIndex,\SecondBasisIndex<\SpectralRank$,
the products
$\ProjectedBasisFunction_{\SequenceIndex,\BasisIndex}
 \ProjectedBasisFunction_{\SequenceIndex,\SecondBasisIndex}$
converge uniformly on
$\AmbientCarrier$
to
$\ProjectedBasisFunction_\BasisIndex\ProjectedBasisFunction_\SecondBasisIndex$.
Since
$\ProjectedBasisFunction_\BasisIndex|_\BaseCarrier=\Eigenfunction_\BasisIndex$
and the prescribed basis is orthonormal,
\[
 \int_\AmbientCarrier
 \ProjectedBasisFunction_\BasisIndex(\IntegrationPoint)\ProjectedBasisFunction_\SecondBasisIndex(\IntegrationPoint)
 \,\IntegrationDifferential\bar\mu(\IntegrationPoint)
 =\KroneckerSymbol_{\BasisIndex\SecondBasisIndex}.
\]
For every sufficiently large
$\SequenceIndex$
and every pair of integers
$0\leq\BasisIndex,\SecondBasisIndex<\SpectralRank$,
define the uniform product error by
\begin{equation}\label{eq:gram-uniform-product-error}
 a_{\SequenceIndex,\BasisIndex,\SecondBasisIndex}
 =\norm{\ProjectedBasisFunction_{\SequenceIndex,\BasisIndex}
       \ProjectedBasisFunction_{\SequenceIndex,\SecondBasisIndex}
       -\ProjectedBasisFunction_\BasisIndex\ProjectedBasisFunction_\SecondBasisIndex}_\infty.
\end{equation}
For the same indices,
define the integration error for the fixed limiting product by
\begin{equation}\label{eq:gram-weak-integration-error}
 b_{\SequenceIndex,\BasisIndex,\SecondBasisIndex}
 =\left|\int_\AmbientCarrier
       \ProjectedBasisFunction_\BasisIndex(\IntegrationPoint)\ProjectedBasisFunction_\SecondBasisIndex(\IntegrationPoint)
       \,\IntegrationDifferential\bar\mu_\SequenceIndex(\IntegrationPoint)
       -\int_\AmbientCarrier
       \ProjectedBasisFunction_\BasisIndex(\IntegrationPoint)\ProjectedBasisFunction_\SecondBasisIndex(\IntegrationPoint)
       \,\IntegrationDifferential\bar\mu(\IntegrationPoint)\right|.
\end{equation}
Since
$\bar\mu_\SequenceIndex(\AmbientCarrier)=1$,
\begin{equation}\label{eq:gram-entry-error}
 \left|(\GramMatrix_\SequenceIndex)_{\BasisIndex\SecondBasisIndex}
             -\KroneckerSymbol_{\BasisIndex\SecondBasisIndex}\right|
 \leq a_{\SequenceIndex,\BasisIndex,\SecondBasisIndex}
       +b_{\SequenceIndex,\BasisIndex,\SecondBasisIndex}.
\end{equation}
For every fixed
$0\leq\BasisIndex,\SecondBasisIndex<\SpectralRank$,
uniform convergence of the products implies
$a_{\SequenceIndex,\BasisIndex,\SecondBasisIndex}\to0$
in \eqref{eq:gram-uniform-product-error},
and weak convergence of
$\bar\mu_\SequenceIndex$
to
$\bar\mu$
implies
$b_{\SequenceIndex,\BasisIndex,\SecondBasisIndex}\to0$
in \eqref{eq:gram-weak-integration-error}.
By \eqref{eq:gram-entry-error},
\[
 (\GramMatrix _\SequenceIndex )_{\BasisIndex \SecondBasisIndex }\to\int_\AmbientCarrier  \ProjectedBasisFunction _\BasisIndex(\IntegrationPoint) \ProjectedBasisFunction _\SecondBasisIndex(\IntegrationPoint) \,\IntegrationDifferential \bar\mu(\IntegrationPoint)=\KroneckerSymbol _{\BasisIndex \SecondBasisIndex }.
\]
Thus
$\GramMatrix _\SequenceIndex \to \IdentityMatrix _\SpectralRank $,
so choose
$\SequenceIndex_0\in\NonnegativeIntegers$
such that for every
$\SequenceIndex\geq\SequenceIndex_0$
the matrix
$\GramMatrix_\SequenceIndex$
is positive definite.

\ProofStep{Normalization by inverse square roots.}\label{step:spectral-basis-normalization}
\paraid{p-dc0f659821c3}
We prove that the inverse square roots of the Gram matrices converge to the identity.
\Cref{thm:matrix-square-root} defines
$\GramMatrix_\SequenceIndex^{-1/2}$
by replacing each positive eigenvalue with its inverse square root in
an orthonormal eigenbasis.
For a real
$\SpectralRank\times\SpectralRank$
matrix
$\MatrixOperator$,
write its Euclidean operator norm as
\[
 \norm{\MatrixOperator}_\OperatorNorm
 =\sup_{\substack{\CoefficientVector\in\RealNumbers^\SpectralRank\\
                  \norm{\CoefficientVector}_2=1}}
       \norm{\MatrixOperator\CoefficientVector}_2.
\]
For every sufficiently large
$\SequenceIndex\in\NonnegativeIntegers$,
put
\[
 \MatrixError_\SequenceIndex
 =\norm{\GramMatrix_\SequenceIndex-\IdentityMatrix_\SpectralRank}_\OperatorNorm.
\]
Entrywise convergence  implies
$\MatrixError_\SequenceIndex\to0$.
Increase
$\SequenceIndex_0$
if necessary so that for every
$\SequenceIndex\geq\SequenceIndex_0$,
\[
 0\leq\MatrixError_\SequenceIndex<1.
\]
Fix
$\SequenceIndex\geq\SequenceIndex_0$.
Since
$\GramMatrix_\SequenceIndex$
is real symmetric,
\Cref{thm:matrix-square-root} implies that every
$\Eigenvalue\in\spec(\GramMatrix_\SequenceIndex)$
is real and admits an eigenvector
$\CoefficientVector\in\RealNumbers^\SpectralRank$
with
$\norm{\CoefficientVector}_2=1$.
For such
$\Eigenvalue$
and
$\CoefficientVector$,
we have
\[
 |\Eigenvalue-1|
 =\norm{(\GramMatrix_\SequenceIndex-\IdentityMatrix_\SpectralRank)
         \CoefficientVector}_2
 \leq\norm{\GramMatrix_\SequenceIndex-\IdentityMatrix_\SpectralRank}_\OperatorNorm
      \norm{\CoefficientVector}_2
 =\MatrixError_\SequenceIndex.
\]
Since
$\Eigenvalue\in\RealNumbers$
and
$\MatrixError_\SequenceIndex<1$,
it follows that
\[
 \spec(\GramMatrix_\SequenceIndex)
 \subset[1-\MatrixError_\SequenceIndex,1+\MatrixError_\SequenceIndex]
 \subset(0,\infty).
\]
For these indices
$\SequenceIndex$,
\Cref{thm:matrix-square-root} implies
\[
 \begin{aligned}
 \norm{\GramMatrix_\SequenceIndex^{-1/2}-\IdentityMatrix_\SpectralRank}_\OperatorNorm
 &=\max_{\Eigenvalue\in\spec(\GramMatrix_\SequenceIndex)}
       |\Eigenvalue^{-1/2}-1|
 \leq\max_{1-\MatrixError_\SequenceIndex\leq\SecondParameter\leq1+\MatrixError_\SequenceIndex}
       |\SecondParameter^{-1/2}-1|\\
 &=\max_{1-\MatrixError_\SequenceIndex\leq\SecondParameter\leq1+\MatrixError_\SequenceIndex}
       \frac{|\SecondParameter-1|}
       {\sqrt{\SecondParameter}(1+\sqrt{\SecondParameter})}\\
 &\leq
 \frac{\MatrixError_\SequenceIndex}
 {\sqrt{1-\MatrixError_\SequenceIndex}
  (1+\sqrt{1-\MatrixError_\SequenceIndex})}
 \longrightarrow0.
 \end{aligned}
\]
Thus
\[
 \GramMatrix_\SequenceIndex^{-1/2}
 \xrightarrow[\SequenceIndex\to\infty]{\norm{\cdot}_\OperatorNorm}
 \IdentityMatrix_\SpectralRank.
\]
For every
$\SequenceIndex\geq\SequenceIndex_0$
and
$\IntegrationPoint\in\AmbientCarrier$,
we define the row map
\[
 \ProjectedBasisRow_\SequenceIndex\colon\AmbientCarrier\to\RealNumbers^{1\times\SpectralRank}
\]
by
\[
 \ProjectedBasisRow_\SequenceIndex(\IntegrationPoint)
 =(\ProjectedBasisFunction_{\SequenceIndex,0}(\IntegrationPoint),\ldots,
   \ProjectedBasisFunction_{\SequenceIndex,\SpectralRank-1}(\IntegrationPoint)).
\]
We multiply this row vector by the matrix
$\GramMatrix_\SequenceIndex^{-1/2}$
on the right and define
\begin{equation}\label{eq:normalized-basis-row}
 \bigl(\widetilde\Eigenfunction_{\SequenceIndex,0}(\IntegrationPoint),\ldots,
       \widetilde\Eigenfunction_{\SequenceIndex,\SpectralRank-1}(\IntegrationPoint)\bigr)
 =\ProjectedBasisRow_\SequenceIndex(\IntegrationPoint)\GramMatrix_\SequenceIndex^{-1/2}.
\end{equation}
For every integer
$0\leq\BasisIndex<\SpectralRank$,
the components in \eqref{eq:normalized-basis-row} satisfy
\[
 \widetilde\Eigenfunction_{\SequenceIndex,\BasisIndex}(\IntegrationPoint)
 =\sum_{\SecondBasisIndex=0}^{\SpectralRank-1}
 \ProjectedBasisFunction_{\SequenceIndex,\SecondBasisIndex}(\IntegrationPoint)
 (\GramMatrix_\SequenceIndex^{-1/2})_{\SecondBasisIndex\BasisIndex},
\]
and
\[
 \Eigenfunction_{\SequenceIndex,\BasisIndex}
 =\widetilde\Eigenfunction_{\SequenceIndex,\BasisIndex}|_{\BaseCarrier_\SequenceIndex}.
\]
Since
$\GramMatrix_\SequenceIndex^{-1/2}$
is symmetric,
the Gram matrix of these restrictions is
\begin{equation}\label{eq:normalized-basis-gram}
 \begin{aligned}
 &\int_\AmbientCarrier
   \bigl(\ProjectedBasisRow_\SequenceIndex(\IntegrationPoint)
         \GramMatrix_\SequenceIndex^{-1/2}\bigr)^{\TransposeSymbol}
   \bigl(\ProjectedBasisRow_\SequenceIndex(\IntegrationPoint)
         \GramMatrix_\SequenceIndex^{-1/2}\bigr)
   \,\IntegrationDifferential\bar\mu_\SequenceIndex(\IntegrationPoint)\\
 &=(\GramMatrix_\SequenceIndex^{-1/2})^{\TransposeSymbol}
   \GramMatrix_\SequenceIndex\GramMatrix_\SequenceIndex^{-1/2}\\
 &=\GramMatrix_\SequenceIndex^{-1/2}\GramMatrix_\SequenceIndex
   \GramMatrix_\SequenceIndex^{-1/2}=\IdentityMatrix_\SpectralRank.
 \end{aligned}
\end{equation}
By \eqref{eq:projection-restriction},
the functions
$\Eigenfunction_{\SequenceIndex,\BasisIndex}$
belong to
$\SpectralSubspace{\BaseCarrier_\SequenceIndex}{\SpectralCutoff}$.
Equation \eqref{eq:normalized-basis-gram} proves orthonormality,
and their number equals
$\dim\SpectralSubspace{\BaseCarrier_\SequenceIndex}{\SpectralCutoff}$,
so they form an orthonormal basis of
$\SpectralSubspace{\BaseCarrier_\SequenceIndex}{\SpectralCutoff}$.
This proves \ref{item:spectral-basis-orthonormality}.

\ProofStep{Continuous extensions and convergence.}\label{step:spectral-basis-extensions}
\paraid{p-9a01e68ec997}
For every
$\SequenceIndex\geq\SequenceIndex_0$
and every integer
$0\leq\BasisIndex<\SpectralRank$,
put
$\BasisExtension_{\SequenceIndex,\BasisIndex}
 =\widetilde\Eigenfunction_{\SequenceIndex,\BasisIndex}$
and
$\BasisExtension_\BasisIndex=\ProjectedBasisFunction_\BasisIndex$.
The projections preserve real functions,
so the construction in Step~\ref{step:spectral-basis-projection} and the finite linear combinations in
\eqref{eq:normalized-basis-row} show that these functions are real and continuous.
Their restriction identities follow from Step~\ref{step:spectral-basis-projection} and the definition of
$\Eigenfunction_{\SequenceIndex,\BasisIndex}$
in Step~\ref{step:spectral-basis-normalization}.
This proves \ref{item:spectral-basis-extensions}.
We next prove \ref{item:spectral-basis-uniform-convergence}.
The uniform convergence of the projected basis functions in Step~\ref{step:spectral-basis-projection},
the convergence
$\GramMatrix_\SequenceIndex^{-1/2}\to\IdentityMatrix_\SpectralRank$,
and \eqref{eq:normalized-basis-row} imply that,
for every integer
$0\leq\BasisIndex<\SpectralRank$,
\[
 \norm{\widetilde\Eigenfunction_{\SequenceIndex,\BasisIndex}
       -\ProjectedBasisFunction_\BasisIndex}_\infty\longrightarrow0.
\]
Since there are finitely many basis indices,
it follows that
\[
 \max_{0\leq\BasisIndex<\SpectralRank}
 \norm{\widetilde\Eigenfunction_{\SequenceIndex,\BasisIndex}
       -\ProjectedBasisFunction_\BasisIndex}_\infty\longrightarrow0.
\]
By the definitions of the extensions,
\[
 \max_{0\leq\BasisIndex<\SpectralRank}
 \norm{\BasisExtension_{\SequenceIndex,\BasisIndex}
       -\BasisExtension_\BasisIndex}_\infty
 =\max_{0\leq\BasisIndex<\SpectralRank}
 \norm{\widetilde\Eigenfunction_{\SequenceIndex,\BasisIndex}
       -\ProjectedBasisFunction_\BasisIndex}_\infty
 \longrightarrow0,
\]
which proves \ref{item:spectral-basis-uniform-convergence}.
Finally,
let
$\Correspondence_\SequenceIndex\subset\BaseCarrier_\SequenceIndex\times\BaseCarrier$
be a sequence of correspondences satisfying the hypothesis in
\ref{item:spectral-basis-correspondence-convergence}.
For every
$\SequenceIndex\geq\SequenceIndex_0$,
put
\[
 \CorrespondenceDisplacement_\SequenceIndex
 =\sup_{(\BasePoint_\SequenceIndex,\BasePoint)\in\Correspondence_\SequenceIndex}
       \AmbientMetric(\BasePoint_\SequenceIndex,\BasePoint).
\]
For every integer
$0\leq\BasisIndex<\SpectralRank$,
the restriction identities in \ref{item:spectral-basis-extensions} imply
\begin{equation}\label{eq:basis-alignment-error}
 \sup_{(\BasePoint_\SequenceIndex,\BasePoint)\in\Correspondence_\SequenceIndex}|\Eigenfunction _{\SequenceIndex ,\BasisIndex }(\BasePoint _\SequenceIndex )-\Eigenfunction _\BasisIndex (\BasePoint )|
 \leq\norm{\widetilde \Eigenfunction _{\SequenceIndex ,\BasisIndex }-\ProjectedBasisFunction _\BasisIndex }_\infty
       +\sup_{\substack{\IntegrationPoint ,\IntegrationPoint _0\in \AmbientCarrier \\\AmbientMetric (\IntegrationPoint ,\IntegrationPoint _0)\leq\CorrespondenceDisplacement _\SequenceIndex }}
                  |\ProjectedBasisFunction _\BasisIndex (\IntegrationPoint )-\ProjectedBasisFunction _\BasisIndex (\IntegrationPoint _0)|.
\end{equation}
The maximum over the basis indices of the norm term in
\eqref{eq:basis-alignment-error} tends to
$0$
by
\ref{item:spectral-basis-uniform-convergence}
and the definitions of the extensions.
For each
$0\leq\BasisIndex<\SpectralRank$,
uniform continuity of
$\ProjectedBasisFunction_\BasisIndex$
on the compact space $\AmbientCarrier$ and
$\CorrespondenceDisplacement_\SequenceIndex\to0$
imply the following convergence,
where the maximum tends to
$0$
because there are finitely many basis indices.
\[
 \max_{0\leq\BasisIndex<\SpectralRank}
 \sup_{\substack{\IntegrationPoint,\IntegrationPoint_0\in\AmbientCarrier\\
                  \AmbientMetric(\IntegrationPoint,\IntegrationPoint_0)
                  \leq\CorrespondenceDisplacement_\SequenceIndex}}
 \left|\ProjectedBasisFunction_\BasisIndex(\IntegrationPoint)
       -\ProjectedBasisFunction_\BasisIndex(\IntegrationPoint_0)\right|
 \longrightarrow0.
\]
Taking the maximum over the basis indices in
\eqref{eq:basis-alignment-error} now yields
\[
 \max_{0\leq\BasisIndex<\SpectralRank}
 \sup_{(\BasePoint_\SequenceIndex,\BasePoint)\in\Correspondence_\SequenceIndex}
 \left|\Eigenfunction_{\SequenceIndex,\BasisIndex}(\BasePoint_\SequenceIndex)
       -\Eigenfunction_\BasisIndex(\BasePoint)\right|
 \longrightarrow0,
\]
which proves \ref{item:spectral-basis-correspondence-convergence}.
\end{proof}

\subsection{Local maps to normed spaces}\label{subsec:spectral-local-models}
\paraid{p-160c1dd745e3}
We now turn the spectral subspace into coordinates in
$\RealNumbers^\ModelDimension $.
A choice of orthonormal basis determines the coordinates,
but the resulting pseudometric is independent of this choice.

\paraid{p-643ccba9258d}
We use the orthogonal group
\[
 \OrthogonalGroup(\ModelDimension )=\{\OrthogonalChange \in\RealNumbers^{\ModelDimension \times \ModelDimension }\mid \OrthogonalChange ^{\TransposeSymbol }\OrthogonalChange =\IdentityMatrix _\ModelDimension \},
\]
acting on column vectors in
$\RealNumbers^\ModelDimension $.
Thus
$\OrthogonalChange $
preserves the Euclidean norm and
$\OrthogonalChange ^{-1}=\OrthogonalChange ^{\TransposeSymbol }$.

\paraid{p-f848e928ba37}
We consider the equivalence class of a coordinate map and its norm,
so that a change of orthonormal basis does not change the model.

\begin{definition}\label{def:normed-coordinate-class}
\paraid{p-49a6a71b7246}
Let
$\ComparisonCarrier$
be a nonempty compact space and let
$\ModelDimension\geq1$.
Write
\[
 \FiniteNormSpace{\ModelDimension}
 =\{\NormSymbol\colon\RealNumbers^\ModelDimension\to[0,\infty)
       \mid\NormSymbol\text{ is a norm}\}.
\]
The group
$\OrthogonalGroup(\ModelDimension)$
acts on
$\ContinuousFunctions(\ComparisonCarrier,\RealNumbers^\ModelDimension)
 \times\FiniteNormSpace{\ModelDimension}$
by
\begin{equation}\label{eq:frame-action}
 \OrthogonalChange\cdot(\FeatureCoordinates,\NormSymbol)
 =(\OrthogonalChange\circ\FeatureCoordinates,
   \NormSymbol\circ\OrthogonalChange^{-1}).
\end{equation}
We call an element of the quotient set
\begin{equation}\label{eq:normed-coordinate-class-space}
 \CoordinateClassSpace{\ModelDimension}{\ComparisonCarrier}
 =\bigl(\ContinuousFunctions(\ComparisonCarrier,\RealNumbers^\ModelDimension)
       \times\FiniteNormSpace{\ModelDimension}\bigr)
       /\OrthogonalGroup(\ModelDimension)
\end{equation}
an \emph{$\ModelDimension$-dimensional normed coordinate class on
$\ComparisonCarrier$}.
The class of a pair
$(\FeatureCoordinates,\NormSymbol)$
is denoted by
$[\FeatureCoordinates,\NormSymbol]_{\OrthogonalGroup(\ModelDimension)}$.
\end{definition}

\paraid{p-f230f9fc55e9}
For an integer
$\ModelDimension\geq1$,
a nonempty compact space
$\ComparisonCarrier$,
a continuous map
$\FeatureCoordinates\colon\ComparisonCarrier\to\RealNumbers^\ModelDimension$,
a norm
$\NormSymbol\in\FiniteNormSpace{\ModelDimension}$,
and points
$\ComparisonPoint,\IntegrationPoint\in\ComparisonCarrier$,
set
\[
 \Pseudometric(\ComparisonPoint,\IntegrationPoint)
 =\NormSymbol(\FeatureCoordinates(\ComparisonPoint)-\FeatureCoordinates(\IntegrationPoint)).
\]
This formula defines a continuous pseudometric on
$\ComparisonCarrier$.
The pseudometric is unchanged by \eqref{eq:frame-action}.
Thus it depends only on the normed coordinate class.
Coordinate classes are defined on concrete carriers.
To compare classes on different carriers,
we use the compatibility with input isometries in \Cref{thm:local-models}.
That theorem constructs such a class for each compact metric space
in a neighborhood of a fixed center.
Representing coordinate pairs can be chosen to converge along each convergent sequence of spaces.
We use their linear coordinates to combine approximations by a partition of unity
in Section~\ref{sec:ar}.
Figure~\ref{fig:orthogonal-coordinate-change} illustrates the invariance
under \eqref{eq:frame-action}.

\begin{figure}[H]
\centering
\begin{tikzpicture}[>=Stealth,font=\small]
\node at (0,1.65)
 {$(\FeatureCoordinates,\NormSymbol)$};
\node at (5.8,1.65)
 {$(\OrthogonalChange\circ\FeatureCoordinates,
     \NormSymbol\circ\OrthogonalChange^{-1})$};
\begin{scope}
 \filldraw[fill=black!4] (0,0) ellipse (1.45 and 0.88);
 \draw[->,black!35] (-1.8,0) -- (1.8,0);
 \draw[->,black!35] (0,-1.15) -- (0,1.2);
 \foreach \x/\y in {-0.8/-0.3,-0.55/0.42,0.12/-0.55,0.48/0.3,0.92/-0.15}
  \fill (\x,\y) circle (1.4pt);
 \draw[black!55] (0.48,0.3) -- (0.85,1.03);
 \node at (0.85,1.25)
  {$\FeatureCoordinates(\ComparisonPoint)$};
 \draw[black!55] (-0.8,-0.3) -- (-0.95,-0.9);
 \node at (-0.95,-1.12)
  {$\FeatureCoordinates(\IntegrationPoint)$};
\end{scope}
\begin{scope}[xshift=5.8cm]
 \begin{scope}[rotate=30]
  \filldraw[fill=black!4] (0,0) ellipse (1.45 and 0.88);
 \end{scope}
 \draw[->,black!35] (-1.8,0) -- (1.8,0);
 \draw[->,black!35] (0,-1.15) -- (0,1.2);
 \begin{scope}[rotate=30]
  \foreach \x/\y in {-0.8/-0.3,-0.55/0.42,0.12/-0.55,0.48/0.3,0.92/-0.15}
   \fill (\x,\y) circle (1.4pt);
  \coordinate (rotatedy) at (0.48,0.3);
  \coordinate (rotatedz) at (-0.8,-0.3);
 \end{scope}
 \draw[black!55] (rotatedy) -- (0.4,1.03);
 \node at (0.4,1.25)
  {$\OrthogonalChange\FeatureCoordinates(\ComparisonPoint)$};
 \draw[black!55] (rotatedz) -- (-1,-0.92);
 \node at (-1,-1.14)
  {$\OrthogonalChange\FeatureCoordinates(\IntegrationPoint)$};
\end{scope}
\draw[->] (2.05,0.15) -- node[above]
 {$\OrthogonalChange\in\OrthogonalGroup(2)$} (3.75,0.15);
\node at (0,-1.55)
 {$\{\FirstVector\mid\NormSymbol(\FirstVector)\leq1\}$};
\node at (5.8,-1.55)
 {$\{\FirstVector\mid
     (\NormSymbol\circ\OrthogonalChange^{-1})(\FirstVector)\leq1\}$};
\node at (2.9,-2.35)
 {$(\NormSymbol\circ\OrthogonalChange^{-1})
     \bigl(\OrthogonalChange\FeatureCoordinates(\ComparisonPoint)
           -\OrthogonalChange\FeatureCoordinates(\IntegrationPoint)\bigr)
   =\NormSymbol\bigl(\FeatureCoordinates(\ComparisonPoint)
                     -\FeatureCoordinates(\IntegrationPoint)\bigr)$};
\end{tikzpicture}
\caption{An orthogonal change of coordinates acts on both the coordinate
image and the norm unit ball.
The dots represent a finite coordinate image in dimension
$2$,
and the shaded regions represent the two unit balls.
Their simultaneous rotation leaves the induced pseudometric unchanged.}
\label{fig:orthogonal-coordinate-change}
\end{figure}


\begin{lemma}\label{lem:coordinate-pseudometric-estimate}
\paraid{p-29c8b5705912}
Let
$\BaseCarrier,\ComparisonCarrier$
be nonempty compact spaces,
let
$\Correspondence\subset\BaseCarrier\times\ComparisonCarrier$
be a correspondence,
and let
$\ModelDimension\geq1$.
Let
$\FeatureCoordinates_0\colon\BaseCarrier\to\RealNumbers^\ModelDimension$
and
$\FeatureCoordinates_1\colon\ComparisonCarrier\to\RealNumbers^\ModelDimension$
be continuous,
and let
$\NormSymbol_0,\NormSymbol_1$
be norms on
$\RealNumbers^\ModelDimension$.
Denote the induced pseudometrics by
$\Pseudometric_0$
and
$\Pseudometric_1$,
respectively.
Put
\[
 \CoefficientBound=\sup_{\BasePoint\in\BaseCarrier}
       \norm{\FeatureCoordinates_0(\BasePoint)}_2,
\]
\[
 \CoordinateComparisonError
 =\sup_{(\BasePoint,\ComparisonPoint)\in\Correspondence}
       \norm{\FeatureCoordinates_0(\BasePoint)-\FeatureCoordinates_1(\ComparisonPoint)}_2,
\]
\[
 \NormComparisonError
 =\sup_{\substack{\CoefficientVector\in\RealNumbers^\ModelDimension\\\norm{\CoefficientVector}_2=1}}
       |\NormSymbol_0(\CoefficientVector)-\NormSymbol_1(\CoefficientVector)|,
\]
and
\[
 \NormComparisonBound
 =\max_{\substack{\CoefficientVector\in\RealNumbers^\ModelDimension\\\norm{\CoefficientVector}_2=1}}\NormSymbol_1(\CoefficientVector).
\]
Then
\begin{equation}\label{eq:coordinate-pseudometric-estimate}
 \PseudometricError_\Correspondence(\Pseudometric_0,\Pseudometric_1)
 \leq2\CoefficientBound\NormComparisonError
       +2\NormComparisonBound\CoordinateComparisonError.
\end{equation}
\end{lemma}

\begin{proof}
\paraid{p-9d7ad7bd9f69}
Fix
$(\BasePoint_0,\ComparisonPoint_0),(\BasePoint_1,\ComparisonPoint_1)\in\Correspondence$
and set
\[
 \FirstVector=\FeatureCoordinates_0(\BasePoint_0)-\FeatureCoordinates_0(\BasePoint_1),
\]
and
\[
 \SecondVector=\FeatureCoordinates_1(\ComparisonPoint_0)-\FeatureCoordinates_1(\ComparisonPoint_1).
\]
Then
$\norm{\FirstVector}_2\leq2\CoefficientBound$
and
$\norm{\FirstVector-\SecondVector}_2\leq2\CoordinateComparisonError$.
Homogeneity and the norm triangle inequality imply
\[
 \begin{aligned}
 |\NormSymbol_0(\FirstVector)-\NormSymbol_1(\SecondVector)|
 &\leq|\NormSymbol_0(\FirstVector)-\NormSymbol_1(\FirstVector)|
       +\NormSymbol_1(\FirstVector-\SecondVector)\\
 &\leq\NormComparisonError\norm{\FirstVector}_2
       +\NormComparisonBound\norm{\FirstVector-\SecondVector}_2\\
 &\leq2\CoefficientBound\NormComparisonError
       +2\NormComparisonBound\CoordinateComparisonError.
 \end{aligned}
\]
Taking the supremum over the two pairs proves
\eqref{eq:coordinate-pseudometric-estimate}.
\end{proof}

\paraid{p-03fcd50a02a7}
When a function is defined on GH isometry classes, its value at
$\MetricSpace{\ComparisonCarrier}{\ComparisonMetric}$ denotes its value at
$[\ComparisonCarrier,\ComparisonMetric]$, the isometry class of that representative.
In particular, this convention applies to the error functions.

\begin{theorem}\label{thm:local-models}
\paraid{p-9559de6c1eef}
Let
$\MetricSpace{\BaseCarrier }{\MetricSymbol }\in\GHSpace$
be a nonempty compact metric space and let
$\LocalTolerance >0$.
Then there exist an open subset
$\ModelNeighborhood\subset\GHSpace$
with
$\MetricSpace{\BaseCarrier }{\MetricSymbol }\in\ModelNeighborhood$
and an integer
$\ModelDimension \geq1$
and a continuous function
$\LocalError\colon\ModelNeighborhood\to[0,\infty)$,
together with the following simultaneous choices.
For each concrete compact metric space
$\MetricSpace{\ComparisonCarrier}{\ComparisonMetric}$
whose isometry class belongs to
$\ModelNeighborhood$,
choose a class
\[
 \LocalModelClass{\ComparisonCarrier}{\ComparisonMetric}
 \in\CoordinateClassSpace{\ModelDimension}{\ComparisonCarrier}.
\]
These choices satisfy the following properties.
\begin{enumerate}[label=\textup{(A\arabic*)},ref=\textup{(A\arabic*)}]
\item\label{item:model-class}
\paraid{p-578e284641e7}
For each such compact metric space
$\MetricSpace{\ComparisonCarrier}{\ComparisonMetric}$,
the assigned class satisfies
\[
 \LocalModelClass{\ComparisonCarrier}{\ComparisonMetric}
 \in\CoordinateClassSpace{\ModelDimension}{\ComparisonCarrier}.
\]
In \ref{item:model-bound} and \ref{item:model-pseudometric}--\ref{item:model-center-error},
$(\FeatureCoordinates_\ComparisonCarrier,\NormSymbol_\ComparisonCarrier)$
denotes an arbitrary coordinate pair in this class.
The bound in \ref{item:model-bound} and the pseudometric and error assertions
in \ref{item:model-pseudometric}--\ref{item:model-center-error} hold for every coordinate pair in this class.
\item\label{item:model-equivariance}
\paraid{p-38960d79969a}
For every isometry
$\InputIsometry\colon\MetricSpace{\ComparisonCarrier}{\ComparisonMetric}
 \to\MetricSpace{\ComparisonCarrier_0}{\ComparisonMetric_0}$
between compact metric spaces whose isometry classes belong to
$\ModelNeighborhood$
and every choice of coordinate pairs
$(\FeatureCoordinates_\ComparisonCarrier,\NormSymbol_\ComparisonCarrier)$
and
$(\FeatureCoordinates_{\ComparisonCarrier_0},\NormSymbol_{\ComparisonCarrier_0})$
of the two assigned coordinate classes,
there exists
$\OrthogonalChange\in\OrthogonalGroup(\ModelDimension)$
such that
\begin{equation}\label{eq:input-isometry-1}
 \FeatureCoordinates_{\ComparisonCarrier_0}\circ\InputIsometry
 =\OrthogonalChange\circ\FeatureCoordinates_\ComparisonCarrier,
\end{equation}
and
\begin{equation}\label{eq:input-isometry-2}
 \NormSymbol_{\ComparisonCarrier_0}=\NormSymbol_\ComparisonCarrier\circ\OrthogonalChange^{-1}.
\end{equation}
\item\label{item:model-bound}
\paraid{p-ac13929fbe66}
For every
$\ComparisonPoint \in \ComparisonCarrier $,
\[
 \NormSymbol _\ComparisonCarrier (\FeatureCoordinates _\ComparisonCarrier (\ComparisonPoint ))\leq2\diam\MetricSpace{\ComparisonCarrier }{\ComparisonMetric }.
\]
\item\label{item:model-continuity}
\paraid{p-e36eef01289e}
Let
$\{\MetricSpace{\ComparisonCarrier_\SequenceIndex}{\ComparisonMetric_\SequenceIndex}\}_{\SequenceIndex\in\NonnegativeIntegers}$
be a sequence of compact metric spaces whose isometry classes belong to
$\ModelNeighborhood$,
and let
$\MetricSpace{\ComparisonCarrier}{\ComparisonMetric}$
be a compact metric space whose isometry class belongs to
$\ModelNeighborhood$.
Assume that these spaces have isometric embeddings into a compact
$\MetricSpace{\AmbientCarrier }{\AmbientMetric }$
with
$\HausdorffDistance{\AmbientMetric }(\ComparisonCarrier _\SequenceIndex ,\ComparisonCarrier )\to0$.
Then there exist a coordinate pair
$(\FeatureCoordinates_\ComparisonCarrier,\NormSymbol_\ComparisonCarrier)$
and a family of coordinate pairs
$\{(\FeatureCoordinates_{\ComparisonCarrier_\SequenceIndex},\NormSymbol_{\ComparisonCarrier_\SequenceIndex})\}_{\SequenceIndex\in\NonnegativeIntegers}$
such that
\[
 (\FeatureCoordinates_\ComparisonCarrier,\NormSymbol_\ComparisonCarrier)
 \in\LocalModelClass{\ComparisonCarrier}{\ComparisonMetric},
\]
and
\[
 (\FeatureCoordinates_{\ComparisonCarrier_\SequenceIndex},\NormSymbol_{\ComparisonCarrier_\SequenceIndex})
 \in\LocalModelClass{\ComparisonCarrier_\SequenceIndex}{\ComparisonMetric_\SequenceIndex},
\]
and
\[
 \sup_{\substack{\CoefficientVector\in\RealNumbers^\ModelDimension\\\norm{\CoefficientVector}_2=1}}|\NormSymbol _{\ComparisonCarrier _\SequenceIndex }(\CoefficientVector )-\NormSymbol _\ComparisonCarrier (\CoefficientVector )|\to0.
\]
These coordinate pairs can be chosen using only the prescribed sequence
and ambient embeddings, independently of any sequence of correspondences with vanishing ambient displacement.
For every sequence of correspondences
$\Correspondence _\SequenceIndex \subset \ComparisonCarrier _\SequenceIndex \times \ComparisonCarrier $
satisfying
\[
 \sup_{(\ComparisonPoint _\SequenceIndex ,\ComparisonPoint )\in \Correspondence _\SequenceIndex }\AmbientMetric (\ComparisonPoint _\SequenceIndex ,\ComparisonPoint )\to0,
\]
the coordinate maps
$\FeatureCoordinates_{\ComparisonCarrier_\SequenceIndex}$
and
$\FeatureCoordinates_\ComparisonCarrier$
satisfy
\begin{equation}\label{eq:local-coordinate-convergence}
 \sup_{(\ComparisonPoint_\SequenceIndex,\ComparisonPoint)\in\Correspondence_\SequenceIndex}
 \norm{\FeatureCoordinates_{\ComparisonCarrier_\SequenceIndex}(\ComparisonPoint_\SequenceIndex)
 -\FeatureCoordinates_\ComparisonCarrier(\ComparisonPoint)}_2\to0.
\end{equation}
\item\label{item:model-pseudometric}
\paraid{p-29f1150553e1}
For
$\ComparisonPoint,\IntegrationPoint\in\ComparisonCarrier$,
the formula
\begin{equation}\label{eq:local-model-pseudometric}
 \Pseudometric _\ComparisonCarrier (\ComparisonPoint ,\IntegrationPoint )=\NormSymbol _\ComparisonCarrier (\FeatureCoordinates _\ComparisonCarrier (\ComparisonPoint )-\FeatureCoordinates _\ComparisonCarrier (\IntegrationPoint ))
\end{equation}
defines a continuous pseudometric on
$\ComparisonCarrier $.
\item\label{item:model-pseudometric-invariance}
\paraid{p-b7eec021787e}
For every isometry
$\InputIsometry\colon\MetricSpace{\ComparisonCarrier}{\ComparisonMetric}
 \to\MetricSpace{\ComparisonCarrier_0}{\ComparisonMetric_0}$
between compact metric spaces whose isometry classes belong to
$\ModelNeighborhood$
and every
$\ComparisonPoint,\IntegrationPoint\in\ComparisonCarrier$,
\[
 \Pseudometric_{\ComparisonCarrier_0}(\InputIsometry(\ComparisonPoint),\InputIsometry(\IntegrationPoint))
 =\Pseudometric_\ComparisonCarrier(\ComparisonPoint,\IntegrationPoint).
\]
\item\label{item:model-error-definition}
\paraid{p-702b7e8de607}
For each compact metric representative
$\MetricSpace{\ComparisonCarrier}{\ComparisonMetric}$
of a class in
$\ModelNeighborhood$,
the error function satisfies
\[
 \LocalError \MetricSpace{\ComparisonCarrier }{\ComparisonMetric }=\sup_{\ComparisonPoint ,\IntegrationPoint \in \ComparisonCarrier }|\Pseudometric _\ComparisonCarrier (\ComparisonPoint ,\IntegrationPoint )-\ComparisonMetric (\ComparisonPoint ,\IntegrationPoint )|.
\]
By \ref{item:model-pseudometric-invariance}, this value is independent
of the representative and of the coordinate pair.
\item\label{item:model-error-continuity}
\paraid{p-685c2b6e8303}
The error function
$\LocalError$
is continuous.
\item\label{item:model-center-error}
\paraid{p-f46126c20706}
At the center, it satisfies
$\LocalError \MetricSpace{\BaseCarrier }{\MetricSymbol }<\LocalTolerance $.
\end{enumerate}
\end{theorem}

\paraid{p-fc5c3b7ac001}
Write
$\mathcal L$
for the choice rule that associates to each concrete compact metric space
$\MetricSpace{\ComparisonCarrier}{\ComparisonMetric}$
whose class belongs to
$\ModelNeighborhood$
the class
\[
 \mathcal L_{\ComparisonCarrier,\ComparisonMetric}
 =\LocalModelClass{\ComparisonCarrier}{\ComparisonMetric}
 \in\CoordinateClassSpace{\ModelDimension}{\ComparisonCarrier}.
\]
We call the data
\[
 \mathfrak M=(\ModelNeighborhood,\ModelDimension,\mathcal L,\LocalError)
\]
a \emph{local model}.
Condition \ref{item:model-equivariance} describes how these values are transported by isometries.
Condition \ref{item:model-continuity} asserts the existence of a convergent family of coordinate pairs
for each prescribed sequence and its ambient embeddings.
These pairs represent the classes of
$\mathfrak M$.
By conditions \ref{item:model-pseudometric-invariance}--\ref{item:model-error-continuity},
the rule
\[
 \MetricSpace{\ComparisonCarrier}{\ComparisonMetric}
 \longmapsto\MetricQuotient{\ComparisonCarrier}{\Pseudometric_\ComparisonCarrier}
\]
induces a well-defined map
$\ModelNeighborhood\to\GHSpace$.
The error function
$\LocalError$
is independently well defined on isometry classes by
\ref{item:model-error-definition}.
The map
$\FeatureCoordinates _\ComparisonCarrier $
need not be injective,
which is why
$\Pseudometric _\ComparisonCarrier $
is a pseudometric.

\begin{proof}
\paraid{p-6d8a3f3aa620}
For the singleton center and each
$\MetricSpace{\ComparisonCarrier}{\ComparisonMetric}\in\GHSpace$,
set
$\ModelDimension =1$,
$\NormSymbol _\ComparisonCarrier =|\cdot|$,
and
$\FeatureCoordinates _\ComparisonCarrier =0$.
Define
$\LocalModelClass{\ComparisonCarrier}{\ComparisonMetric}
 =[(0,|\cdot|)]_{\OrthogonalGroup(1)}$.
Then
$\Pseudometric _\ComparisonCarrier =0$
and
$\LocalError \MetricSpace{\ComparisonCarrier }{\ComparisonMetric }=\diam\MetricSpace{\ComparisonCarrier }{\ComparisonMetric }$.
The error is continuous because the diameter is continuous for
$\GHDistance$.
For a prescribed
$\LocalTolerance>0$,
we may take its domain to be the open neighborhood
\[
 \ModelNeighborhood
 =\{\MetricSpace{\ComparisonCarrier}{\ComparisonMetric}\in\GHSpace
       \mid\diam\MetricSpace{\ComparisonCarrier}{\ComparisonMetric}<\LocalTolerance\}
 \ni\SingletonSpace.
\]
On this domain the zero maps and fixed norm satisfy
\ref{item:model-bound}--\ref{item:model-center-error},
with the error equal to the diameter.
The coordinate maps
$\FeatureCoordinates_\ComparisonCarrier$
vanish and the norm
$\NormSymbol_\ComparisonCarrier$
is fixed,
so the coordinate and isometry identities also hold.
This proves the theorem when the center is a singleton.

\paraid{p-08d995d825ac}
Assume henceforth that the center is not a singleton.
We construct the local data from a single positive real number
$\SpectralCutoff$,
prove their continuity first on each carrier and then as the carrier varies,
and finally verify invariance and the error estimate.

\ProofStep{A neighborhood with constant spectral dimension.}\label{step:local-model-spectral-dimension}
\paraid{p-3a9b99318187}
Choose
$\Radius>0$
sufficiently small.
Then choose
$2\leq\IntegrabilityExponent<\infty$
sufficiently large.
Finally,
choose
$\SmallError>0$
sufficiently small and a corresponding positive real number
$\SpectralCutoff $
for
$\MetricSpace{\BaseCarrier }{\MetricSymbol }$
so that
\[
\begin{gathered}
 \norm{\Pseudometric_\BaseCarrier-\MetricSymbol}_\infty
 \leq2\SmallError
      +\DiameterBound(1-\BallMassBound_\Radius^{1/\IntegrabilityExponent})
      +2\Radius
 <\LocalTolerance,
\\
 \SpectralCutoff,-\SpectralCutoff
 \notin\spec(\DistanceOperator{\BaseCarrier}{\MetricSymbol}),
 \qquad
 \dim\SpectralSubspace{\BaseCarrier}{\SpectralCutoff}
 =\ModelDimension\geq1.
\end{gathered}
\]
Define
\[
 \ModelNeighborhood=
 \{\MetricSpace{\ComparisonCarrier }{\ComparisonMetric }\in\GHSpace\mid
     \SpectralCutoff ,-\SpectralCutoff \notin\spec(\DistanceOperator{\ComparisonCarrier}{\ComparisonMetric}),\
     \dim\SpectralSubspace{\ComparisonCarrier }{\SpectralCutoff }=\ModelDimension \}.
\]
Then
$\MetricSpace{\BaseCarrier}{\MetricSymbol}\in\ModelNeighborhood$.
We prove that
$\ModelNeighborhood$
is open.
Let
$\MetricSpace{\ComparisonCarrier_\SequenceIndex}{\ComparisonMetric_\SequenceIndex}$
converge to
$\MetricSpace{\ComparisonCarrier}{\ComparisonMetric}\in\ModelNeighborhood$
in
$\GHSpace$.
By \Cref{lem:common-gh-embedding},
we realize the sequence and its limit as compact subsets of a common compact metric space
$\MetricSpace{\AmbientCarrier}{\AmbientMetric}$
with
\[
 \HausdorffDistance{\AmbientMetric}
 (\ComparisonCarrier_\SequenceIndex,\ComparisonCarrier)\to0.
\]
\Cref{thm:invariant-measures} implies
\[
 \SelectedMeasure{\ComparisonCarrier_\SequenceIndex}{\ComparisonMetric_\SequenceIndex}
 \to\SelectedMeasure{\ComparisonCarrier}{\ComparisonMetric}
 \quad\text{weakly on }\AmbientCarrier.
\]
\Cref{lem:spectral-projections} now yields
$\MetricSpace{\ComparisonCarrier_\SequenceIndex}{\ComparisonMetric_\SequenceIndex}
 \in\ModelNeighborhood$
for all sufficiently large
$\SequenceIndex$.
Since
$\GHSpace$
is metrizable,
this proves openness.

\ProofStep{Norms and continuous coordinate maps.}\label{step:local-model-coordinates}
\paraid{p-5e1efa60ac95}
For
$\MetricSpace{\ComparisonCarrier }{\ComparisonMetric }\in\ModelNeighborhood$,
choose a real orthonormal basis
$\BasisFunction_0,\ldots,\BasisFunction_{\ModelDimension-1}$
of
$\SpectralSubspace{\ComparisonCarrier }{\SpectralCutoff }$.
Define the functional
$\NormSymbol_\ComparisonCarrier\colon\RealNumbers^\ModelDimension\to[0,\infty)$
as follows.
Use
$\IntegrationPoint\in\ComparisonCarrier$
as the integration variable in
\eqref{eq:local-norm} and \eqref{eq:best-approximation-objective}.
For every
$\CoefficientVector\in\RealNumbers^\ModelDimension$,
put
\begin{equation}\label{eq:local-norm}
 \NormSymbol _\ComparisonCarrier (\CoefficientVector )=\left(\int_\ComparisonCarrier \left|\sum_{\BasisIndex=0}^{\ModelDimension-1}  \CoefficientVector _\BasisIndex \BasisFunction _\BasisIndex (\IntegrationPoint )\right|^\IntegrabilityExponent
                                     \IntegrationDifferential \SelectedMeasure{\ComparisonCarrier}{\ComparisonMetric}(\IntegrationPoint )\right)^{1/\IntegrabilityExponent }.
\end{equation}
Since
$\IntegrabilityExponent \geq2$
and the measure
$\SelectedMeasure{\ComparisonCarrier}{\ComparisonMetric}$
is a probability,
orthonormality implies for every
$\CoefficientVector\in\RealNumbers^\ModelDimension$
\begin{equation}\label{eq:local-norm-euclidean-lower-bound}
 \NormSymbol _\ComparisonCarrier (\CoefficientVector )\geq\left\|\sum_{\BasisIndex=0}^{\ModelDimension-1}  \CoefficientVector _\BasisIndex \BasisFunction _\BasisIndex \right\|_2=\norm{\CoefficientVector }_2.
\end{equation}
The functional is homogeneous and satisfies the triangle inequality by the
$\LebesgueSpace^\IntegrabilityExponent$-norm properties.
The lower bound in \eqref{eq:local-norm-euclidean-lower-bound} proves nondegeneracy,
so
$\NormSymbol_\ComparisonCarrier$ is a norm.

\paraid{p-8cbc0bb85af6}
Define the objective function
$\ApproximationObjective_\ComparisonCarrier\colon\ComparisonCarrier\times\RealNumbers^\ModelDimension\to[0,\infty)$
as follows.
For every
$\ComparisonPoint\in\ComparisonCarrier$
and
$\CoefficientVector\in\RealNumbers^\ModelDimension$,
put
\begin{equation}\label{eq:best-approximation-objective}
 \ApproximationObjective _\ComparisonCarrier (\ComparisonPoint ,\CoefficientVector )
 =\int_\ComparisonCarrier \left|\ComparisonMetric (\ComparisonPoint ,\IntegrationPoint )-\sum_{\BasisIndex=0}^{\ModelDimension-1}  \CoefficientVector _\BasisIndex \BasisFunction _\BasisIndex (\IntegrationPoint )\right|^\IntegrabilityExponent
                                      \IntegrationDifferential \SelectedMeasure{\ComparisonCarrier}{\ComparisonMetric}(\IntegrationPoint ).
\end{equation}
For each
$\ComparisonPoint\in\ComparisonCarrier$,
define
$\FeatureCoordinates_\ComparisonCarrier(\ComparisonPoint)\in\RealNumbers^\ModelDimension$
to be the unique minimizer of
$\CoefficientVector\mapsto\ApproximationObjective_\ComparisonCarrier(\ComparisonPoint,\CoefficientVector)$.
These minimizers define a map
$\FeatureCoordinates_\ComparisonCarrier\colon\ComparisonCarrier\to\RealNumbers^\ModelDimension$.
Existence and uniqueness follow from
\Cref{thm:best-approximation},
applied to the finite-dimensional spectral subspace
$\SpectralSubspace{\ComparisonCarrier}{\SpectralCutoff}$,
viewed as a subspace of
$\LebesgueSpace^\IntegrabilityExponent(\ComparisonCarrier,
\SelectedMeasure{\ComparisonCarrier}{\ComparisonMetric})$.
For each
$\ComparisonPoint\in\ComparisonCarrier$,
define the function
$\BestApproximationFunction_{\ComparisonCarrier,\ComparisonPoint}\colon\ComparisonCarrier\to\RealNumbers$
by setting, for every
$\IntegrationPoint\in\ComparisonCarrier$,
\begin{equation}\label{eq:local-best-approximation-function}
 \BestApproximationFunction_{\ComparisonCarrier,\ComparisonPoint}(\IntegrationPoint)
 =\sum_{\BasisIndex=0}^{\ModelDimension-1}
 (\FeatureCoordinates_\ComparisonCarrier(\ComparisonPoint))_\BasisIndex
 \BasisFunction_\BasisIndex(\IntegrationPoint).
\end{equation}
By \eqref{eq:best-approximation-objective},
$\BestApproximationFunction_{\ComparisonCarrier,\ComparisonPoint}$
is the best approximation of
$\ComparisonMetric(\ComparisonPoint,\cdot)$
in
$\SpectralSubspace{\ComparisonCarrier}{\SpectralCutoff}$.
The components of the vector
$\FeatureCoordinates_\ComparisonCarrier(\ComparisonPoint)$
are the coefficients of this best approximation in the chosen orthonormal basis
$\{\BasisFunction_\BasisIndex\}_{0\leq\BasisIndex<\ModelDimension}$.
Figure~\ref{fig:local-coordinates} displays this construction of
$\FeatureCoordinates_\ComparisonCarrier$.
\begin{figure}[H]
\centering
\begin{tikzpicture}[>=Stealth,node distance=12mm,font=\small]
\node (profile) {$\ComparisonMetric(\ComparisonPoint,\cdot)
 \in\LebesgueSpace^\IntegrabilityExponent
 (\ComparisonCarrier,\SelectedMeasure{\ComparisonCarrier}{\ComparisonMetric})$};
\node[below=of profile] (approximation) {$
 \displaystyle\sum_{\BasisIndex=0}^{\ModelDimension-1}
 (\FeatureCoordinates_\ComparisonCarrier(\ComparisonPoint))_\BasisIndex
 \BasisFunction_\BasisIndex
 \in\SpectralSubspace{\ComparisonCarrier}{\SpectralCutoff}$};
\node[below=of approximation] (coordinates) {$
 \FeatureCoordinates_\ComparisonCarrier(\ComparisonPoint)
 \in(\RealNumbers^\ModelDimension,\NormSymbol_\ComparisonCarrier)$};
\draw[->] (profile) -- node[right]
 {best approximation in $\LebesgueSpace^\IntegrabilityExponent$} (approximation);
\draw[->] (approximation) -- node[right]
 {coordinates in $\{\BasisFunction_\BasisIndex\}_{0\leq\BasisIndex<\ModelDimension}$} (coordinates);
\end{tikzpicture}
\caption{The best approximation of
$\ComparisonMetric(\ComparisonPoint,\cdot)$
in
$\SpectralSubspace{\ComparisonCarrier}{\SpectralCutoff}$
has coordinate vector
$\FeatureCoordinates_\ComparisonCarrier(\ComparisonPoint)$
in the chosen
$\LebesgueSpace^2$-orthonormal
basis.
This vector is the minimizer in \eqref{eq:best-approximation-objective}.
With the
$\LebesgueSpace^\IntegrabilityExponent$
norm on
$\SpectralSubspace{\ComparisonCarrier}{\SpectralCutoff}$,
the coordinate map in the second arrow is a linear isometry onto
$(\RealNumbers^\ModelDimension,\NormSymbol_\ComparisonCarrier)$
by \eqref{eq:local-norm}.}
\label{fig:local-coordinates}
\end{figure}


\paraid{p-aa449f741605}
For each
$\ComparisonPoint\in\ComparisonCarrier$,
write
$\ComparisonMetric_\ComparisonPoint=\ComparisonMetric(\ComparisonPoint,\cdot)
 \in\ContinuousFunctions(\ComparisonCarrier)$
for the distance profile at
$\ComparisonPoint$.
Since
$\SelectedMeasure{\ComparisonCarrier}{\ComparisonMetric}$
is a probability and
$\ComparisonMetric(\ComparisonPoint,\IntegrationPoint)\leq
\diam\MetricSpace{\ComparisonCarrier}{\ComparisonMetric}$
for every
$\IntegrationPoint\in\ComparisonCarrier$,
we have
\[
 \norm{\ComparisonMetric_\ComparisonPoint}_\IntegrabilityExponent^{\IntegrabilityExponent}
 =\int_\ComparisonCarrier
   |\ComparisonMetric(\ComparisonPoint,\IntegrationPoint)|^{\IntegrabilityExponent}
   \,\IntegrationDifferential\SelectedMeasure{\ComparisonCarrier}{\ComparisonMetric}(\IntegrationPoint)
 \leq\diam\MetricSpace{\ComparisonCarrier}{\ComparisonMetric}^{\IntegrabilityExponent}
   \int_\ComparisonCarrier
   \,\IntegrationDifferential\SelectedMeasure{\ComparisonCarrier}{\ComparisonMetric}(\IntegrationPoint)
 =\diam\MetricSpace{\ComparisonCarrier}{\ComparisonMetric}^{\IntegrabilityExponent}.
\]
By minimality of
$\FeatureCoordinates_\ComparisonCarrier(\ComparisonPoint)$,
\[
 \norm{\ComparisonMetric_\ComparisonPoint
       -\BestApproximationFunction_{\ComparisonCarrier,\ComparisonPoint}}_\IntegrabilityExponent
 \leq\norm{\ComparisonMetric_\ComparisonPoint}_\IntegrabilityExponent.
\]
Therefore, the triangle inequality and the definition of
$\NormSymbol_\ComparisonCarrier$
yield
\[
 \NormSymbol_\ComparisonCarrier
   (\FeatureCoordinates_\ComparisonCarrier(\ComparisonPoint))
 =\norm{\BestApproximationFunction_{\ComparisonCarrier,\ComparisonPoint}}_\IntegrabilityExponent
 \leq\norm{\ComparisonMetric_\ComparisonPoint}_\IntegrabilityExponent
    +\norm{\ComparisonMetric_\ComparisonPoint
           -\BestApproximationFunction_{\ComparisonCarrier,\ComparisonPoint}}_\IntegrabilityExponent
 \leq2\diam\MetricSpace{\ComparisonCarrier}{\ComparisonMetric}.
\]
Together with the lower bound in
\eqref{eq:local-norm-euclidean-lower-bound},
this proves
\begin{equation}\label{eq:coefficient-bounds}
 \norm{\FeatureCoordinates _\ComparisonCarrier (\ComparisonPoint )}_2\leq \NormSymbol _\ComparisonCarrier (\FeatureCoordinates _\ComparisonCarrier (\ComparisonPoint ))\leq2\diam\MetricSpace{\ComparisonCarrier }{\ComparisonMetric }.
\end{equation}
This proves \ref{item:model-bound}.
In particular,
the minimizers stay in a fixed compact Euclidean ball when the diameters
are bounded.

\paraid{p-29af57cad7b7}
For
$\CoefficientBound>0$,
write
\[
 \CoefficientBall_\CoefficientBound
 =\{\CoefficientVector\in\RealNumbers^\ModelDimension
      \mid\norm{\CoefficientVector}_2\leq\CoefficientBound\}.
\]
For fixed
$\ComparisonCarrier$,
the function
$\ApproximationObjective_\ComparisonCarrier$
is continuous.
If
$\ComparisonPoint_\SubsequenceIndex\to\ComparisonPoint$,
continuity on the product of the compact space
$\ComparisonCarrier$
and each
$\CoefficientBall_\CoefficientBound$
implies
\[
 \sup_{\CoefficientVector\in\CoefficientBall_\CoefficientBound}
 |\ApproximationObjective_\ComparisonCarrier(\ComparisonPoint_\SubsequenceIndex,\CoefficientVector)
  -\ApproximationObjective_\ComparisonCarrier(\ComparisonPoint,\CoefficientVector)|\to0.
\]
The bound \eqref{eq:coefficient-bounds} and
\Cref{lem:minimizer-continuity} prove continuity of
$\FeatureCoordinates_\ComparisonCarrier$.

\ProofStep{Convergence of the objectives and norms.}\label{step:local-model-norm-convergence}
\paraid{p-4cd36224c006}
Assume that
$\HausdorffDistance{\AmbientMetric}(\ComparisonCarrier_\SequenceIndex,\ComparisonCarrier)\to0$
in the stated compact ambient space.
Write
$\ComparisonMetric_\SequenceIndex=\AmbientMetric|_{\ComparisonCarrier_\SequenceIndex\times\ComparisonCarrier_\SequenceIndex}$
and
$\ComparisonMetric=\AmbientMetric|_{\ComparisonCarrier\times\ComparisonCarrier}$.
Let $\iota_\SequenceIndex\colon\ComparisonCarrier_\SequenceIndex\hookrightarrow\AmbientCarrier$
and $\iota\colon\ComparisonCarrier\hookrightarrow\AmbientCarrier$ be the inclusions.
Define $\bar\mu_\SequenceIndex=(\iota_\SequenceIndex)_*\SelectedMeasure{\ComparisonCarrier_\SequenceIndex}{\ComparisonMetric_\SequenceIndex}$
and $\bar\mu=\iota_*\SelectedMeasure{\ComparisonCarrier}{\ComparisonMetric}$.
By \Cref{thm:invariant-measures},
\[
 \bar\mu_\SequenceIndex\to\bar\mu
 \quad\text{weakly on }\AmbientCarrier.
\]
Apply \Cref{lem:spectral-continuity} to the basis
$\{\BasisFunction_\BasisIndex\}_{0\leq\BasisIndex<\ModelDimension}$
of
$\SpectralSubspace{\ComparisonCarrier}{\SpectralCutoff}$
chosen in Step~\ref{step:local-model-coordinates}.
We obtain real orthonormal bases
$\{\BasisFunction_{\SequenceIndex,\BasisIndex}\}_{0\leq\BasisIndex<\ModelDimension}$
of
$\SpectralSubspace{\ComparisonCarrier_\SequenceIndex}{\SpectralCutoff}$
and
continuous functions
$\ExtendedBasisFunction_{\SequenceIndex,\BasisIndex},\ExtendedBasisFunction_\BasisIndex
 \colon\AmbientCarrier\to\RealNumbers$
such that
\[
 \ExtendedBasisFunction_{\SequenceIndex,\BasisIndex}|_{\ComparisonCarrier_\SequenceIndex}
 =\BasisFunction_{\SequenceIndex,\BasisIndex},
\]
and
\[
 \ExtendedBasisFunction_\BasisIndex|_\ComparisonCarrier=\BasisFunction_\BasisIndex,
\]
and
\[
 \max_{0\leq\BasisIndex<\ModelDimension}
 \norm{\ExtendedBasisFunction_{\SequenceIndex,\BasisIndex}-\ExtendedBasisFunction_\BasisIndex}_\infty\to0.
\]
Define
$\ApproximationObjective_\SequenceIndex,\ApproximationObjective
 \colon\AmbientCarrier\times\RealNumbers^\ModelDimension\to[0,\infty)$
as follows.
In both integrals below,
$\SecondVector$
ranges over
$\AmbientCarrier$.
For every sufficiently large
$\SequenceIndex\in\NonnegativeIntegers$,
every
$\IntegrationPoint\in\AmbientCarrier$,
and every
$\CoefficientVector\in\RealNumbers^\ModelDimension$,
put
\begin{equation}\label{eq:ambient-approximation-objectives-1}
 \ApproximationObjective _\SequenceIndex (\IntegrationPoint ,\CoefficientVector )
 =\int_\AmbientCarrier \left|\AmbientMetric (\IntegrationPoint ,\SecondVector )-\sum_{\BasisIndex=0}^{\ModelDimension-1}  \CoefficientVector _\BasisIndex \ExtendedBasisFunction _{\SequenceIndex ,\BasisIndex }(\SecondVector )\right|^\IntegrabilityExponent
                                     \IntegrationDifferential \bar\mu_\SequenceIndex(\SecondVector ).
\end{equation}
For every
$\IntegrationPoint\in\AmbientCarrier$
and
$\CoefficientVector\in\RealNumbers^\ModelDimension$,
put
\begin{equation}\label{eq:ambient-approximation-objectives-2}
 \ApproximationObjective(\IntegrationPoint,\CoefficientVector)
 =\int_\AmbientCarrier\left|\AmbientMetric(\IntegrationPoint,\SecondVector)
       -\sum_{\BasisIndex=0}^{\ModelDimension-1}
        \CoefficientVector_\BasisIndex\ExtendedBasisFunction_\BasisIndex(\SecondVector)
       \right|^\IntegrabilityExponent
       \IntegrationDifferential\bar\mu(\SecondVector).
\end{equation}
By \eqref{eq:best-approximation-objective},
\[
 \ApproximationObjective_\SequenceIndex
 |_{\ComparisonCarrier_\SequenceIndex\times\RealNumbers^\ModelDimension}
 =\ApproximationObjective_{\ComparisonCarrier_\SequenceIndex},
\]
and
\[
 \ApproximationObjective|_{\ComparisonCarrier\times\RealNumbers^\ModelDimension}
 =\ApproximationObjective_\ComparisonCarrier.
\]
For every
$\CoefficientBound >0$,
\begin{equation}\label{eq:uniform-objectives}
 \sup_{\substack{\IntegrationPoint\in\AmbientCarrier\\\CoefficientVector\in\RealNumbers^\ModelDimension\\\norm{\CoefficientVector}_2\leq\CoefficientBound}}
       |\ApproximationObjective _\SequenceIndex (\IntegrationPoint ,\CoefficientVector )-\ApproximationObjective (\IntegrationPoint ,\CoefficientVector )|\to0.
\end{equation}
Apply \Cref{lem:parameter-integrals} with compact parameter space
$\AmbientCarrier\times\{\CoefficientVector\in\RealNumbers^\ModelDimension
 \mid\norm{\CoefficientVector}_2\leq\CoefficientBound\}$.
Uniform convergence of the extended basis functions
$\ExtendedBasisFunction_{\SequenceIndex,\BasisIndex}$
to
$\ExtendedBasisFunction_\BasisIndex$
implies uniform convergence of the integrands in \eqref{eq:ambient-approximation-objectives-1} and \eqref{eq:ambient-approximation-objectives-2}
on this product with
$\AmbientCarrier$.
For the norms
$\NormSymbol_{\ComparisonCarrier_\SequenceIndex},\NormSymbol_\ComparisonCarrier$,
apply the same lemma with the Euclidean unit sphere
$\UnitSphere^{\ModelDimension-1}$
as parameter space and integrands
$\left|\sum_{\BasisIndex=0}^{\ModelDimension-1}\CoefficientVector_\BasisIndex
 \ExtendedBasisFunction_{\SequenceIndex,\BasisIndex}(\SecondVector)\right|^\IntegrabilityExponent$.
It follows that
\begin{equation}\label{eq:local-norm-power-convergence}
 \sup_{\substack{\CoefficientVector\in\RealNumbers^\ModelDimension\\\norm{\CoefficientVector}_2=1}}
 |\NormSymbol_{\ComparisonCarrier_\SequenceIndex}(\CoefficientVector)^\IntegrabilityExponent
 -\NormSymbol_\ComparisonCarrier(\CoefficientVector)^\IntegrabilityExponent|\to0.
\end{equation}
For nonnegative real numbers
$\FirstLength,\SecondLength$
and
$\IntegrabilityExponent\geq1$,
\begin{equation}\label{eq:root-holder-bound}
 |\FirstLength^{1/\IntegrabilityExponent}-\SecondLength^{1/\IntegrabilityExponent}|
 \leq|\FirstLength-\SecondLength|^{1/\IntegrabilityExponent}.
\end{equation}
Apply \eqref{eq:root-holder-bound} with
$\FirstLength=\NormSymbol_{\ComparisonCarrier_\SequenceIndex}(\CoefficientVector)^\IntegrabilityExponent$
and
$\SecondLength=\NormSymbol_\ComparisonCarrier(\CoefficientVector)^\IntegrabilityExponent$.
Taking the supremum and using \eqref{eq:local-norm-power-convergence},
we obtain
\begin{equation}\label{eq:local-norm-convergence}
 \begin{aligned}
 &\sup_{\substack{\CoefficientVector\in\RealNumbers^\ModelDimension\\\norm{\CoefficientVector}_2=1}}
 |\NormSymbol_{\ComparisonCarrier_\SequenceIndex}(\CoefficientVector)
       -\NormSymbol_\ComparisonCarrier(\CoefficientVector)|\\
 &\quad\leq
 \left(\sup_{\substack{\CoefficientVector\in\RealNumbers^\ModelDimension\\\norm{\CoefficientVector}_2=1}}
 |\NormSymbol_{\ComparisonCarrier_\SequenceIndex}(\CoefficientVector)^\IntegrabilityExponent
       -\NormSymbol_\ComparisonCarrier(\CoefficientVector)^\IntegrabilityExponent|
 \right)^{1/\IntegrabilityExponent}
 \longrightarrow0.
 \end{aligned}
\end{equation}

\ProofStep{Uniform convergence of the coordinate maps.}\label{step:local-model-coordinate-convergence}
\paraid{p-5f47a359eb40}
We now prove the second convergence in \ref{item:model-continuity}
for the bases
$\{\BasisFunction_{\SequenceIndex,\BasisIndex}\}_{0\leq\BasisIndex<\ModelDimension}$
and
$\{\BasisFunction_\BasisIndex\}_{0\leq\BasisIndex<\ModelDimension}$
chosen in Step~\ref{step:local-model-norm-convergence}.
Fix an arbitrary sequence of correspondences
$\{\Correspondence_\SequenceIndex\}_{\SequenceIndex\in\NonnegativeIntegers}$,
where
$\Correspondence_\SequenceIndex\subset\ComparisonCarrier_\SequenceIndex\times\ComparisonCarrier$
for every
$\SequenceIndex\in\NonnegativeIntegers$,
satisfying
\[
 \sup_{(\ComparisonPoint_\SequenceIndex,\ComparisonPoint)\in\Correspondence_\SequenceIndex}
 \AmbientMetric(\ComparisonPoint_\SequenceIndex,\ComparisonPoint)\to0.
\]
For the sake of contradiction,
suppose that \eqref{eq:local-coordinate-convergence} fails.
Choose
$\ErrorTolerance>0$,
a strictly increasing sequence of indices
$n_k$,
and for each
$k\in\NonnegativeIntegers$,
a pair
$(\ComparisonPoint _{n_k} ,\ComparisonPoint _{0,{n_k} })\in \Correspondence _{n_k} $
with
\begin{equation}\label{eq:nonconvergent-coordinate-separation}
 \norm{\FeatureCoordinates _{\ComparisonCarrier _{n_k} }(\ComparisonPoint _{n_k} )-\FeatureCoordinates _\ComparisonCarrier (\ComparisonPoint _{0,{n_k} })}_2\geq\ErrorTolerance .
\end{equation}
By compactness of
$\ComparisonCarrier$
and \eqref{eq:coefficient-bounds},
pass to a further subsequence and denote its indices again by
$n_k$.
For some
$\ComparisonPoint\in\ComparisonCarrier$
and
$\CoefficientVector\in\RealNumbers^\ModelDimension$,
we then obtain
\[
 \ComparisonPoint_{0,{n_k}}\to\ComparisonPoint,
\]
and
\[
 \FeatureCoordinates_{\ComparisonCarrier_{n_k}}(\ComparisonPoint_{n_k})
 \to\CoefficientVector.
\]
The condition on the correspondences implies
\[
 \AmbientMetric(\ComparisonPoint_{n_k},\ComparisonPoint)
 \leq\AmbientMetric(\ComparisonPoint_{n_k},\ComparisonPoint_{0,{n_k}})
      +\AmbientMetric(\ComparisonPoint_{0,{n_k}},\ComparisonPoint)
 \longrightarrow0.
\]
Fix an arbitrary
$\CoefficientBound>0$.
For every
$k\in\NonnegativeIntegers$,
define the uniform objective error by
\begin{equation}\label{eq:local-uniform-objective-error}
 a_k=\sup_{\substack{\IntegrationPoint\in\AmbientCarrier\\
                  \CoefficientVector\in\CoefficientBall_\CoefficientBound}}
 |\ApproximationObjective_{n_k}(\IntegrationPoint,\CoefficientVector)
  -\ApproximationObjective(\IntegrationPoint,\CoefficientVector)|.
\end{equation}
For every
$k\in\NonnegativeIntegers$,
define the error due to the moving evaluation point by
\begin{equation}\label{eq:local-moving-point-error}
 b_k=\sup_{\CoefficientVector\in\CoefficientBall_\CoefficientBound}
 |\ApproximationObjective(\ComparisonPoint_{n_k},\CoefficientVector)
  -\ApproximationObjective(\ComparisonPoint,\CoefficientVector)|.
\end{equation}
By \eqref{eq:uniform-objectives},
$a_k\to0$.
Uniform continuity of
$\ApproximationObjective$
on the compact set
$\AmbientCarrier\times\CoefficientBall_\CoefficientBound$
and convergence of
$\ComparisonPoint_{n_k}$
to
$\ComparisonPoint$
imply
$b_k\to0$.
Thus the triangle inequality yields
\begin{equation}\label{eq:local-moving-objective-convergence}
 \sup_{\CoefficientVector\in\CoefficientBall_\CoefficientBound}
 |\ApproximationObjective_{n_k}(\ComparisonPoint_{n_k},\CoefficientVector)
  -\ApproximationObjective(\ComparisonPoint,\CoefficientVector)|
 \leq a_k+b_k\longrightarrow0.
\end{equation}
The convergence in \eqref{eq:local-moving-objective-convergence} holds for every fixed
$\CoefficientBound>0$,
so the objectives converge uniformly on every compact subset of
$\RealNumbers^\ModelDimension$.
For every $k$, the vector
$\FeatureCoordinates_{\ComparisonCarrier_{n_k}}(\ComparisonPoint_{n_k})$
minimizes $\ApproximationObjective_{n_k}(\ComparisonPoint_{n_k},\cdot)$.
By \eqref{eq:coefficient-bounds},
choose
$\CoefficientBound>0$
large enough that
$\CoefficientBall_\CoefficientBound$
contains all these minimizers.
Equation \eqref{eq:local-moving-objective-convergence} proves uniform convergence of the objectives on that ball.
The limit objective has the unique minimizer
$\FeatureCoordinates_\ComparisonCarrier(\ComparisonPoint)$.
\Cref{lem:minimizer-continuity} therefore implies
$\CoefficientVector=\FeatureCoordinates_\ComparisonCarrier(\ComparisonPoint)$.
But continuity of
$\FeatureCoordinates _\ComparisonCarrier $
also implies
$\FeatureCoordinates _\ComparisonCarrier (\ComparisonPoint _{0,{n_k} })\to \FeatureCoordinates _\ComparisonCarrier (\ComparisonPoint )$,
contradicting \eqref{eq:nonconvergent-coordinate-separation}.
This proves \ref{item:model-continuity}.

\ProofStep{Changes of basis and isometries.}\label{step:local-model-isometries}
\paraid{p-67fdf3abc74d}
Let
$\OrthogonalChange\in\OrthogonalGroup(\ModelDimension)$
and let
$\{\widehat\BasisFunction_\BasisIndex\}_{0\leq\BasisIndex<\ModelDimension}$
be a second orthonormal basis such that for every
$0\leq\BasisIndex<\ModelDimension$
we have
\[
 \widehat\BasisFunction_\BasisIndex
 =\sum_{\SubsequenceIndex=0}^{\ModelDimension-1}
   \OrthogonalChange_{\BasisIndex\SubsequenceIndex}\BasisFunction_\SubsequenceIndex.
\]
For every
$\ComparisonPoint\in\ComparisonCarrier$
and
$\CoefficientVector\in\RealNumbers^\ModelDimension$,
the coordinate map
$\widehat\FeatureCoordinates_\ComparisonCarrier$
and the norm
$\widehat\NormSymbol_\ComparisonCarrier$
defined using the second basis satisfy
\[
 \widehat \FeatureCoordinates _\ComparisonCarrier (\ComparisonPoint )=\OrthogonalChange  \FeatureCoordinates _\ComparisonCarrier (\ComparisonPoint ),
\]
and
\[
 \widehat \NormSymbol _\ComparisonCarrier (\CoefficientVector )=\NormSymbol _\ComparisonCarrier (\OrthogonalChange ^{-1}\CoefficientVector ).
\]
Consequently,
$[\FeatureCoordinates_\ComparisonCarrier,\NormSymbol_\ComparisonCarrier]_{\OrthogonalGroup(\ModelDimension)}$
is independent of the orthonormal basis used in Step~\ref{step:local-model-coordinates}.
Define
\[
 \LocalModelClass{\ComparisonCarrier}{\ComparisonMetric}
 =[\FeatureCoordinates_\ComparisonCarrier,\NormSymbol_\ComparisonCarrier]_{\OrthogonalGroup(\ModelDimension)}.
\]
The coordinate pairs
$(\FeatureCoordinates_{\ComparisonCarrier_\SequenceIndex},
  \NormSymbol_{\ComparisonCarrier_\SequenceIndex})$
and
$(\FeatureCoordinates_\ComparisonCarrier,\NormSymbol_\ComparisonCarrier)$
constructed using the bases from Step~\ref{step:local-model-norm-convergence} represent
$\LocalModelClass{\ComparisonCarrier_\SequenceIndex}{\ComparisonMetric_\SequenceIndex}$
and
$\LocalModelClass{\ComparisonCarrier}{\ComparisonMetric}$,
respectively.
Their convergence in \eqref{eq:local-norm-convergence}
and \eqref{eq:local-coordinate-convergence}
proves \ref{item:model-continuity}.
For an isometry
$\InputIsometry \colon\MetricSpace{\ComparisonCarrier }{\ComparisonMetric }\to\MetricSpace{\ComparisonCarrier _0}{\ComparisonMetric _0}$,
define the pullback map
\[
 \IsometryPullback\colon
 \LebesgueSpace^2(\ComparisonCarrier_0,\SelectedMeasure{\ComparisonCarrier_0}{\ComparisonMetric_0})
 \to
 \LebesgueSpace^2(\ComparisonCarrier,\SelectedMeasure{\ComparisonCarrier}{\ComparisonMetric})
\]
as follows.
For every
$\HilbertFunction\in\LebesgueSpace^2(\ComparisonCarrier_0,\SelectedMeasure{\ComparisonCarrier_0}{\ComparisonMetric_0})$,
set
\[
 \IsometryPullback\HilbertFunction=\HilbertFunction\circ\InputIsometry.
\]
The map
$\IsometryPullback$
is unitary because
$\InputIsometry \Pushforward\SelectedMeasure{\ComparisonCarrier}{\ComparisonMetric}=\SelectedMeasure{\ComparisonCarrier _0}{\ComparisonMetric _0}$.
It intertwines the distance operators,
\[
 \DistanceOperator{\ComparisonCarrier}{\ComparisonMetric}(\HilbertFunction \circ \InputIsometry )
       =(\DistanceOperator{\ComparisonCarrier _0}{\ComparisonMetric _0}\HilbertFunction )\circ \InputIsometry .
\]
Thus
\[
 \IsometryPullback\bigl(\SpectralSubspace{\ComparisonCarrier_0}{\SpectralCutoff}\bigr)
 =\SpectralSubspace{\ComparisonCarrier}{\SpectralCutoff}.
\]
Let
$\{\BasisFunction_{0,\BasisIndex}\}_{0\leq\BasisIndex<\ModelDimension}$
be the orthonormal basis chosen on
$\ComparisonCarrier_0$.
Since both
$\{\BasisFunction_{0,\BasisIndex}\circ\InputIsometry\}_{\BasisIndex}$
and
$\{\BasisFunction_\BasisIndex\}_{\BasisIndex}$
are orthonormal bases of
$\SpectralSubspace{\ComparisonCarrier}{\SpectralCutoff}$,
there is
$\OrthogonalChange\in\OrthogonalGroup(\ModelDimension)$
such that for every
$0\leq\BasisIndex<\ModelDimension$
we have
\begin{equation}\label{eq:isometry-pullback-basis}
 \BasisFunction_{0,\BasisIndex}\circ\InputIsometry
 =\sum_{\SecondBasisIndex=0}^{\ModelDimension-1}
   \OrthogonalChange_{\BasisIndex\SecondBasisIndex}\BasisFunction_\SecondBasisIndex.
\end{equation}
Using \eqref{eq:isometry-pullback-basis} and
$\InputIsometry\Pushforward\SelectedMeasure{\ComparisonCarrier}{\ComparisonMetric}
 =\SelectedMeasure{\ComparisonCarrier_0}{\ComparisonMetric_0}$,
we obtain,
for every
$\ComparisonPoint\in\ComparisonCarrier$
and
$\CoefficientVector\in\RealNumbers^\ModelDimension$,
\begin{equation}\label{eq:isometry-objective-transform-1}
 \NormSymbol_{\ComparisonCarrier_0}(\CoefficientVector)
 =\NormSymbol_\ComparisonCarrier(\OrthogonalChange^{-1}\CoefficientVector),
\end{equation}
and
\begin{equation}\label{eq:isometry-objective-transform-2}
 \ApproximationObjective_{\ComparisonCarrier_0}
   (\InputIsometry(\ComparisonPoint),\CoefficientVector)
 =\ApproximationObjective_\ComparisonCarrier
   (\ComparisonPoint,\OrthogonalChange^{-1}\CoefficientVector).
\end{equation}
The identity \eqref{eq:isometry-objective-transform-2}
and uniqueness of the minimizers imply
\[
 \FeatureCoordinates_{\ComparisonCarrier_0}(\InputIsometry(\ComparisonPoint))
 =\OrthogonalChange\FeatureCoordinates_\ComparisonCarrier(\ComparisonPoint).
\]
This proves both identities in \eqref{eq:input-isometry-1} and \eqref{eq:input-isometry-2}.
Consequently,
for every
$\ComparisonPoint,\IntegrationPoint\in\ComparisonCarrier$,
\[
 \begin{aligned}
 \Pseudometric_{\ComparisonCarrier_0}
   (\InputIsometry(\ComparisonPoint),\InputIsometry(\IntegrationPoint))
 &=\NormSymbol_{\ComparisonCarrier_0}
   \bigl(\OrthogonalChange(\FeatureCoordinates_\ComparisonCarrier(\ComparisonPoint)
                         -\FeatureCoordinates_\ComparisonCarrier(\IntegrationPoint))\bigr)\\
 &=\Pseudometric_\ComparisonCarrier(\ComparisonPoint,\IntegrationPoint).
 \end{aligned}
\]
Thus the pseudometric
$\Pseudometric_\ComparisonCarrier$
is independent of the coordinates and is invariant under isometries of
$\MetricSpace{\ComparisonCarrier}{\ComparisonMetric}$.

\ProofStep{Continuity of the approximation error.}\label{step:local-model-error-continuity}
\paraid{p-2d2c5d189dec}
Equation \eqref{eq:local-model-pseudometric} defines a continuous pseudometric
because it pulls back the distance of a normed space by the continuous map
$\FeatureCoordinates_\ComparisonCarrier$.
Step~\ref{step:local-model-isometries} proves its invariance.
Set
\[
 \NormComparisonError_\SequenceIndex
 =\sup_{\substack{\CoefficientVector\in\RealNumbers^\ModelDimension\\\norm{\CoefficientVector}_2=1}}
       |\NormSymbol_{\ComparisonCarrier_\SequenceIndex}(\CoefficientVector)
          -\NormSymbol_\ComparisonCarrier(\CoefficientVector)|,
\]
and
\[
 \CoordinateComparisonError_\SequenceIndex
 =\sup_{(\ComparisonPoint_\SequenceIndex,\ComparisonPoint)\in\Correspondence_\SequenceIndex}
       \norm{\FeatureCoordinates_{\ComparisonCarrier_\SequenceIndex}(\ComparisonPoint_\SequenceIndex)
               -\FeatureCoordinates_\ComparisonCarrier(\ComparisonPoint)}_2.
\]
Steps~\ref{step:local-model-norm-convergence} and~\ref{step:local-model-coordinate-convergence} show that both numbers tend to
$0$.
Choose
$\CoefficientBound>0$
which bounds all the coordinate norms in \eqref{eq:coefficient-bounds}
along the sequence and on the limit space.
Such a bound exists because the diameters converge.
By \Cref{lem:coordinate-pseudometric-estimate},
\begin{equation}\label{eq:local-norm-pair-estimate}
 \PseudometricError_{\Correspondence_\SequenceIndex}
   (\Pseudometric_{\ComparisonCarrier_\SequenceIndex},\Pseudometric_\ComparisonCarrier)
 \leq2\CoefficientBound\NormComparisonError_\SequenceIndex
       +2\CoordinateComparisonError_\SequenceIndex
          \max_{\substack{\CoefficientVector\in\RealNumbers^\ModelDimension\\\norm{\CoefficientVector}_2=1}}\NormSymbol_\ComparisonCarrier(\CoefficientVector).
\end{equation}
The right-hand side tends to
$0$.
Thus
\begin{equation}\label{eq:local-pseudometric-convergence}
 \PseudometricError_{\Correspondence_\SequenceIndex}
 (\Pseudometric_{\ComparisonCarrier_\SequenceIndex},\Pseudometric_\ComparisonCarrier)\to0.
\end{equation}
If
$\AmbientDisplacement _\SequenceIndex =\sup_{(\ComparisonPoint_\SequenceIndex,\ComparisonPoint)\in\Correspondence_\SequenceIndex}\AmbientMetric (\ComparisonPoint _\SequenceIndex ,\ComparisonPoint )$,
then for every
$(\ComparisonPoint_\SequenceIndex,\ComparisonPoint),
 (\IntegrationPoint_\SequenceIndex,\IntegrationPoint)\in\Correspondence_\SequenceIndex$,
\[
 |\ComparisonMetric _\SequenceIndex(\ComparisonPoint _\SequenceIndex ,\IntegrationPoint _\SequenceIndex )-\ComparisonMetric (\ComparisonPoint ,\IntegrationPoint )|\leq2\AmbientDisplacement _\SequenceIndex .
\]
By \eqref{eq:pseudometric-uniform-comparison}
and \eqref{eq:local-pseudometric-convergence},
we obtain
\[
 |\LocalError \MetricSpace{\ComparisonCarrier _\SequenceIndex }{\ComparisonMetric _\SequenceIndex}-\LocalError \MetricSpace{\ComparisonCarrier }{\ComparisonMetric }|
 \leq\PseudometricError_{\Correspondence_\SequenceIndex}
 (\Pseudometric_{\ComparisonCarrier_\SequenceIndex},\Pseudometric_\ComparisonCarrier)
 +2\AmbientDisplacement _\SequenceIndex \to0.
\]
The bound at
$\MetricSpace{\BaseCarrier }{\MetricSymbol }$
was arranged in Step~\ref{step:local-model-spectral-dimension} by \eqref{eq:finite-spectral-error}.
This proves \ref{item:model-pseudometric}--\ref{item:model-center-error}.
\end{proof}

\paraid{p-ec88a14da3e5}
The construction also represents the isometries of each input space by
orthogonal matrices.

\begin{remark}\label{rem:local-isometry-representation}
\paraid{p-a393c40d9c5a}
Fix a local model from \Cref{thm:local-models} and a representative
$(\FeatureCoordinates_\ComparisonCarrier,\NormSymbol_\ComparisonCarrier)$
constructed from the real orthonormal basis
$\{\BasisFunction_\BasisIndex\}_{0\leq\BasisIndex<\ModelDimension}$
of
$\SpectralSubspace{\ComparisonCarrier}{\SpectralCutoff}$.
Then the construction induces a continuous group homomorphism
\begin{equation}\label{eq:local-isometry-representation-1}
 \IsometryRepresentation{\ComparisonCarrier}{\ComparisonMetric}
 \colon\Isom\MetricSpace{\ComparisonCarrier}{\ComparisonMetric}
       \to\OrthogonalGroup(\ModelDimension),
\end{equation}
with the following matrix entries.
For every
$\FirstIsometry\in\Isom\MetricSpace{\ComparisonCarrier}{\ComparisonMetric}$
and integers
$0\leq\BasisIndex,\SecondBasisIndex<\ModelDimension$,
set
\begin{equation}\label{eq:local-isometry-representation-2}
 \bigl(\IsometryRepresentation{\ComparisonCarrier}{\ComparisonMetric}(\FirstIsometry)\bigr)_{\BasisIndex\SecondBasisIndex}
 =\int_\ComparisonCarrier
   \BasisFunction_\BasisIndex(\FirstIsometry\ComparisonPoint)
   \BasisFunction_\SecondBasisIndex(\ComparisonPoint)
   \,\IntegrationDifferential\SelectedMeasure{\ComparisonCarrier}{\ComparisonMetric}(\ComparisonPoint).
\end{equation}
Indeed,
Step~\ref{step:local-model-isometries} of the proof shows that composition with
$\FirstIsometry$
preserves the spectral subspace and its inner product.
Hence
\begin{equation}\label{eq:isometry-basis-action}
 \BasisFunction_\BasisIndex\circ\FirstIsometry
 =\sum_{\SecondBasisIndex=0}^{\ModelDimension-1}
  \bigl(\IsometryRepresentation{\ComparisonCarrier}{\ComparisonMetric}(\FirstIsometry)\bigr)_{\BasisIndex\SecondBasisIndex}
  \BasisFunction_\SecondBasisIndex.
\end{equation}
For every
$\FirstIsometry,\SecondIsometry\in\Isom\MetricSpace{\ComparisonCarrier}{\ComparisonMetric}$,
applying \eqref{eq:isometry-basis-action} successively to
$\FirstIsometry$
and
$\SecondIsometry$
and using uniqueness of basis coefficients proves
\[
 \IsometryRepresentation{\ComparisonCarrier}{\ComparisonMetric}(\FirstIsometry\circ\SecondIsometry)
 =\IsometryRepresentation{\ComparisonCarrier}{\ComparisonMetric}(\FirstIsometry)
  \IsometryRepresentation{\ComparisonCarrier}{\ComparisonMetric}(\SecondIsometry).
\]
The matrix entries in \eqref{eq:local-isometry-representation-2} are continuous
for the uniform topology on the isometry group.
By \eqref{eq:input-isometry-1},
\eqref{eq:input-isometry-2} and uniqueness of the best approximations,
we obtain the equivariance identity \eqref{eq:local-model-equivariance-1}
and the norm preservation identity \eqref{eq:local-model-equivariance-2}.
For every
$\FirstIsometry\in\Isom\MetricSpace{\ComparisonCarrier}{\ComparisonMetric}$
and
$\ComparisonPoint\in\ComparisonCarrier$,
\begin{equation}\label{eq:local-model-equivariance-1}
 \FeatureCoordinates_\ComparisonCarrier(\FirstIsometry\ComparisonPoint)
 =\IsometryRepresentation{\ComparisonCarrier}{\ComparisonMetric}(\FirstIsometry)
   \FeatureCoordinates_\ComparisonCarrier(\ComparisonPoint).
\end{equation}
For every
$\FirstIsometry\in\Isom\MetricSpace{\ComparisonCarrier}{\ComparisonMetric}$
and
$\CoefficientVector\in\RealNumbers^\ModelDimension$,
\begin{equation}\label{eq:local-model-equivariance-2}
 \NormSymbol_\ComparisonCarrier
   \bigl(\IsometryRepresentation{\ComparisonCarrier}{\ComparisonMetric}(\FirstIsometry)\CoefficientVector\bigr)
 =\NormSymbol_\ComparisonCarrier(\CoefficientVector).
\end{equation}
A change of basis by
$\OrthogonalChange\in\OrthogonalGroup(\ModelDimension)$
replaces this representation by
$\FirstIsometry\mapsto\OrthogonalChange
 \IsometryRepresentation{\ComparisonCarrier}{\ComparisonMetric}(\FirstIsometry)
 \OrthogonalChange^{-1}$.
For the constant model used at a singleton center,
take the trivial homomorphism into
$\OrthogonalGroup(1)$.
For a compact group
$\ActingGroup$,
the Peter--Weyl theorem asserts that the matrix coefficients of finite-dimensional
continuous unitary representations have dense linear span in
$\ContinuousFunctions(\ActingGroup,\ComplexNumbers)$
\cite[Chapter~6, Section~2, pp.~165--167]{DiestelSpalsbury2014}.
In that proof,
convolution operators commute with left translations,
and their nonzero eigenspaces carry finite-dimensional unitary representations.
Here the distance operator commutes with the action of
$\Isom\MetricSpace{\ComparisonCarrier}{\ComparisonMetric}$,
so its invariant spectral subspace defines
$\IsometryRepresentation{\ComparisonCarrier}{\ComparisonMetric}$.
\Cref{lem:spectral-continuity} permits comparison of these subspaces as the compact
metric space varies.
\end{remark}

\begin{remark}\label{rem:mds-comparison}
\paraid{p-1384428ed1f6}
For compact metric measure spaces,
Kroshnin,
Stepanov,
and Trevisan
\cite[arXiv v2, Theorem~5.4 and Remark~5.5]{KroshninStepanovTrevisan2022}
prove stability of multidimensional scaling coordinates in
$\LebesgueSpace^2$
using optimal couplings for the Gromov--Wasserstein distance of order four
and minimizing over orthogonal transformations,
assuming convergence for that distance and a gap
after the last retained positive eigenvalue.
\Cref{thm:local-models} uses best
$\LebesgueSpace^\IntegrabilityExponent$
approximations of distance functions and varying norms to approximate
the metric uniformly.
\end{remark}

% Retired paragraph IDs from former Subsection 8.5; see SAVED-MATERIALS.md.
% \paraid{p-aac7ee003672}
% \paraid{p-39b097fcf654}
% \paraid{p-0c5e017edead}
% \paraid{p-09fe39859c20}
% \paraid{p-f72bfed56be6}
% \paraid{p-977d89fb2970}
% \paraid{p-9d78f01958f0}
% \paraid{p-146a1d37efcc}
% \paraid{p-2c2e9f81b579}
% \paraid{p-4ced780973fc}
% \paraid{p-ba88bcbd74fe}
% \paraid{p-1c044364faca}
% \paraid{p-65c05f32842b}
% \paraid{p-ee56c954b4dc}
% \paraid{p-b8072b65a2b8}
% \paraid{p-48feee0285c6}
% \paraid{p-20f8b1991114}
% \paraid{p-2480d2fdfab2}
% \paraid{p-5de9462f1fdb}
% \paraid{p-34542b375466}
% \paraid{p-0720c8effdbc}
% \paraid{p-1e670d2fae30}
% \paraid{p-a2653c1926fd}
% Retired paragraph IDs from former Subsection 8.6; see SAVED-MATERIALS.md.
% \paraid{p-b13e5b3c82a4}
% \paraid{p-8f0b4175df73}
% \paraid{p-e2bafde4496d}
% \paraid{p-b75d4bf8c5ae}
% \paraid{p-c4288955bcdb}

\clearpage
\part{The topology of Gromov--Hausdorff space III: The Gromov--Hausdorff space is an absolute retract}\label{part:topology}
\begin{quote}
\small
\noindent\textbf{Abstract.}
\paraid{p-fedf4665176f}
We prove that the Gromov--Hausdorff space is an absolute retract for all metrizable spaces. Starting from continuous finite-dimensional local approximations, we represent the varying norms in fixed Banach spaces and combine the local maps using a partition of unity. We use hyperspaces modulo compact groups to construct absolute retracts through which we approximate the identity map with arbitrarily small homotopy tracks. By Hanner's domination theorem, we obtain the absolute neighborhood retract property. Since the Gromov--Hausdorff space is contractible, we conclude that the Gromov--Hausdorff space is an absolute retract.
\par\smallskip
\noindent\textbf{Keywords.} Gromov--Hausdorff space, absolute retract, absolute extensor, hyperspace, approximation.

\par\smallskip
\noindent\textbf{2020 Mathematics Subject Classification.} Primary 54C55; Secondary 54E35, 54B20.

\end{quote}

\section{Introduction to Part III}\label{sec:intro-topology}
\paraid{p-61c139e831ec}
The Gromov--Hausdorff space is a global analogue of a hyperspace.
For a fixed metric space,
the hyperspace of its nonempty compact subsets is equipped with the
Hausdorff distance.
In the Gromov--Hausdorff space
$\GHSpace$,
we consider the isometry classes of all nonempty compact metric spaces
and allow the ambient metric space to vary.
The ordinary unpointed
Gromov--Hausdorff distance
$\GHDistance$
is the infimum of the Hausdorff distances between isometric copies
in all common ambient metric spaces.



\paraid{p-f1859d6a167a}
Curtis
\cite[Theorem~1.6]{Curtis1980}
proves,
in particular,
that the hyperspace of nonempty compact subsets of a Banach space,
equipped with the Hausdorff metric of its norm,
is an absolute retract and hence an absolute neighborhood retract.
This result for a fixed ambient space motivates the analogous question
for the global space
$\GHSpace$.
Antonyan
\cite[p.~2]{Antonyan2020Euclidean}
asks whether
$\GHSpace$
is an absolute retract.
In this part,
we answer this question affirmatively by proving the following theorem.

\begin{theorem}\label{thm:part-iii-main}
\paraid{p-11b6e24f7baa}
The space
$\GHSpace$
is an absolute retract for all metrizable spaces.
\end{theorem}



\paraid{p-4df38d2e7f6c}
As a related result,
for each positive integer
$\FiniteSize$,
Antonyan identifies the subspace represented by compact subsets of
$\RealNumbers^\FiniteSize$
with the Hilbert cube minus a point
\cite[Corollary~5.3]{Antonyan2021Euclidean}.
Hyperspaces also appear in the study of Gromov--Hausdorff geodesics.
M\'emoli and Wan
\cite[arXiv v2, Theorem~1]{MemoliWan2023}
prove that every Gromov--Hausdorff geodesic can be realized as a Hausdorff
geodesic of compact subsets of one compact metric space.

\paraid{p-1cc29943ea5e}
For metrics on a fixed metrizable space,
the author proved an interpolation theorem with a uniform error bound
for prescribed metrics on a discrete family of closed subsets
\cite[Theorem~1.1]{Ishiki2026Interpolation}.
In this paper,
the domain of an extension problem maps into
$\GHSpace$,
so the underlying compact spaces vary with the parameter.
The coordinate pairs in the local models can be chosen to satisfy
the norm and coordinate convergence in \ref{item:model-continuity}
of \Cref{thm:local-models}.
We preserve their norms when placing the models in Banach spaces,
so that the partition of unity controls the error in all pairwise distances.


\paraid{p-402bb42dd02f}
To prove that
$\GHSpace$
is an ANR,
we use Hanner's domination theorem
(\Cref{thm:hanner-domination}).
Roughly speaking,
this theorem states that a separable metric space is an ANR if its identity
map can be deformed to a map factoring through an ANR by a homotopy
whose point tracks have arbitrarily small diameter.
For
$\GHSpace$,
we construct the auxiliary ANR as the quotient of the hyperspace of a
Banach space by a compact group action.
The Banach space and the group are obtained from the local models,
and the absolute retract property of this quotient follows from
Antonyan's results on equivariant hyperspaces and orbit spaces.

\paraid{p-433e89812f55}
Fix a continuous error function
$\ErrorControl\colon\GHSpace\to(0,\infty)$.
Choose local models whose errors are less than
$\ErrorControl$
on their domains,
and choose a countable locally finite partition of unity
$\{\PartitionWeight_\CoverIndex\}_{\CoverIndex\in\NonnegativeIntegers}$
subordinate to these domains.
For each
$\CoverIndex\in\NonnegativeIntegers$,
fix a local model of dimension
$\ModelDimension_\CoverIndex\geq1$
on an open set
$\ModelDomain_\CoverIndex$
containing
$\supp\PartitionWeight_\CoverIndex$.
For every
$\CoverIndex\in\NonnegativeIntegers$
and every compact metric representative
$\MetricSpace{\ComparisonCarrier}{\ComparisonMetric}$
whose isometry class belongs to
$\ModelDomain_\CoverIndex$,
choose a coordinate pair for the local model and write its point map and norm as
\[
 \FeatureCoordinates_{\CoverIndex,\ComparisonCarrier}\colon
 \ComparisonCarrier\to\RealNumbers^{\ModelDimension_\CoverIndex},
 \qquad
 \NormSymbol_{\CoverIndex,\ComparisonCarrier}\colon
 \RealNumbers^{\ModelDimension_\CoverIndex}\to[0,\infty).
\]

\paraid{p-b1d3e3825e91}
We first represent the varying local norms in fixed Banach spaces.
For each integer
$\ModelDimension\geq1$,
let
$\UnitSphere^{\ModelDimension-1}$
be the Euclidean unit sphere and set
\[
 \BanachAmbientSpace_\ModelDimension
 =\ContinuousFunctions(\UnitSphere^{\ModelDimension-1},\RealNumbers),
\]
with the uniform norm.
For every integer
$\ModelDimension\geq1$
and every norm
$\NormSymbol$
on
$\RealNumbers^\ModelDimension$,
let
$\NormSymbol\DualNorm$
denote its dual norm with respect to the Euclidean inner product.
\Cref{lem:dual-representation} constructs the linear isometric embedding
\[
 \NormEmbedding_\NormSymbol\colon
 (\RealNumbers^\ModelDimension,\NormSymbol)
 \to\BanachAmbientSpace_\ModelDimension
\]
defined,
for
$\FirstVector\in\RealNumbers^\ModelDimension$
and
$\SpherePoint\in\UnitSphere^{\ModelDimension-1}$,
by
\[
 \NormEmbedding_\NormSymbol(\FirstVector)(\SpherePoint)
 =\frac{\langle\FirstVector,\SpherePoint\rangle}
       {\NormSymbol\DualNorm(\SpherePoint)}.
\]
We take the
$\SummableNorm$
direct sum of these spaces over the local models and let the product
of the corresponding orthogonal groups act on it.
Thus we set
\[
 \BanachAmbientSpace
 =\left(\bigoplus_{\CoverIndex\in\NonnegativeIntegers}
        \BanachAmbientSpace_{\ModelDimension_\CoverIndex}\right)_{\SummableNorm},
\]
and
\[
 \ActingGroup
 =\prod_{\CoverIndex\in\NonnegativeIntegers}
   \OrthogonalGroup(\ModelDimension_\CoverIndex).
\]
Each orthogonal group acts on the corresponding space of functions
by precomposition with inverse orthogonal transformations of the sphere.
The resulting action of the compact metrizable group
$\ActingGroup$
on the separable Banach space
$\BanachAmbientSpace$
is continuous and consists of linear isometries.
Equation~\eqref{eq:dual-equivariance} makes the embeddings compatible
with orthogonal changes of the local norms and coordinates.

\paraid{p-b04298598a10}
Let
$\CompactHyperspace(\BanachAmbientSpace)$
be the space of nonempty compact subsets of
$\BanachAmbientSpace$,
with the Hausdorff metric induced by its norm,
and set
\[
 \AuxiliaryAR=\CompactHyperspace(\BanachAmbientSpace)/\ActingGroup.
\]
The absolute retract property of
$\AuxiliaryAR$
follows from two results of Antonyan.
Since
$\BanachAmbientSpace$
is connected and locally continuum-connected,
his equivariant hyperspace theorem
\cite[Proposition~3.1]{Antonyan2003West},
with the corrigendum \cite{Antonyan2006Correction},
implies that
$\CompactHyperspace(\BanachAmbientSpace)$
is an equivariant absolute retract,
or a
$\ActingGroup$-AR
(\Cref{thm:equivariant-hyperspace}).
His theorem on orbit spaces
\cite[Theorem~8 and Corollary~1]{Antonyan1990}
then implies that
$\AuxiliaryAR$
is an AR because
$\ActingGroup$
is compact with a countable basis
(\Cref{thm:orbit-retract}).
\Cref{lem:hyperspace-orbit-ar} also proves that
$\AuxiliaryAR$
is separable and metrizable.
In particular,
$\AuxiliaryAR$
is the ANR through which we will approximate the identity of
$\GHSpace$.

\paraid{p-8e3de8ba284c}
We combine the local point maps using the partition of unity.
For each compact metric representative
$\MetricSpace{\ComparisonCarrier}{\ComparisonMetric}$,
define
$\JointFeatureMap_\ComparisonCarrier\colon
\ComparisonCarrier\to\BanachAmbientSpace$
as follows.
For every
$\CoverIndex\in\NonnegativeIntegers$
and
$\ComparisonPoint\in\ComparisonCarrier$,
set
\[
 \bigl(\JointFeatureMap_\ComparisonCarrier(\ComparisonPoint)\bigr)_\CoverIndex
 =
 \begin{cases}
 \PartitionWeight_\CoverIndex\MetricSpace{\ComparisonCarrier}{\ComparisonMetric}
 \NormEmbedding_{\NormSymbol_{\CoverIndex,\ComparisonCarrier}}
   (\FeatureCoordinates_{\CoverIndex,\ComparisonCarrier}(\ComparisonPoint))
 &\text{if }\PartitionWeight_\CoverIndex
   \MetricSpace{\ComparisonCarrier}{\ComparisonMetric}>0,\\
 0&\text{otherwise}.
 \end{cases}
\]
For fixed
$\MetricSpace{\ComparisonCarrier}{\ComparisonMetric}$,
only finitely many coordinate maps are nonzero.
Thus
$\JointFeatureMap_\ComparisonCarrier$
is continuous and its image
$\CompactFeatureImage_\ComparisonCarrier
 =\JointFeatureMap_\ComparisonCarrier(\ComparisonCarrier)$
is nonempty and compact.
Changes of local bases and isometric representatives change
$\CompactFeatureImage_\ComparisonCarrier$
by elements of
$\ActingGroup$.
Thus its orbit is well defined by the input isometry class.
This construction applies to both trivial and nontrivial isometry groups,
although spaces with trivial isometry group are dense in
$\GHSpace$
by Rouyer's results
\cite[arXiv v1, Theorems~2 and~4]{Rouyer2011}.
Writing
$\OrbitProjection\colon\CompactHyperspace(\BanachAmbientSpace)\to\AuxiliaryAR$
for the orbit map,
we construct in \Cref{thm:global-domination} continuous maps
\[
 \GHSpace\xrightarrow{\ \ApproximationMap\ }
 \AuxiliaryAR
 \xrightarrow{\ \RealizationMap\ }\GHSpace.
\]
For every compact metric representative
$\MetricSpace{\ComparisonCarrier}{\ComparisonMetric}$,
the first map satisfies
\[
 \ApproximationMap\MetricSpace{\ComparisonCarrier}{\ComparisonMetric}
 =\OrbitProjection(\CompactFeatureImage_\ComparisonCarrier).
\]
The second map sends the orbit of a nonempty compact subset
$K\subset\BanachAmbientSpace$
to its isometry class with the metric induced by the norm of
$\BanachAmbientSpace$.

\paraid{p-5e91fc971616}
\Cref{thm:global-domination} also constructs a continuous homotopy
\[
 \ApproximationHomotopy\colon\GHSpace\times\UnitInterval\to\GHSpace.
\]
For every compact metric representative
$\MetricSpace{\ComparisonCarrier}{\ComparisonMetric}$,
its endpoints satisfy
\[
 \ApproximationHomotopy(\MetricSpace{\ComparisonCarrier}{\ComparisonMetric},0)
 =\MetricSpace{\ComparisonCarrier}{\ComparisonMetric},
 \qquad
 \ApproximationHomotopy(\MetricSpace{\ComparisonCarrier}{\ComparisonMetric},1)
 =\RealizationMap\ApproximationMap
   \MetricSpace{\ComparisonCarrier}{\ComparisonMetric}.
\]
For every compact metric representative
$\MetricSpace{\ComparisonCarrier}{\ComparisonMetric}$,
the point track satisfies
\begin{equation}\label{eq:part-iii-domination-3}
 \diam_{\GHDistance}
 \bigl\{\ApproximationHomotopy(\MetricSpace{\ComparisonCarrier}{\ComparisonMetric},t)
       \bigm| t\in\UnitInterval\bigr\}
 <\ErrorControl\MetricSpace{\ComparisonCarrier}{\ComparisonMetric}.
\end{equation}
For each
$\ErrorTolerance>0$,
we apply this construction with
$\ErrorControl\equiv\ErrorTolerance$.
Hanner's domination theorem
(\Cref{thm:hanner-domination})
and \eqref{eq:part-iii-domination-3} then imply that
$\GHSpace$
is an ANR for all separable metrizable spaces.
\Cref{thm:anr-category} extends this conclusion to all metrizable spaces.
Finally,
scaling the metrics contracts
$\GHSpace$
to the one-point space.
The ANR/ANE and AR/AE equivalences in \Cref{thm:retract-extensor},
together with \Cref{thm:contractible-ane},
therefore imply the absolute retract property in
\Cref{thm:part-iii-main}.

\medskip
\noindent\textbf{Dependence on the other parts.}
\paraid{p-065600b3665d}
\Cref{thm:local-models} is the input from Part~\ref{part:spectral}.
The continuity assertion \ref{item:model-continuity}
and the isometry identities \eqref{eq:input-isometry-1} and \eqref{eq:input-isometry-2}
in \Cref{thm:local-models} are used in the global construction.
The dependence on Part~\ref{part:measures} is through \Cref{thm:local-models}.
Part~\ref{part:hilbert} uses the absolute retract property established in
this part together with an independently constructed discrete
approximation to identify
$\GHSpace$
with Hilbert space.

\medskip
\noindent\textbf{Organization.}
\paraid{p-fc2cd8290ce7}
In Section~\ref{sec:prelim-topology},
we recall the equivariant hyperspace results and the domination and
extension theorems used in the proof.
In Subsection~\ref{subsec:ar-norm-representations},
we represent varying finite-dimensional norms isometrically
in a fixed Banach space,
with compatibility under orthogonal changes of coordinates.
In Subsection~\ref{subsec:ar-hyperspace-orbits},
we construct absolute retracts from hyperspaces modulo
compact group actions and compare their orbit metrics with the
Gromov--Hausdorff distance.
In Subsection~\ref{subsec:ar-global-approximation},
we use a partition of unity to combine local models into
global approximations with small homotopy tracks.
In Subsection~\ref{subsec:ar-extension},
we apply Hanner's
\Cref{thm:hanner-domination,thm:anr-category}
and use contractibility to obtain extension from
closed subsets of arbitrary metrizable spaces,
proving \Cref{thm:absolute-extensor}.

\medskip
\noindent\textbf{Conventions and notation.}
\paraid{p-97bafc9921fd}
All compact metric spaces are nonempty,
and Banach and Hilbert spaces are real.
We identify compact metric spaces with their isometry classes in
$\GHSpace$.
An isometry is a surjective isometric embedding.
For a metric space
$\MetricSpace{\BaseCarrier}{\MetricSymbol}$,
we write
$\CompactHyperspace(\BaseCarrier)$
for its nonempty compact subsets,
equipped with the Hausdorff metric
$\HausdorffDistance{\MetricSymbol}$.
For a compact space with an isometric embedding into an ambient metric space,
we identify the compact space with its image and use the restricted ambient metric.
The notation
$\IdentityMap$
denotes the identity map on its specified domain.
Sequences are indexed by
$\NonnegativeIntegers=\{0,1,2,\ldots\}$.


\section{Preliminaries}\label{sec:prelim-topology}
\paraid{p-0a5e363880c8}
We recall the equivariant hyperspace results and the domination and
extension theorems used in this part.
The terms AR,
ANR,
AE,
and ANE have the meanings fixed in Definition~\ref{def:extension}.

\subsection{Equivariant hyperspaces and orbit spaces}\label{subsec:ar-preliminaries-hyperspaces}
\paraid{p-44f297e51ff0}
A metrizable
$\ActingGroup$-space
is a metrizable space with a continuous action of the topological group
$\ActingGroup$.
For metrizable
$\ActingGroup$-spaces
$\BaseCarrier$ and
$\ComparisonCarrier$,
a continuous map
$\Mapping:\BaseCarrier\to\ComparisonCarrier$
is equivariant if, for every
$\GroupElement\in\ActingGroup$
and
$\BasePoint\in\BaseCarrier$,
\[
 \Mapping(\GroupElement\cdot\BasePoint)
 =\GroupElement\cdot\Mapping(\BasePoint).
\]
A metrizable
$\ActingGroup$-space
$\BaseCarrier$
is a
$\ActingGroup$-AR
if every closed equivariant embedding
$\FirstEmbedding:\BaseCarrier\to\ComparisonCarrier$
into a metrizable
$\ActingGroup$-space
$\ComparisonCarrier$
admits a continuous equivariant retraction
$\EquivariantRetraction:\ComparisonCarrier\to\BaseCarrier$
such that
\[
 \EquivariantRetraction\circ\FirstEmbedding=\IdentityMap_\BaseCarrier.
\]


A space
$\BanachAmbientSpace$
is \emph{locally continuum-connected} if for every
$\BasePoint\in\BanachAmbientSpace$
and every open set
$\OpenSubset\subset\BanachAmbientSpace$
with
$\BasePoint\in\OpenSubset$,
there is an open set
$\InnerNeighborhood$
such that
\[
 \BasePoint\in\InnerNeighborhood\subset\OpenSubset,
 \qquad
 \forall\ComparisonPoint\in\InnerNeighborhood\quad
 \exists\JoiningContinuum\subset\BanachAmbientSpace\quad
 \begin{cases}
 \JoiningContinuum\text{ is compact and connected},\\
 \{\BasePoint,\ComparisonPoint\}\subset\JoiningContinuum\subset\OpenSubset.
 \end{cases}
\]
In a Banach space we can take
$\InnerNeighborhood$
to be a sufficiently small open ball and use the segments
$\JoiningContinuum=\{(1-\HomotopyTime)\BasePoint+\HomotopyTime\ComparisonPoint
 \mid\HomotopyTime\in\UnitInterval\}$.

\paraid{p-c196bd37c296}
For an action of a group $\ActingGroup$ on a space $\BanachAmbientSpace$,
the induced action on $\CompactHyperspace(\BanachAmbientSpace)$ is defined as follows.
For every $\GroupElement\in\ActingGroup$ and
$\FirstCompactSet\in\CompactHyperspace(\BanachAmbientSpace)$, set
\[
 \GroupElement\FirstCompactSet
 =\{\GroupElement\FirstVector\mid\FirstVector\in\FirstCompactSet\}.
\]
The following hyperspace theorem is
\cite[Proposition~3.1]{Antonyan2003West},
with the corrigendum \cite{Antonyan2006Correction}.
It is also stated in \cite[Theorem~2.3]{Antonyan2020Euclidean}.

\begin{theorem}[{\cite[Proposition~3.1]{Antonyan2003West}}]\label{thm:equivariant-hyperspace}
\paraid{p-d6a3d16c8e27}
Let a compact group
$\ActingGroup$
act continuously on a nonempty connected,
locally continuum-connected metrizable space
$\BanachAmbientSpace$.
Then
$\CompactHyperspace(\BanachAmbientSpace)$,
with its Vietoris topology and this induced action,
is a
$\ActingGroup$-AR.
\end{theorem}

\begin{theorem}[{\cite[Theorem~8 and Corollary~1]{Antonyan1990}}]\label{thm:orbit-retract}
\paraid{p-687bedbec7c3}
Let
$\ActingGroup$
be a compact group with a countable basis and let
$\BanachAmbientSpace$
be a metrizable
$\ActingGroup$-AR.
Then its orbit space
$\BanachAmbientSpace/\ActingGroup$
is an AR for
all metrizable spaces.
\end{theorem}

\subsection{Domination and absolute extension}\label{subsec:ar-preliminaries-domination}
\paraid{p-e539d8246cc3}
We first use domination by ANRs to obtain the ANR property for
all separable
metrizable spaces.
For a contractible space,
we obtain the AE property.
Roughly speaking,
the next theorem states that a space is an ANR if,
for every
$\ErrorTolerance>0$,
its identity map can be deformed to a map factoring through an ANR
by a homotopy whose point tracks have diameter less than
$\ErrorTolerance$.

\begin{theorem}[{\cite[Theorem~7.2(b)]{Hanner1951}}]\label{thm:hanner-domination}
\paraid{p-91db7d801e6b}
Let
$\MetricSpace{\RecognitionTarget}{\TargetMetric}$
be a separable metric space.
Assume that for every
$\ErrorTolerance>0$
there exist a separable metrizable ANR
$\AuxiliaryAR$
and continuous maps
\[
 \ApproximationMap\colon\RecognitionTarget\to\AuxiliaryAR,
 \quad\RealizationMap\colon\AuxiliaryAR\to\RecognitionTarget,
 \quad\ApproximationHomotopy\colon\RecognitionTarget\times\UnitInterval\to\RecognitionTarget
\]
such that for every
$\DomainPoint\in\RecognitionTarget$,
\[
 \ApproximationHomotopy(\DomainPoint,0)=\DomainPoint,
 \quad\ApproximationHomotopy(\DomainPoint,1)=\RealizationMap\ApproximationMap(\DomainPoint),
 \quad
 \diam_{\TargetMetric}\{\ApproximationHomotopy(\DomainPoint,\HomotopyTime)
 \mid\HomotopyTime\in\UnitInterval\}<\ErrorTolerance.
\]
Then
$\RecognitionTarget$
is an ANR
for all  separable metrizable spaces.
\end{theorem}

\begin{theorem}[{\cite[Theorem~13.4]{Hanner1952}}]\label{thm:anr-category}
\paraid{p-acd75c965ce7}
A separable metrizable space which is an ANR for all separable metrizable
spaces is an ANR for all metrizable spaces.
\end{theorem}

\begin{theorem}[{\cite[Theorem~12.3]{Hanner1952}}]\label{thm:contractible-ane}
\paraid{p-741437522ce9}
A nonempty contractible metrizable ANE for
all metrizable spaces is an AE
for all metrizable spaces.
\end{theorem}


\section{Absolute extension by global approximation}\label{sec:ar}
\paraid{p-ec939c6eac66}
We now combine the local maps of \Cref{thm:local-models}
into a global approximation through an absolute retract.
A partition of unity permits arbitrarily many local models to be active,
with only local finiteness required.
Using this construction,
we prove the absolute extension property without
a dimension restriction on the extension domain.


\subsection{Isometric representations of varying norms}\label{subsec:ar-norm-representations}
\paraid{p-8059e6be8ce0}
The norms
$\NormSymbol_\ComparisonCarrier$
in the local models vary with the compact metric space
$\MetricSpace{\ComparisonCarrier}{\ComparisonMetric}$.
We first represent all norms of a fixed finite dimension in one Banach
space in a way compatible with orthogonal changes of coordinates.

\paraid{p-e433ffb7555a}
For
$\ModelDimension \geq1$,
write
\[
 \UnitSphere^{\ModelDimension -1}=\{\SpherePoint \in\RealNumbers^\ModelDimension \mid\norm{\SpherePoint }_2=1\},
\]
and
\[
 \BanachAmbientSpace _\ModelDimension =\ContinuousFunctions(\UnitSphere^{\ModelDimension -1},\RealNumbers).
\]
We equip
$\BanachAmbientSpace _\ModelDimension $
with the uniform norm.
For a norm
$\NormSymbol $
on
$\RealNumbers^\ModelDimension $,
define its dual norm on
$\RealNumbers^\ModelDimension$
as follows.
For every
$\SecondVector\in\RealNumbers^\ModelDimension$,
set
\[
 \NormSymbol\DualNorm(\SecondVector)
 =\sup\{|\langle\FirstVector,\SecondVector\rangle|
       \mid\FirstVector\in\RealNumbers^\ModelDimension,
            \NormSymbol(\FirstVector)\leq1\}.
\]
The pairing is the Euclidean inner product.
For every nonzero
$\SecondVector\in\RealNumbers^\ModelDimension$,
we have
$\NormSymbol\DualNorm(\SecondVector)>0$.
In particular,
the restriction of the dual norm to
$\UnitSphere^{\ModelDimension-1}$
is positive.
Figure~\ref{fig:dual-norm-embedding} illustrates the construction in
\Cref{lem:dual-representation} in dimension
$2$.

\begin{figure}[H]
\centering
\begin{tikzpicture}[>=Stealth,font=\small]
\node at (0,1.65) {radial normalization};
\filldraw[fill=black!4] (0,0) ellipse (1.55 and 0.85);
\draw[dashed] (0,0) circle (1);
\draw[->,black!35] (-1.85,0) -- (1.9,0);
\draw[->,black!35] (0,-1.1) -- (0,1.25);
\pgfmathsetmacro{\radialfactor}{1/sqrt((cos(10)/1.55)^2+(sin(10)/0.85)^2)}
\coordinate (spherepoint) at ({cos(10)},{sin(10)});
\coordinate (dualpoint) at ({\radialfactor*cos(10)},{\radialfactor*sin(10)});
\draw[->] (0,0) -- (dualpoint);
\fill (spherepoint) circle (1.3pt);
\fill (dualpoint) circle (1.3pt);
\draw[black!55] (spherepoint) -- (0.65,0.3);
\node at (0.65,0.46) {$\SpherePoint$};
\draw[black!55] (dualpoint) -- (1.5,0.58);
\node at (1.5,0.98)
 {$\dfrac{\SpherePoint}{\NormSymbol\DualNorm(\SpherePoint)}$};
\node[below left,inner sep=2pt] at (0,0) {$0$};
\node at (-0.85,0.95) {$\UnitSphere^1$};
\node at (0,-1.4)
 {$\{\SecondVector\mid\NormSymbol\DualNorm(\SecondVector)=1\}$};
\draw[->] (2.45,0.05) -- node[above,align=center]
 {pair with\\$\FirstVector$} (3.45,0.05);
\begin{scope}[xshift=4.25cm]
 \node at (2.1,1.65)
  {$\NormEmbedding_\NormSymbol(\FirstVector)(\SpherePoint(\theta))$};
 \draw[->,black!50] (-0.15,0) -- (4.5,0) node[right] {$\theta$};
 \draw[->,black!50] (0,-1.35) -- (0,1.35);
 \draw[densely dotted,black!45] (0,1.1) -- (4.2,1.1);
 \draw[densely dotted,black!45] (0,-1.1) -- (4.2,-1.1);
 \draw[thick,smooth,samples=121,domain=0:360,variable=\angle]
  plot ({4.2*\angle/360},
        {1.1*cos(\angle)/(1.55*sqrt((cos(\angle)/1.55)^2+(sin(\angle)/0.85)^2))});
 \node[left,inner sep=3pt] at (0,1.1) {$\NormSymbol(\FirstVector)$};
 \node[left,inner sep=3pt] at (0,-1.1) {$-\NormSymbol(\FirstVector)$};
 \foreach \x/\tick in {0/0,2.1/\pi,4.2/2\pi}
 {
  \draw (\x,0.05) -- (\x,-0.05);
  \node[below,inner sep=3pt] at (\x,0) {$\tick$};
 }
\end{scope}
\node at (3.25,-2.15)
 {$\NormEmbedding_\NormSymbol(\FirstVector)(\SpherePoint)
   =\left\langle\FirstVector,
      \dfrac{\SpherePoint}{\NormSymbol\DualNorm(\SpherePoint)}\right\rangle,
   \qquad
   \norm{\NormEmbedding_\NormSymbol(\FirstVector)}_\infty
   =\NormSymbol(\FirstVector)$};
\end{tikzpicture}
\caption{The dual unit sphere is parametrized radially by the fixed
Euclidean sphere.
In the illustrated case,
the dual unit sphere is an ellipse,
$\FirstVector=(1,0)$,
and
$\SpherePoint(\theta)=(\cos\theta,\sin\theta)$.
Pairing with
$\FirstVector$
produces a continuous function on
$\UnitSphere^1$
whose uniform norm is
$\NormSymbol(\FirstVector)$.
Changing the norm changes the function while its domain remains fixed.}
\label{fig:dual-norm-embedding}
\end{figure}


\paraid{p-6b2bac5b7bed}
The following proposition is a specific version of the Banach--Mazur theorem.
Its linear isometry is obtained from the dual norm formula
\cite[Proposition~11.21]{Muscat2024}.

\begin{proposition}\label{lem:dual-representation}
\paraid{p-53660d2864b0}
Let
$\NormSymbol $
be a norm on
$\RealNumbers^\ModelDimension $.
For
$\FirstVector\in\RealNumbers^\ModelDimension$
and
$\SpherePoint\in\UnitSphere^{\ModelDimension-1}$,
define
\[
 \NormEmbedding _\NormSymbol \colon(\RealNumbers^\ModelDimension ,\NormSymbol )\to(\BanachAmbientSpace _\ModelDimension ,\norm{\cdot}_\infty)
\]
by
\begin{equation}\label{eq:dual-embedding}
 \NormEmbedding _\NormSymbol (\FirstVector )(\SpherePoint )=\frac{\langle \FirstVector ,\SpherePoint \rangle}{\NormSymbol \DualNorm(\SpherePoint )}.
\end{equation}
Then
$\NormEmbedding_\NormSymbol$
is a linear isometric embedding.
Thus,
for every
$\FirstVector,\SecondVector\in\RealNumbers^\ModelDimension$,
we have
\begin{equation}\label{eq:dual-isometry}
 \norm{\NormEmbedding _\NormSymbol (\FirstVector )-\NormEmbedding _\NormSymbol (\SecondVector )}_\infty=\NormSymbol (\FirstVector -\SecondVector ).
\end{equation}
\end{proposition}

\begin{proof}
\paraid{p-7ee99d206db2}
The formula in \eqref{eq:dual-embedding} defines a continuous function
because the dual norm is continuous and positive on the compact sphere.
It is linear in
$\FirstVector$.
For every
$\SpherePoint\in\UnitSphere^{\ModelDimension-1}$,
the definition of the dual norm gives
\[
 \frac{|\langle\FirstVector,\SpherePoint\rangle|}
      {\NormSymbol\DualNorm(\SpherePoint)}
 \leq\NormSymbol(\FirstVector).
\]
Taking the supremum over
$\UnitSphere^{\ModelDimension-1}$
shows that
\begin{equation}\label{eq:dual-embedding-upper-bound}
 \sup_{\SpherePoint\in\UnitSphere^{\ModelDimension-1}}
 \frac{|\langle\FirstVector,\SpherePoint\rangle|}
      {\NormSymbol\DualNorm(\SpherePoint)}
 \leq\NormSymbol(\FirstVector).
\end{equation}
If
$\FirstVector\ne0$,
define
$\SubspaceFunctional\colon\LinearSpan\{\FirstVector\}\to\RealNumbers$
 by
 $\SubspaceFunctional(\SpanCoefficient\FirstVector)=\SpanCoefficient\NormSymbol(\FirstVector)$.
Its operator norm is
\[
 \sup_{\substack{\HahnBanachTestVector\in\LinearSpan\{\FirstVector\}\\\HahnBanachTestVector\ne0}}
 \frac{|\SubspaceFunctional(\HahnBanachTestVector)|}{\NormSymbol(\HahnBanachTestVector)}=1.
\]
By the Hahn--Banach theorem
\cite[Theorem~11.19 and Proposition~11.20]{Muscat2024},
the functional
$\SubspaceFunctional$
extends to a functional
$\HahnBanachExtension\colon\RealNumbers^\ModelDimension\to\RealNumbers$
such that
\[
 \sup_{\HahnBanachTestVector\in\RealNumbers^\ModelDimension\setminus\{0\}}
 \frac{|\HahnBanachExtension(\HahnBanachTestVector)|}{\NormSymbol(\HahnBanachTestVector)}=1.
\]
Write
$\HahnBanachExtension(\HahnBanachTestVector)=\langle\HahnBanachTestVector,\NormingVector\rangle$.
Then
$\NormSymbol\DualNorm(\NormingVector)=1$
and
$\langle\FirstVector,\NormingVector\rangle=\NormSymbol(\FirstVector)$.
For every
$\HomogeneityScalar\in\RealNumbers\setminus\{0\}$,
we have
\[
 \frac{|\langle\FirstVector,\HomogeneityScalar\SpherePoint\rangle|}
      {\NormSymbol\DualNorm(\HomogeneityScalar\SpherePoint)}
 =\frac{|\HomogeneityScalar|\,|\langle\FirstVector,\SpherePoint\rangle|}
        {|\HomogeneityScalar|\,\NormSymbol\DualNorm(\SpherePoint)}
 =\frac{|\langle\FirstVector,\SpherePoint\rangle|}
        {\NormSymbol\DualNorm(\SpherePoint)}.
\]
Set
\[
 \NormalizedNormingVector
 =\frac{\NormingVector}{\norm{\NormingVector}_2}
 \in\UnitSphere^{\ModelDimension-1}.
\]
Then
\begin{equation}\label{eq:dual-embedding-lower-bound}
 \sup_{\SpherePoint\in\UnitSphere^{\ModelDimension-1}}
 \frac{|\langle\FirstVector,\SpherePoint\rangle|}
      {\NormSymbol\DualNorm(\SpherePoint)}
 \geq
 \frac{|\langle\FirstVector,\NormalizedNormingVector\rangle|}
      {\NormSymbol\DualNorm(\NormalizedNormingVector)}
 =\frac{|\langle\FirstVector,\NormingVector\rangle|}
        {\NormSymbol\DualNorm(\NormingVector)}
 =\NormSymbol(\FirstVector).
\end{equation}
The upper bound in
\eqref{eq:dual-embedding-upper-bound}
and the lower bound in
\eqref{eq:dual-embedding-lower-bound}
together with
\eqref{eq:dual-embedding},
 for
$\FirstVector\ne0$,
we have
\[
 \norm{\NormEmbedding_\NormSymbol(\FirstVector)}_\infty
 =\sup_{\SpherePoint\in\UnitSphere^{\ModelDimension-1}}
       \frac{|\langle\FirstVector,\SpherePoint\rangle|}
            {\NormSymbol\DualNorm(\SpherePoint)}
 =\NormSymbol(\FirstVector).
\]
This identity is immediate when
$\FirstVector=0$.
Linearity now proves \eqref{eq:dual-isometry}.
\end{proof}

\begin{proposition}\label{lem:dual-representation-continuity}
\paraid{p-e6ffdb8e19db}
Let
$\NormSymbol$
be a norm on
$\RealNumbers^\ModelDimension$.
If norms
$\NormSymbol _\SequenceIndex $
converge uniformly to
$\NormSymbol $
on
$\UnitSphere^{\ModelDimension -1}$
and
$\FirstVector _\SequenceIndex \to \FirstVector $
in
$\RealNumbers^\ModelDimension $,
then
\[
 \norm{\NormEmbedding _{\NormSymbol _\SequenceIndex }(\FirstVector _\SequenceIndex )-\NormEmbedding _\NormSymbol (\FirstVector )}_\infty\to0.
\]
\end{proposition}

\begin{proof}
\paraid{p-d51c2824eb6d}
Set
\[
 \MinimumNormValue =\min_{\UnitSphere^{\ModelDimension -1}}\NormSymbol >0.
\]
Set
\[
 \ErrorTolerance _\SequenceIndex =\sup_{\UnitSphere^{\ModelDimension -1}}|\NormSymbol _\SequenceIndex -\NormSymbol |.
\]
Set
\[
 \RelativeNormError _\SequenceIndex =\ErrorTolerance _\SequenceIndex /\MinimumNormValue .
\]
Since
$\RelativeNormError_\SequenceIndex\to0$,
choose
$n_0\in\NonnegativeIntegers$
such that
$0\leq\RelativeNormError_\SequenceIndex<1$
for every
$\SequenceIndex\geq n_0$.
For every
$\SequenceIndex\geq n_0$
and every
$\FirstVector\in\RealNumbers^\ModelDimension$,
homogeneity implies
\[
 (1-\RelativeNormError _\SequenceIndex )\NormSymbol (\FirstVector )\leq \NormSymbol _\SequenceIndex (\FirstVector )\leq(1+\RelativeNormError _\SequenceIndex )\NormSymbol (\FirstVector ).
\]
For every
$\SequenceIndex\geq n_0$
and every
$\SpherePoint\in\UnitSphere^{\ModelDimension-1}$,
the inclusions of unit balls imply
\[
 (1+\RelativeNormError _\SequenceIndex )^{-1}\NormSymbol \DualNorm(\SpherePoint )\leq \NormSymbol _\SequenceIndex \DualNorm(\SpherePoint )
                       \leq(1-\RelativeNormError _\SequenceIndex )^{-1}\NormSymbol \DualNorm(\SpherePoint ).
\]
Thus the dual norms
$\NormSymbol_\SequenceIndex\DualNorm$
converge uniformly to
$\NormSymbol\DualNorm$
on the sphere
$\UnitSphere^{\ModelDimension-1}$
and have a common positive lower bound there.
Consequently,
\begin{equation}\label{eq:dual-embedding-continuity}
 \begin{split}
 &\norm{\NormEmbedding _{\NormSymbol _\SequenceIndex }(\FirstVector _\SequenceIndex )-\NormEmbedding _\NormSymbol (\FirstVector )}_\infty
 \leq\\ &
 \frac{\norm{\FirstVector _\SequenceIndex -\FirstVector }_2}{\min_{\UnitSphere^{\ModelDimension -1}}\NormSymbol _\SequenceIndex \DualNorm}+\norm{\FirstVector }_2\sup_{\UnitSphere^{\ModelDimension -1}}
                       |(\NormSymbol _\SequenceIndex \DualNorm)^{-1}-(\NormSymbol \DualNorm)^{-1}|
 \longrightarrow0.
 \end{split}
\end{equation}
Applying \eqref{eq:dual-embedding-continuity} to fixed vectors,
we also obtain uniform convergence on each bounded set.
For every
$\CoefficientBound>0$,
\begin{equation}\label{eq:dual-embedding-uniform-on-bounded}
 \sup_{\substack{\FirstVector\in\RealNumbers^\ModelDimension\\\norm{\FirstVector}_2\leq\CoefficientBound}}
 \norm{\NormEmbedding_{\NormSymbol_\SequenceIndex}(\FirstVector)
       -\NormEmbedding_\NormSymbol(\FirstVector)}_\infty
 \leq\CoefficientBound
 \sup_{\UnitSphere^{\ModelDimension-1}}
 |(\NormSymbol_\SequenceIndex\DualNorm)^{-1}
   -(\NormSymbol\DualNorm)^{-1}|
 \longrightarrow0.
\end{equation}

\end{proof}

\begin{proposition}\label{lem:dual-representation-equivariance}
\paraid{p-ecf17de1513a}
Let
$\NormSymbol$
be a norm on
$\RealNumbers^\ModelDimension$.
For
$\OrthogonalChange\in\OrthogonalGroup(\ModelDimension)$,
$\SphereFunction\in\BanachAmbientSpace_\ModelDimension$,
and
$\SpherePoint\in\UnitSphere^{\ModelDimension-1}$,
define
\[
 (\OrthogonalChange\SphereFunction)(\SpherePoint)
 =\SphereFunction(\OrthogonalChange^{-1}\SpherePoint).
\]
This defines a continuous action of
$\OrthogonalGroup(\ModelDimension )$
on
$\BanachAmbientSpace _\ModelDimension $
by linear isometries.
For every
$\OrthogonalChange\in\OrthogonalGroup(\ModelDimension)$
and
$\FirstVector\in\RealNumbers^\ModelDimension$,
we have
\begin{equation}\label{eq:dual-equivariance}
 \NormEmbedding _{\NormSymbol \circ \OrthogonalChange ^{-1}}(\OrthogonalChange \FirstVector )=\OrthogonalChange \NormEmbedding _\NormSymbol (\FirstVector ).
\end{equation}
\end{proposition}

\begin{proof}
\paraid{p-2ab06945f322}
Orthogonal transformations preserve the Euclidean sphere,
so
$\norm{\OrthogonalChange \SphereFunction }_\infty=\norm{\SphereFunction }_\infty$.
If
$\OrthogonalChange _\SequenceIndex \to \OrthogonalChange $
and
$\SphereFunction _\SequenceIndex \to \SphereFunction $
uniformly,
then
\[
 \norm{\OrthogonalChange _\SequenceIndex \SphereFunction _\SequenceIndex -\OrthogonalChange \SphereFunction }_\infty
 \leq\norm{\SphereFunction _\SequenceIndex -\SphereFunction }_\infty+
      \sup_{\SpherePoint \in\UnitSphere^{\ModelDimension -1}}|\SphereFunction (\OrthogonalChange _\SequenceIndex ^{-1}\SpherePoint )-\SphereFunction (\OrthogonalChange ^{-1}\SpherePoint )|\to0
\]
by uniform continuity of
$\SphereFunction $.
Finally,
substituting
$\FirstVector =\OrthogonalChange \SecondVector $
in the definition of the dual norm
$\NormSymbol\DualNorm$
implies
\[
 (\NormSymbol \circ \OrthogonalChange ^{-1})\DualNorm(\SpherePoint )=\NormSymbol \DualNorm(\OrthogonalChange ^{-1}\SpherePoint ).
\]
Substituting this identity into \eqref{eq:dual-embedding},
we obtain
\[
 \NormEmbedding _{\NormSymbol \circ \OrthogonalChange ^{-1}}(\OrthogonalChange \FirstVector )(\SpherePoint )
 =\frac{\langle \OrthogonalChange \FirstVector ,\SpherePoint \rangle}{\NormSymbol \DualNorm(\OrthogonalChange ^{-1}\SpherePoint )}
 =\NormEmbedding _\NormSymbol (\FirstVector )(\OrthogonalChange ^{-1}\SpherePoint ),
\]
which proves \eqref{eq:dual-equivariance}.
\end{proof}

\subsection{Absolute retracts from hyperspace orbits}\label{subsec:ar-hyperspace-orbits}
\paraid{p-bb0d31d746c9}
Equation \eqref{eq:dual-equivariance} converts a change of coordinates into an
isometry of a fixed Banach space.
The global approximation will use compact subsets of such a space
modulo this action.
We apply \Cref{thm:equivariant-hyperspace,thm:orbit-retract} to prove that the resulting hyperspace orbit is
an absolute retract.

\begin{theorem}\label{lem:hyperspace-orbit-ar}
\paraid{p-4a30396f26c4}
Let
$(\BanachAmbientSpace ,\norm{\cdot}_\BanachAmbientSpace )$
be a separable Banach space,
and let a compact metrizable group
$\ActingGroup $
act continuously on
$\BanachAmbientSpace $
by linear isometries.
Equip $\BanachAmbientSpace$ with its induced metric
$\BanachMetric\colon\BanachAmbientSpace\times\BanachAmbientSpace\to[0,\infty)$.
For every $\FirstVector,\SecondVector\in\BanachAmbientSpace$, it satisfies
$\BanachMetric(\FirstVector,\SecondVector)=\norm{\FirstVector-\SecondVector}_\BanachAmbientSpace$.
Then the orbit space
$\CompactHyperspace(\BanachAmbientSpace )/\ActingGroup $
is a separable metrizable AR.
For nonempty compact subsets
$\FirstCompactSet,\SecondCompactSet\subset\BanachAmbientSpace$,
the formula
\begin{equation}\label{eq:orbit-metric}
  \OrbitMetric([\FirstCompactSet ],[\SecondCompactSet ])=\min_{\GroupElement \in \ActingGroup }\HausdorffDistance{\BanachMetric }(\FirstCompactSet ,\GroupElement \SecondCompactSet )
\end{equation}
defines a metric that induces the quotient topology.
\end{theorem}

\begin{proof}
\textbf{Step 1. The absolute retract property.}
\paraid{p-55569123ecd5}
The induced action on
$\CompactHyperspace(\BanachAmbientSpace )$
is continuous.
To see this,
assume that
$\GroupElement _\SequenceIndex \to \GroupElement $
and
$\HausdorffDistance{\BanachMetric }(\FirstCompactSet _\SequenceIndex ,\FirstCompactSet )\to0$.
Then
\[
 \HausdorffDistance{\BanachMetric }(\GroupElement _\SequenceIndex \FirstCompactSet _\SequenceIndex ,\GroupElement \FirstCompactSet )
 \leq\HausdorffDistance{\BanachMetric }(\FirstCompactSet _\SequenceIndex ,\FirstCompactSet )+\sup_{\CompactSetPoint \in \FirstCompactSet }\norm{\GroupElement _\SequenceIndex \CompactSetPoint -\GroupElement \CompactSetPoint }_\BanachAmbientSpace \to0.
\]
The last supremum tends to
$0$
because the action is continuous and
$\FirstCompactSet $
is compact.
Each group element acts isometrically on the hyperspace
$\CompactHyperspace(\BanachAmbientSpace)$.
By \Cref{thm:compact-hyperspace-topology},
the Hausdorff distance induces the Vietoris topology on
$\CompactHyperspace(\BanachAmbientSpace)$.
\Cref{thm:equivariant-hyperspace} therefore applies because
$\BanachAmbientSpace $
is connected and locally continuum-connected.
Thus
$\CompactHyperspace(\BanachAmbientSpace )$
is a
$\ActingGroup $-AR.
The compact metrizable group
$\ActingGroup $
has a countable basis,
so \Cref{thm:orbit-retract} implies that
$\CompactHyperspace(\BanachAmbientSpace )/\ActingGroup $
is an AR in the metrizable category.

\textbf{Step 2. The orbit metric and quotient topology.}
\paraid{p-154b529a9133}
Let
$\FirstCompactSet,\SecondCompactSet\subset\BanachAmbientSpace$
be nonempty compact subsets.
The minimum in \eqref{eq:orbit-metric} exists because
$\GroupElement \mapsto\HausdorffDistance{\BanachMetric }(\FirstCompactSet ,\GroupElement \SecondCompactSet )$
is continuous on the compact group
$\ActingGroup$.

\paraid{p-fe18da324b94}
The Hausdorff distance on
$\CompactHyperspace(\BanachAmbientSpace)$
is
$\ActingGroup$-invariant,
and every orbit is compact and hence closed.
The argument in
\cite[proof of Theorem~4.3.4, p.~319]{Palais1961}
(see also \cite[p.~4, (2.1) and the preceding paragraph]{Antonyan2020Euclidean})
shows that \eqref{eq:orbit-metric} defines a metric inducing the quotient topology.
Let
$\OrbitProjection \colon\CompactHyperspace(\BanachAmbientSpace )\to\CompactHyperspace(\BanachAmbientSpace )/\ActingGroup $
be the orbit map.

\paraid{p-1e603606de4c}
Since
$\BanachAmbientSpace$
is separable,
\Cref{thm:compact-hyperspace-topology} implies that
$\CompactHyperspace(\BanachAmbientSpace)$
is separable.
Its continuous image under
$\OrbitProjection$
is therefore separable as well.
\end{proof}

\begin{lemma}\label{lem:hyperspace-orbit-realization}
\paraid{p-0459ff3d5b20}
Let
$(\BanachAmbientSpace ,\norm{\cdot}_\BanachAmbientSpace )$
be a separable Banach space,
and let a compact metrizable group
$\ActingGroup $
act continuously on
$\BanachAmbientSpace $
by linear isometries.
Equip $\BanachAmbientSpace$ with its induced metric
$\BanachMetric\colon\BanachAmbientSpace\times\BanachAmbientSpace\to[0,\infty)$.
For every $\FirstVector,\SecondVector\in\BanachAmbientSpace$, it satisfies
$\BanachMetric(\FirstVector,\SecondVector)=\norm{\FirstVector-\SecondVector}_\BanachAmbientSpace$.
Equip
$\CompactHyperspace(\BanachAmbientSpace)/\ActingGroup$
with the metric in \eqref{eq:orbit-metric}.
Then the map
\[
 \RealizationMap\colon\CompactHyperspace(\BanachAmbientSpace )/\ActingGroup \to\GHSpace
\]
defined by
\[
 \RealizationMap([\FirstCompactSet ])=\MetricSpace{\FirstCompactSet }{\RestrictedMetric{\BanachMetric }{\FirstCompactSet }},
\]
is well-defined and
$1$-Lipschitz.
\end{lemma}

\begin{proof}
\paraid{p-0ad0a6a4a958}
For each
$\GroupElement \in \ActingGroup $,
the map
$\SecondCompactSet \to \GroupElement \SecondCompactSet $
obtained by restricting
$\GroupElement$
is an isometry.
Consequently
$\RealizationMap$
is independent of the representative of an orbit.
For every
$\GroupElement\in\ActingGroup$,
we also have
\[
 \begin{aligned}
 \GHDistance(\RealizationMap([\FirstCompactSet ]),\RealizationMap([\SecondCompactSet ]))
 &=\GHDistance(\MetricSpace{\FirstCompactSet }{\BanachMetric |_{\FirstCompactSet \times \FirstCompactSet }},
         \MetricSpace{\GroupElement \SecondCompactSet }{\BanachMetric |_{\GroupElement \SecondCompactSet \times \GroupElement \SecondCompactSet }})\\
 &\leq\HausdorffDistance{\BanachMetric }(\FirstCompactSet ,\GroupElement \SecondCompactSet ).
 \end{aligned}
\]
Taking the minimum over
$\GroupElement\in\ActingGroup$
proves that
$\RealizationMap$
is
$1$-Lipschitz.
\end{proof}

\begin{remark}\label{rem:peano-hyperspace}
\paraid{p-15a612ab8c6c}
A classical hyperspace theorem also relates compact subsets to the Hilbert cube.
Let
$(X,d)$
be a nondegenerate Peano continuum,
that is,
a compact connected locally connected metric space with at least two points.
Curtis and Schori \cite[Theorem~3.2]{CurtisSchori1978}
prove that
\[
 \bigl(\CompactHyperspace(X),\HausdorffDistance{d}\bigr)
 \text{ is homeomorphic to }\HilbertCube=[0,1]^{\NonnegativeIntegers}.
\]
This hyperspace consists of all nonempty compact subsets of the fixed space
$X$.
The space in \Cref{lem:hyperspace-orbit-ar}
instead consists of compact subsets of a Banach space modulo a compact group.
\end{remark}

\subsection{Global approximation with small homotopy tracks}\label{subsec:ar-global-approximation}
\paraid{p-441d3863d070}
In the next theorem,
we combine the local models to construct approximations
through the auxiliary AR
$\AuxiliaryAR=\CompactHyperspace(\BanachAmbientSpace)/\ActingGroup$.
We choose this orbit space for the prescribed error function.
For maps on $\GHSpace$, evaluation at $\MetricSpace{\ComparisonCarrier}{\ComparisonMetric}$
means evaluation at its isometry class $[\ComparisonCarrier,\ComparisonMetric]$.
This convention applies in particular to $\ModelError$ and $\ErrorControl$.

\begin{theorem}\label{thm:global-domination}
\paraid{p-c8631fd0a7ae}
Let
$\ErrorControl \colon \GHSpace\to(0,\infty)$
be continuous.
Then there exist a separable metrizable AR
$\AuxiliaryAR$,
continuous maps
$\ApproximationMap\colon \GHSpace\to \AuxiliaryAR$
and
$\RealizationMap\colon \AuxiliaryAR\to\GHSpace$,
a continuous function
$\ModelError \colon \GHSpace\to[0,\infty)$,
and a continuous homotopy
$\ApproximationHomotopy \colon \GHSpace\times\UnitInterval\to\GHSpace$
with the following properties.
For every
$\MetricSpace{\ComparisonCarrier}{\ComparisonMetric}\in\GHSpace$,
we have
\[
 \ModelError\MetricSpace{\ComparisonCarrier}{\ComparisonMetric}
 <\ErrorControl\MetricSpace{\ComparisonCarrier}{\ComparisonMetric},
\]
and the endpoint identities
\begin{equation}\label{eq:global-domination-1}
 \ApproximationHomotopy (\MetricSpace{\ComparisonCarrier }{\ComparisonMetric },0)=\MetricSpace{\ComparisonCarrier }{\ComparisonMetric },
\end{equation}
and
\begin{equation}\label{eq:global-domination-2}
 \ApproximationHomotopy (\MetricSpace{\ComparisonCarrier }{\ComparisonMetric },1)=\RealizationMap\ApproximationMap\MetricSpace{\ComparisonCarrier }{\ComparisonMetric }.
\end{equation}
For every
$\MetricSpace{\ComparisonCarrier}{\ComparisonMetric}\in\GHSpace$
and every
$\FirstTime,\HomotopyTime\in\UnitInterval$,
we have
\begin{equation}\label{eq:global-track-control}
  \GHDistance(\ApproximationHomotopy (\MetricSpace{\ComparisonCarrier }{\ComparisonMetric },\FirstTime ),\ApproximationHomotopy (\MetricSpace{\ComparisonCarrier }{\ComparisonMetric },\HomotopyTime ))\leq\frac{\abs{\FirstTime -\HomotopyTime }}2 \ModelError \MetricSpace{\ComparisonCarrier }{\ComparisonMetric }.
\end{equation}
\end{theorem}

\begin{proof}
\paraid{p-665396df3bb9}
We construct the local cover and its Banach space,
define the global map,
and then prove continuity and the homotopy estimates.

\ProofStep{Local models and a partition of unity.}\label{step:global-local-models}
\paraid{p-45dee78ed9f8}
For each center
$\MetricSpace{\BaseCarrier}{\MetricSymbol}\in\GHSpace$,
apply \Cref{thm:local-models} with a positive tolerance smaller than
$\ErrorControl\MetricSpace{\BaseCarrier}{\MetricSymbol}$.
By continuity of the model error and of
$\ErrorControl$,
shrink the model domain so that for every space
$\MetricSpace{\ComparisonCarrier}{\ComparisonMetric}$
in that domain the model error is less than
$\ErrorControl\MetricSpace{\ComparisonCarrier}{\ComparisonMetric}$.
Let $\OpenCover$ be the family of these domains.
Then $\OpenCover$ is an open cover of
$\GHSpace$.

\paraid{p-d2972dd59fa1}
Since
$\GHSpace$
is separable and metrizable,
choose a countable locally finite partition of unity
$\{\PartitionWeight _\CoverIndex \}_{\CoverIndex \in\NonnegativeIntegers}$
subordinate to $\OpenCover$
(\Cref{thm:metric-partition-unity}).
For each index
$\CoverIndex$,
fix one of the local models whose domain
$\ModelDomain_\CoverIndex$
contains
$\supp\PartitionWeight_\CoverIndex$.
Different indices may use the same local model.
The supports form a locally finite family.
For each
$\CoverIndex\in\NonnegativeIntegers$,
write the selected local model as
\[
 \mathfrak M_\CoverIndex
 =(\ModelDomain_\CoverIndex,\ModelDimension_\CoverIndex,
   \mathcal L_\CoverIndex,\ModelError_\CoverIndex).
\]
Thus
$\ModelDimension_\CoverIndex\geq1$
is its dimension,
$\mathcal L_\CoverIndex$
is its coordinate-class assignment,
and
$\ModelError_\CoverIndex\colon\ModelDomain_\CoverIndex\to[0,\infty)$
is its error function.
For a compact metric space
$\MetricSpace{\ComparisonCarrier}{\ComparisonMetric}$
whose isometry class belongs to
$\ModelDomain_\CoverIndex$,
choose a pair in
$\mathcal L_\CoverIndex(\ComparisonCarrier,\ComparisonMetric)$.
Write its coordinate map as
$\FeatureCoordinates_{\CoverIndex,\ComparisonCarrier}\colon\ComparisonCarrier\to\RealNumbers^{\ModelDimension_\CoverIndex}$
and its norm as
$\NormSymbol_{\CoverIndex,\ComparisonCarrier}\colon\RealNumbers^{\ModelDimension_\CoverIndex}\to[0,\infty)$.
For
$\ComparisonPoint,\IntegrationPoint\in\ComparisonCarrier$,
write
\[
 \ModelPseudometric_{\CoverIndex,\ComparisonCarrier}(\ComparisonPoint,\IntegrationPoint)
 =\NormSymbol_{\CoverIndex,\ComparisonCarrier}
   (\FeatureCoordinates_{\CoverIndex,\ComparisonCarrier}(\ComparisonPoint)
    -\FeatureCoordinates_{\CoverIndex,\ComparisonCarrier}(\IntegrationPoint)).
\]
The model error satisfies
\[
 \ModelError_\CoverIndex\MetricSpace{\ComparisonCarrier}{\ComparisonMetric}
 =\norm{\ComparisonMetric-\ModelPseudometric_{\CoverIndex,\ComparisonCarrier}}_\infty
 <\ErrorControl\MetricSpace{\ComparisonCarrier}{\ComparisonMetric}.
\]
The zero-weight blocks will be defined separately in Step~\ref{step:global-map}.

\ProofStep{A Banach space with a compact group action.}\label{step:global-orbit-space}
\paraid{p-ab63bc101601}
Form the separable Banach space and compact metrizable group
\begin{equation}\label{eq:banach-sum-group-1}
 \BanachAmbientSpace =\left(\bigoplus_{\CoverIndex \in\NonnegativeIntegers}\BanachAmbientSpace _{\ModelDimension _\CoverIndex }\right)_{\SummableNorm},
\end{equation}
and
\begin{equation}\label{eq:banach-sum-group-2}
 \ActingGroup =\prod_{\CoverIndex \in\NonnegativeIntegers}\OrthogonalGroup(\ModelDimension _\CoverIndex ).
\end{equation}
The norm in
$\BanachAmbientSpace $
is
$\norm{\{\BlockVector _\CoverIndex \}_{\CoverIndex \in\NonnegativeIntegers}}_\BanachAmbientSpace =\sum_{\CoverIndex\in\NonnegativeIntegers} \norm{\BlockVector _\CoverIndex }_\infty$.
Define the induced metric
\[
 \BanachMetric\colon\BanachAmbientSpace\times\BanachAmbientSpace\to[0,\infty)
\]
as follows.
For every
$\FirstVector,\SecondVector\in\BanachAmbientSpace$,
set
\[
 \BanachMetric(\FirstVector,\SecondVector)=\norm{\FirstVector-\SecondVector}_\BanachAmbientSpace.
\]
The group
$\ActingGroup$
acts on each summand by linear isometries.
We prove that the action map
\[
 \ActingGroup\times\BanachAmbientSpace\to\BanachAmbientSpace,
 \qquad
 (\GroupElement,\DirectSumVector)\longmapsto\GroupElement\DirectSumVector
\]
is continuous.
Indeed,
if
$\GroupElement _\SequenceIndex \to \GroupElement $
in the product group and
$\DirectSumVector =\{\BlockVector _\CoverIndex \}\in \BanachAmbientSpace $,
then for each
$\TruncationSize \geq1$,
\[
 \norm{\GroupElement _\SequenceIndex \DirectSumVector -\GroupElement \DirectSumVector }_\BanachAmbientSpace
 \leq\sum_{\CoverIndex =0}^\TruncationSize \norm{(\GroupElement _\SequenceIndex )_\CoverIndex \BlockVector _\CoverIndex -\GroupElement _\CoverIndex \BlockVector _\CoverIndex }_\infty
        +2\sum_{\CoverIndex >\TruncationSize }\norm{\BlockVector _\CoverIndex }_\infty.
\]
First choose
$\TruncationSize $
to control the tail,
and then use continuity on the finitely many remaining blocks.
For varying
$\ApproximatingDirectSumVector _\SequenceIndex \to \DirectSumVector $,
the inequality
\[
 \norm{\GroupElement _\SequenceIndex \ApproximatingDirectSumVector _\SequenceIndex -\GroupElement \DirectSumVector }_\BanachAmbientSpace \leq\norm{\ApproximatingDirectSumVector _\SequenceIndex -\DirectSumVector }_\BanachAmbientSpace +\norm{\GroupElement _\SequenceIndex \DirectSumVector -\GroupElement \DirectSumVector }_\BanachAmbientSpace\longrightarrow0
\]
proves continuity of the action map on the product
$\ActingGroup\times\BanachAmbientSpace$.

\ProofStep{The global map and its error.}\label{step:global-map}
\paraid{p-4b598342fd3f}
Fix a compact metric space
$\MetricSpace{\ComparisonCarrier}{\ComparisonMetric}$
representing an input class in
$\GHSpace$.
Define the finite set of active indices for this input by
\[
 \mathcal I(\ComparisonCarrier,\ComparisonMetric)
 =\{\CoverIndex\in\NonnegativeIntegers\mid
 \PartitionWeight_\CoverIndex\MetricSpace{\ComparisonCarrier}{\ComparisonMetric}>0\}.
\]
Finiteness follows from local finiteness of the partition of unity.
For each index with
$\PartitionWeight_\CoverIndex\MetricSpace{\ComparisonCarrier}{\ComparisonMetric}>0$,
choose a representative coordinate pair
$(\FeatureCoordinates_{\CoverIndex,\ComparisonCarrier},\NormSymbol_{\CoverIndex,\ComparisonCarrier})$
from the coordinate class of the assigned local model.
For every $\CoverIndex\in\NonnegativeIntegers$ and
$\ComparisonPoint\in\ComparisonCarrier$, define
\[
 B_{\CoverIndex,\ComparisonCarrier}(\ComparisonPoint)=
 \begin{cases}
 \PartitionWeight_\CoverIndex\MetricSpace{\ComparisonCarrier}{\ComparisonMetric}
 \NormEmbedding_{\NormSymbol_{\CoverIndex,\ComparisonCarrier}}
 (\FeatureCoordinates_{\CoverIndex,\ComparisonCarrier}(\ComparisonPoint))
 &\text{if }\PartitionWeight_\CoverIndex\MetricSpace{\ComparisonCarrier}{\ComparisonMetric}>0,\\
 0&\text{otherwise}.
 \end{cases}
\]
The zero in the second case is the zero vector of
$\BanachAmbientSpace_{\ModelDimension_\CoverIndex}$.
Using these blocks,
define the map
$\JointFeatureMap_\ComparisonCarrier\colon\ComparisonCarrier\to\BanachAmbientSpace$
as follows.
For every
$\ComparisonPoint\in\ComparisonCarrier$,
put
\begin{equation}\label{eq:joint-feature-map-1}
 \JointFeatureMap_\ComparisonCarrier(\ComparisonPoint)=\{B_{\CoverIndex,\ComparisonCarrier}(\ComparisonPoint)\}_{\CoverIndex\in\NonnegativeIntegers},
\end{equation}
and denote its image by
\begin{equation}\label{eq:joint-feature-map-2}
 \CompactFeatureImage _\ComparisonCarrier =\JointFeatureMap _\ComparisonCarrier (\ComparisonCarrier ).
\end{equation}
For every point of
$\ComparisonCarrier$,
all nonzero blocks have indices in the same finite set
$\mathcal I(\ComparisonCarrier,\ComparisonMetric)$.
Thus
$\JointFeatureMap _\ComparisonCarrier $
is continuous and
$\CompactFeatureImage _\ComparisonCarrier $
is nonempty and compact.
Equation \eqref{eq:dual-isometry} and \ref{item:model-bound} of \Cref{thm:local-models}
imply that for every
$\ComparisonPoint\in\ComparisonCarrier$
\begin{equation}\label{eq:global-feature-bound}
  \norm{\JointFeatureMap _\ComparisonCarrier (\ComparisonPoint )}_\BanachAmbientSpace \leq2\diam\MetricSpace{\ComparisonCarrier }{\ComparisonMetric }.
\end{equation}

\paraid{p-b86cdaebafd5}
If the local coordinate pairs are changed as in \eqref{eq:frame-action} by
$\GroupElement=\{\OrthogonalChange_\CoverIndex\}_\CoverIndex\in\ActingGroup$,
then \eqref{eq:dual-equivariance} transforms the joint point map into
$\GroupElement\circ\JointFeatureMap_\ComparisonCarrier$
and its image into
$\GroupElement\CompactFeatureImage_\ComparisonCarrier$.
If
$\InputIsometry\colon\MetricSpace{\ComparisonCarrier}{\ComparisonMetric}
 \to\MetricSpace{\ComparisonCarrier_0}{\ComparisonMetric_0}$
is an isometry,
then
\eqref{eq:input-isometry-1},
\eqref{eq:input-isometry-2},
and \eqref{eq:dual-equivariance}
imply that for some
$\GroupElement\in\ActingGroup$,
\[
 \JointFeatureMap_{\ComparisonCarrier_0}\circ\InputIsometry
 =\GroupElement\circ\JointFeatureMap_\ComparisonCarrier.
\]
Thus both changes leave the orbit of the compact image unchanged.
Consequently,
we can define
\[
 \ApproximationMap\colon \GHSpace\to \AuxiliaryAR=\CompactHyperspace(\BanachAmbientSpace )/\ActingGroup
\]
by
\[
 \ApproximationMap\MetricSpace{\ComparisonCarrier }{\ComparisonMetric }=[\CompactFeatureImage _\ComparisonCarrier ].
\]
The value is independent of the representative of the input isometry class.
By \Cref{lem:hyperspace-orbit-ar},
$\AuxiliaryAR$
is a separable metrizable AR.
By \Cref{lem:hyperspace-orbit-realization},
we obtain a
$1$-Lipschitz map
$\RealizationMap\colon \AuxiliaryAR\to\GHSpace$.

\paraid{p-1264f3c8b4b3}
Figure~\ref{fig:global-models} shows how the local maps are combined.
\begin{figure}[H]
\centering
\begin{tikzpicture}[>=Stealth,font=\small]
\node (input) at (0,0) {$\ComparisonPoint\in\ComparisonCarrier$};
\node (first) at (4,1) {$\PartitionWeight_{\CoverIndex_0}\MetricSpace{\ComparisonCarrier}{\ComparisonMetric}
 \NormEmbedding_{\NormSymbol_{\CoverIndex_0,\ComparisonCarrier}}
 (\FeatureCoordinates_{\CoverIndex_0,\ComparisonCarrier}(\ComparisonPoint))$};
\node (second) at (4,-1) {$\PartitionWeight_{\CoverIndex_1}\MetricSpace{\ComparisonCarrier}{\ComparisonMetric}
 \NormEmbedding_{\NormSymbol_{\CoverIndex_1,\ComparisonCarrier}}
 (\FeatureCoordinates_{\CoverIndex_1,\ComparisonCarrier}(\ComparisonPoint))$};
\node at (4,0) {$\vdots$};
\node (joint) at (8,0) {$\JointFeatureMap_\ComparisonCarrier(\ComparisonPoint)\in\BanachAmbientSpace$};
\draw[->] (input)--(first);
\draw[->] (input)--(second);
\draw[->] (first)--(joint);
\draw[->] (second)--(joint);
\end{tikzpicture}
\par\medskip
\begin{tikzcd}[column sep=large]
 \GHSpace \arrow[r,"\ApproximationMap"] &
 \AuxiliaryAR=\CompactHyperspace(\BanachAmbientSpace)/\ActingGroup
 \arrow[r,"\RealizationMap"] & \GHSpace
\end{tikzcd}
\[
 \RealizationMap\ApproximationMap\simeq\IdentityMap_\GHSpace
 \quad\text{via }\ApproximationHomotopy.
\]
\caption{For a fixed compact metric space
$\MetricSpace{\ComparisonCarrier}{\ComparisonMetric}$,
the active blocks use the same point
$\ComparisonPoint$.
Their joint image is
$\CompactFeatureImage_\ComparisonCarrier=\JointFeatureMap_\ComparisonCarrier(\ComparisonCarrier)$.
Changing the chosen coordinate pairs as in \eqref{eq:frame-action} acts on this image by
$\ActingGroup$,
so
$\ApproximationMap\MetricSpace{\ComparisonCarrier}{\ComparisonMetric}
 =[\CompactFeatureImage_\ComparisonCarrier]$.
The displayed composite is joined to the identity by the homotopy in Step~\ref{step:global-homotopy}
with the track estimate \eqref{eq:global-track-control}.}
\label{fig:global-models}
\end{figure}


\paraid{p-b45630ebc16d}
For
$\ComparisonPoint,\IntegrationPoint\in\ComparisonCarrier$,
define the pseudometric induced by
$\JointFeatureMap_\ComparisonCarrier$
by
\begin{equation}\label{eq:global-pseudometric}
  \ModelPseudometric _\ComparisonCarrier (\ComparisonPoint ,\IntegrationPoint )=\norm{\JointFeatureMap _\ComparisonCarrier (\ComparisonPoint )-\JointFeatureMap _\ComparisonCarrier (\IntegrationPoint )}_\BanachAmbientSpace.
\end{equation}
By \eqref{eq:dual-isometry},
we have
\begin{equation}\label{eq:global-pseudometric-sum}
 \ModelPseudometric_\ComparisonCarrier(\ComparisonPoint,\IntegrationPoint)
 =\sum_{\substack{\CoverIndex\in\NonnegativeIntegers\\\PartitionWeight_\CoverIndex\MetricSpace{\ComparisonCarrier}{\ComparisonMetric}>0}} \PartitionWeight _\CoverIndex \MetricSpace{\ComparisonCarrier }{\ComparisonMetric }\ModelPseudometric _{\CoverIndex ,\ComparisonCarrier }(\ComparisonPoint ,\IntegrationPoint ).
\end{equation}
Every block in \eqref{eq:joint-feature-map-1} is evaluated at the same point
$\ComparisonPoint\in\ComparisonCarrier$.
Define the error by
\begin{equation}\label{eq:global-error}
  \ModelError \MetricSpace{\ComparisonCarrier }{\ComparisonMetric }=\norm{\ComparisonMetric -\ModelPseudometric _\ComparisonCarrier }_\infty.
\end{equation}
Since the weights sum to one,
\eqref{eq:global-pseudometric-sum} implies
\begin{equation}\label{eq:global-error-bound}
 \ModelError\MetricSpace{\ComparisonCarrier}{\ComparisonMetric}
 \leq
 \sum_{\CoverIndex\in \mathcal I(\ComparisonCarrier,\ComparisonMetric)}\PartitionWeight _\CoverIndex \MetricSpace{\ComparisonCarrier }{\ComparisonMetric }\ModelError _\CoverIndex \MetricSpace{\ComparisonCarrier }{\ComparisonMetric }
 <
 \sum_{\CoverIndex\in \mathcal I(\ComparisonCarrier,\ComparisonMetric)}\PartitionWeight _\CoverIndex \MetricSpace{\ComparisonCarrier }{\ComparisonMetric }
 \ErrorControl \MetricSpace{\ComparisonCarrier }{\ComparisonMetric }
=\ErrorControl \MetricSpace{\ComparisonCarrier }{\ComparisonMetric }.
\end{equation}





\ProofStep{Continuity along common correspondences.}\label{step:global-continuity}
\paraid{p-613bbdb06ba1}
We next prove continuity of
$\ApproximationMap$
and
$\ModelError$
by estimating the maps
$\JointFeatureMap_\ComparisonCarrier$
along common correspondences.
Local finiteness reduces the argument to finitely many blocks.
We treat separately the blocks whose weights have positive limits
and those whose weights tend to
$0$.
We will also use the estimate \eqref{eq:global-alignment} to prove continuity of
$\ApproximationHomotopy$.

\ProofSubstep{A common correspondence and finitely many blocks.}\label{step:global-common-correspondence}
\paraid{p-527721648b5d}
Let
$\MetricSpace{\ComparisonCarrier _\SequenceIndex }{\ComparisonMetric _\SequenceIndex}\to\MetricSpace{\ComparisonCarrier }{\ComparisonMetric }$
in
$\GHSpace$.
By \Cref{lem:common-gh-embedding},
we may regard all the carriers as compact subsets of one compact metric space
$\MetricSpace{\AmbientCarrier}{\AmbientMetric}$
such that
$\HausdorffDistance{\AmbientMetric}
 (\ComparisonCarrier_\SequenceIndex,\ComparisonCarrier)\to0$.
Choose correspondences
$\Correspondence _\SequenceIndex \subset \ComparisonCarrier _\SequenceIndex \times \ComparisonCarrier $
such that
\[
 \AmbientDisplacement_\SequenceIndex
 =\sup_{(\ComparisonPoint_\SequenceIndex,\ComparisonPoint)
          \in\Correspondence_\SequenceIndex}
   \AmbientMetric(\ComparisonPoint_\SequenceIndex,\ComparisonPoint)
 \longrightarrow0.
\]
We use these same correspondences for every block.
Local finiteness of the supports provides an open neighborhood
$\OpenNeighborhood$
of
$\MetricSpace{\ComparisonCarrier}{\ComparisonMetric}$
such that
\[
 \ActiveModelIndices
 =\{\CoverIndex\in\NonnegativeIntegers\mid
       \OpenNeighborhood\cap\supp\PartitionWeight_\CoverIndex
       \ne\emptyset\}
 \quad\text{is finite}.
\]
For all sufficiently large
$\SequenceIndex$,
we have
$\MetricSpace{\ComparisonCarrier_\SequenceIndex}{\ComparisonMetric_\SequenceIndex}
 \in\OpenNeighborhood$.
Thus,
for every
$\CoverIndex\in\NonnegativeIntegers\setminus\ActiveModelIndices$,
\[
 \PartitionWeight_\CoverIndex\MetricSpace{\ComparisonCarrier}{\ComparisonMetric}
 =\PartitionWeight_\CoverIndex
   \MetricSpace{\ComparisonCarrier_\SequenceIndex}{\ComparisonMetric_\SequenceIndex}
 =0.
\]
Partition the finite set
$\ActiveModelIndices$
according to the limiting weights by putting
\[
 I_+=\{\CoverIndex\in\ActiveModelIndices\mid
 \PartitionWeight_\CoverIndex\MetricSpace{\ComparisonCarrier}{\ComparisonMetric}>0\}
\]
and
$I_-=\ActiveModelIndices\setminus I_+$.
Namely,
\[
 I_-=\{\CoverIndex\in\ActiveModelIndices\mid
 \PartitionWeight_\CoverIndex\MetricSpace{\ComparisonCarrier}{\ComparisonMetric}=0\}.
\]

\ProofSubstep{Blocks with positive limiting weights.}\label{step:global-positive-blocks}
\paraid{p-5b0d3e37abfc}
For each
$\CoverIndex\in I_+$,
continuity of
$\PartitionWeight_\CoverIndex$
implies
\[
 \PartitionWeight_\CoverIndex
 \MetricSpace{\ComparisonCarrier_\SequenceIndex}{\ComparisonMetric_\SequenceIndex}
 \longrightarrow
 \PartitionWeight_\CoverIndex\MetricSpace{\ComparisonCarrier}{\ComparisonMetric}>0
 \quad\text{as }\SequenceIndex\to\infty.
\]
Since
$I_+$
is finite,
we ignore finitely many initial terms so that
$\PartitionWeight_\CoverIndex
 \MetricSpace{\ComparisonCarrier_\SequenceIndex}{\ComparisonMetric_\SequenceIndex}>0$
for every
$\CoverIndex\in I_+$
and every remaining index
$\SequenceIndex$.
Thus,
for each
$\CoverIndex\in I_+$,
the corresponding local model is defined at both
$\MetricSpace{\ComparisonCarrier_\SequenceIndex}{\ComparisonMetric_\SequenceIndex}$
and
$\MetricSpace{\ComparisonCarrier}{\ComparisonMetric}$.
We now examine the continuity of the factor
$\NormEmbedding_{\NormSymbol_{\CoverIndex,\ComparisonCarrier}}
 \circ\FeatureCoordinates_{\CoverIndex,\ComparisonCarrier}$
in the definition of
$B_{\CoverIndex,\ComparisonCarrier}$
as both the norm and the coordinate map vary.
By \ref{item:model-continuity} of \Cref{thm:local-models},
choose coordinate pairs for each
$\CoverIndex\in I_+$
such that
\begin{equation}\label{eq:global-block-norm-convergence}
 \sup_{\SpherePoint\in\UnitSphere^{\ModelDimension_\CoverIndex-1}}
 \bigl|\NormSymbol_{\CoverIndex,\ComparisonCarrier_\SequenceIndex}(\SpherePoint)
       -\NormSymbol_{\CoverIndex,\ComparisonCarrier}(\SpherePoint)\bigr|
 \longrightarrow0
\end{equation}
and
\begin{equation}\label{eq:global-block-coordinate-convergence}
 \sup_{(\ComparisonPoint_\SequenceIndex,\ComparisonPoint)
          \in\Correspondence_\SequenceIndex}
 \norm{\FeatureCoordinates_{\CoverIndex,\ComparisonCarrier_\SequenceIndex}
          (\ComparisonPoint_\SequenceIndex)
       -\FeatureCoordinates_{\CoverIndex,\ComparisonCarrier}
          (\ComparisonPoint)}_2
\longrightarrow0.
\end{equation}
Fix
$\CoverIndex\in I_+$.
For every sufficiently large
$\SequenceIndex$,
the supremum in \eqref{eq:global-block-coordinate-convergence}
is at most
$1$.
Fix such an index
$\SequenceIndex$
and take an arbitrary point
$\ComparisonPoint_\SequenceIndex\in\ComparisonCarrier_\SequenceIndex$.
Since the first coordinate projection
$\Correspondence_\SequenceIndex\to\ComparisonCarrier_\SequenceIndex$
is surjective,
choose
$\ComparisonPoint\in\ComparisonCarrier$
such that
$(\ComparisonPoint_\SequenceIndex,\ComparisonPoint)\in\Correspondence_\SequenceIndex$.
Applying the triangle inequality to this pair,
we obtain
\begin{equation}\label{eq:global-block-pointwise-bound}
 \begin{aligned}
 \norm{\FeatureCoordinates_{\CoverIndex,\ComparisonCarrier_\SequenceIndex}
       (\ComparisonPoint_\SequenceIndex)}_2
 &\leq
 \norm{\FeatureCoordinates_{\CoverIndex,\ComparisonCarrier_\SequenceIndex}
       (\ComparisonPoint_\SequenceIndex)
       -\FeatureCoordinates_{\CoverIndex,\ComparisonCarrier}(\ComparisonPoint)}_2
 +\norm{\FeatureCoordinates_{\CoverIndex,\ComparisonCarrier}(\ComparisonPoint)}_2\\
 &\leq1+\max_{\IntegrationPoint\in\ComparisonCarrier}
       \norm{\FeatureCoordinates_{\CoverIndex,\ComparisonCarrier}(\IntegrationPoint)}_2.
 \end{aligned}
\end{equation}
Taking the supremum over all
$\ComparisonPoint_\SequenceIndex\in\ComparisonCarrier_\SequenceIndex$
in \eqref{eq:global-block-pointwise-bound},
we obtain
\begin{equation}\label{eq:global-block-coordinate-bound}
 \sup_{\ComparisonPoint_\SequenceIndex\in\ComparisonCarrier_\SequenceIndex}
 \norm{\FeatureCoordinates_{\CoverIndex,\ComparisonCarrier_\SequenceIndex}
       (\ComparisonPoint_\SequenceIndex)}_2
 \leq1+\max_{\IntegrationPoint\in\ComparisonCarrier}
       \norm{\FeatureCoordinates_{\CoverIndex,\ComparisonCarrier}(\IntegrationPoint)}_2.
\end{equation}
Put
\[
 \Radius_\CoverIndex
 =1+\max_{\IntegrationPoint\in\ComparisonCarrier}
       \norm{\FeatureCoordinates_{\CoverIndex,\ComparisonCarrier}(\IntegrationPoint)}_2
 <\infty.
\]
The finiteness follows from continuity and compactness.
By \eqref{eq:global-block-coordinate-bound},
for all sufficiently large
$\SequenceIndex$,
the images of
$\FeatureCoordinates_{\CoverIndex,\ComparisonCarrier_\SequenceIndex}$
and
$\FeatureCoordinates_{\CoverIndex,\ComparisonCarrier}$
lie in the closed Euclidean ball
\[
 \overline{\MetricBall}(0,\Radius_\CoverIndex;\norm{\cdot}_2)
 =\{\FirstVector\in\RealNumbers^{\ModelDimension_\CoverIndex}\mid
     \norm{\FirstVector}_2\leq\Radius_\CoverIndex\}.
\]
For
$(\ComparisonPoint_\SequenceIndex,\ComparisonPoint)
 \in\Correspondence_\SequenceIndex$,
we separate the change of norm from the change of coordinates by the estimate
\begin{equation}\label{eq:global-block-embedding-comparison}
 \begin{aligned}
 &\Bigl\|\NormEmbedding_{\NormSymbol_{\CoverIndex,\ComparisonCarrier_\SequenceIndex}}
   \bigl(\FeatureCoordinates_{\CoverIndex,\ComparisonCarrier_\SequenceIndex}
      (\ComparisonPoint_\SequenceIndex)\bigr)
 -\NormEmbedding_{\NormSymbol_{\CoverIndex,\ComparisonCarrier}}
   \bigl(\FeatureCoordinates_{\CoverIndex,\ComparisonCarrier}
      (\ComparisonPoint)\bigr)\Bigr\|_\infty\\
 &\quad\leq
 \Bigl\|\bigl(\NormEmbedding_{\NormSymbol_{\CoverIndex,\ComparisonCarrier_\SequenceIndex}}
             -\NormEmbedding_{\NormSymbol_{\CoverIndex,\ComparisonCarrier}}\bigr)
   \bigl(\FeatureCoordinates_{\CoverIndex,\ComparisonCarrier_\SequenceIndex}
      (\ComparisonPoint_\SequenceIndex)\bigr)\Bigr\|_\infty
 +
 \NormSymbol_{\CoverIndex,\ComparisonCarrier}
   \bigl(\FeatureCoordinates_{\CoverIndex,\ComparisonCarrier_\SequenceIndex}
      (\ComparisonPoint_\SequenceIndex)
       -\FeatureCoordinates_{\CoverIndex,\ComparisonCarrier}(\ComparisonPoint)\bigr)\\
 &\quad\leq
 \sup_{\FirstVector\in\overline{\MetricBall}(0,\Radius_\CoverIndex;\norm{\cdot}_2)}
 \Bigl\|\bigl(\NormEmbedding_{\NormSymbol_{\CoverIndex,\ComparisonCarrier_\SequenceIndex}}
             -\NormEmbedding_{\NormSymbol_{\CoverIndex,\ComparisonCarrier}}\bigr)
   (\FirstVector)\Bigr\|_\infty
 +
 \NormSymbol_{\CoverIndex,\ComparisonCarrier}
   \bigl(\FeatureCoordinates_{\CoverIndex,\ComparisonCarrier_\SequenceIndex}
      (\ComparisonPoint_\SequenceIndex)
       -\FeatureCoordinates_{\CoverIndex,\ComparisonCarrier}(\ComparisonPoint)\bigr).
 \end{aligned}
\end{equation}
The first inequality in \eqref{eq:global-block-embedding-comparison}
follows from the triangle inequality and \eqref{eq:dual-isometry}.
The second follows from \eqref{eq:global-block-coordinate-bound}.
By \eqref{eq:global-block-norm-convergence}
and \eqref{eq:dual-embedding-uniform-on-bounded},
we obtain
\begin{equation}\label{eq:global-block-embedding-uniform-convergence}
 \sup_{\FirstVector\in\overline{\MetricBall}(0,\Radius_\CoverIndex;\norm{\cdot}_2)}
 \Bigl\|\bigl(\NormEmbedding_{\NormSymbol_{\CoverIndex,\ComparisonCarrier_\SequenceIndex}}
             -\NormEmbedding_{\NormSymbol_{\CoverIndex,\ComparisonCarrier}}\bigr)
   (\FirstVector)\Bigr\|_\infty
 \longrightarrow0.
\end{equation}
Since the fixed norm
$\NormSymbol_{\CoverIndex,\ComparisonCarrier}$
is continuous on the compact unit sphere,
the constant
\[
 \NormComparisonBound_\CoverIndex
 =\max_{\SpherePoint\in\UnitSphere^{\ModelDimension_\CoverIndex-1}}
       \NormSymbol_{\CoverIndex,\ComparisonCarrier}(\SpherePoint)
\]
is finite.
For every
$\FirstVector\in\RealNumbers^{\ModelDimension_\CoverIndex}$,
homogeneity implies
\[
 \NormSymbol_{\CoverIndex,\ComparisonCarrier}(\FirstVector)
 \leq\NormComparisonBound_\CoverIndex\norm{\FirstVector}_2.
\]
Thus,
\eqref{eq:global-block-coordinate-convergence} implies
\begin{equation}\label{eq:global-block-coordinate-norm-convergence}
 \begin{aligned}
 &\sup_{(\ComparisonPoint_\SequenceIndex,\ComparisonPoint)
          \in\Correspondence_\SequenceIndex}
 \NormSymbol_{\CoverIndex,\ComparisonCarrier}
   \bigl(\FeatureCoordinates_{\CoverIndex,\ComparisonCarrier_\SequenceIndex}
      (\ComparisonPoint_\SequenceIndex)
       -\FeatureCoordinates_{\CoverIndex,\ComparisonCarrier}(\ComparisonPoint)\bigr)\\
 &\quad\leq\NormComparisonBound_\CoverIndex
 \sup_{(\ComparisonPoint_\SequenceIndex,\ComparisonPoint)
          \in\Correspondence_\SequenceIndex}
 \norm{\FeatureCoordinates_{\CoverIndex,\ComparisonCarrier_\SequenceIndex}
          (\ComparisonPoint_\SequenceIndex)
       -\FeatureCoordinates_{\CoverIndex,\ComparisonCarrier}
          (\ComparisonPoint)}_2
 \longrightarrow0.
 \end{aligned}
\end{equation}
By \eqref{eq:global-block-embedding-comparison},
\eqref{eq:global-block-embedding-uniform-convergence},
and \eqref{eq:global-block-coordinate-norm-convergence},
for every
$\CoverIndex\in I_+$,
the embedded coordinate maps are uniformly bounded and converge
uniformly along
$\Correspondence_\SequenceIndex$.
By continuity of the weights
$\PartitionWeight_\CoverIndex$
and the definition of the blocks
$B_{\CoverIndex,\ComparisonCarrier_\SequenceIndex}$
and
$B_{\CoverIndex,\ComparisonCarrier}$
in Step~\ref{step:global-map},
we obtain
\begin{equation}\label{eq:global-positive-weight-block}
 \sup_{(\ComparisonPoint_\SequenceIndex,\ComparisonPoint)
          \in\Correspondence_\SequenceIndex}
 \norm{B_{\CoverIndex,\ComparisonCarrier_\SequenceIndex}
           (\ComparisonPoint_\SequenceIndex)
       -B_{\CoverIndex,\ComparisonCarrier}(\ComparisonPoint)}_\infty
 \longrightarrow0.
\end{equation}

\ProofSubstep{Blocks with vanishing weights.}\label{step:global-vanishing-blocks}
\paraid{p-50b148b589f3}
Fix
$\CoverIndex\in I_-$.
By continuity of
$\PartitionWeight_\CoverIndex$,
we have
$\PartitionWeight_\CoverIndex
 \MetricSpace{\ComparisonCarrier_\SequenceIndex}{\ComparisonMetric_\SequenceIndex}
 \to\PartitionWeight_\CoverIndex\MetricSpace{\ComparisonCarrier}{\ComparisonMetric}=0$.
The definition in Step~\ref{step:global-map} implies
$B_{\CoverIndex,\ComparisonCarrier}=0$.
If the weight at
$\MetricSpace{\ComparisonCarrier_\SequenceIndex}{\ComparisonMetric_\SequenceIndex}$
is zero,
then the corresponding block is zero by its definition in Step~\ref{step:global-map}.
If the weight is positive,
then the local model is defined,
and \ref{item:model-bound} of \Cref{thm:local-models}
and \eqref{eq:dual-isometry} apply.
In both cases,
\begin{equation}\label{eq:vanishing-weight-block}
 \sup_{\ComparisonPoint_\SequenceIndex\in\ComparisonCarrier_\SequenceIndex}
 \norm{B_{\CoverIndex,\ComparisonCarrier_\SequenceIndex}(\ComparisonPoint_\SequenceIndex)}_\infty
 \leq2\PartitionWeight _\CoverIndex \MetricSpace{\ComparisonCarrier _\SequenceIndex }{\ComparisonMetric _\SequenceIndex}\diam\MetricSpace{\ComparisonCarrier _\SequenceIndex }{\ComparisonMetric _\SequenceIndex}\to0.
\end{equation}
Indeed,
the diameters are bounded along the convergent sequence.
Thus these blocks converge uniformly to
$0$
for any choices of coordinate pairs.

\ProofSubstep{Continuity of the global map and its error.}\label{step:global-map-continuity}
\paraid{p-f79a92fe8ba3}
Use the coordinate pairs chosen in Step~\ref{step:global-positive-blocks} on
$I_+$
and arbitrary coordinate pairs for the blocks with positive weights in
$I_-$.
Put
\[
 \FeatureDisplacement _\SequenceIndex =\sup_{(\ComparisonPoint _\SequenceIndex ,\ComparisonPoint )\in \Correspondence _\SequenceIndex }\norm{\JointFeatureMap _{\ComparisonCarrier _\SequenceIndex }(\ComparisonPoint _\SequenceIndex )-\JointFeatureMap _\ComparisonCarrier (\ComparisonPoint )}_\BanachAmbientSpace ,
\]
so that
$\FeatureDisplacement_\SequenceIndex$
measures the displacement of the joint feature maps along
$\Correspondence_\SequenceIndex$.
All blocks outside
$\ActiveModelIndices$
vanish for sufficiently large
$\SequenceIndex$.
Since the norm of
$\BanachAmbientSpace$
is the sum of the block norms and
$\ActiveModelIndices$
is finite,
\eqref{eq:global-positive-weight-block} and \eqref{eq:vanishing-weight-block}
imply
\begin{equation}\label{eq:global-alignment}
 \FeatureDisplacement_\SequenceIndex
 \leq\sum_{\CoverIndex\in\ActiveModelIndices}
 \sup_{(\ComparisonPoint_\SequenceIndex,\ComparisonPoint)
          \in\Correspondence_\SequenceIndex}
 \norm{B_{\CoverIndex,\ComparisonCarrier_\SequenceIndex}
           (\ComparisonPoint_\SequenceIndex)
       -B_{\CoverIndex,\ComparisonCarrier}(\ComparisonPoint)}_\infty
\longrightarrow0.
\end{equation}
Since both coordinate projections of
$\Correspondence_\SequenceIndex$
are onto,
the definition of
$\FeatureDisplacement_\SequenceIndex$
implies
\[
 \HausdorffDistance{\BanachMetric}
 (\CompactFeatureImage_{\ComparisonCarrier_\SequenceIndex},
  \CompactFeatureImage_\ComparisonCarrier)
 \leq\FeatureDisplacement_\SequenceIndex\longrightarrow0.
\]
Here the representatives
$\CompactFeatureImage_{\ComparisonCarrier_\SequenceIndex}$
and
$\CompactFeatureImage_\ComparisonCarrier$
are constructed using,
for each
$\CoverIndex\in I_+$,
the coordinate pairs guaranteed by
\ref{item:model-continuity} of \Cref{thm:local-models}
to satisfy \eqref{eq:global-block-norm-convergence}
and \eqref{eq:global-block-coordinate-convergence}.
Taking the identity element of
$\ActingGroup$
in \eqref{eq:orbit-metric},
we obtain
\[
 \OrbitMetric([\CompactFeatureImage_{\ComparisonCarrier_\SequenceIndex}],
             [\CompactFeatureImage_\ComparisonCarrier])
 \leq\HausdorffDistance{\BanachMetric}
 (\CompactFeatureImage_{\ComparisonCarrier_\SequenceIndex},
  \CompactFeatureImage_\ComparisonCarrier)
 \longrightarrow0.
\]
This proves continuity of
$\ApproximationMap$.

\paraid{p-39150d57f5ed}
We finally prove continuity of
$\ModelError$.
For
$(\ComparisonPoint _\SequenceIndex ,\ComparisonPoint ),(\IntegrationPoint _\SequenceIndex ,\IntegrationPoint )\in \Correspondence _\SequenceIndex $,
the triangle inequalities in
$\BanachAmbientSpace$
and
$\MetricSpace{\AmbientCarrier}{\AmbientMetric}$
imply
\[
 |\ModelPseudometric _{\ComparisonCarrier _\SequenceIndex }(\ComparisonPoint _\SequenceIndex ,\IntegrationPoint _\SequenceIndex )-\ModelPseudometric _\ComparisonCarrier (\ComparisonPoint ,\IntegrationPoint )|\leq2\FeatureDisplacement _\SequenceIndex ,
\]
and
\[
 |\ComparisonMetric _\SequenceIndex(\ComparisonPoint _\SequenceIndex ,\IntegrationPoint _\SequenceIndex )-\ComparisonMetric (\ComparisonPoint ,\IntegrationPoint )|\leq2\AmbientDisplacement _\SequenceIndex .
\]
By \eqref{eq:pseudometric-uniform-comparison},
we obtain
\[
 |\ModelError \MetricSpace{\ComparisonCarrier _\SequenceIndex }{\ComparisonMetric _\SequenceIndex}-\ModelError \MetricSpace{\ComparisonCarrier }{\ComparisonMetric }|
 \leq2\FeatureDisplacement _\SequenceIndex +2\AmbientDisplacement _\SequenceIndex \longrightarrow0.
\]
This proves continuity of
$\ModelError $
in \eqref{eq:global-error}.

\ProofStep{The interpolating homotopy.}\label{step:global-homotopy}
\paraid{p-35ac6ff9ee97}
We interpolate the original metric and the model pseudometric to construct a continuous map
$\ApproximationHomotopy\colon\GHSpace\times\UnitInterval\to\GHSpace$
with the endpoint identities in \eqref{eq:global-domination-1} and \eqref{eq:global-domination-2}
and the estimate in \eqref{eq:global-track-control}.

\ProofSubstep{Definition and compactness.}\label{step:global-homotopy-definition}
\paraid{p-132d78a0f66f}
For
$\MetricSpace{\ComparisonCarrier}{\ComparisonMetric}\in\GHSpace$
and
$\HomotopyTime\in\UnitInterval$,
put
\begin{equation}\label{eq:interpolating-pseudometrics-1}
\InterpolatingPseudometric_{\ComparisonCarrier,\HomotopyTime}
 =(1-\HomotopyTime)\ComparisonMetric
       +\HomotopyTime\ModelPseudometric_\ComparisonCarrier,
\end{equation}
and
\begin{equation}\label{eq:interpolating-pseudometrics-2}
\ApproximationHomotopy(\MetricSpace{\ComparisonCarrier}{\ComparisonMetric},\HomotopyTime)
 =\MetricQuotient{\ComparisonCarrier}
       {\InterpolatingPseudometric_{\ComparisonCarrier,\HomotopyTime}}.
\end{equation}
Both
$\ComparisonMetric$
and
$\ModelPseudometric_\ComparisonCarrier$
are continuous pseudometrics on the compact space
$\ComparisonCarrier$.
By the definition,
$\InterpolatingPseudometric_{\ComparisonCarrier,\HomotopyTime}$
is a metric
generating the same topology on
$Y$
when
$\HomotopyTime<1$.

\paraid{p-7cd5915365f0}
By \eqref{eq:global-pseudometric-sum}
and the invariance under \eqref{eq:frame-action},
$\ModelPseudometric_\ComparisonCarrier$
is independent of the coordinate pairs chosen in the local models.
Moreover,
by \eqref{eq:input-isometry-1} and \eqref{eq:input-isometry-2},
every isometry
$\InputIsometry\colon\MetricSpace{\ComparisonCarrier}{\ComparisonMetric}
 \to\MetricSpace{\ComparisonCarrier_0}{\ComparisonMetric_0}$
satisfies
\[
 \ModelPseudometric_{\ComparisonCarrier_0}
   (\InputIsometry(\ComparisonPoint),\InputIsometry(\IntegrationPoint))
 =\ModelPseudometric_\ComparisonCarrier(\ComparisonPoint,\IntegrationPoint)
 \qquad(\ComparisonPoint,\IntegrationPoint\in\ComparisonCarrier).
\]
Thus
$\ApproximationHomotopy$
is well-defined on
$\GHSpace\times\UnitInterval$.

\ProofSubstep{The two endpoints.}\label{step:global-homotopy-endpoints}
\paraid{p-76f372559d9c}
We next prove \eqref{eq:global-domination-1} and \eqref{eq:global-domination-2}.
By \eqref{eq:interpolating-pseudometrics-1},
we obtain
$\InterpolatingPseudometric_{\ComparisonCarrier,0}=\ComparisonMetric$,
and
$\InterpolatingPseudometric_{\ComparisonCarrier,1}
       =\ModelPseudometric_\ComparisonCarrier$.
At the second endpoint,
for every
$\ComparisonPoint,\IntegrationPoint\in\ComparisonCarrier$,
\eqref{eq:global-pseudometric}
implies
that
$\ModelPseudometric_\ComparisonCarrier(\ComparisonPoint,\IntegrationPoint)=0$
if and only if
$\JointFeatureMap_\ComparisonCarrier(\ComparisonPoint)
       =\JointFeatureMap_\ComparisonCarrier(\IntegrationPoint)$.
Since
$\CompactFeatureImage_\ComparisonCarrier
 =\JointFeatureMap_\ComparisonCarrier(\ComparisonCarrier)$,
the induced map
\begin{equation}\label{eq:quotient-feature-isometry}
 \MetricQuotient{\ComparisonCarrier}{\ModelPseudometric_\ComparisonCarrier}
 \to\MetricSpace{\CompactFeatureImage_\ComparisonCarrier}
       {\RestrictedMetric{\BanachMetric}{\CompactFeatureImage_\ComparisonCarrier}}
\end{equation}
is defined as follows.
For each
$\ComparisonPoint\in\ComparisonCarrier$,
send
$[\ComparisonPoint]$
to
$\JointFeatureMap_\ComparisonCarrier(\ComparisonPoint)$.
It is a bijection.
For every
$\ComparisonPoint,\IntegrationPoint\in\ComparisonCarrier$,
the identity
\[
 \ModelPseudometric_\ComparisonCarrier(\ComparisonPoint,\IntegrationPoint)
 =\norm{\JointFeatureMap_\ComparisonCarrier(\ComparisonPoint)
        -\JointFeatureMap_\ComparisonCarrier(\IntegrationPoint)}_\BanachAmbientSpace
\]
shows that this map is an isometry.
By \eqref{eq:interpolating-pseudometrics-2},
the isometry
\eqref{eq:quotient-feature-isometry},
and the definitions of
$\ApproximationMap$
and
$\RealizationMap$,
we obtain the following equalities in
$\GHSpace$.
\[
 \ApproximationHomotopy(\MetricSpace{\ComparisonCarrier}{\ComparisonMetric},0)
  =\MetricSpace{\ComparisonCarrier}{\ComparisonMetric},
\]
and
\[
 \begin{aligned}
 \ApproximationHomotopy(\MetricSpace{\ComparisonCarrier}{\ComparisonMetric},1)
 &=\MetricQuotient{\ComparisonCarrier}{\ModelPseudometric_\ComparisonCarrier}=\MetricSpace{\CompactFeatureImage_\ComparisonCarrier}
       {\RestrictedMetric{\BanachMetric}{\CompactFeatureImage_\ComparisonCarrier}}
  =\RealizationMap\ApproximationMap\MetricSpace{\ComparisonCarrier}{\ComparisonMetric}.
 \end{aligned}
\]

\ProofSubstep{Variation of the time parameter.}\label{step:global-homotopy-time}
\paraid{p-8853a418623b}
Fix
$\MetricSpace{\ComparisonCarrier}{\ComparisonMetric}$.
To prove \eqref{eq:global-track-control},
let
$\FirstTime,\HomotopyTime\in\UnitInterval$.
By \eqref{eq:global-error},
\[
 \InterpolatingPseudometric_{\ComparisonCarrier,\FirstTime}
       -\InterpolatingPseudometric_{\ComparisonCarrier,\HomotopyTime}
 =(\FirstTime-\HomotopyTime)
       (\ModelPseudometric_\ComparisonCarrier-\ComparisonMetric),
\]
and
\[
 \norm{\InterpolatingPseudometric_{\ComparisonCarrier,\FirstTime}
       -\InterpolatingPseudometric_{\ComparisonCarrier,\HomotopyTime}}_\infty
 =|\FirstTime-\HomotopyTime|
       \ModelError\MetricSpace{\ComparisonCarrier}{\ComparisonMetric}.
\]
Applying \eqref{eq:pseudometric-bound} on the common carrier
$\ComparisonCarrier$,
we obtain
\[
 \begin{aligned}
 &\GHDistance\bigl(
   \ApproximationHomotopy(\MetricSpace{\ComparisonCarrier}{\ComparisonMetric},\FirstTime),
   \ApproximationHomotopy(\MetricSpace{\ComparisonCarrier}{\ComparisonMetric},\HomotopyTime)
   \bigr)\\
 &\qquad\leq\tfrac12
       \norm{\InterpolatingPseudometric_{\ComparisonCarrier,\FirstTime}
       -\InterpolatingPseudometric_{\ComparisonCarrier,\HomotopyTime}}_\infty
  =\tfrac12|\FirstTime-\HomotopyTime|
       \ModelError\MetricSpace{\ComparisonCarrier}{\ComparisonMetric}.
 \end{aligned}
\]
This is \eqref{eq:global-track-control}.
Figure~\ref{fig:pseudometric-homotopy} illustrates the interpolation
and the metric quotient at its endpoint.

\begin{figure}[H]
\centering
\begin{tikzpicture}[>=Stealth,font=\small]
\node at (0,2.15) {$\HomotopyTime=0$};
\node at (3.8,2.15) {$\HomotopyTime=\tfrac12$};
\node at (7.6,2.15) {$\HomotopyTime=1$};
\coordinate (leftzero) at (-0.375,0);
\coordinate (leftone) at (0.375,0);
\coordinate (lefttwo) at (0,1.4524);
\draw (leftzero) -- node[left] {$2$} (lefttwo)
 -- node[right] {$2$} (leftone)
 -- node[below] {$1$} (leftzero);
\foreach \point in {leftzero,leftone,lefttwo}
 \fill (\point) circle (1.4pt);
\node[below left] at (leftzero) {$\ComparisonPoint_0$};
\node[below right] at (leftone) {$\ComparisonPoint_1$};
\node[above] at (lefttwo) {$\ComparisonPoint_2$};
\begin{scope}[xshift=3.8cm]
 \coordinate (middlezero) at (-0.1875,0);
 \coordinate (middleone) at (0.1875,0);
 \coordinate (middletwo) at (0,1.4882);
 \draw (middlezero) -- node[left] {$2$} (middletwo)
 -- node[right] {$2$} (middleone)
 -- node[below] {$\tfrac12$} (middlezero);
 \foreach \point in {middlezero,middleone,middletwo}
  \fill (\point) circle (1.4pt);
 \node[below left] at (middlezero) {$\ComparisonPoint_0$};
 \node[below right] at (middleone) {$\ComparisonPoint_1$};
 \node[above] at (middletwo) {$\ComparisonPoint_2$};
\end{scope}
\begin{scope}[xshift=7.6cm]
 \draw (0,0) -- node[right] {$2$} (0,1.5);
 \fill (0,0) circle (1.4pt);
 \fill (0,1.5) circle (1.4pt);
 \node[below] at (0,0) {$[\ComparisonPoint_0]=[\ComparisonPoint_1]$};
 \node[above] at (0,1.5) {$[\ComparisonPoint_2]$};
\end{scope}
\draw[->] (1.15,0.7) -- (2.6,0.7);
\draw[->] (4.9,0.7) -- (6.4,0.7);
\node at (3.8,-1)
 {$\InterpolatingPseudometric_{\ComparisonCarrier,\HomotopyTime}
   =(1-\HomotopyTime)\ComparisonMetric
     +\HomotopyTime\ModelPseudometric_\ComparisonCarrier$};
\node at (3.8,-1.85)
 {$\begin{aligned}
  &\GHDistance\bigl(
   \ApproximationHomotopy(\MetricSpace{\ComparisonCarrier}{\ComparisonMetric},\FirstTime),
   \ApproximationHomotopy(\MetricSpace{\ComparisonCarrier}{\ComparisonMetric},\HomotopyTime)
   \bigr)\\
  &\qquad\leq\tfrac12|\FirstTime-\HomotopyTime|
       \ModelError\MetricSpace{\ComparisonCarrier}{\ComparisonMetric}
  \end{aligned}$};
\end{tikzpicture}
\caption{Interpolation on a three-point carrier.
The edge labels are distances for
$\InterpolatingPseudometric_{\ComparisonCarrier,\HomotopyTime}$.
The distance between
$\ComparisonPoint_0$
and
$\ComparisonPoint_1$
decreases from
$1$
to
$0$,
while their distances to
$\ComparisonPoint_2$
remain
$2$.
Distinct points can be identified at the endpoint.
The general track estimate \eqref{eq:global-track-control}
also holds there.}
\label{fig:pseudometric-homotopy}
\end{figure}


\ProofSubstep{Continuity when the space and time both vary.}\label{step:global-homotopy-continuity}
\paraid{p-f7c37d1cb847}
Let
\[
 \MetricSpace{\ComparisonCarrier_\SequenceIndex}{\ComparisonMetric_\SequenceIndex}
       \longrightarrow\MetricSpace{\ComparisonCarrier}{\ComparisonMetric}
       \quad\text{in }\GHSpace,
\]
and
\[
 \HomotopyTime_\SequenceIndex\longrightarrow\HomotopyTime
       \quad\text{in }\UnitInterval.
\]
We prove that
\[
 \GHDistance\bigl(
   \ApproximationHomotopy(
       \MetricSpace{\ComparisonCarrier_\SequenceIndex}{\ComparisonMetric_\SequenceIndex},
       \HomotopyTime_\SequenceIndex),
   \ApproximationHomotopy(\MetricSpace{\ComparisonCarrier}{\ComparisonMetric},\HomotopyTime)
   \bigr)\longrightarrow0.
\]
Use the correspondences
$\Correspondence_\SequenceIndex$
and the numbers
$\AmbientDisplacement_\SequenceIndex,\FeatureDisplacement_\SequenceIndex\to0$
from Step~\ref{step:global-continuity}.
First
we consider the case of
$\HomotopyTime_\SequenceIndex=\HomotopyTime$.
For
$(\ComparisonPoint_\SequenceIndex,\ComparisonPoint),
  (\IntegrationPoint_\SequenceIndex,\IntegrationPoint)\in\Correspondence_\SequenceIndex$,
the estimates in Step~\ref{step:global-continuity} imply
\[
 \begin{aligned}
 &\bigl|\InterpolatingPseudometric_{\ComparisonCarrier_\SequenceIndex,\HomotopyTime}
       (\ComparisonPoint_\SequenceIndex,\IntegrationPoint_\SequenceIndex)
       -\InterpolatingPseudometric_{\ComparisonCarrier,\HomotopyTime}
       (\ComparisonPoint,\IntegrationPoint)\bigr|\\
 &\qquad\leq(1-\HomotopyTime)
       |\ComparisonMetric_\SequenceIndex(\ComparisonPoint_\SequenceIndex,\IntegrationPoint_\SequenceIndex)
         -\ComparisonMetric(\ComparisonPoint,\IntegrationPoint)|
       +\HomotopyTime
       |\ModelPseudometric_{\ComparisonCarrier_\SequenceIndex}
          (\ComparisonPoint_\SequenceIndex,\IntegrationPoint_\SequenceIndex)
         -\ModelPseudometric_\ComparisonCarrier(\ComparisonPoint,\IntegrationPoint)|\\
 &\qquad\leq2(1-\HomotopyTime)\AmbientDisplacement_\SequenceIndex
       +2\HomotopyTime\FeatureDisplacement_\SequenceIndex.
 \end{aligned}
\]
At
$\HomotopyTime=1$,
the induced pseudometrics may identify distinct points,
and we use their metric quotients.
\Cref{lem:pseudometric-comparison} therefore applies also at this endpoint and yields
\begin{equation}\label{eq:fixed-time-homotopy-convergence}
 \begin{aligned}
 &\GHDistance\bigl(
   \ApproximationHomotopy(
       \MetricSpace{\ComparisonCarrier_\SequenceIndex}{\ComparisonMetric_\SequenceIndex},\HomotopyTime),
   \ApproximationHomotopy(\MetricSpace{\ComparisonCarrier}{\ComparisonMetric},\HomotopyTime)
   \bigr)\leq(1-\HomotopyTime)\AmbientDisplacement_\SequenceIndex
       +\HomotopyTime\FeatureDisplacement_\SequenceIndex\longrightarrow0.
 \end{aligned}
\end{equation}

\paraid{p-6a7506bb68e5}
We combine the fixed-time convergence with the track estimate to compare the values at
$\HomotopyTime_\SequenceIndex$
and
$\HomotopyTime$.
The continuity
of
$\ModelError$
proved in Step~\ref{step:global-continuity} implies
\[
 \ModelError\MetricSpace{\ComparisonCarrier_\SequenceIndex}{\ComparisonMetric_\SequenceIndex}
       \longrightarrow\ModelError\MetricSpace{\ComparisonCarrier}{\ComparisonMetric}.
\]
By the triangle inequality and \eqref{eq:global-track-control},
\[
 \begin{aligned}
 &\GHDistance\bigl(
   \ApproximationHomotopy(
       \MetricSpace{\ComparisonCarrier_\SequenceIndex}{\ComparisonMetric_\SequenceIndex},
       \HomotopyTime_\SequenceIndex),
   \ApproximationHomotopy(\MetricSpace{\ComparisonCarrier}{\ComparisonMetric},\HomotopyTime)
   \bigr)\\
 &\qquad\leq\GHDistance\bigl(
   \ApproximationHomotopy(
       \MetricSpace{\ComparisonCarrier_\SequenceIndex}{\ComparisonMetric_\SequenceIndex},
       \HomotopyTime_\SequenceIndex),
   \ApproximationHomotopy(
       \MetricSpace{\ComparisonCarrier_\SequenceIndex}{\ComparisonMetric_\SequenceIndex},\HomotopyTime)
   \bigr)+\GHDistance\bigl(
   \ApproximationHomotopy(
       \MetricSpace{\ComparisonCarrier_\SequenceIndex}{\ComparisonMetric_\SequenceIndex},\HomotopyTime),
   \ApproximationHomotopy(\MetricSpace{\ComparisonCarrier}{\ComparisonMetric},\HomotopyTime)
   \bigr)\\
 &\qquad\leq\tfrac12|\HomotopyTime_\SequenceIndex-\HomotopyTime|
       \ModelError\MetricSpace{\ComparisonCarrier_\SequenceIndex}{\ComparisonMetric_\SequenceIndex}
       +(1-\HomotopyTime)\AmbientDisplacement_\SequenceIndex
       +\HomotopyTime\FeatureDisplacement_\SequenceIndex\longrightarrow0.
 \end{aligned}
\]
Since
$\GHSpace\times\UnitInterval$
is metrizable,
this proves continuity of
$\ApproximationHomotopy\colon\GHSpace\times\UnitInterval\to\GHSpace$.
\end{proof}

\subsection{Absolute retracts}\label{subsec:ar-extension}
\paraid{p-32b0b613231b}
We now apply the domination and extension theorems from
Subsection~\ref{subsec:ar-preliminaries-domination}.

\begin{theorem}\label{thm:absolute-extensor}
\paraid{p-8ec15f8bbe5b}
The space
$\GHSpace$
is an ANR for all metrizable spaces.
In particular,
the space
$\GHSpace$ is also  an AR
for all metrizable spaces.
\end{theorem}

\begin{proof}
\paraid{p-2a571e30ef7f}
Apply \Cref{thm:global-domination} with
$\ErrorControl \equiv\ErrorTolerance $
for each
$\ErrorTolerance >0$.
We obtain continuous maps
$\ApproximationMap\colon\GHSpace\to\AuxiliaryAR$
and
$\RealizationMap\colon\AuxiliaryAR\to\GHSpace$
and a homotopy
$\ApproximationHomotopy$
from
$\IdentityMap_{\GHSpace}$
to
$\RealizationMap\circ\ApproximationMap$.
By \eqref{eq:global-track-control},
for every
$\MetricSpace{\ComparisonCarrier}{\ComparisonMetric}\in\GHSpace$,
we have
\[
 \diam_{\GHDistance}
 \{\ApproximationHomotopy(\MetricSpace{\ComparisonCarrier}{\ComparisonMetric},\HomotopyTime)
   \mid\HomotopyTime\in\UnitInterval\}
 \leq\tfrac12\ModelError\MetricSpace{\ComparisonCarrier}{\ComparisonMetric}
 <\tfrac12\ErrorTolerance.
\]
By \Cref{thm:hanner-domination},
$\GHSpace$
is an ANR
for all  separable metrizable spaces.

\paraid{p-6516cafeac15}
By \Cref{thm:anr-category},
the space
$\GHSpace$
is an ANR for all metrizable spaces.
Thus
\Cref{thm:retract-extensor}
implies
that
$\GHSpace$
is an ANE for all metrizable spaces.

\paraid{p-6d9be79c3189}
\Cref{lem:scaling-contraction} proves that
$\GHSpace$
is contractible by scaling every distance to
$0$.
Thus
 \Cref{thm:contractible-ane}
 shows that
the contractible ANE
$\GHSpace$
is an AE for all metrizable spaces.
Finally,
using  \Cref{thm:retract-extensor},
we conclude that
$\GHSpace$
is an AR for all metrizable spaces.
\end{proof}


\begin{proof}[Proof of \Cref{thm:part-iii-main}]
\paraid{p-0b3fa815dd6f}
This is a direct consequence of \Cref{thm:absolute-extensor}.
\end{proof}


\clearpage
\part{The topology of Gromov--Hausdorff space IV: The Gromov--Hausdorff space is homeomorphic to the  Hilbert space}\label{part:hilbert}
\begin{quote}
\small
\noindent\textbf{Abstract.}
\paraid{p-4a904c863731}
We prove that the Gromov--Hausdorff space  is homeomorphic to the separable infinite-dimensional Hilbert space, by combining a discrete approximation theorem with the absolute retract property from Part~\ref{part:topology}, completeness, separability, and Toru\'nczyk's characterization.

\par\smallskip
\noindent\textbf{Keywords.} Gromov--Hausdorff space, Hilbert space, discrete approximation property, metric product.

\par\smallskip
\noindent\textbf{2020 Mathematics Subject Classification.} Primary 57N20; Secondary 54E35, 54F45.

\end{quote}

\section{Introduction to Part IV}\label{sec:intro-hilbert}
\paraid{p-96eb53c66867}
In Part~\ref{part:topology},
we proved that the Gromov--Hausdorff space
is an absolute retract.
We now prove the discrete approximation property needed to identify
this space with Hilbert space.
Let
$\GHSpace$
denote the space of isometry classes of nonempty compact metric spaces
with the ordinary unpointed Gromov--Hausdorff distance
$\GHDistance$.
Antonyan asks whether
$\GHSpace$
is homeomorphic to the separable
infinite-dimensional Hilbert space
\cite[p.~2]{Antonyan2020Euclidean}.
For each positive integer
$\FiniteSize $,
Antonyan identifies the subspace represented by compact subsets of
$\RealNumbers^\FiniteSize $
with the Hilbert cube minus a point
\cite[Corollary~5.3]{Antonyan2021Euclidean}.

\paraid{p-707b7043674a}
Previous work has explored the topology of $\GHSpace$ through explicit
families and embeddings.
The author constructed families of branching Gromov--Hausdorff geodesics
continuously parametrized by the Hilbert cube
\cite[Theorem~1.3]{Ishiki2022Branching}.
The author also constructed topological embeddings of arbitrary compact metrizable
spaces into the subsets of
$\GHSpace$
consisting of continua
\cite[Theorem~1.1]{Ishiki2022Continua},
spaces with prescribed dimensions
\cite[Theorem~1.3]{Ishiki2023FractalDimensions},
and compact metric trees
\cite[Theorem~1.1]{Ishiki2023MetricTrees}.
In each construction,
finitely many distinct values of the embedding may be prescribed.
Byakuno
\cite[arXiv v1, Corollary~1.4]{Byakuno2026Embedding}
proves that countable products of closed intervals with positive summable lengths,
equipped with the supremum metric,
admit isometric embeddings into
$\GHSpace$.

\paraid{p-1084d3b137e6}
We also recall the dimension of the subspace of finite metric spaces.
Nakajima,
Yamauchi,
and Zava
\cite[arXiv v1, Theorem~A]{NakajimaYamauchiZava2025}
prove that the subspace of spaces with at most
$\FiniteSize$
points has topological dimension
$\FiniteSize(\FiniteSize-1)/2$.
Nakajima and Shioya
\cite[arXiv v1, Main Theorem~1.2]{NakajimaShioyaCompactification}
construct a compact metrizable space containing a dense topological copy of
$\GHSpace$.

\paraid{p-c46ed027b219}
A related Hilbert space classification concerns the Urysohn universal
metric space.
Uspenskij
\cite[arXiv v1, Theorem~2.1]{Uspenskij2004Urysohn}
proves that this space is homeomorphic to
$\SeparableHilbertSpace$.

\paraid{p-6d7aa84e3e0d}
Our approximation theorem concerns a countable family of prescribed maps
from compact metrizable domains.
We approximate all these maps so that every point of
$\GHSpace$
has a neighborhood meeting at most one of their images.
The construction uses products with finite equilateral spaces,
independently of Parts~\ref{part:measures}--\ref{part:topology}.
We combine this discrete approximation with the absolute retract property
to prove the following theorem.

\begin{theorem}\label{thm:part-iv-main}
\paraid{p-13f85da9e5b7}
The space
$\GHSpace$
is homeomorphic to the real Hilbert space
$\SeparableHilbertSpace$.
\end{theorem}

\paraid{p-add73a912728}
The classification concerns only topology and does not preserve the geometric
information carried by the Gromov--Hausdorff distance.
Ivanov and Tuzhilin
\cite[Main Theorem]{IvanovTuzhilin2019Isometries}
prove that every surjective isometry from
$\GHSpace$ to itself
is the identity.
Zava
\cite[arXiv v3, Theorem~B]{Zava2025Coarse}
proves that the subspace of finite metric spaces cannot be coarsely
embedded into any Hilbert space.
Consequently,
a homeomorphism in \Cref{thm:part-iv-main} cannot be bi-Lipschitz.

\paraid{p-60f77024d745}
For a sequence of compact metrizable domains
$\{\ApproximationDomain_\FamilyIndex\}_{\FamilyIndex\in\NonnegativeIntegers}$,
continuous maps
$\InputMap_\FamilyIndex\colon\ApproximationDomain_\FamilyIndex\to\GHSpace$,
and an open cover
$\OpenCover$
of
$\GHSpace$,
we construct continuous maps
$\ApproximatingMap_\FamilyIndex\colon\ApproximationDomain_\FamilyIndex\to\GHSpace$
such that
\[
 \forall\FamilyIndex\in\NonnegativeIntegers\
 \forall\DomainPoint\in\ApproximationDomain_\FamilyIndex\
 \exists U\in\OpenCover,\qquad
 \{\InputMap_\FamilyIndex(\DomainPoint),
   \ApproximatingMap_\FamilyIndex(\DomainPoint)\}\subset U,
\]
and
\[
 \{\ApproximatingMap_\FamilyIndex(\ApproximationDomain_\FamilyIndex)
   \}_{\FamilyIndex\in\NonnegativeIntegers}
 \text{ is discrete in }\GHSpace.
\]
The first condition means that the maps are
$\OpenCover$-close.
The second means that every point of
$\GHSpace$
has a neighborhood meeting at most one image.
This is \Cref{thm:discrete-approximation}.
Taking
$\ApproximationDomain_\FamilyIndex=\HilbertCube
 =\UnitInterval^{\NonnegativeIntegers}$
for every
$\FamilyIndex$
yields the approximation
condition in Toru\'nczyk's characterization,
stated in \Cref{thm:torunczyk}
and taken from
\cite[p.~248, assertion~(i)]{Torunczyk1981},
with the correction in
\cite[Section~C]{Torunczyk1985}.
Together with completeness,
separability,
and the absolute retract property,
this condition proves \Cref{thm:part-iv-main}.

\paraid{p-01c02735fab9}
The main theorem also implies that every point has a basis of open neighborhoods
that strongly deformation retract onto that point.
We do not know whether these neighborhoods can be chosen to be balls for
$\GHDistance$.

\medskip
\noindent\textbf{Dependence on the other parts.}
\paraid{p-518fca37991d}
\Cref{thm:discrete-approximation} uses the common preliminaries
and the metric constructions recalled in Section~\ref{sec:prelim-hilbert}.
The absolute retract result \Cref{thm:part-iii-main} enters only when we
apply Toru\'nczyk's characterization in
Subsection~\ref{subsec:approximation-hilbert-recognition}.
Thus the proof of the classification combines the approximation in this
part with the result of Parts~\ref{part:measures}--\ref{part:topology}.

\medskip
\noindent\textbf{Organization.}
\paraid{p-c7e99d59b5f0}
In Section~\ref{sec:prelim-hilbert},
we define discrete families and state the Hilbert space characterization.
In Subsection~\ref{subsec:approximation-products},
we recall the Gromov--Hausdorff estimate for metric products.
In Subsection~\ref{subsec:approximation-discrete},
we construct approximations of maps from countably many
compact domains whose images form a discrete family,
proving \Cref{thm:discrete-approximation}.
In Subsection~\ref{subsec:approximation-hilbert-recognition},
we apply Toru\'nczyk's characterization to prove
\Cref{thm:part-iv-main} and obtain strong local contractions.
In \Cref{sec:questions},
we conclude with questions about related spaces
and the geometry of Gromov--Hausdorff balls.

\medskip
\noindent\textbf{Conventions and notation.}
\paraid{p-bb6d5f435441}
All compact metric spaces are nonempty,
and Banach and Hilbert spaces are real.
We identify compact metric spaces with their isometry classes in
$\GHSpace$.
An isometry is a surjective isometric embedding.
For a metric space
$\MetricSpace{\BaseCarrier}{\MetricSymbol}$,
we write
$\CompactHyperspace(\BaseCarrier)$
for its nonempty compact subsets,
equipped with the Hausdorff metric
$\HausdorffDistance{\MetricSymbol}$.
For a compact space with an isometric embedding into an ambient metric space,
we identify the compact space with its image and use the restricted ambient metric.
The notation
$\IdentityMap$
denotes the identity map on its specified domain.
Sequences are indexed by
$\NonnegativeIntegers=\{0,1,2,\ldots\}$.


\section{Preliminaries}\label{sec:prelim-hilbert}
\paraid{p-6c5ffdcadadc}
We define discrete families and approximation with respect to an open cover.
We recall metric products and state Toru\'nczyk's characterization of Hilbert space.

\subsection{Discrete families and close maps}\label{subsec:hilbert-preliminaries-discrete}
\paraid{p-0e8bda81dbda}
To state the approximation condition in \Cref{thm:torunczyk},
we recall discreteness of a family of subsets and closeness of maps
with respect to an open cover.

\begin{definition}\label{def:discrete-family}
\paraid{p-360939bb32bd}
A family
$\{\ClosedDomain _\FamilyIndex \}_{\FamilyIndex \in \DomainIndexSet }$
of subsets of a space
$\RecognitionTarget $
is \emph{discrete in}
$\RecognitionTarget $
if every point of
$\RecognitionTarget $
has a neighborhood meeting at most one member of the family
$\{\ClosedDomain_\FamilyIndex\}_{\FamilyIndex\in\DomainIndexSet}$.
We call the family
$\{\ClosedDomain_\FamilyIndex\}_{\FamilyIndex\in\DomainIndexSet}$
\emph{locally finite} if every point of
$\RecognitionTarget$
has a neighborhood meeting only finitely many members.
These requirements apply also at points outside the union.
For an open cover
$\OpenCover$
of
$\RecognitionTarget $,
two maps
$\InputMap ,\ApproximatingMap \colon \CommonMapDomain \to \RecognitionTarget $
are
\emph{$\OpenCover$-close}
if for each
$\DomainPoint \in \CommonMapDomain $
there exists
$U\in\OpenCover$
such that
\[
 \{\InputMap(\DomainPoint),\ApproximatingMap(\DomainPoint)\}\subset U.
\]
\end{definition}

\subsection{Metric products}\label{subsec:approximation-products}
\paraid{p-94e60cbaf8be}
We recall a Gromov--Hausdorff estimate for metric products
used in the approximation construction.
For compact metric spaces
$\MetricSpace{\FirstFactor }{\FirstFactorMetric }$
and
$\MetricSpace{\SecondFactor }{\SecondFactorMetric }$,
write
$\FirstFactorMetric \vee \SecondFactorMetric $
for the max metric on
$\FirstFactor \times \SecondFactor $.
For every
$\FirstFactorPoint,\OtherFirstFactorPoint\in\FirstFactor$
and
$\SecondFactorPoint,\OtherSecondFactorPoint\in\SecondFactor$,
put
\[
  (\FirstFactorMetric \vee \SecondFactorMetric )((\FirstFactorPoint ,\SecondFactorPoint ),(\OtherFirstFactorPoint ,\OtherSecondFactorPoint ))
  =\max\{\FirstFactorMetric (\FirstFactorPoint ,\OtherFirstFactorPoint ),\SecondFactorMetric (\SecondFactorPoint ,\OtherSecondFactorPoint )\}.
\]

\begin{lemma}\label{lem:metric-products}
\paraid{p-0e0f1f1764e1}
Let
$\MetricSpace{\FirstFactor }{\FirstFactorMetric }$,
$\MetricSpace{\SecondFactor }{\SecondFactorMetric }$,
$\MetricSpace{\ComparisonFirstFactor }{\ComparisonFirstFactorMetric }$,
and
$\MetricSpace{\ComparisonSecondFactor }{\ComparisonSecondFactorMetric }$
be nonempty compact metric spaces.
Then
\begin{equation}\label{eq:max-product}
 \begin{aligned}
  &\GHDistance(\MetricSpace{\FirstFactor \times \SecondFactor }{\FirstFactorMetric \vee \SecondFactorMetric },
        \MetricSpace{\ComparisonFirstFactor \times \ComparisonSecondFactor }{\ComparisonFirstFactorMetric \vee \ComparisonSecondFactorMetric })\\
  &\qquad\leq\max\{\GHDistance(\MetricSpace{\FirstFactor }{\FirstFactorMetric },\MetricSpace{\ComparisonFirstFactor }{\ComparisonFirstFactorMetric }),
                       \GHDistance(\MetricSpace{\SecondFactor }{\SecondFactorMetric },\MetricSpace{\ComparisonSecondFactor }{\ComparisonSecondFactorMetric })\}.
 \end{aligned}
\end{equation}
\end{lemma}

\begin{proof}
\paraid{p-dd7692535fb1}
Let
$\Correspondence_\FirstFactor\subset
 \FirstFactor\times\ComparisonFirstFactor$
and
$\Correspondence_\SecondFactor\subset
 \SecondFactor\times\ComparisonSecondFactor$
be correspondences.
By rearranging coordinates,
we regard their product
$\Correspondence_\times
 =\Correspondence_\FirstFactor\times\Correspondence_\SecondFactor$
as a correspondence between
$\FirstFactor\times\SecondFactor$
and
$\ComparisonFirstFactor\times\ComparisonSecondFactor$.
For real numbers
$a,b,c,d$,
we have
\[
 |\max\{a,b\}-\max\{c,d\}|
 \leq\max\{|a-c|,|b-d|\}.
\]
Hence
$\dis\Correspondence_\times
 \leq\max\{\dis\Correspondence_\FirstFactor,
            \dis\Correspondence_\SecondFactor\}$.
Taking the infimum over the two factor correspondences proves
\eqref{eq:max-product} by \eqref{eq:correspondence}.
\end{proof}

% \paraid{p-4fe5350ab6f2}
% \paraid{p-990fe723df47}
% \paraid{p-593e7df8a9f5}
% \paraid{p-cba33d789d31}
% \paraid{p-539467a269ef}
% \paraid{p-117a415627d1}
% \paraid{p-de5425fa9621}
% \paraid{p-9eac5ece7654}
% \paraid{p-0596666bd9f9}

\subsection{A characterization of Hilbert space}\label{subsec:hilbert-preliminaries-recognition}
\paraid{p-a8af9823b28c}
We use the following characterization of Hilbert space.
It is Toru\'nczyk's theorem
\cite[p.~248, condition~(i)]{Torunczyk1981},
with its correction
\cite[Section~C]{Torunczyk1985}.

\begin{theorem}[{Toru\'nczyk, \cite[p.~248, condition~(i)]{Torunczyk1981}, \cite[Section~C]{Torunczyk1985}}]\label{thm:torunczyk}
\paraid{p-e19e6ccab7f7}
Let
$\MetricSpace{\RecognitionTarget }{\TargetMetric }$
be a nonempty complete separable metric AR,
and let
$\HilbertCube =\UnitInterval^{\NonnegativeIntegers}$
be the Hilbert cube.
Equip
$\NonnegativeIntegers$
with the discrete topology.
Then
$\RecognitionTarget $
is homeomorphic to the real Hilbert space
$\SeparableHilbertSpace$
if and only if,
for every continuous map
$\InputMap \colon \NonnegativeIntegers\times \HilbertCube \to \RecognitionTarget $
and every open cover
$\OpenCover$
of
$\RecognitionTarget $,
there exists a continuous map
$\ApproximatingMap \colon \NonnegativeIntegers\times \HilbertCube \to \RecognitionTarget $
that is
$\OpenCover$-close to
$\InputMap$
and whose image family
$\{\ApproximatingMap (\{\FamilyIndex \}\times \HilbertCube )\}_{\FamilyIndex \in\NonnegativeIntegers}$
is discrete in
$\RecognitionTarget $.
\end{theorem}


\section{Discrete approximation and Hilbert space topology}\label{sec:approximation}
\paraid{p-3d2b7bf36880}
To apply Toru\'nczyk's characterization of Hilbert space,
we prove the discrete approximation property for maps from countably many Hilbert cubes.

\subsection{Discrete approximation for compact domains}\label{subsec:approximation-discrete}
\paraid{p-39f5655ab8f2}
We now construct the approximating maps and prove discreteness at
every point of the ambient Gromov--Hausdorff space.

\paraid{p-db4cf11c2a9d}
We first choose a positive Lipschitz scale adapted to an open cover.

\begin{lemma}\label{lem:cover-adapted-scale}
\paraid{p-90d3aa515184}
Let
$\OpenCover$
be an open cover of
$\GHSpace$.
Then there exists a
$1/16$-Lipschitz function
$\CoverScale\colon\GHSpace\to(0,1/16]$
with the following property.
For every
$\MetricSpace{\BaseCarrier}{\MetricSymbol},
 \MetricSpace{\ComparisonCarrier}{\ComparisonMetric}\in\GHSpace$,
if
\[
 \GHDistance(\MetricSpace{\BaseCarrier}{\MetricSymbol},
             \MetricSpace{\ComparisonCarrier}{\ComparisonMetric})
 <\CoverScale\MetricSpace{\BaseCarrier}{\MetricSymbol},
\]
then there exists
$U\in\OpenCover$
with
\[
 \{\MetricSpace{\BaseCarrier}{\MetricSymbol},
   \MetricSpace{\ComparisonCarrier}{\ComparisonMetric}\}\subset U.
\]
\end{lemma}

\begin{proof}
\paraid{p-8da60d9b543c}
Use the convention that distance to the empty set is
$\infty$.
For every
$\MetricSpace{\BaseCarrier}{\MetricSymbol}\in\GHSpace$,
put
\begin{equation}\label{eq:cover-scale-2}
 \CoverScale\MetricSpace{\BaseCarrier}{\MetricSymbol}
 =\frac1{16}\sup_{\OpenSubset\in\OpenCover}
   \min\left\{1,\dist_{\GHDistance}
       \bigl(\MetricSpace{\BaseCarrier}{\MetricSymbol},\GHSpace\setminus\OpenSubset\bigr)\right\}.
\end{equation}
Each truncated distance function is
$1$-Lipschitz,
and their supremum is positive because
$\OpenCover$
is an open cover.
Thus,
for every
$\MetricSpace{\BaseCarrier}{\MetricSymbol},
 \MetricSpace{\ComparisonCarrier}{\ComparisonMetric}\in\GHSpace$,
we have
\begin{equation}\label{eq:scale-lipschitz-1}
 0<\CoverScale\MetricSpace{\BaseCarrier}{\MetricSymbol}\leq\frac1{16},
\end{equation}
and
\begin{equation}\label{eq:scale-lipschitz-2}
 \abs{\CoverScale\MetricSpace{\BaseCarrier}{\MetricSymbol}
      -\CoverScale\MetricSpace{\ComparisonCarrier}{\ComparisonMetric}}
 \leq\frac1{16}\GHDistance(\MetricSpace{\BaseCarrier}{\MetricSymbol},
                          \MetricSpace{\ComparisonCarrier}{\ComparisonMetric}).
\end{equation}
Now let
$\MetricSpace{\BaseCarrier}{\MetricSymbol},
 \MetricSpace{\ComparisonCarrier}{\ComparisonMetric}\in\GHSpace$
satisfy
$\GHDistance(\MetricSpace{\BaseCarrier}{\MetricSymbol},
             \MetricSpace{\ComparisonCarrier}{\ComparisonMetric})
 <\CoverScale\MetricSpace{\BaseCarrier}{\MetricSymbol}$.
By \eqref{eq:cover-scale-2},
there exists
$\mathscr O\in\OpenCover$
such that
\[
 \GHDistance(\MetricSpace{\BaseCarrier}{\MetricSymbol},
             \MetricSpace{\ComparisonCarrier}{\ComparisonMetric})
 <\CoverScale\MetricSpace{\BaseCarrier}{\MetricSymbol}
 <\dist_{\GHDistance}
   \bigl(\MetricSpace{\BaseCarrier}{\MetricSymbol},\GHSpace\setminus\mathscr O\bigr).
\]
Both spaces therefore belong to
$\mathscr O$.
\end{proof}

\begin{theorem}\label{thm:discrete-approximation}
\paraid{p-07c6df8fb234}
Let
$\{\ApproximationDomain _\FamilyIndex \}_{\FamilyIndex \in\NonnegativeIntegers}$
be a sequence of compact metrizable spaces.
For each
$\FamilyIndex\in\NonnegativeIntegers$,
let
$\InputMap _\FamilyIndex \colon \ApproximationDomain _\FamilyIndex \to\GHSpace$
be continuous.
Let
$\OpenCover$
be an open cover of
$\GHSpace$.
Then there exist continuous maps
$\ApproximatingMap _\FamilyIndex \colon \ApproximationDomain _\FamilyIndex \to\GHSpace$
such that each
$\ApproximatingMap _\FamilyIndex $
is
$\OpenCover$-close
to
$\InputMap _\FamilyIndex $
and the family
$\{\ApproximatingMap _\FamilyIndex (\ApproximationDomain _\FamilyIndex )\}_{\FamilyIndex \in\NonnegativeIntegers}$
is discrete in
$\GHSpace$.
\end{theorem}

\begin{proof}
\paraid{p-bce29ad7f159}
We first dispose of the empty domains.
For each empty domain
$\ApproximationDomain_\FamilyIndex$,
choose
$\ApproximatingMap_\FamilyIndex$
to be the unique empty map.
Empty images do not affect closeness or discreteness,
so the assertion holds if every domain
$\ApproximationDomain_\FamilyIndex$
is empty.
We now consider only the nonempty domains and put
$\DomainIndexSet
 =\{\FamilyIndex\in\NonnegativeIntegers\mid
    \ApproximationDomain_\FamilyIndex\ne\emptyset\}$.
For each
$\FamilyIndex\in\DomainIndexSet$,
we define
a positive integer
$\EquilateralSize_{\FamilyIndex}$
recursively,
and
for
$\DomainPoint\in\ApproximationDomain_{\FamilyIndex}$,
we construct the approximating map
$\ApproximatingMap_{\FamilyIndex}(z)$
by taking the product of each value
$\InputMap_{\FamilyIndex}(\DomainPoint)$
with finite equilateral spaces of cardinality
$\EquilateralSize_{\FamilyIndex}$.
We then prove that the resulting family of images is locally finite and discrete.

\ProofStep{Product approximations.}\label{step:discrete-products}
\paraid{p-021b6255b9ed}
By \Cref{lem:cover-adapted-scale},
choose
$\CoverScale\colon\GHSpace\to(0,1/16]$
satisfying \eqref{eq:scale-lipschitz-2} and the cover condition in that lemma.

\paraid{p-71bd13b39f7a}
For an integer
$\EquilateralSize\geq2$,
let
$\EquilateralCarrier_{\EquilateralSize}=\{0,\ldots,\EquilateralSize-1\}$.
For
$\ProductScale>0$,
let
$\EquilateralMetric_{\EquilateralSize,\ProductScale}$
be the equilateral metric on
$\EquilateralCarrier_{\EquilateralSize}$
with off-diagonal distance
$\ProductScale$.
For each integer
$\EquilateralSize\geq2$,
define
$\ProductApproximation_{\EquilateralSize}\colon\GHSpace\to\GHSpace$
as follows.
For
$\MetricSpace{\BaseCarrier}{\MetricSymbol}\in\GHSpace$,
put
\begin{equation}\label{eq:product-approximation-2}
 \ProductScale=\CoverScale\MetricSpace{\BaseCarrier}{\MetricSymbol}
\end{equation}
and
\begin{equation}\label{eq:product-approximation-1}
 \ProductApproximation_{\EquilateralSize}\MetricSpace{\BaseCarrier}{\MetricSymbol}
 =\MetricSpace{\BaseCarrier\times\EquilateralCarrier_{\EquilateralSize}}
              {\MetricSymbol\vee\EquilateralMetric_{\EquilateralSize,\ProductScale}}.
\end{equation}
The graph of the projection onto
$\BaseCarrier$
is a correspondence of distortion at most
$\ProductScale$.
Hence
we have
\begin{equation}\label{eq:product-error}
 \GHDistance(\ProductApproximation_{\EquilateralSize}\MetricSpace{\BaseCarrier}{\MetricSymbol},
             \MetricSpace{\BaseCarrier}{\MetricSymbol})
 \leq\frac12\CoverScale\MetricSpace{\BaseCarrier}{\MetricSymbol}.
\end{equation}

\paraid{p-7490df62137b}
For an integer
$\EquilateralSize\geq2$
and spaces
$\MetricSpace{\BaseCarrier}{\MetricSymbol},
 \MetricSpace{\ComparisonCarrier}{\ComparisonMetric}\in\GHSpace$,
pair equal labels in the two equilateral factors.
As in the proof of
\cite[Proposition~5.3]{Ishiki2022Branching},
the product correspondence and \eqref{eq:scale-lipschitz-2} imply
\begin{equation}\label{eq:equilateral-product-lipschitz}
 \begin{aligned}
 &\GHDistance(\ProductApproximation_{\EquilateralSize}\MetricSpace{\BaseCarrier}{\MetricSymbol},
              \ProductApproximation_{\EquilateralSize}\MetricSpace{\ComparisonCarrier}{\ComparisonMetric})\\
 &\quad\leq\max\left\{
     \GHDistance(\MetricSpace{\BaseCarrier}{\MetricSymbol},
                 \MetricSpace{\ComparisonCarrier}{\ComparisonMetric}),
     \frac12\left|\CoverScale\MetricSpace{\BaseCarrier}{\MetricSymbol}
                 -\CoverScale\MetricSpace{\ComparisonCarrier}{\ComparisonMetric}\right|
                  \right\}\\
 &\quad\leq\GHDistance(\MetricSpace{\BaseCarrier}{\MetricSymbol},
                      \MetricSpace{\ComparisonCarrier}{\ComparisonMetric}).
 \end{aligned}
\end{equation}
Thus every map
$\ProductApproximation_{\EquilateralSize}$
is
$1$-Lipschitz.

\ProofStep{Choice of the cardinalities and separation of the images.}\label{step:discrete-separation}
\paraid{p-04700bc9ef5a}
We choose the positive integers
$\EquilateralSize_{\FamilyIndex}$
and the maps
$\ApproximatingMap_{\FamilyIndex}$
recursively over
$\FamilyIndex\in\DomainIndexSet$.
Fix
$\FamilyIndex\in\DomainIndexSet$
and assume that the maps with smaller indices have been constructed
and have compact images.
Put
\begin{equation}\label{eq:previous-approximation-images}
 \PreviousApproximationImages_{\FamilyIndex}
 =\bigcup_{\substack{k\in\DomainIndexSet\\k<\FamilyIndex}}
       \ApproximatingMap_k(\ApproximationDomain_k).
\end{equation}
This set is compact,
and
$\PreviousApproximationImages_0=\emptyset$.
Since
$\ApproximationDomain_{\FamilyIndex}$
is nonempty and compact,
the number
\begin{equation}\label{eq:minimum-product-scale}
 \ProductScale_{\FamilyIndex}
 =\min_{\DomainPoint\in\ApproximationDomain_{\FamilyIndex}}
       \CoverScale(\InputMap_{\FamilyIndex}(\DomainPoint))
\end{equation}
is positive.

\paraid{p-5bd0a1e059c4}
Compactness of
$\PreviousApproximationImages_{\FamilyIndex}$
and Corollary~\ref{cor:uniform-separated-sets}
imply that there is an integer
$\PackingBound_{\FamilyIndex}\geq0$
such that every
$\ProductScale_{\FamilyIndex}/2$-separated
subset of every space in
$\PreviousApproximationImages_{\FamilyIndex}$
has cardinality at most
$\PackingBound_{\FamilyIndex}$.
If
$\PreviousApproximationImages_{\FamilyIndex}=\emptyset$,
take
$\PackingBound_{\FamilyIndex}=0$.

\paraid{p-82b288935ffb}
Choose an integer
$\EquilateralSize_{\FamilyIndex}$
larger than all previously chosen cardinalities and satisfying
\begin{equation}\label{eq:equilateral-cardinality-choice}
 \EquilateralSize_{\FamilyIndex}
 >\max\{\FamilyIndex+1,\PackingBound_{\FamilyIndex}\}.
\end{equation}
Define
$\ApproximatingMap_{\FamilyIndex}\colon\ApproximationDomain_{\FamilyIndex}\to\GHSpace$
by
\begin{equation}\label{eq:discrete-maps}
 \ApproximatingMap_{\FamilyIndex}
 =\ProductApproximation_{\EquilateralSize_{\FamilyIndex}}\circ\InputMap_{\FamilyIndex}.
\end{equation}
The map
$\ApproximatingMap_{\FamilyIndex}$
is continuous by \eqref{eq:equilateral-product-lipschitz},
and its image is compact.
For every
$\DomainPoint\in\ApproximationDomain_{\FamilyIndex}$,
equation \eqref{eq:product-error} implies
\begin{equation}\label{eq:approximation-error}
 \GHDistance(\ApproximatingMap_{\FamilyIndex}(\DomainPoint),
             \InputMap_{\FamilyIndex}(\DomainPoint))
 \leq\frac12\CoverScale(\InputMap_{\FamilyIndex}(\DomainPoint))
 <\CoverScale(\InputMap_{\FamilyIndex}(\DomainPoint)).
\end{equation}
Thus
$\ApproximatingMap_{\FamilyIndex}$
and
$\InputMap_{\FamilyIndex}$
are
$\OpenCover$-close.

\paraid{p-dbde286ff7a5}
Every product in
$\ApproximatingMap_{\FamilyIndex}(\ApproximationDomain_{\FamilyIndex})$
contains a copy of its equilateral factor,
obtained by fixing a point in its first factor.
By \eqref{eq:minimum-product-scale},
this copy has
$\EquilateralSize_{\FamilyIndex}$
points separated by at least
$\ProductScale_{\FamilyIndex}$.
If the GH distance from such a product to a space in
$\PreviousApproximationImages_{\FamilyIndex}$
were less than
$\ProductScale_{\FamilyIndex}/4$,
a correspondence of distortion less than
$\ProductScale_{\FamilyIndex}/2$
would send these points to distinct points separated by more than
$\ProductScale_{\FamilyIndex}/2$
in that space.
This would imply
$\EquilateralSize_{\FamilyIndex}\leq\PackingBound_{\FamilyIndex}$,
contrary to \eqref{eq:equilateral-cardinality-choice}.
Consequently,
for every
$\DomainPoint\in\ApproximationDomain_{\FamilyIndex}$
and every
$\MetricSpace{\ComparisonCarrier}{\ComparisonMetric}\in\PreviousApproximationImages_{\FamilyIndex}$,
we have
\begin{equation}\label{eq:recursive-image-separation}
 \GHDistance(\ApproximatingMap_{\FamilyIndex}(\DomainPoint),
             \MetricSpace{\ComparisonCarrier}{\ComparisonMetric})
 \geq\frac{\ProductScale_{\FamilyIndex}}4.
\end{equation}
This completes the recursive construction and proves that the images are pairwise disjoint.

\paraid{p-4a46362ad0ff}
Figure~\ref{fig:discrete-approximation} shows how the scales and cardinalities
control the approximation and the separation of the images.
\begin{figure}[H]
\centering
\begin{tikzpicture}[>=Stealth,font=\small]
\node at (0,1.6) {$\BaseCarrier$};
\draw (0,0) ellipse (1.15 and 0.4);
\fill (-0.55,0.04) circle (1.1pt);
\fill (0.5,-0.12) circle (1.1pt);
\fill (0.05,0.08) circle (1.5pt);
\node[above,inner sep=3pt] at (0.05,0.08) {$\BasePoint$};
\draw[->] (1.5,0) -- node[above]
 {$\times\EquilateralCarrier_3$} (3.1,0);
\node at (4.65,1.6) {$\BaseCarrier\times\EquilateralCarrier_3$};
\fill[black!8] (4.5,-1.15) rectangle (4.9,1.3);
\foreach \level/\index in {0.95/0,0/1,-0.95/2}
{
 \draw (4.65,\level) ellipse (1.15 and 0.4);
 \fill (4.1,{\level+0.04}) circle (1.1pt);
 \fill (5.15,{\level-0.12}) circle (1.1pt);
 \fill (4.7,{\level+0.08}) circle (1.5pt);
 \node[right] at (5.95,\level)
  {$\BaseCarrier\times\{\index\}$};
}
\node[align=center] at (4.7,-1.95)
 {$\{\BasePoint\}\times\EquilateralCarrier_3$\\
  pairwise distance $\ProductScale$};
\end{tikzpicture}
\par\medskip
\begin{tikzpicture}[
 >=Stealth,
 font=\small,
 role/.style={draw,align=center,text width=3.4cm,
              minimum height=2.1cm,inner sep=5pt,anchor=north}
]
\node at (4,0.9)
 {$\MetricSpace{\ClosedDomain}{\SubspaceMetric}
   =\ProductApproximation_{\EquilateralSize_{\FamilyIndex}}
      \MetricSpace{\BaseCarrier}{\MetricSymbol},
   \qquad
   \ProductScale=\CoverScale\MetricSpace{\BaseCarrier}{\MetricSymbol}$};
\node (formula) at (4,0)
 {$\MetricSpace{\ClosedDomain}{\SubspaceMetric}
   =\MetricSpace{\BaseCarrier\times\EquilateralCarrier_{\EquilateralSize_{\FamilyIndex}}}
                {\MetricSymbol\vee\EquilateralMetric_{\EquilateralSize_{\FamilyIndex},\ProductScale}}$};
\node[role] (approximation) at (0,-1.7)
 {small scale $\ProductScale$\\
  $\begin{gathered}
    \GHDistance(\MetricSpace{\ClosedDomain}{\SubspaceMetric},
                 \MetricSpace{\BaseCarrier}{\MetricSymbol})\\
    \leq\ProductScale/2
   \end{gathered}$\\
  approximation};
\node[role] (separation) at (4,-1.7)
 {$\EquilateralSize_{\FamilyIndex}>\PackingBound_{\FamilyIndex}$\\
  recursive choice\\
  disjoint images};
\node[role] (packing) at (8,-1.7)
 {$\EquilateralSize_{\FamilyIndex}\to\infty$\\
  positive scale bound\\
  no accumulation};
\draw[->,shorten <=2pt]
 ([xshift=-1.8cm]formula.south) .. controls +(0,-0.55) and +(0,0.55) .. (approximation.north);
\draw[->,shorten <=2pt]
 (formula.south) -- (separation.north);
\draw[->,shorten <=2pt]
 ([xshift=1.8cm]formula.south) .. controls +(0,-0.55) and +(0,0.55) .. (packing.north);
\end{tikzpicture}
\caption{The upper diagram shows three copies of
$\BaseCarrier$
in a product with an equilateral space.
The marked fiber has
$\EquilateralSize$
points at pairwise distance
$\ProductScale$
for a general factor
$\EquilateralCarrier_\EquilateralSize$.
The scales and cardinalities of the equilateral factors control the
 discrete approximation.
The bound
$\PackingBound_{\FamilyIndex}$
applies to subsets separated by
$\ProductScale_{\FamilyIndex}/2$
in the previously constructed images,
where
$\ProductScale_{\FamilyIndex}$
is the minimum input scale in \eqref{eq:minimum-product-scale}.
The recursive choice \eqref{eq:equilateral-cardinality-choice}
separates the images.
Along a convergent sequence of approximating spaces,
the scales have a positive lower bound,
so the divergence of the cardinalities contradicts
Corollary~\ref{cor:uniform-separated-sets}.}
\label{fig:discrete-approximation}
\end{figure}


\ProofStep{Local finiteness.}\label{step:discrete-local-finiteness}
\paraid{p-86b6e1b5bb9e}
For the sake of contradiction,
suppose that the family
$\{\ApproximatingMap_{\FamilyIndex}(\ApproximationDomain_{\FamilyIndex})\}_{\FamilyIndex\in\DomainIndexSet}$
is not locally finite at some
$\MetricSpace{\LimitCarrier}{\LimitMetric}\in\GHSpace$.
Then every neighborhood of this point meets infinitely many images.
For each
$\SequenceIndex\in\NonnegativeIntegers$,
choose an index
$\FamilyIndex_{\SequenceIndex}$
larger than all previously chosen indices and a point
$\DomainPoint_{\SequenceIndex}\in\ApproximationDomain_{\FamilyIndex_{\SequenceIndex}}$
whose image lies in the GH ball of radius
$2^{-\SequenceIndex}$
about
$\MetricSpace{\LimitCarrier}{\LimitMetric}$.
For each
$\SequenceIndex\in\NonnegativeIntegers$,
put
$\InputMap_{\FamilyIndex_{\SequenceIndex}}(\DomainPoint_{\SequenceIndex})
=\MetricSpace{\BaseCarrier_{\SequenceIndex}}{\MetricSymbol_{\SequenceIndex}}$,
and
$\ApproximatingMap_{\FamilyIndex_{\SequenceIndex}}(\DomainPoint_{\SequenceIndex})
=\MetricSpace{\ComparisonCarrier_{\SequenceIndex}}{\ComparisonMetric_{\SequenceIndex}}$,
given by \eqref{eq:product-approximation-1} and \eqref{eq:discrete-maps}.
Then
\begin{equation}\label{eq:equilateral-image-convergence}
 \MetricSpace{\ComparisonCarrier_{\SequenceIndex}}{\ComparisonMetric_{\SequenceIndex}}
 \to\MetricSpace{\LimitCarrier}{\LimitMetric}.
\end{equation}
For each
$\SequenceIndex\in\NonnegativeIntegers$,
put
$\CoverScale_{\SequenceIndex}
 =\CoverScale\MetricSpace{\BaseCarrier_{\SequenceIndex}}{\MetricSymbol_{\SequenceIndex}}$.
For every
$\SequenceIndex\in\NonnegativeIntegers$,
equations \eqref{eq:scale-lipschitz-2} and \eqref{eq:approximation-error} imply
\begin{equation}\label{eq:scale-lower-bound}
 \begin{aligned}
 \CoverScale\MetricSpace{\LimitCarrier}{\LimitMetric}
 &\leq\CoverScale_{\SequenceIndex}
      +\frac1{16}\GHDistance(
         \MetricSpace{\BaseCarrier_{\SequenceIndex}}{\MetricSymbol_{\SequenceIndex}},
         \MetricSpace{\LimitCarrier}{\LimitMetric})\\
 &\leq\CoverScale_{\SequenceIndex}
      +\frac1{16}\Bigl(\GHDistance(
         \MetricSpace{\BaseCarrier_{\SequenceIndex}}{\MetricSymbol_{\SequenceIndex}},
         \MetricSpace{\ComparisonCarrier_{\SequenceIndex}}{\ComparisonMetric_{\SequenceIndex}})
         +\GHDistance(
         \MetricSpace{\ComparisonCarrier_{\SequenceIndex}}{\ComparisonMetric_{\SequenceIndex}},
         \MetricSpace{\LimitCarrier}{\LimitMetric})\Bigr)\\
 &\leq\CoverScale_{\SequenceIndex}
      +\frac1{16}\Bigl(\frac12\CoverScale_{\SequenceIndex}
         +\GHDistance(
         \MetricSpace{\ComparisonCarrier_{\SequenceIndex}}{\ComparisonMetric_{\SequenceIndex}},
         \MetricSpace{\LimitCarrier}{\LimitMetric})\Bigr)\\
 &\leq\frac{33}{32}\CoverScale_{\SequenceIndex}
      +\frac1{16}\GHDistance(
         \MetricSpace{\ComparisonCarrier_{\SequenceIndex}}{\ComparisonMetric_{\SequenceIndex}},
         \MetricSpace{\LimitCarrier}{\LimitMetric}).
 \end{aligned}
\end{equation}
Put
$\SeparationScale=\frac12\CoverScale\MetricSpace{\LimitCarrier}{\LimitMetric}>0$.
By \eqref{eq:equilateral-image-convergence} and \eqref{eq:scale-lower-bound},
for all sufficiently large
$\SequenceIndex$
we have
\begin{equation}\label{eq:equilateral-positive-scale}
 \CoverScale_{\SequenceIndex}\geq\SeparationScale.
\end{equation}

\paraid{p-866cf4c6571b}
For every
$\SequenceIndex\in\NonnegativeIntegers$,
choose
$\BasePoint_{\SequenceIndex}\in\BaseCarrier_{\SequenceIndex}$
and define
\[
 \SeparatedSet_{\SequenceIndex}
 =\{(\BasePoint_{\SequenceIndex},\EquilateralPoint)\mid
     \EquilateralPoint\in\EquilateralCarrier_{\EquilateralSize_{\FamilyIndex_{\SequenceIndex}}}\}
 \subset\ComparisonCarrier_{\SequenceIndex}.
\]
For all sufficiently large
$\SequenceIndex$
and all distinct
$\EquilateralPoint_0,\EquilateralPoint_1
 \in\EquilateralCarrier_{\EquilateralSize_{\FamilyIndex_{\SequenceIndex}}}$,
equation \eqref{eq:equilateral-positive-scale} implies
\[
 \ComparisonMetric_{\SequenceIndex}
 ((\BasePoint_{\SequenceIndex},\EquilateralPoint_0),
  (\BasePoint_{\SequenceIndex},\EquilateralPoint_1))
 =\max\{0,\CoverScale_{\SequenceIndex}\}
 =\CoverScale_{\SequenceIndex}\geq\SeparationScale.
\]
Thus
$\SeparatedSet_{\SequenceIndex}$
is
$\SeparationScale$-separated
for all sufficiently large
$\SequenceIndex$.
By \eqref{eq:equilateral-image-convergence}
and Corollary~\ref{cor:uniform-separated-sets},
the cardinalities
$\Card(\SeparatedSet_{\SequenceIndex})$
are uniformly bounded.
This contradicts \eqref{eq:equilateral-cardinality-choice},
since
\[
 \Card(\SeparatedSet_{\SequenceIndex})
 =\EquilateralSize_{\FamilyIndex_{\SequenceIndex}}
 >\FamilyIndex_{\SequenceIndex}+1\longrightarrow\infty.
\]

\paraid{p-eb900edd6e7c}
Therefore the family
$\{\ApproximatingMap_{\FamilyIndex}(\ApproximationDomain_{\FamilyIndex})\}_{\FamilyIndex\in\DomainIndexSet}$
is locally finite.

\ProofStep{Discreteness.}\label{step:discrete-discreteness}
\paraid{p-0f7257809982}
A locally finite family of pairwise disjoint closed sets is discrete.
By Steps~\ref{step:discrete-separation} and~\ref{step:discrete-local-finiteness},
this fact applies to our compact images.
Local finiteness is automatic when
$\DomainIndexSet$
is finite.
\end{proof}

\subsection{Proof of Main Theorem}\label{subsec:approximation-hilbert-recognition}
\paraid{p-0fb4af10c68c}
We apply \Cref{thm:torunczyk} to
$\GHSpace$
using the AR property from Part~\ref{part:topology}
and the discrete approximation in \Cref{thm:discrete-approximation}.

\begin{proof}[Proof of \Cref{thm:part-iv-main}]
\paraid{p-a8eece7b38b0}
By \Cref{lem:gh-basics},
$\GHSpace$
is complete and separable,
and by \Cref{thm:part-iii-main},
it is a nonempty AR.
\Cref{thm:discrete-approximation} with
$\ApproximationDomain_\FamilyIndex=\HilbertCube$
for every
$\FamilyIndex\in\NonnegativeIntegers$
establishes the approximation condition in \Cref{thm:torunczyk}.
Thus
\Cref{thm:torunczyk} implies that
$\GHSpace$
is homeomorphic to
$\SeparableHilbertSpace$.
\end{proof}

\paraid{p-f04f1370b85d}
The proof of \Cref{thm:part-iv-main} also applies to the subspace
of spaces of diameter one.

\begin{corollary}\label{cor:unit-diameter-hilbert-space}
\paraid{p-673b86d6e522}
The subspace
\[
 \UnitDiameterGHSpace
 =\{\MetricSpace{\BaseCarrier}{\MetricSymbol}\in\GHSpace
     \mid\diam\MetricSpace{\BaseCarrier}{\MetricSymbol}=1\}
\]
of
$\GHSpace$
is homeomorphic to the real Hilbert space
$\SeparableHilbertSpace$.
\end{corollary}

\begin{proof}
\paraid{p-a048b3017172}
Since
$\UnitDiameterGHSpace$
is a nonempty closed subspace of
$\GHSpace$,
the space
$\UnitDiameterGHSpace$
is also  complete and separable.
Define
$\DiameterNormalization\colon\GHSpace\setminus\{\SingletonSpace\}\to\UnitDiameterGHSpace$
as follows.
For every
$\MetricSpace{\BaseCarrier}{\MetricSymbol}\in\GHSpace\setminus\{\SingletonSpace\}$,
put
\[
 \DiameterNormalization\MetricSpace{\BaseCarrier}{\MetricSymbol}
 =\MetricSpace{\BaseCarrier}
   {\MetricSymbol/\diam\MetricSpace{\BaseCarrier}{\MetricSymbol}}.
\]
By \Cref{lem:scaling-contraction},
this is a continuous retraction.
Namely,
$\UnitDiameterGHSpace$
is a retract of an open subset of the AR
$\GHSpace$
from \Cref{thm:absolute-extensor}.
This means that
$\UnitDiameterGHSpace$ is
also an ANR.

\paraid{p-e1df20d1720f}
Fix a two-point metric space
$\MetricSpace{D}{\delta}$
of diameter one.
Define
$\ContractionHomotopy\colon\UnitDiameterGHSpace\times\UnitInterval\to\UnitDiameterGHSpace$
as follows.
For every
$\MetricSpace{\BaseCarrier}{\MetricSymbol}\in\UnitDiameterGHSpace$
and every
$\HomotopyTime\in\UnitInterval$,
put
\[
 \ContractionHomotopy(\MetricSpace{\BaseCarrier}{\MetricSymbol},\HomotopyTime)
 =\DiameterNormalization\bigl(
    \MetricQuotient{\BaseCarrier}{(1-\HomotopyTime)\MetricSymbol}
    \times\MetricQuotient{D}{\HomotopyTime\delta}\bigr),
\]
using the maximum metric on the product,
whose diameter is
$\max\{1-\HomotopyTime,\HomotopyTime\}>0$.
\Cref{lem:scaling-contraction,lem:metric-products}
show that
$\ContractionHomotopy$
is continuous and contracts
$\UnitDiameterGHSpace$
to
$\MetricSpace{D}{\delta}$.
Hence
$\UnitDiameterGHSpace$
is an AR by \Cref{thm:contractible-ane,thm:retract-extensor}.

\paraid{p-fa3dc2186b0d}
Let
$\OpenCover_{0}$ be
 a family of  open sets of
 $\GHSpace$
 covering
 $\UnitDiameterGHSpace$
and consider
the open cover
$\OpenCover=\OpenCover_{0}\cup\{\GHSpace\setminus\UnitDiameterGHSpace\}$
of
$\GHSpace$.
Apply the construction in the proof of
\Cref{thm:discrete-approximation}
to maps into
$\UnitDiameterGHSpace$
and the cover
$\OpenCover$.
For every
$\MetricSpace{\BaseCarrier}{\MetricSymbol}\in\UnitDiameterGHSpace$
and every integer
$\EquilateralSize\geq2$,
the maps in Step~\ref{step:discrete-products} satisfy
$ \diam\bigl(\ProductApproximation_\EquilateralSize
             \MetricSpace{\BaseCarrier}{\MetricSymbol}\bigr)
 =\max\{1,\CoverScale\MetricSpace{\BaseCarrier}{\MetricSymbol}\}=1$
because
$\CoverScale\leq1/16$.
Thus the recursive choices and the discreteness argument in Steps~\ref{step:discrete-separation}--\ref{step:discrete-discreteness}
remain valid in
$\UnitDiameterGHSpace$.
The arguments on
closeness to the given cover and discreteness
are still valid for  this subspace.
\Cref{thm:torunczyk} now applies as in the proof of
\Cref{thm:part-iv-main}.
\end{proof}

\begin{remark}\label{rem:banach-homeomorphism}
\paraid{p-70d10e1e63b3}
Kadets proved that all separable infinite-dimensional Banach spaces are
homeomorphic \cite[Theorem]{Kadets1967}.
Consequently,
\Cref{thm:part-iv-main} implies that
$\GHSpace$
is homeomorphic to every real separable infinite-dimensional Banach space,
including
$\ContinuousFunctions([0,1],\RealNumbers)$
with the uniform norm and
$\ell^p$
for every
$1\leq p<\infty$.
\end{remark}

\paraid{p-6349dead9826}
The main theorem also implies strong local contractions.

\begin{corollary}\label{cor:strong-local-contractions}
\paraid{p-82c01acb09b2}
Every point
$\MetricSpace{\BaseCarrier }{\MetricSymbol }\in \GHSpace$
has a basis of open neighborhoods which strongly deformation retract
onto
$\{\MetricSpace{\BaseCarrier }{\MetricSymbol }\}$.
\end{corollary}

\begin{remark}\label{rem:nonarchimedean-gh}
\paraid{p-1f95d530606f}
The author determined the topology of the non-Archimedean
Gromov--Hausdorff space by identifying it isometrically with a function space
\cite[Theorems~1.1 and~1.3(3)--(4)]{Ishiki2024NonseparableUrysohn}.
Let
$\UltrametricGHSpace$
be the set of isometry classes of nonempty compact ultrametric spaces.
Its non-Archimedean Gromov--Hausdorff distance
$\NonArchimedeanGHDistance$
is the infimum of the Hausdorff distances between isometric images
in common ultrametric spaces.
Let
$\UrysohnFunctionModel$
consist of all functions
$f\colon[0,\infty)\to\NonnegativeIntegers$
such that
\[
 f(0)=0,
\]
and for every
$\varepsilon>0$,
\[
 \Card\bigl(\{r\geq\varepsilon\mid f(r)\ne0\}\bigr)<\infty.
\]
Equip
$\UrysohnFunctionModel$
with the ultrametric
\[
 \UrysohnFunctionDistance(f,g)
 =
 \begin{cases}
  \max\{r\in[0,\infty)\mid f(r)\ne g(r)\},&f\ne g,\\
  0,&f=g.
 \end{cases}
\]
Then the cited theorems imply that
\[
 (\UltrametricGHSpace,\NonArchimedeanGHDistance)
 \text{ is isometric to }
 (\UrysohnFunctionModel,\UrysohnFunctionDistance).
\]
\end{remark}


\section{Questions}\label{sec:questions}
\paraid{p-23b92a7e59a6}
We conclude with questions about related spaces and the geometry of GH balls.
We use the mm-isomorphism classes and box topology of
\Cref{def:box-distance}.

\begin{question}\label{question:mm-hilbert-space}
\paraid{p-854913de1991}
Is the space of mm-isomorphism classes of metric measure spaces,
equipped with the box topology,
homeomorphic to the real Hilbert space
$\SeparableHilbertSpace$?
\end{question}

\paraid{p-84551a16a309}
We recall the concentration topology to distinguish it from the box topology.
For measurable real functions
$f,g$
on
$\BoxParameterInterval$,
identified
up to equality almost everywhere,
define the \emph{Ky Fan distance} by
\[
 \KyFanDistance(f,g)=\inf\{\ErrorTolerance>0\mid
 \BoxLebesgueMeasure(\{s\in\BoxParameterInterval\mid
 |f(s)-g(s)|>\ErrorTolerance\})\leq\ErrorTolerance\}.
\]
Let
$\LipschitzObservables(X)$
be the real
$1$-Lipschitz
functions on
$X$.
For a parameter
$\BoxFirstParameter$
of
$X$,
put
$\BoxFirstParameter^*\LipschitzObservables(X)
=\{f\circ\BoxFirstParameter\mid f\in\LipschitzObservables(X)\}$.
Let
$\BoxFirstParameter$
and
$\BoxSecondParameter$
range over parameters of
$X$
and
$Y$,
respectively.
The \emph{observable distance} is
\[
 \ConcentrationDistance(X,Y)
 =\inf_{\substack{\BoxFirstParameter\text{ a parameter of }X\\\BoxSecondParameter\text{ a parameter of }Y}}
 \HausdorffDistance{\KyFanDistance}
 \bigl(\BoxFirstParameter^*\LipschitzObservables(X),
       \BoxSecondParameter^*\LipschitzObservables(Y)\bigr).
\]
For nonempty sets
$A,B$
of measurable real functions on
$\BoxParameterInterval$,
we use
\[
 \HausdorffDistance{\KyFanDistance}(A,B)
 =\max\left\{
   \sup_{f\in A}\inf_{g\in B}\KyFanDistance(f,g),
   \sup_{g\in B}\inf_{f\in A}\KyFanDistance(f,g)
 \right\}.
\]
This definition does not require compactness of
$A$
or
$B$.
It is finite because
$\KyFanDistance\leq1$.
The metric
$\ConcentrationDistance$
on
$\MMClassSpace$
induces the
\emph{concentration topology}
\cite[arXiv v1, Definitions~2.4, 2.13 and Theorem~2.14]{KazukawaNakajimaShioya2024}.

\paraid{p-f5124bb9fbcd}
For the concentration topology,
the analogous question has a negative answer.
Indeed,
Kazukawa,
Nakajima,
and Shioya prove that
$(\MMClassSpace,\ConcentrationDistance)$
is not a Baire space
\cite[arXiv v1, Theorem~1.3]{KazukawaNakajimaShioya2024},
whereas
$\SeparableHilbertSpace$
is a Baire space.

\paraid{p-8aae9293fb4b}
Nakajima and Shioya construct a compact metrizable space containing a dense
topological copy of
$\GHSpace$
\cite[arXiv v1, Main Theorem~1.2]{NakajimaShioyaCompactification}.
This leads to the following question.

\begin{question}\label{question:gh-compactification-hilbert-cube}
\paraid{p-ae64e7766f6b}
Is the compactification of
$\GHSpace$
constructed by Nakajima and Shioya homeomorphic to the Hilbert cube
$\HilbertCube$?
\end{question}

\paraid{p-35809a8f88d6}
The neighborhood basis in Corollary~\ref{cor:strong-local-contractions}
need not consist of GH balls.
We use the notation
$\GHball(\MetricSpace{\BaseCarrier}{\MetricSymbol},\Radius)$
for open Gromov--Hausdorff balls introduced in
\Cref{sec:preliminaries}.

\begin{question}\label{question:gh-balls-contractible}
\paraid{p-bcf67eb18b1e}
For every
$\MetricSpace{\BaseCarrier}{\MetricSymbol}\in\GHSpace$,
does there exist
$\Radius_0>0$
such that for every
$0<\Radius<\Radius_0$,
the ball
$\GHball(\MetricSpace{\BaseCarrier}{\MetricSymbol},\Radius)$
is contractible?
More strongly,
can each such ball be strongly deformation retracted onto
$\{\MetricSpace{\BaseCarrier}{\MetricSymbol}\}$?
\end{question}

\paraid{p-22ea8b8a03bb}
One can also ask for a contraction whose image lies in a larger ball
with a controlled radius.

\begin{question}\label{question:gh-linear-local-contraction}
\paraid{p-e1b22ae40390}
For every
$\MetricSpace{\BaseCarrier}{\MetricSymbol}\in\GHSpace$,
do there exist
$\ContractionFactor\geq1$
and
$\Radius_0>0$
such that for every
$0<\Radius<\Radius_0$,
there is a continuous map
$\ContractionHomotopy\colon
 \GHball(\MetricSpace{\BaseCarrier}{\MetricSymbol},\Radius)\times\UnitInterval
 \to
 \GHball(\MetricSpace{\BaseCarrier}{\MetricSymbol},\ContractionFactor\Radius)$
such that for every
$\MetricSpace{\ComparisonCarrier}{\ComparisonMetric}\in\GHball(\MetricSpace{\BaseCarrier}{\MetricSymbol},\Radius)$,
we have
$\ContractionHomotopy(\MetricSpace{\ComparisonCarrier}{\ComparisonMetric},0)=\MetricSpace{\ComparisonCarrier}{\ComparisonMetric}$,
and
$\ContractionHomotopy(\MetricSpace{\ComparisonCarrier}{\ComparisonMetric},1)=\MetricSpace{\BaseCarrier}{\MetricSymbol}$,
and for every
$\HomotopyTime\in\UnitInterval$,
$\ContractionHomotopy(\MetricSpace{\BaseCarrier}{\MetricSymbol},\HomotopyTime)=\MetricSpace{\BaseCarrier}{\MetricSymbol}$?
Can
$\ContractionFactor$
be chosen independently of
$\MetricSpace{\BaseCarrier}{\MetricSymbol}$,
with
$\Radius_0$
still allowed to depend on
$\MetricSpace{\BaseCarrier}{\MetricSymbol}$?
\end{question}

\paraid{p-c13b61a9c8ad}
To define the final compactification,
we use
$\MMClassSpace$
and
$\BoxDistance$
from \Cref{def:box-distance}.
For
$X,Y\in\MMClassSpace$,
choose representatives
$(A,d_A,\mu_A)$
and
$(B,d_B,\mu_B)$,
respectively.
Write
$Y\prec X$
if there exists a
$1$-Lipschitz map
$f\colon\supp\mu_A\to\supp\mu_B$
with
$f_*(\mu_A|_{\supp\mu_A})=\mu_B|_{\supp\mu_B}$.
This condition is independent of the representatives.
A pyramid is a nonempty box-closed subset
$\Pyramid\subset\MMClassSpace$
such that for every
$X\in\Pyramid$
and every
$Y\in\MMClassSpace$
with
$Y\prec X$,
we have
$Y\in\Pyramid$,
and for every
$X,Y\in\Pyramid$,
there exists
$Z\in\Pyramid$
such that
$X\prec Z$,
and
$Y\prec Z$
\cite[arXiv v1, Definitions~2.2 and~2.16]{KazukawaNakajimaShioya2024}.
Let
$\PyramidSpace$
be the set of pyramids.
Let
$\{\Pyramid_n\}_{n\in\NonnegativeIntegers}$
be a sequence in
$\PyramidSpace$
and let
$\Pyramid\in\PyramidSpace$.
For
$X\in\MMClassSpace$,
put
$\BoxDistance(X,\Pyramid_n)
 =\inf_{Y\in\Pyramid_n}\BoxDistance(X,Y)$.
The weak topology on
$\PyramidSpace$
is metrizable and is characterized by
$\Pyramid_n\to\Pyramid$
if and only if both of the following conditions hold.
For every
$X\in\Pyramid$,
$\BoxDistance(X,\Pyramid_n)\longrightarrow0$,
and for every
$X\in\MMClassSpace\setminus\Pyramid$,
$\liminf_{n\to\infty}\BoxDistance(X,\Pyramid_n)>0$.
The assignment
$X\mapsto\{Y\in\MMClassSpace\mid Y\prec X\}$
identifies the concentration space
$(\MMClassSpace,\ConcentrationDistance)$
with a dense subspace of the compact space
$\PyramidSpace$
\cite[arXiv v1, Definition~2.17 and Theorem~2.18]{KazukawaNakajimaShioya2024}.
Kazukawa,
Nakajima,
and Shioya prove that
$\PyramidSpace$
is contractible and locally path connected
\cite[arXiv v1, Theorems~1.4 and~1.6]{KazukawaNakajimaShioya2024}.
This is a different compactification from the one in
Question~\ref{question:gh-compactification-hilbert-cube}.

\begin{question}\label{question:pyramids-hilbert-cube}
\paraid{p-78e27525b559}
Is
$\PyramidSpace$
with the weak topology homeomorphic to the Hilbert cube
$\HilbertCube$?
\end{question}

\paraid{p-89bc50316911}
Corollary~\ref{cor:ghp-retract} embeds the Gromov--Hausdorff space
$\GHSpace$
as a retract of the measured space
$\MeasuredSpaces$.
We ask whether
$\MeasuredSpaces$
has the same homeomorphism type as
$\GHSpace$.

\paraid{p-734d1158d30a}
We equip
$\MeasuredSpaces$
with the Gromov--Hausdorff--Prokhorov distance
$\GHPDistance$
of \Cref{def:ghp-distance}.
The probabilities need not have full support.

\begin{question}\label{question:ghp-hilbert-space}
\paraid{p-42d66cf6a0c5}
Is
$(\MeasuredSpaces,\GHPDistance)$
homeomorphic to the real Hilbert space
$\SeparableHilbertSpace$?
\end{question}


\printbibliography
\end{document}
